% Generated by build_chapter.py from the GY1985 files. Do not edit by hand; edit the files and rebuild.
\chapter{The academic year 1985--86}\label{ch:1985}
The supplied year is dominated by the passage from string path integrals to
determinant lines and from oscillator calculations to Virasoro representation
theory. Feigin's sequence makes the partition combinatorics, boson--fermion
correspondence, vertex operators and critical dimension part of one programme;
Knizhnik and Beilinson then expose its geometric content on moduli of curves.
The companion covers all 26 requested catalogue records, with 260 mathematical
items. Record 665 is a heading and has no mathematical item. The records are
ordered by the manuscript, including 646 before 645 and 653 before 652.

\begin{center}
\begin{tabular}{@{}llr@{}}
\toprule
Class & Meaning & Items\\
\midrule
A & Soon published by the speaker or Gelfand & 4\\
B & Developed into established results by others & 14\\
C & Partly developed, with a specified remaining question & 2\\
D & No development found after an adequate search & 0\\
E & Already known or circulating before the talk & 217\\
F & Exact assertion or fate not identifiable & 23\\
\midrule
 & Total & 260\\
\bottomrule
\end{tabular}
\end{center}

These counts describe the mathematical occurrences, not the originality or
quality of the speakers. A single construction can recur in several talks.
Much of the seminar is deliberate exposition of recent work. First circulation
matters: Witten's ten-dimensional construction was already a May 1985 preprint,
and Connes's relevant Bourbaki lecture was given in June 1985, despite their
1986 publication dates. Conversely, a named announcement is not automatically
the complete later theorem: Vigneras's construction and the full local
\(GSp_4\) correspondence have different scopes
\citep[\S11 and Main Theorem]{gantakeda2011local}.

The most direct subsequent publications include Kapranov--Manin's twistor
construction, Ryzhikova's decision theorem for subspace-lattice identities,
and Beilinson--Manin's Mumford-form treatment
\citep[\S3]{kapranovmanin1986twistor}
\citep{ryzhikova1987decidability,beilinsonmanin1986mumford}.
Arnold's Abelian-integral question belongs to the first-period form of the
infinitesimal Hilbert-sixteenth problem; the later uniform bound solves that
precise question, not every version of the broader problem
\citep[Theorem~2]{binyamininovikovyakovenko2010zeros}.
The 1976 and 1984 Schrödinger-operator constructions in Novikov's talk were
already published. Their zero-energy Manakov evolution must be separated from
an all-energy Lax assertion and from the last, underspecified scattering slogan.

Two C items remain. The unrestricted arithmetic classification of the
automorphic spectrum, GP-1985-04, has major solutions for specified families,
but the notebook chooses no group or packet that would support a bounded
contribution here \citep[\S4]{arthur2018functoriality}.
GP-1985-42 concerns realization-space universality with both simpliciality and
dimension four fixed. General four-polytope universality and simplicial
universality in growing dimension do not settle that question
\citep[Main Theorem]{richtergebertziegler1995realization}
\citep[Theorem~2 and Conjecture~4]{adiprasitopadrol2017universality}.
The development proves an explicit continuous insertion section and contraction
for any prescribed history of facet stackings on a fixed simplex. This is a
certificate for a restricted, already understood family, with exact rational
checks; it does not identify every realization of the final combinatorial type
with that history or prove universality. The latest opened explicit open-status
evidence is from 2024, with a bounded search through November 2026 documented
separately. No D item is manufactured from a failed search or vague remark.

The most consequential corrections are often errors faithfully copied from the
notebook. Examples include missing oscillator normalization factors, the
hyperelliptic exception to products of one-forms, a reversed Beltrami wedge
sign, swapped Novikov--Veselov derivatives, and a Quillen-norm formula on the
wrong dual determinant line. The ledger distinguishes these emendations from
translation errors and uncertain readings. The appendix supplies complete
proofs or exact checks, including a smooth explicit counterexample to the
unrestricted no-Lax statement. Flatness of a determinant metric is also kept
distinct from a canonical global trivialization or a unique scale.

The previous index supplies earlier Baker-function, Clifford, lattice,
partition, theta and Virasoro occurrences; the reducibility question reuses
GP-1982-03. Contextual resemblance alone does not merge problem identifiers.
The supplied earlier index ends with the 1982 academic year, so no unseen
1983--84 coverage is implied. All 61 supplied scans are accounted for: the
last two, P001258--P001259, begin separate Kontsevich and Morozov talks, outside
the requested 26-record list. Their itemized supplement preserves the evidence
without misattribution. The conflicting catalogue and handwritten dates are
recorded explicitly; a plausible corrected date is not presented as a direct
reading.

\medskip
\noindent\textit{Reading the chapter.} Each item separates the notebook,
the precise reconstruction, the publication match and the dated status.
``Verified'' in the bibliography means that the cited URL or DOI was opened;
the item field and source audit distinguish reading the result from checking
metadata. Full-text access failures and unresolved identifications remain
visible. The problem index, corrections and reproducible checks follow the
occurrence index.

\GStalk{641}{1985-09-09}{M. M. Kapranov}{P001199-P001200}{Baker functions, spinors, and twistor transforms}
The talk joins the construction of line bundles from exponential principal parts to twistor descriptions of Yang--Mills equations. The manuscript distinguishes the newly generalized Baker-function construction from material inspired by Witten, Clifford algebra, and the four-dimensional Penrose--Ward correspondence. A. Goncharov supplies the spinor explanation. The ten-dimensional projective diagram is too compressed to support an additional theorem. Earlier related exposition occurs at 1975-374.4, 1975-379.7, 1977-432.15, 1978-455.6, and 1982-592.7; these are connections, not evidence for a recurring open problem.

\begin{GSitem}{1985-641.1}{A}{Baker functions with an exponential singularity}
\GSfield{Kind; source}{definition; \GSpage{P001199}}
\GSfield{The notes}{The heading is \ru{Функция Бейкера произв. мн-зия}; the notes write \(\varphi=e^{q(z)}\varphi _0\), with \(\varphi _0\) holomorphic, and require meromorphic behavior elsewhere. The English ``arbitrary set'' mistranslates the abbreviated Russian for manifold.}
\GSfield{Statement}{Let \(X\) be a complex manifold, \(D\) an effective Cartier divisor, and \(q\in H^0(X,\mathcal O_X(D)/\mathcal O_X)\). Choose local meromorphic lifts \(q_i\). A scalar Baker function of exponential type \(q\) is meromorphic on \(X\setminus |D|\), with \(e^{-q_i}\varphi\) meromorphically extendible across \(D\) on each chart. The notebook's holomorphic condition is a more restrictive local version. Allowed poles and normalization must be specified before asserting existence or uniqueness.}
\GSfield{Reconstruction}{The global gluing data and the distinction between holomorphic and meromorphic prefactors are supplied. No existence theorem for arbitrary prescribed data is attributed to the talk.}
\GSfield{Fate}{\citet[\S3.3]{kapranovmanin1986twistor} give this construction; the paper was received on 10 October 1985 and printed in 1986. Its abstract and full text were checked.}
\GSfield{Present status}{As of November 2026: Established construction; existence and dimensions require the specified line bundle and divisor data. This is not an unsolved existence claim.}
\GSfield{Assessment}{literal statement: correct with missing hypotheses; publication match: exact; verification: read the result; development: none; confidence: high}
\end{GSitem}
\begin{GSitem}{1985-641.2}{E}{Exponential principal parts and the Picard group}
\GSfield{Kind; source}{construction; \GSpage{P001199}}
\GSfield{The notes}{The displayed chain is \(H^0(\mathcal O(D)/\mathcal O)\xrightarrow{\delta}H^1(\mathcal O)\to H^1(\mathcal O^*)\).}
\GSfield{Statement}{For the preceding \(X,D\), the connecting map of \(0\to\mathcal O_X\to\mathcal O_X(D)\to\mathcal O_X(D)/\mathcal O_X\to0\), followed by the sheaf exponential, associates to \(q\) a line bundle with transition functions \(\exp(q_i-q_j)\). Changing the lifts changes this cocycle by a coboundary. This construction produces the image of these principal parts, not necessarily every element of \(\operatorname{Pic}X\).}
\GSfield{Reconstruction}{The cocycle and the limitation on surjectivity make explicit what the arrows mean.}
\GSfield{Fate}{\citet[\S3.1--2]{kapranovmanin1986twistor} explicitly describe the cohomological linearization as known, while employing it in their 1985--86 construction.}
\GSfield{Present status}{As of November 2026: Standard sheaf-cohomological construction; the generalized Baker-function application is the adjacent A item.}
\GSfield{Assessment}{literal statement: correct with missing hypotheses; publication match: exact; verification: read the result; development: none; confidence: high}
\end{GSitem}
\begin{GSitem}{1985-641.3}{E}{Super light rays of dimension one even and eight odd}
\GSfield{Kind; source}{construction; \GSpage{P001199}}
\GSfield{The notes}{``Now something inspired by Witten'': in \(\mathbb C^{10/16}\), consider lines of dimension \(1/8\). The even space has a nondegenerate symmetric form, and its chosen line is isotropic.}
\GSfield{Statement}{Write the complex supertranslation space as \(V^{10|16}\), whose odd space is a half-spin module \(S^+\) for \(\operatorname{Spin}(10,\mathbb C)\). A nonzero null vector \(v\) determines an eight-dimensional subspace \(\ker(c(v):S^+\to S^-)\). The corresponding supertranslation subsupergroup has dimension \(1|8\); its translates are the super light rays. Their bracket is the spin-equivariant symmetric map \(\operatorname{Sym}^2 S^+\to V\).}
\GSfield{Reconstruction}{The slash is interpreted as superdimension, not division. Clifford multiplication and the supertranslation bracket complete the omitted definition; ordinary affine linear subspaces alone do not encode this bracket.}
\GSfield{Fate}{Witten's preprint circulated in May 1985, as recorded in \citet[reference 52; \S4.2]{kapranovmanin1986twistor}; \citet{witten1986twistor} appeared in March 1986. Witten metadata and the explicit construction in Kapranov--Manin were read.}
\GSfield{Present status}{As of November 2026: An established supertwistor construction already circulating before this talk; the 1986 print date does not make it a new September announcement.}
\GSfield{Assessment}{literal statement: correct with missing hypotheses; publication match: partial; verification: read the result; development: none; confidence: high}
\end{GSitem}
\begin{GSitem}{1985-641.4}{E}{Spin(10), the isotropic-line stabilizer, and two eight-dimensional spinors}
\GSfield{Kind; source}{claim; \GSpage{P001199}; A. Goncharov}
\GSfield{The notes}{A. Goncharov explains that the sixteen-dimensional space is the spinor representation of \(SO(10)\), and that the stabilizer of an isotropic line leads to \(SO(8)\) and two invariant eight-dimensional subspaces.}
\GSfield{Statement}{The linear half-spin representation belongs to \(\operatorname{Spin}(10,\mathbb C)\). For an isotropic line \(\ell\subset V\), the quotient \(\ell^\perp/\ell\) is nondegenerate of dimension eight. Restriction to the semisimple factor \(\operatorname{Spin}(8,\mathbb C)\) of a chosen Levi subgroup of the stabilizer gives \(S^+|_{\operatorname{Spin}(8)}\simeq S_8^+\oplus S_8^-\), with both summands eight-dimensional. A Levi splitting must not be confused with a splitting preserved by the entire parabolic.}
\GSfield{Reconstruction}{The covering group and the Levi qualification repair the shorthand; no canonical pair of complementary subbundles under the full parabolic is asserted.}
\GSfield{Fate}{The classical Clifford-module framework is in \citet[\S\S3--5]{atiyahbottshapiro1964clifford}; the eight-dimensional null kernel used here is explicit in \citet[\S4.2]{kapranovmanin1986twistor}.}
\GSfield{Present status}{As of November 2026: Standard representation theory. The literal \(SO(10)\) assertion requires the spin double cover, or a projective representation.}
\GSfield{Assessment}{literal statement: false as written; publication match: partial; verification: read the result; development: none; confidence: high}
\end{GSitem}
\begin{GSitem}{1985-641.5}{E}{The complex Clifford algebra in ten dimensions}
\GSfield{Kind; source}{theorem; \GSpage{P001200}}
\GSfield{The notes}{The generators satisfy \(e_i e_k+e_k e_i=2\delta_{ik}I\); the notes give algebra dimension \(2^{10}\) and matrices of size \(2^5\). A marginal vector formula omits the factor two.}
\GSfield{Statement}{For a ten-dimensional complex vector space with nondegenerate symmetric form \(b\), use the convention \(uv+vu=2b(u,v)1\). Then \(\operatorname{Cl}(V,b)\cong M_{32}(\mathbb C)\), with vector-space dimension \(1024\). Its irreducible module has dimension \(32\), and the even subalgebra preserves its two half-spin summands of dimension \(16\).}
\GSfield{Reconstruction}{The even/odd distinction explains why the preceding superspace uses sixteen spinors whereas the full Clifford module has dimension thirty-two.}
\GSfield{Fate}{This is the classical complex Clifford classification; see \citet[\S4, Table 1]{atiyahbottshapiro1964clifford}, published in 1964.}
\GSfield{Present status}{As of November 2026: Established; only the normalization inconsistency in the margin needs correction.}
\GSfield{Assessment}{literal statement: correct with missing hypotheses; publication match: exact; verification: read the result; development: none; confidence: high}
\end{GSitem}
\begin{GSitem}{1985-641.6}{E}{Inner automorphisms yield a projective orthogonal action}
\GSfield{Kind; source}{claim; \GSpage{P001200}}
\GSfield{The notes}{An orthogonal transformation acts on the Clifford algebra; its automorphism is inner, and the notes conclude that this produces a representation of the orthogonal group.}
\GSfield{Statement}{The orthogonal action on \(M_{32}(\mathbb C)\) gives a homomorphism \(O(10,\mathbb C)\to\operatorname{PGL}_{32}(\mathbb C)\). Implementing matrices are determined only up to scalars, so innerness alone supplies a projective representation. Passing to \(\operatorname{Pin}(10)\), or to \(\operatorname{Spin}(10)\) above \(SO(10)\), supplies the usual linear spin representation.}
\GSfield{Reconstruction}{The scalar ambiguity is essential and is absent from the notebook's inference.}
\GSfield{Fate}{\citet[\S3]{atiyahbottshapiro1964clifford} construct the spin double cover and its Clifford-module action. This predates the talk by two decades.}
\GSfield{Present status}{As of November 2026: The projective statement is correct; the unqualified ordinary representation of \(SO(10)\) is false.}
\GSfield{Assessment}{literal statement: false as written; publication match: partial; verification: read the result; development: annotated; confidence: high}
\end{GSitem}
\begin{GSitem}{1985-641.7}{E}{Flatness along super light rays and ten-dimensional Yang--Mills constraints}
\GSfield{Kind; source}{claim; \GSpage{P001200}}
\GSfield{The notes}{The curvature is said to vanish on the indicated superlines, equivalently giving a generalized self-dual Yang--Mills equation.}
\GSfield{Statement}{For a holomorphic connection on ten-dimensional supertranslation space, flatness on every \(1|8\) super light ray means that the pullback curvature vanishes on the tangent distribution of each ray. Witten's correspondence relates the appropriate supersymmetric Yang--Mills constraints to bundles on the space of such rays that are trivial on the incidence fibers. This is a supersymmetric integrability statement, not the ordinary four-dimensional equation \(F=\pm *F\) transplanted to ten dimensions.}
\GSfield{Reconstruction}{The notes do not list the superconnection constraints or bundle triviality hypotheses. Accordingly the correspondence is identified, but a complete component-field theorem is not reconstructed as something stated in this talk.}
\GSfield{Fate}{The pre-September circulation of Witten's construction is documented above; see \citet[Introduction (f), \S4]{kapranovmanin1986twistor} and \citet{witten1986twistor}.}
\GSfield{Present status}{As of November 2026: The identifiable correspondence is established. The notebook is insufficient to specify exactly which component equations the speaker imposed.}
\GSfield{Assessment}{literal statement: correct with missing hypotheses; publication match: partial; verification: read the result; development: none; confidence: medium}
\end{GSitem}
\begin{GSitem}{1985-641.8}{F}{The twelve-dimensional projective diagram}
\GSfield{Kind; source}{remark; \GSpage{P001200}}
\GSfield{The notes}{A diagram places \(\mathbb C^{12}\setminus\{0\}\), a projective space with an overwritten superscript, and \(Q^{10}\) together; the last phrase mentions a covering supermanifold.}
\GSfield{Statement}{One elementary repair is forced: \((\mathbb C^{12}\setminus\{0\})/\mathbb C^*=\mathbb P^{11}\). A nonsingular quadric hypersurface therein has dimension ten. The manuscript does not determine the arrows, covering, superdimension, or intended relation between this quadric and the eight-dimensional quadric of null directions in \(\mathbb C^{10}\).}
\GSfield{Reconstruction}{The dimension check is ours. Identifying the diagram with a particular compactification or supertwistor space would require data not present on the page.}
\GSfield{Fate}{The explicit ten-dimensional correspondence in \citet[\S4.2]{kapranovmanin1986twistor} was compared, but it does not establish the identity of every arrow in this separate diagram.}
\GSfield{Present status}{As of November 2026: Unresolved identification of a diagram, not an open theorem. The projective superscript is retained as an ambiguity in the corrections ledger.}
\GSfield{Assessment}{literal statement: not determinable; publication match: thematic only; verification: unresolved lead; development: none; confidence: low}
\end{GSitem}
\begin{GSitem}{1985-641.9}{E}{Lines in projective three-space and null two-planes}
\GSfield{Kind; source}{construction; \GSpage{P001200}}
\GSfield{The notes}{The Radon--Penrose discussion identifies lines in \(\mathbb P^3\) with \(\operatorname{Gr}(2,4)\), and draws the two families of null two-planes in a tangent space.}
\GSfield{Statement}{For \(W\in\operatorname{Gr}(2,\mathbb C^4)\), \(T_W\operatorname{Gr}(2,4)\simeq\operatorname{Hom}(W,\mathbb C^4/W)\). The determinant is a quadratic form valued in a one-dimensional line, hence defines a conformal class. Its rank-one cone has two rulings by two-dimensional totally null linear spaces: maps with a fixed one-dimensional kernel, or maps with image in a fixed line. These give the two types of maximal null surfaces.}
\GSfield{Reconstruction}{The tangent-space description makes the diagram precise and explains why only a conformal class, rather than a preferred metric, is intrinsic.}
\GSfield{Fate}{This is the classical twistor incidence geometry used in \citet[Introduction (a)]{kapranovmanin1986twistor}; the four-dimensional instanton correspondence was published by \citet{atiyahward1977instantons} in 1977.}
\GSfield{Present status}{As of November 2026: Standard geometry; no claim of a new classification is present.}
\GSfield{Assessment}{literal statement: correct with missing hypotheses; publication match: partial; verification: read the result; development: none; confidence: high}
\end{GSitem}
\begin{GSitem}{1985-641.10}{E}{Rank two and two relative-position parameters}
\GSfield{Kind; source}{remark; \GSpage{P001200}}
\GSfield{The notes}{The Grassmannian is said to have rank two: pairs of its points require two numbers, whereas pairs of points on a sphere require one.}
\GSfield{Statement}{Equip \(\mathbb C^4\) with a positive Hermitian form and consider the compact homogeneous space \(U(4)/(U(2)\times U(2))\). A pair of two-planes is classified under simultaneous unitary transformations by its two principal angles \(0\leq\theta_1\leq\theta_2\leq\pi/2\), obtained as singular values \(\cos\theta_i\) of orthogonal projection from one plane to the other. This is the rank-two statement. Under arbitrary complex linear transformations, this continuous two-parameter assertion would be false.}
\GSfield{Reconstruction}{The compact metric and acting group are supplied; they are needed to make the comparison with a sphere correct. The classification follows directly from singular-value decomposition.}
\GSfield{Fate}{The notebook gives an elementary homogeneous-space observation, not a problem. The modern statement above is a direct linear-algebraic verification, rather than a publication match inferred from the twistor bibliography.}
\GSfield{Present status}{As of November 2026: Correct in the compact Hermitian interpretation; the source does not specify that interpretation.}
\GSfield{Assessment}{literal statement: correct with missing hypotheses; publication match: none found; verification: indirect; development: none; confidence: medium}
\end{GSitem}
\begin{GSitem}{1985-641.11}{E}{Curvature on alpha-surfaces and the Penrose--Ward transform}
\GSfield{Kind; source}{theorem; \GSpage{P001200}}
\GSfield{The notes}{Vanishing curvature on one family of null surfaces is equated with self-duality. The connection then gives a bundle, without a connection, on the parameter space of those surfaces.}
\GSfield{Statement}{On a suitable complex four-dimensional conformally flat domain, the vanishing of the curvature restricted to every alpha-plane is equivalent to the vanishing of one of its two Hodge components. Which component is called self-dual depends on the chosen orientation. If the relevant alpha-surfaces and incidence fibers satisfy the usual local simple-connectivity/triviality conditions, horizontal sections along them define a holomorphic bundle on twistor space, trivial on each twistor line. Conversely such a bundle reconstructs the connection up to gauge.}
\GSfield{Reconstruction}{Locality, triviality on twistor lines, and the orientation convention are necessary additions. Flatness on each alpha-surface does not imply vanishing of all curvature.}
\GSfield{Fate}{The correspondence predates the talk: \citet{atiyahward1977instantons}, with its complex formulation reviewed in \citet[Introduction (a)]{kapranovmanin1986twistor}. The latter exposition was read; the former publisher record was checked.}
\GSfield{Present status}{As of November 2026: Established correspondence. The statement here is local complex geometry; real instantons require additional reality and global conditions.}
\GSfield{Assessment}{literal statement: correct with missing hypotheses; publication match: partial; verification: read the result; development: none; confidence: high}
\end{GSitem}

\GStalk{642}{1985-09-09}{V. I. Arnold}{P001201}{Zeros of Abelian integrals and Petrov nonoscillation}
Arnold asks for bounds on zeros of Abelian integrals and reports Petrov's sharp cubic-Hamiltonian theorem. The page then moves through Chebyshev spaces, Picard--Fuchs equations, vanishing periods, and two Jacobian-algebra examples. The compactness and degree qualifications are essential. A general linear-independence sentence is false literally; it is not used to justify Petrov's theorem.

\begin{GSitem}{1985-642.1}{B}{Uniform bounds for zeros of polynomial Abelian integrals}
\GSfield{Kind; source}{problem; \GSpage{P001201}}
\GSfield{The notes}{For \(I(c)=\oint_{H=c}(P\,dx+Q\,dy)\), the opening question asks how many zeros it can have.}
\GSfield{Statement}{Given degree bounds \(\deg H\leq d\) and \(\max(\deg P,\deg Q)\leq m\), seek a finite bound depending only on \((d,m)\) for the total isolated zeros, with multiplicity, of the real period branches over all regular compact ovals of \(H\). Identically zero branches contribute no isolated zeros. Integration over an entire noncompact real level set is not intended.}
\GSfield{Reconstruction}{Degree bounds and regular oval branches are supplied to make the notebook's question definite. No sharp asymptotic conjecture, or claim concerning higher return-map variations, is imputed to Arnold.}
\GSfield{Fate}{\citet[Theorem 2 and \S1.2]{binyamininovikovyakovenko2010zeros} give a uniform double-exponential degree bound. The preprint appeared in August 2008; publication followed in 2010.}
\GSfield{Present status}{As of November 2026: The finite-uniform-bound question is solved in this precise sense. The cited theorem does not settle Hilbert's full sixteenth problem, cycles born from singular levels, or perturbations whose first Abelian integral vanishes identically.}
\GSfield{Assessment}{literal statement: correct with missing hypotheses; publication match: exact; verification: read the result; development: none; confidence: high; problem \hyperref[prob:GP-1985-01]{GP-1985-01}}
\end{GSitem}
\begin{GSitem}{1985-642.2}{B}{Petrov: the cubic Hamiltonian gives an n-dimensional Chebyshev space}
\GSfield{Kind; source}{theorem; \GSpage{P001201}}
\GSfield{The notes}{For \(H=y^2/2+P_3(x)\), \(\deg P,\deg Q\leq n\), Arnold states that the period space has dimension \(n\) and that a nonzero period has fewer than \(n\) zeros. The manuscript's personal-name wording is retained in the correction ledger; the theorem identifies the author as G. S. Petrov.}
\GSfield{Statement}{Let \(n\geq1\), and let \(H=y^2+x^3-x\). On its maximal interval of regular compact ovals, the real vector space \(\{\oint_{\gamma(t)}(P\,dx+Q\,dy):\deg P,\deg Q\leq n\}\) has dimension \(n\); every nonzero element has at most \(n-1\) zeros counted with multiplicity. The bound is attained. A real cubic potential with two distinct real critical points reduces to this model by affine changes and rescaling; without such a real oval annulus the stated real zero-count assertion has no domain.}
\GSfield{Reconstruction}{The normalized model, multiplicities, and nondegenerate real oval interval are supplied. The statement concerns the dimension of the image of the period map, not the much larger space of polynomial one-forms.}
\GSfield{Fate}{This is \citet[Theorem 1 and \S2.1, Russian pp.~46--47]{petrov1986elliptic}. The paper was received on 24 January 1985 and printed in 1986; the September talk therefore reports work already completed, by someone other than its speaker.}
\GSfield{Present status}{As of November 2026: A proved sharp result. Its restriction to cubic potentials must not be promoted to a Chebyshev assertion for all Hamiltonians.}
\GSfield{Assessment}{literal statement: correct with missing hypotheses; publication match: exact; verification: read the result; development: none; confidence: high}
\end{GSitem}
\begin{GSitem}{1985-642.3}{E}{Picard--Fuchs systems and Sturm-type nonoscillation}
\GSfield{Kind; source}{method; \GSpage{P001201}}
\GSfield{The notes}{The period vector is said to satisfy a differential equation, with a margin naming the Gauss--Manin connection. Arnold asks about a relation to oscillation theorems.}
\GSfield{Statement}{Over an interval of regular values of an algebraic family, transport a homology cycle continuously and integrate a basis of relative de Rham classes. These periods satisfy a linear Picard--Fuchs system with meromorphic coefficients. In the cubic elliptic case it reduces to a two-component system; scalar second-order equations and the Riccati equation for a ratio of periods permit zero estimates. Mere satisfaction of a linear ODE does not by itself imply the Chebyshev bound.}
\GSfield{Reconstruction}{The distinction between a general period equation and a special nonoscillation argument is supplied. The marginal question has no separate quantified conjecture.}
\GSfield{Fate}{\citet[\S2.1--2.3]{petrov1986elliptic} explicitly uses Sturm separation and a Riccati reduction. These classical tools were already available; the new sharp estimate is indexed separately.}
\GSfield{Present status}{As of November 2026: The specific connection to oscillation theory is verified in the contemporaneous proof. No unaddressed general theorem is extracted from the margin.}
\GSfield{Assessment}{literal statement: correct with missing hypotheses; publication match: partial; verification: read the result; development: none; confidence: high}
\end{GSitem}
\begin{GSitem}{1985-642.4}{E}{Chebyshev spaces versus arbitrary linearly independent functions}
\GSfield{Kind; source}{definition; \GSpage{P001201}}
\GSfield{The notes}{The page says that for linearly independent \(f_1,\ldots,f_N\), a linear combination has \(N-1\) zeros, and then defines Chebyshev spaces by an upper bound smaller than their dimension.}
\GSfield{Statement}{A real \(N\)-dimensional space of sufficiently differentiable functions on an interval is a Chebyshev space in the multiplicity convention used here if every nonzero function has at most \(N-1\) zeros counted with multiplicity. Linear independence alone gives no such upper bound: \(1\) and \(\sin t\) are independent on \((0,10\pi)\), but \(\sin t\) has nine interior zeros. Conversely, at any prescribed \(N-1\) distinct points, homogeneous linear algebra supplies a nonzero combination vanishing there; it does not give exactly that many zeros.}
\GSfield{Reconstruction}{\ours{} The counterexample and the corrected existential statement are supplied here. The notebook may have compressed this lower-bound observation, but its literal universal reading is false.}
\GSfield{Fate}{The precise multiplicity convention is stated in \citet[\S1]{petrov1986elliptic}; the elementary counterexample independently checks the missing logical distinction.}
\GSfield{Present status}{As of November 2026: The definition is standard. Linear independence and the Chebyshev property must remain separate.}
\GSfield{Assessment}{literal statement: false as written; publication match: partial; verification: read the result; development: annotated; confidence: high}
\end{GSitem}
\begin{GSitem}{1985-642.5}{E}{The kernel of restriction and integration on a level curve}
\GSfield{Kind; source}{construction; \GSpage{P001201}}
\GSfield{The notes}{The notes observe that exact differentials integrate to zero and that multiplication by \(H-c\) does not change the restricted form on \(H=c\).}
\GSfield{Statement}{For a smooth affine level curve \(C_c=\{H=c\}\), every closed cycle \(\gamma\) satisfies \(\int_\gamma dF=0\). Forms divisible by \(H-c\) restrict to zero, as do forms \(R\,dH\). Thus the period map factors through the algebraic de Rham quotient of the fiber. For the whole moving family, one instead uses forms modulo exact forms and multiples of \(dH\), with its \(\mathbb C[H]\)-module structure; \((H-c)\omega\) is not identically period-zero as the parameter varies.}
\GSfield{Reconstruction}{The final distinction repairs a possible confusion between a fixed fiber and the family. The notebook does not claim that periods over one chosen oval detect the entire de Rham quotient.}
\GSfield{Fate}{These are direct restriction and Stokes identities. Their role in polynomial-period reduction is visible in \citet[\S2.1]{petrov1986elliptic}; that paper is cited for the application, not as the origin of the identities.}
\GSfield{Present status}{As of November 2026: Established elementary facts. Additional hypotheses are required for any asserted identification of a relative module with all periods.}
\GSfield{Assessment}{literal statement: correct with missing hypotheses; publication match: partial; verification: read the result; development: none; confidence: high}
\end{GSitem}
\begin{GSitem}{1985-642.6}{E}{Jacobian algebra of x cubed plus y to the fourth}
\GSfield{Kind; source}{construction; \GSpage{P001201}}
\GSfield{The notes}{A displayed quotient for \(f=x^3+y^4\) and a lattice diagram record six monomials.}
\GSfield{Statement}{\oursComputed{} Over \(\mathbb C\), the local Jacobian algebra is \(\mathbb C[[x,y]]/(\partial_x f,\partial_y f)=\mathbb C[[x,y]]/(x^2,y^3)\). Its basis is \(1,y,y^2,x,xy,xy^2\), so its dimension is six. The singularity at the origin is isolated.}
\GSfield{Reconstruction}{The basis and isolatedness are checked directly from the displayed derivatives; no theorem about all Newton diagrams is inferred.}
\GSfield{Fate}{This is a standard Jacobian-algebra computation, in the classical Newton/Milnor setting of \citet{kouchnirenko1976polyedres}. The cited original was opened; the numerical conclusion here is independently verified rather than quoted from it.}
\GSfield{Present status}{As of November 2026: The value six is correct; the accompanying code verifies the standard monomial basis.}
\GSfield{Assessment}{literal statement: correct; publication match: thematic only; verification: indirect; development: annotated; confidence: high}
\end{GSitem}
\begin{GSitem}{1985-642.7}{E}{Jacobian algebra of x squared y plus y to the fourth}
\GSfield{Kind; source}{construction; \GSpage{P001201}}
\GSfield{The notes}{The second lattice sketch is attached to \(f=x^2y+y^4\). It does not explicitly state a dimension.}
\GSfield{Statement}{\oursComputed{} Here the Jacobian ideal is \((2xy,x^2+4y^3)\). The quotient has basis \(1,x,y,y^2,y^3\): \(xy=0\), \(x^2=-4y^3\), and multiplication by \(y\) gives \(y^4=0\). A Groebner basis has leading monomials \(xy,x^2,y^4\), proving independence and dimension five. The common derivative zero is isolated.}
\GSfield{Reconstruction}{The value five is our reconstruction and computation, not an omitted numeral claimed to have been read from the scan.}
\GSfield{Fate}{The example belongs to the same established Jacobian-algebra framework as the preceding item. \citet{kouchnirenko1976polyedres} supplies historical context; the exact basis is checked here directly.}
\GSfield{Present status}{As of November 2026: Computation verified algebraically and by the accompanying code; it is an annotation, not new singularity theory.}
\GSfield{Assessment}{literal statement: correct with missing hypotheses; publication match: thematic only; verification: indirect; development: annotated; confidence: high}
\end{GSitem}

\GStalk{643}{1985-09-09}{A. V. Zelevinsky}{P001202-P001203 (top)}{Jacobi triple product through binary words and partitions}
The correct expanded triple product is followed by an incorrect compressed Pochhammer version. Zelevinsky then gives a combinatorial proof through partitions, eventually constant binary words, a distinguished gap, and two sets of odd distances. The method is attributed in the notebook to Leibenzon; Arnold adds the Young-diagram interpretation. The proof is elementary exposition, not an unsolved partition identity. Related encodings occur at 1976-415.1 and 1981-548.1--2.

\begin{GSitem}{1985-643.1}{E}{Jacobi triple product}
\GSfield{Kind; source}{theorem; \GSpage{P001202}}
\GSfield{The notes}{The initial displayed formula is \(\sum_{r\in\mathbb Z}z^rq^{r^2}=\prod_{k\geq1}(1-q^{2k})(1+zq^{2k-1})(1+z^{-1}q^{2k-1})\).}
\GSfield{Statement}{The identity holds analytically for \(0<|q|<1\) and \(z\in\mathbb C^*\), with locally uniform convergence in \(z\), and formally as a \(q\)-series whose coefficients are Laurent polynomials in \(z\). The case \(q=0\) is obtained directly as the constant term.}
\GSfield{Reconstruction}{Convergence or the formal-series interpretation is supplied. The independent combinatorial construction is indexed in the following items.}
\GSfield{Fate}{The classical identity is recorded in \citet[20.5.9]{nist2026dlmf}; this checked reference is an exposition of a result long predating 1985.}
\GSfield{Present status}{As of November 2026: Established identity; the first display is correct.}
\GSfield{Assessment}{literal statement: correct; publication match: exact; verification: read the result; development: none; confidence: high}
\end{GSitem}
\begin{GSitem}{1985-643.2}{E}{Pochhammer notation and the three incorrect arguments}
\GSfield{Kind; source}{definition; \GSpage{P001202}}
\GSfield{The notes}{The notebook defines \((a;q)_\infty=\prod_{n\geq0}(1-aq^n)\) but rewrites the product as \((1;q^2)_\infty(-zq^{-1};q^2)_\infty(-z^{-1}q^{-1};q^2)_\infty\). The scan confirms the literal \(1\) and both negative exponents.}
\GSfield{Statement}{With the stated zero-based convention, the correct expression is \((q^2;q^2)_\infty(-zq;q^2)_\infty(-z^{-1}q;q^2)_\infty\). In particular \((1;q^2)_\infty=0\), so the recorded compressed formula cannot equal the theta series. This is an error in the notebook, not an error introduced by the transcription.}
\GSfield{Reconstruction}{The correction follows by substituting the definition into each of the three factors; it changes all three arguments and leaves the expanded product unchanged.}
\GSfield{Fate}{The standard notation and product are checked against \citet[17.8.1 and 20.5.9]{nist2026dlmf}.}
\GSfield{Present status}{As of November 2026: The emended expression is correct. The diplomatic reading remains documented separately.}
\GSfield{Assessment}{literal statement: false as written; publication match: exact; verification: read the result; development: annotated; confidence: high}
\end{GSitem}
\begin{GSitem}{1985-643.3}{E}{The reference to Ramanujan summation}
\GSfield{Kind; source}{claim; \GSpage{P001202}}
\GSfield{The notes}{An emphatic line says the triple product follows from Ramanujan's identity, without writing that identity.}
\GSfield{Statement}{A natural precise candidate is Ramanujan's bilateral summation \( {}_1\psi_1(a;b;q,z)=(q,b/a,az,q/(az);q)_\infty/(b,q/a,z,b/(az);q)_\infty\), valid away from poles for \(0<|q|<1\), \(|b/a|<|z|<1\). Here \({}_1\psi_1=\sum_{n\in\mathbb Z}(a;q)_n z^n/(b;q)_n\). A specialization and limit give the triple product.}
\GSfield{Reconstruction}{Identification with \({}_1\psi_1\) is plausible but not uniquely forced by the unnamed reference. The notes contain no limiting argument; none is attributed to Zelevinsky.}
\GSfield{Fate}{The summation and its triple-product specialization are the classical formulas \citet[17.8.1--2]{nist2026dlmf}.}
\GSfield{Present status}{As of November 2026: The mathematical relation is established; the intended Ramanujan identity remains a medium-confidence identification.}
\GSfield{Assessment}{literal statement: not determinable; publication match: partial; verification: read the result; development: none; confidence: medium}
\end{GSitem}
\begin{GSitem}{1985-643.4}{E}{Coefficient counting by two strict odd partitions}
\GSfield{Kind; source}{construction; \GSpage{P001202}}
\GSfield{The notes}{After dividing by \(\prod_{k\geq1}(1-q^{2k})\), the coefficient of \(z^rq^{r^2+2n}\) on the left is \(p(n)\). The right is encoded by two decreasing lists \(\alpha,\beta\) of positive odd integers.}
\GSfield{Statement}{Choose two finite strictly decreasing sequences \(\alpha_1>\cdots>\alpha_s>0\) and \(\beta_1>\cdots>\beta_t>0\), all entries odd. Their weight in the right-hand product is \(z^{s-t}q^{\sum\alpha_i+\sum\beta_j}\). Thus the coefficient counts pairs with \(s-t=r\) and total sum \(r^2+2n\). Euler's product on the left gives \(p(n)\). Strictness is required because each factor may be selected at most once.}
\GSfield{Reconstruction}{The product determines the meaning of the overwritten comparison signs in the beta list; it does not establish that every underlying handwritten stroke is a strict sign.}
\GSfield{Fate}{Euler's partition product and the triple product are classical; see \citet[Chapter I, \S1]{macdonald1995symmetric} and \citet[20.5.9]{nist2026dlmf}. The equality of the two counts is justified by the next bijection.}
\GSfield{Present status}{As of November 2026: Correct coefficient interpretation. The strictness qualification is mathematical, and the doubtful scan reading remains marked as such.}
\GSfield{Assessment}{literal statement: correct with missing hypotheses; publication match: partial; verification: read the result; development: none; confidence: high}
\end{GSitem}
\begin{GSitem}{1985-643.5}{E}{Eventually constant binary words and partition size}
\GSfield{Kind; source}{construction; \GSpage{P001202-P001203}}
\GSfield{The notes}{Leibenzon's construction uses bi-infinite zero--one words, zeros sufficiently far to the left and ones sufficiently far to the right. The number of linkages equals the integer being partitioned.}
\GSfield{Statement}{Take binary words with only finitely many inversions \((i,j)\), \(i<j\), \(a_i=1,a_j=0\), and the stated tails, modulo translation of their indexing. Reading a zero as a horizontal step and a one as a vertical step produces the boundary of a finite Young diagram. Inversions correspond to its boxes; consequently translation classes with \(n\) inversions are in bijection with partitions of \(n\).}
\GSfield{Reconstruction}{Taking words modulo translation, or equivalently choosing the unique charge-zero gap, is indispensable: otherwise every partition would have infinitely many shifted representatives.}
\GSfield{Fate}{The partition-boundary encoding is classical and consistent with \citet[Chapter I, \S1]{macdonald1995symmetric}, whose first edition appeared in 1979. The personal attribution to Leibenzon was not independently traced to a publication; that does not turn the elementary bijection into an open problem.}
\GSfield{Present status}{As of November 2026: The construction is valid after specifying the indexing convention. Earlier related beta-number descriptions occur at 1981-548.1--2.}
\GSfield{Assessment}{literal statement: correct with missing hypotheses; publication match: partial; verification: indirect; development: annotated; confidence: medium}
\end{GSitem}
\begin{GSitem}{1985-643.6}{E}{Gap charge and the two odd-distance sets}
\GSfield{Kind; source}{construction; \GSpage{P001203 (top)}}
\GSfield{The notes}{Each gap is assigned the number of ones to its left minus the number of zeros to its right. The resulting integers are said to run through \(\mathbb Z\). The odd-distance lists read off from a gap give the two partitions.}
\GSfield{Statement}{\ours{} For a fixed word class and gap, put \(r=s-t\), where \(s\) counts left-hand ones and \(t\) right-hand zeros. Moving the gap one place right increases \(r\) by one, regardless of the intervening bit; the tail conditions make the values unbounded both ways. Therefore each integer charge selects exactly one gap. Distances \(1,3,5,\ldots\) from that gap to left ones and right zeros define \(\alpha,\beta\), and conversely these finite sets reconstruct the word with its gap.}
\GSfield{Reconstruction}{The monotonicity proof and inversion of the encoding are explicit reconstructions of the notebook's compressed diagrams.}
\GSfield{Fate}{This is the elementary mechanism of the recorded proof; the standard Frobenius/partition framework appears in \citet[Chapter I, \S1]{macdonald1995symmetric}. No exact independent publication of this particular seminar attribution was found.}
\GSfield{Present status}{As of November 2026: Verified combinatorial bijection. It supplies the missing uniqueness needed by the coefficient count.}
\GSfield{Assessment}{literal statement: correct with missing hypotheses; publication match: partial; verification: indirect; development: annotated; confidence: high}
\end{GSitem}
\begin{GSitem}{1985-643.7}{E}{Young-diagram area equals half the odd sum minus charge squared}
\GSfield{Kind; source}{claim; \GSpage{P001203 (top)}; V. I. Arnold}
\GSfield{The notes}{Arnold draws the zero/right, one/up path and interprets the word as a Young-diagram boundary. The resulting relation is \(\sum\alpha_i+\sum\beta_j=r^2+2n\). A coordinate label \(x-y\) depends on the direction chosen for the vertical axis.}
\GSfield{Statement}{\ours{} Write \(\alpha_i=2a_i+1\), \(\beta_j=2b_j+1\). Counting inversions separately on the left, on the right, and across the gap gives \(n=\sum_i a_i+\sum_j b_j-s(s-1)/2-t(t-1)/2+st\). Hence \(\sum\alpha_i+\sum\beta_j=2n+(s-t)^2\). Together with the unique charge gap, this proves the coefficient equality and therefore the triple product.}
\GSfield{Reconstruction}{\oursComputed{} The accompanying independent finite enumerations check the inversion identity, the odd-part counts, and the product coefficients. The general inversion proof is the one supplied above. For an upward-positive coordinate axis, the monotonically changing gap coordinate is \(x+y\); \(x-y\) works with the downward-positive Young-diagram convention. No definite transcription error is asserted solely from the sketch.}
\GSfield{Fate}{The conclusion is the classical identity in \citet[20.5.9]{nist2026dlmf}; the finite counting verification included with this chapter checks the reconstructed bijection independently.}
\GSfield{Present status}{As of November 2026: The proof is complete under the explicit conventions. It is an annotation of an old identity, not a C/D development.}
\GSfield{Assessment}{literal statement: correct with missing hypotheses; publication match: partial; verification: indirect; development: annotated; confidence: high}
\end{GSitem}

\GStalk{644}{1985-09-16}{A. V. Zelevinsky, with A. A. Beilinson}{P001203-P001204 (middle)}{Reports on Vigneras, triple products, and trace Paley--Wiener}
The 16 September heading starts below the completed partition argument on P001203. It reports Vigneras's work and the Piatetski-Shapiro--Rallis triple product, then sketches the trace transform and unramified families. The separate Beilinson heading in the middle of P001204 belongs to this same meeting and concerns the quadratic Hermitian space and theta correspondence; it is not the 23 September Beilinson talk. Vigneras's intervening explanation at the top of P001204 has its own record 646.

\begin{GSitem}{1985-644.1}{B}{Vigneras: from GL(2) over a quadratic extension to GSp(4)}
\GSfield{Kind; source}{remark; \GSpage{P001203}}
\GSfield{The notes}{The first report is titled \ru{От \(GL_2(\mathbb Q(\sqrt d))\) к \(GSp_4(\mathbb Q)\)}. It mentions a lifting, its local counterpart, and Langlands functoriality.}
\GSfield{Statement}{The identifiable subject is transfer of suitable automorphic representations over a quadratic extension \(E/F\) toward symplectic similitude representations over \(F\). Here \(GSp_4(F)=\{g:g^tJg=\nu(g)J\}\). Locally the relevant four-dimensional induced parameters must preserve an alternating form up to a character. The notes give no central-character, cuspidality, or nonvanishing hypotheses, so they do not state a universal lift for every representation of \(GL_2(E)\).}
\GSfield{Reconstruction}{The qualifications describe the scope that any precise lifting theorem needs, not missing claims supplied to the speaker.}
\GSfield{Fate}{The report matches the title of \citet{vigneras1986correspondances}, published in 1986; its author bibliography was checked. \citet[\S11]{gantakeda2011local} identify and revisit the quadratic-extension case of that work.}
\GSfield{Present status}{As of November 2026: A specific historical construction is identified. Its full global hypotheses were not read in the inaccessible original, so no exact global theorem is claimed here.}
\GSfield{Assessment}{literal statement: not determinable; publication match: partial; verification: publisher metadata checked; development: none; confidence: medium}
\end{GSitem}
\begin{GSitem}{1985-644.2}{B}{The degree-eight triple-product L-function}
\GSfield{Kind; source}{remark; \GSpage{P001203}}
\GSfield{The notes}{The second report names Rallis and Piatetski-Shapiro and writes \(L(s,\pi_1\times\pi_2\times\pi_3)\), for three representations of \(GL_2\).}
\GSfield{Statement}{For cuspidal automorphic representations \(\pi_i\) of \(GL_2(\mathbb A_F)\), the unramified local factor uses the tensor product of their two-dimensional Satake parameters: \(L_v(s)=\det(1-q_v^{-s}(t_{1,v}\otimes t_{2,v}\otimes t_{3,v}))^{-1}\), of degree eight. An integral representation then addresses continuation and a functional equation, with compatible central characters in its initial setup. The page records an announcement, not a statement of all analytic hypotheses.}
\GSfield{Reconstruction}{The degree-eight factor is the exact meaning of the product notation; the completed ramified and archimedean factors are not reconstructed from the one-line report.}
\GSfield{Fate}{\citet[Introduction and \S3]{piatetskishapirorallis1987rankin} is the exact publication, received January 1986 and printed in 1987. Its introduction also records Garrett's communication of his integral construction in April 1985; this work already circulated before print.}
\GSfield{Present status}{As of November 2026: The triple-product theory is established. This item is not evidence that arbitrary automorphic tensor-product functoriality was proved at the seminar.}
\GSfield{Assessment}{literal statement: correct with missing hypotheses; publication match: partial; verification: read the result; development: none; confidence: high}
\end{GSitem}
\begin{GSitem}{1985-644.3}{E}{Hecke functions and trace characters}
\GSfield{Kind; source}{construction; \GSpage{P001203}}
\GSfield{The notes}{A discussion associated with Jacquet asks how to use the trace formula. For locally constant compactly supported \(f\), it writes \(\pi(f)=\int_G f(g)\pi(g)\,dg\) and seeks to describe \(\pi\mapsto\operatorname{tr}\pi(f)\).}
\GSfield{Statement}{Let \(F\) be a nonarchimedean local field of characteristic zero, \(G=\mathbf G(F)\) for a connected reductive group, and fix Haar measure. The Hecke algebra \(\mathcal H(G)=C_c^\infty(G)\) has convolution product. For a smooth admissible representation of finite length, \(\pi(f)\) has finite rank, so its trace exists. Additivity in exact sequences makes this a linear form on the Grothendieck group \(R(G)\). The problem is the image of \(\mathcal H(G)\to R(G)^*\), not a Fourier transform on a naive manifold of irreducibles.}
\GSfield{Reconstruction}{Admissibility, convolution, and the Grothendieck-group domain supply the missing analytic and algebraic hypotheses.}
\GSfield{Fate}{The construction is stated in \citet[\S1.1--1.2]{bernsteindelignekazhdan1986trace}; it was standard before the new image theorem. The latter is separately recorded at 1985-646.3.}
\GSfield{Present status}{As of November 2026: Standard character theory. The name written as ``Jacket'' is bibliographically identified as Jacquet, without altering the diplomatic record.}
\GSfield{Assessment}{literal statement: correct with missing hypotheses; publication match: exact; verification: read the result; development: none; confidence: high}
\end{GSitem}
\begin{GSitem}{1985-644.4}{E}{Parabolic induction and unramified-character families}
\GSfield{Kind; source}{construction; \GSpage{P001203}}
\GSfield{The notes}{The notes describe irreducible representations through induced representations from \(P=MU\), calling \(M\) the semisimple part, and varying a cuspidal representation by characters.}
\GSfield{Statement}{Here \(M\) is a reductive Levi factor, which generally has a center. Every irreducible smooth representation of \(G\) is a subquotient of normalized parabolic induction \(i_P^G(\rho)\) from a supercuspidal representation of a Levi, unique up to conjugacy as cuspidal support. Twisting \(\rho\) by unramified characters \(\chi\in X_{\rm nr}(M)\) produces algebraic families; neither all induced representations nor all their constituents are irreducible uniformly in \(\chi\).}
\GSfield{Reconstruction}{``Semisimple part'' is corrected to Levi factor, and subquotient replaces any implied equality with an induced representation.}
\GSfield{Fate}{\citet[\S2.1]{bernsteindelignekazhdan1986trace} recall this established Bernstein theory, citing the earlier central-algebra work.}
\GSfield{Present status}{As of November 2026: Established support theory. Reducibility loci explain why its parameter variety is not literally the irreducible dual.}
\GSfield{Assessment}{literal statement: false as written; publication match: partial; verification: read the result; development: none; confidence: high}
\end{GSitem}
\begin{GSitem}{1985-644.5}{F}{An unfinished comparison with the real and complex cases}
\GSfield{Kind; source}{remark; \GSpage{P001203}}
\GSfield{The notes}{The final short line mentions the real case, possibly the complex case, and a difficult abbreviated word followed by ``on all Cartans''.}
\GSfield{Statement}{No definite theorem can be extracted: the line does not specify a function space, transform, growth condition, or the meaning of the abbreviation. Real-group Paley--Wiener theorems would be relevant background, but do not identify this statement.}
\GSfield{Reconstruction}{No completion of the missing predicate is attempted.}
\GSfield{Fate}{Comparison with the precisely stated nonarchimedean theorem in \citet[\S1.2]{bernsteindelignekazhdan1986trace} and the primary Arthur publication list did not resolve the shorthand. The literature's thematic resemblance is not treated as an exact match.}
\GSfield{Present status}{As of November 2026: Unidentified fragment, with no determinate open question to develop.}
\GSfield{Assessment}{literal statement: not determinable; publication match: none found; verification: unresolved lead; development: none; confidence: low}
\end{GSitem}
\begin{GSitem}{1985-644.6}{B}{Local Langlands packets for GSp(4)}
\GSfield{Kind; source}{problem; \GSpage{P001204 (middle)}; A. A. Beilinson}
\GSfield{The notes}{Under Beilinson's heading an arrow runs from irreducible representations of \(GSp_4\) to Weil--Deligne representations, with finite fibers called \(L\)-packets. The displayed field \(\mathbb Q\) must be read in the local discussion as \(\mathbb Q_p\).}
\GSfield{Statement}{For a nonarchimedean local field \(F\) of characteristic zero, classify \(\operatorname{Irr}GSp_4(F)\) by a surjective finite-to-one map to admissible homomorphisms \(W'_F\to GSp_4(\mathbb C)\), modulo conjugacy, with the expected central-character and local-factor compatibilities. Here \(W'_F\) denotes the Weil--Deligne group. Its fibers are the local \(L\)-packets.}
\GSfield{Reconstruction}{The target dual group, local field, and compatibility requirement make precise the finite-fiber arrow. The notes alone do not establish exhaustion of all irreducibles.}
\GSfield{Fate}{\citet[Main Theorem, pp.~1842--1843]{gantakeda2011local} prove the correspondence for every such \(F\), including residue characteristic two; their preprint dates from June 2007. Their introduction explains the scope and exhaustion gap in the 1986 construction \citep{vigneras1986correspondances}, whose original was bibliographically checked.}
\GSfield{Present status}{As of November 2026: Solved for \(GSp_4\) in the stated characteristic-zero local setting. This is not a claim about an unrestricted global automorphic classification.}
\GSfield{Assessment}{literal statement: correct with missing hypotheses; publication match: exact; verification: read the result; development: none; confidence: high; problem \hyperref[prob:GP-1985-02]{GP-1985-02}}
\end{GSitem}
\begin{GSitem}{1985-644.7}{E}{Quadratic space of Hermitian two-by-two matrices}
\GSfield{Kind; source}{construction; \GSpage{P001204 (middle)}; A. A. Beilinson}
\GSfield{The notes}{The displayed matrices are \(\bigl(\begin{smallmatrix}x_1&x_2\\\bar x_2&x_3\end{smallmatrix}\bigr)\), and the determinant is a rational quadratic form. The invariant condition is abbreviated as \(x=\bar x\), although the display is Hermitian.}
\GSfield{Statement}{Let \(E/F\) be quadratic, \(\operatorname{char}F\ne2\), and let \(V=\{X\in M_2(E):X=\bar X^{\,t}\}\). Then \(\dim_FV=4\) and \(q(X)=\det X\) is a nondegenerate quadratic form over \(F\). The action \(X\mapsto gX\bar g^{\,t}\) of \(GL_2(E)\) satisfies \(q(gX\bar g^{\,t})=N_{E/F}(\det g)q(X)\), defining a homomorphism to \(GO(V,q)\).}
\GSfield{Reconstruction}{\oursComputed{} The missing transpose and the similitude multiplier are supplied by direct matrix calculation. No surjectivity on \(F\)-points is asserted.}
\GSfield{Fate}{The quadratic-extension orthogonal-similitude construction is used in the Vigneras lift \citep{vigneras1986correspondances} and reviewed in \citet[\S11]{gantakeda2011local}; the elementary determinant identity is checked directly here.}
\GSfield{Present status}{As of November 2026: Correct standard construction after the Hermitian condition is repaired.}
\GSfield{Assessment}{literal statement: false as written; publication match: partial; verification: read the result; development: annotated; confidence: high}
\end{GSitem}
\begin{GSitem}{1985-644.8}{B}{Weil representation and the similitude theta correspondence}
\GSfield{Kind; source}{method; \GSpage{P001204 (middle)}; A. A. Beilinson}
\GSfield{The notes}{The construction is said to use the Weil representation of \(GO_4\times GSp_4\).}
\GSfield{Statement}{For a four-dimensional quadratic space \(V\) and four-dimensional symplectic space \(W\) over a characteristic-zero local field, \(V\otimes W\) is symplectic. The usual oscillator construction restricts to the isometry dual pair \(O(V)\times Sp(W)\). With the tensor action, similitudes preserve the symplectic form only on \(R_0=\{(h,g):\nu_V(h)\nu_W(g)=1\}\). Extending the Weil representation to this subgroup and inducing to the relevant product produces the similitude theta representation. It is not a direct restriction of the fixed-character oscillator representation to an arbitrary full product of similitudes.}
\GSfield{Reconstruction}{The compatible-multiplier subgroup and induced extension repair the schematic product notation; cover and character choices are part of the oscillator construction.}
\GSfield{Fate}{The exact similitude formalism was developed beyond the original isometry construction; \citet[\S2, especially p.~1847]{gantakeda2011local} spell out \(R_0\), its extension, and induction. The 1986 report \citep{vigneras1986correspondances} is the contemporaneous application, verified bibliographically.}
\GSfield{Present status}{As of November 2026: Established method. The recorded phrase is informative shorthand, but not a complete definition of a representation.}
\GSfield{Assessment}{literal statement: correct with missing hypotheses; publication match: partial; verification: read the result; development: none; confidence: high}
\end{GSitem}

\GStalk{646}{1985-09-16}{M.-F. Vigneras}{P001204 (top)}{Bernstein support and the trace Paley--Wiener theorem}
Vigneras's short intervention continues the preceding trace-transform discussion. The two essential corrections are that the Bernstein variety parametrizes cuspidal support rather than individual irreducibles, and that the trace Paley--Wiener theorem requires a uniform compact-open invariance bound, in addition to algebraicity along induced families. The item is placed before record 645 because that is their order on the page.

\begin{GSitem}{1985-646.1}{E}{The Bernstein variety is a support space, not the irreducible dual}
\GSfield{Kind; source}{construction; \GSpage{P001204 (top)}}
\GSfield{The notes}{A disjoint union of unramified-character families modulo Weyl actions is equated with \(\operatorname{Irr}G\).}
\GSfield{Statement}{For a reductive \(p\)-adic group, let \(\Omega(G)\) be the set of conjugacy classes of supercuspidal pairs \((M,\rho)\). Its components are unramified-character tori, with finite stabilizer and relative Weyl quotients. Cuspidal support defines a surjective finite-to-one map \(\operatorname{Irr}G\to\Omega(G)\). At reducibility points different irreducible subquotients can have the same support; hence the displayed equality of sets is false in general.}
\GSfield{Reconstruction}{The quotient notation is interpreted as an algebraic support variety. No choice of one constituent is hidden in it.}
\GSfield{Fate}{\citet[\S2.1]{bernsteindelignekazhdan1986trace} recall the earlier Bernstein support theory and state explicitly that the map is finite-to-one.}
\GSfield{Present status}{As of November 2026: Standard theory, with the notebook's equality replaced by the support map.}
\GSfield{Assessment}{literal statement: false as written; publication match: exact; verification: read the result; development: none; confidence: high}
\end{GSitem}
\begin{GSitem}{1985-646.2}{E}{Regular trace functions along induced families}
\GSfield{Kind; source}{claim; \GSpage{P001204 (top)}}
\GSfield{The notes}{The trace is said to be algebraic in the unramified character and to vanish on all but finitely many families.}
\GSfield{Statement}{For \(f\in C_c^\infty(G)\), any Levi \(M\), and irreducible smooth \(\rho\) of \(M\), the function \(\chi\mapsto\operatorname{tr}(i_P^G(\rho\otimes\chi)(f))\) is regular on \(X_{\rm nr}(M)\). One compact open subgroup can be chosen so that \(f\) is bi-invariant; its trace then vanishes on irreducibles having no vectors fixed by that subgroup. Only finitely many Bernstein components can contribute.}
\GSfield{Reconstruction}{The stronger compact-open bound is made explicit; finiteness of components is a consequence, not a substitute for all compatibility conditions.}
\GSfield{Fate}{These necessary conditions are recorded in \citet[\S1.2 and \S2.3]{bernsteindelignekazhdan1986trace}, which treats them as immediate or established preliminaries to its new theorem.}
\GSfield{Present status}{As of November 2026: Correct necessary properties of the trace transform.}
\GSfield{Assessment}{literal statement: correct with missing hypotheses; publication match: exact; verification: read the result; development: none; confidence: high}
\end{GSitem}
\begin{GSitem}{1985-646.3}{B}{The precise trace Paley--Wiener image theorem}
\GSfield{Kind; source}{theorem; \GSpage{P001204 (top)}}
\GSfield{The notes}{The discussion identifies the trace-transform image using algebraicity and finite support in families; the preceding page attributes the theorem to Kazhdan and Bernstein.}
\GSfield{Statement}{A linear form \(\ell:R(G)\to\mathbb C\) is \(\ell([\pi])=\operatorname{tr}\pi(f)\) for some \(f\in\mathcal H(G)\) if and only if: (i) \(\chi\mapsto\ell(i_P^G(\sigma\otimes\chi))\) is regular for every standard Levi \(M\) and \(\sigma\in\operatorname{Irr}M\); and (ii) there is one compact open \(K\) such that \(\ell(E)=0\) whenever \(E^K=0\). The field and reductivity hypotheses are those of 1985-644.3.}
\GSfield{Reconstruction}{The exact Grothendieck-group statement and condition (ii) come from the publication. Merely specifying regular functions on the support variety would lose information at reducibility points and is not this theorem.}
\GSfield{Fate}{\citet[\S1.2, theorem]{bernsteindelignekazhdan1986trace} prove the equivalence. The paper includes Deligne as coauthor, was received on 26 August 1985, and appeared in 1986. It was therefore circulating at the time of the talk.}
\GSfield{Present status}{As of November 2026: Proved and correctly attributable to Bernstein--Deligne--Kazhdan. The seminar's short version omits a material hypothesis and one author.}
\GSfield{Assessment}{literal statement: correct with missing hypotheses; publication match: exact; verification: read the result; development: none; confidence: high}
\end{GSitem}

\GStalk{645}{1985-09-16}{B. L. Feigin}{P001204 (bottom)-P001205}{Worldsheets, area actions, and Polyakov integration}
Feigin describes string worldsheets and the passage from the induced-area functional to integration over a map and an independent metric. The formulas are the established Nambu--Goto and Polyakov constructions. Statements about all particles, a crystal-growth analogy, and an asserted operator/integral correspondence lack enough detail for theorem identification. The fixed-metric Gaussian step is precise; the full measure over surfaces is left formal. The continuation on 23 September begins only after Beilinson's conclusion on P001207.

\begin{GSitem}{1985-645.1}{E}{String histories as maps of surfaces}
\GSfield{Kind; source}{construction; \GSpage{P001204 (bottom)}}
\GSfield{The notes}{A closed string is represented by a circle, and its history sweeps out a surface mapped into space. The notes compare the construction with statistical mechanics on surfaces.}
\GSfield{Statement}{A parametrized closed string history is a map \(X:\Sigma\to\mathbb R^d\), with a cylindrical worldsheet in the simplest free propagation picture. Changes of the topology of \(\Sigma\) represent joining and splitting in perturbative string diagrams. The model requires a chosen target metric and an action; the topological picture alone determines no dynamics.}
\GSfield{Reconstruction}{A worldsheet is the domain of a map and need not be an embedded surface in the target. The cylinder and topology-changing interpretation clarify the drawings.}
\GSfield{Fate}{The world-surface and random-surface formulation is explicit in \citet[p.~207, equations (1)--(2)]{polyakov1981quantum}, received in May and published in July 1981.}
\GSfield{Present status}{As of November 2026: Established formulation; it is background for the action formulas below.}
\GSfield{Assessment}{literal statement: correct with missing hypotheses; publication match: partial; verification: read the result; development: none; confidence: high}
\end{GSitem}
\begin{GSitem}{1985-645.2}{F}{The claim about all particles and a ten-dimensional superspace}
\GSfield{Kind; source}{remark; \GSpage{P001204 (bottom)}}
\GSfield{The notes}{The notes suggest that all particles can be described through maps into a superspace denoted schematically by \(\mathbb R^{10/m}\).}
\GSfield{Statement}{The even dimension ten is recognizable from superstring theory, but the odd dimension \(m\), worldsheet fields, supersymmetry convention, and meaning of ``all particles'' are unspecified. No precise universal equivalence of particle theories is stated.}
\GSfield{Reconstruction}{No identification with a particular ten-dimensional superstring model is imposed.}
\GSfield{Fate}{The bosonic random-surface paper \citet{polyakov1981quantum} and Witten's ten-dimensional twistor paper \citet{witten1986twistor} are adjacent contexts, not exact sources for this unrestricted assertion.}
\GSfield{Present status}{As of November 2026: A programmatic remark whose mathematical content cannot be resolved further from the page; it is not indexed as an open conjecture.}
\GSfield{Assessment}{literal statement: not determinable; publication match: thematic only; verification: unresolved lead; development: none; confidence: low}
\end{GSitem}
\begin{GSitem}{1985-645.3}{F}{The crystal-growth analogy}
\GSfield{Kind; source}{remark; \GSpage{P001204 (bottom)}}
\GSfield{The notes}{A short comparison with crystal growth accompanies the picture of random surfaces.}
\GSfield{Statement}{No lattice, energy, probability measure, boundary condition, or limiting regime is specified. Thus no precise statistical-mechanical surface model or theorem can be reconstructed.}
\GSfield{Reconstruction}{The analogy is retained as an identifiable intervention in the exposition without assigning it a research problem.}
\GSfield{Fate}{Random surfaces do occur in the physical motivation of \citet[p.~207]{polyakov1981quantum}, but this does not identify the notebook's specific crystal-growth comparison.}
\GSfield{Present status}{As of November 2026: Unresolved model identification; no statement of a mathematical limit or universality class is justified.}
\GSfield{Assessment}{literal statement: not determinable; publication match: thematic only; verification: unresolved lead; development: none; confidence: low}
\end{GSitem}
\begin{GSitem}{1985-645.4}{E}{The Nambu--Goto area functional}
\GSfield{Kind; source}{definition; \GSpage{P001205}}
\GSfield{The notes}{The notebook writes \(L=\int(\det g_{ij})^{1/2}\,d\sigma\,d\tau\), with \(g\) induced by \(X\), followed by the explicit square-root formula in first derivatives.}
\GSfield{Statement}{For a smooth immersion \(X:\Sigma\to\mathbb R^d\) with Euclidean target metric, let \(h_{ab}=\partial_aX\cdot\partial_bX\). Its area action is \(S_{NG}[X]=T\int_\Sigma\sqrt{\det h}\,d^2\xi\), \(T>0\). In two coordinates, \(\det h=|X_\sigma|^2|X_\tau|^2-(X_\sigma\cdot X_\tau)^2\). For a timelike worldsheet in a Lorentzian target, the real area density instead uses \(\sqrt{-\det h}\).}
\GSfield{Reconstruction}{The Euclidean convention makes the recorded square root literal; the tension is an optional positive scale absent from the displayed notebook normalization.}
\GSfield{Fate}{The identical Euclidean area formula appears in \citet[equation (3)]{polyakov1981quantum}, as classical geometric background.}
\GSfield{Present status}{As of November 2026: Standard geometric action; no new variational theorem is asserted.}
\GSfield{Assessment}{literal statement: correct with missing hypotheses; publication match: exact; verification: read the result; development: none; confidence: high}
\end{GSitem}
\begin{GSitem}{1985-645.5}{E}{The formal sum over maps}
\GSfield{Kind; source}{definition; \GSpage{P001205}}
\GSfield{The notes}{A Feynman integral is written \(\int e^{\beta L(X)}\,DX\).}
\GSfield{Statement}{This is a formal functional-integral prescription. A Euclidean convention uses \(e^{-S[X]}\), while a real-time convention uses \(e^{iS[X]/\hbar}\). Thus the recorded plus sign is not itself a transcription error: it depends on \(\beta\). To define an actual integral one must specify a regulator, boundary data, and the treatment of reparametrizations; no translation-invariant Lebesgue measure \(DX\) on all smooth maps is supplied.}
\GSfield{Reconstruction}{The distinction between the symbolic prescription and a defined measure is explicit. No rigorous existence of the complete string path integral is attributed to this display.}
\GSfield{Fate}{\citet[equation (2) and the following discussion]{polyakov1981quantum} expressly treats the original surface sum as symbolic and develops a gauge-fixed prescription.}
\GSfield{Present status}{As of November 2026: Established formal notation. Its mathematical definition depends on the chosen finite-dimensional or field-theoretic regularization.}
\GSfield{Assessment}{literal statement: correct with missing hypotheses; publication match: partial; verification: read the result; development: none; confidence: high}
\end{GSitem}
\begin{GSitem}{1985-645.6}{F}{The asserted correspondence of operator and path-integral approaches}
\GSfield{Kind; source}{claim; \GSpage{P001205}}
\GSfield{The notes}{The page says there are approaches by integrals and by ordinary operators, and that a correspondence has now been established.}
\GSfield{Statement}{The sentence does not identify the operator algebra, state space, dimension, anomaly conditions, boundary observables, or theorem of equivalence. A finite-dimensional Gaussian integral/operator correspondence does not establish equality of full interacting string theories.}
\GSfield{Reconstruction}{No particular equivalence theorem is inserted into this unspecified announcement.}
\GSfield{Fate}{The opened Polyakov paper treats a path-integral formalism, but it does not supply a uniquely matching theorem for this sentence. Searches through that immediate bibliography did not identify what precise equivalence was being reported.}
\GSfield{Present status}{As of November 2026: The statement remains unidentifiable at theorem level. It is not an unsolved problem with recoverable quantifiers.}
\GSfield{Assessment}{literal statement: not determinable; publication match: none found; verification: unresolved lead; development: none; confidence: low}
\end{GSitem}
\begin{GSitem}{1985-645.7}{E}{Polyakov action with an independent worldsheet metric}
\GSfield{Kind; source}{definition; \GSpage{P001205}}
\GSfield{The notes}{The new integration variables are a pair \((X,g)\). The displayed action contains a metric-area term and a quadratic derivative term, with the first integration measure suppressed.}
\GSfield{Statement}{For a fixed smooth compact surface \(\Sigma\), independent positive Riemannian metric \(g\), and map \(X:\Sigma\to\mathbb R^d\), write \(S_P[X,g]=\lambda\int_\Sigma d\operatorname{vol}_g+(T/2)\int_\Sigma g^{ab}\partial_aX\cdot\partial_bX\,d\operatorname{vol}_g\). The second term is quadratic in \(X\) for fixed \(g\). With \(\lambda=0\), variation in \(g\) makes \(g\) conformal to the induced metric at nondegenerate critical points; a nonzero cosmological term changes that classical variational statement.}
\GSfield{Reconstruction}{The integration measures and coefficients are normalized explicitly. Classical equivalence with area is not asserted for arbitrary \(\lambda\), nor is quantum equivalence inferred from the change of variables.}
\GSfield{Fate}{See \citet[equations (6)--(7) and (14)]{polyakov1981quantum}; the independent-metric functional and the cosmological weight were published in 1981.}
\GSfield{Present status}{As of November 2026: Established formalism with the stated variational qualification. The notebook's missing measure is a harmless notational omission once supplied.}
\GSfield{Assessment}{literal statement: correct with missing hypotheses; publication match: exact; verification: read the result; development: none; confidence: high}
\end{GSitem}
\begin{GSitem}{1985-645.8}{E}{Integrate the map first: a Gaussian determinant}
\GSfield{Kind; source}{method; \GSpage{P001205}}
\GSfield{The notes}{The closing instruction is to fix \(g_{ij}\) and integrate over \(X\), since the action is quadratic.}
\GSfield{Statement}{On a closed connected Riemannian surface, expand each scalar component of \(X\) in Laplacian eigenfunctions. After removing the constant mode and retaining finitely many positive eigenvalues \(\lambda_1,\ldots,\lambda_N\), the integral is proportional to \(\prod_{j=1}^N\lambda_j^{-d/2}\), with the normalization depending on \(T\). A zeta-regularized continuum notation is \((\det'\Delta_g)^{-d/2}\); the prime records the omitted zero mode. This calculation leaves the metric integral and its gauge factors to be treated.}
\GSfield{Reconstruction}{The finite-mode integral supplies an exact meaning to the Gaussian step. The determinant notation is an explicit modern reconstruction, not a formula written on this page.}
\GSfield{Fate}{The procedure is the integration step following \citet[equation (14)]{polyakov1981quantum}. Its finite-dimensional justification is ordinary Gaussian integration.}
\GSfield{Present status}{As of November 2026: The fixed-metric calculation is standard; it does not by itself construct the remaining measure over metrics.}
\GSfield{Assessment}{literal statement: correct with missing hypotheses; publication match: partial; verification: read the result; development: none; confidence: high}
\end{GSitem}

\GStalk{647}{1985-09-23}{I. M. Gelfand}{P001206 (top)}{A central diagonal section of the four-dimensional cube}
Gelfand poses one elementary geometry problem, explicitly addressed to a ninth-grade audience. The answer is a regular octahedron; a complete coordinate verification is supplied as an annotation. The question is not an unresolved research problem.

\begin{GSitem}{1985-647.1}{E}{Which polyhedron is the diagonal-orthogonal central section of a four-cube?}
\GSfield{Kind; source}{problem; \GSpage{P001206 (top)}}
\GSfield{The notes}{\ru{Задача (для 9 кл.)}: what polyhedron results from cutting a four-dimensional cube through its center by a three-dimensional hyperplane perpendicular to a main diagonal?}
\GSfield{Statement}{\ours{} For the cube \([-1,1]^4\), the section is \(P=\{x:|x_i|\leq1,\ x_1+x_2+x_3+x_4=0\}\). Its vertices are exactly the six permutations of \((1,1,-1,-1)\). Set \(u_1=(1,1,-1,-1)\), \(u_2=(1,-1,1,-1)\), \(u_3=(1,-1,-1,1)\). They are pairwise orthogonal, each of length two, and \(P=\operatorname{conv}\{\pm u_1,\pm u_2,\pm u_3\}\). Hence \(P\) is a regular octahedron.}
\GSfield{Reconstruction}{The supplied proof is as follows: a vertex of the three-dimensional section needs at least three active cube-coordinate bounds; the fourth coordinate is then forced to \(\pm1\), with exactly two plus signs. Orthogonality identifies the resulting cross-polytope. There are twelve equal edges of length \(2\sqrt2\) and eight equilateral triangular faces.}
\GSfield{Fate}{This is an elementary classical-polytope exercise, not an announcement requiring a later research solution. The direct proof and the independently enumerated vertices and edge lengths establish the answer without claiming an unverified historical attribution.}
\GSfield{Present status}{As of November 2026: Completely solved in the stated setting. \oursComputed{} The accompanying independent exact-arithmetic enumeration confirms the vertices, edge lengths, and faces.}
\GSfield{Assessment}{literal statement: correct; publication match: none found; verification: indirect; development: annotated; confidence: high; problem \hyperref[prob:GP-1985-03]{GP-1985-03}}
\end{GSitem}

\GStalk{648}{1985-09-23}{A. A. Beilinson, with I. M. Gelfand interventions}{P001206-P001207 (upper)}{Adeles, automorphic representations, Satake parameters, and L-functions}
Beilinson introduces the automorphic setting for the announced work of Piatetski-Shapiro and Rallis. Gelfand explains p-adic expansions and points to Connes's Bourbaki lecture. The mathematical core passes from adeles and congruence quotients to restricted tensor products, spherical representations, dual groups, and Euler factors. Several handwritten shortcuts are false literally: the unrestricted product of completions, the punctured-sphere descriptions, and the suggestion that local representation lists alone determine the automorphic spectrum. The general occurrence problem is separated from the already established formalism.

\begin{GSitem}{1985-648.1}{E}{The automorphic L2 representation and its spectral decomposition}
\GSfield{Kind; source}{construction; \GSpage{P001206}}
\GSfield{The notes}{The opening task is to describe the action of the adelic group in \(L^2(G(\mathbb A)/G(\mathbb Q))\); the corresponding representations are called automorphic. The heading mentions the new Rallis--Piatetski-Shapiro work on functional equations.}
\GSfield{Statement}{For a connected semisimple group \(G/\mathbb Q\), use the equivalent left-quotient convention \(L^2(G(\mathbb Q)\backslash G(\mathbb A))\), with right translation. Its decomposition includes a discrete Hilbert sum and, in the noncompact case, continuous direct integrals induced from discrete automorphic data on Levi subgroups. The word automorphic also has a broader smooth moderate-growth definition; it is not synonymous with every Hilbert-space subrepresentation. For \(GL_n\), a reductive rather than semisimple example, one fixes a central character or removes the split center to obtain the usual finite-volume setting.}
\GSfield{Reconstruction}{Quotient convention, continuous spectrum, and the reductive example are clarified. The abstract spectral theorem is separated from the arithmetic occurrence problem below.}
\GSfield{Fate}{The pre-1985 theory is summarized precisely in \citet[Theorem 3.3 and p.~28]{arthur1981automorphic}; its lecture setting was 1980 and publication 1981.}
\GSfield{Present status}{As of November 2026: Established spectral formalism. It does not supply a complete arithmetic list of the cuspidal representations.}
\GSfield{Assessment}{literal statement: correct with missing hypotheses; publication match: partial; verification: read the result; development: none; confidence: high}
\end{GSitem}
\begin{GSitem}{1985-648.2}{E}{Gelfand: p-adic expansions and the expansion of one third}
\GSfield{Kind; source}{definition; \GSpage{P001206}; I. M. Gelfand}
\GSfield{The notes}{Gelfand writes finite base-\(p\) expansions, infinite expansions for rational numbers, the \(p\)-adic norm, and \(1/3=1+2+2^3+2^5+\cdots\) for \(p=2\). The coprimality condition in the norm formula is written \((a,m)=1\).}
\GSfield{Statement}{Every nonzero \(x\in\mathbb Q_p\) has a unique expansion \(x=\sum_{j=N}^{\infty}a_jp^j\), where \(N\in\mathbb Z\), \(a_j\in\{0,\ldots,p-1\}\), and \(a_N\ne0\); then \(|x|_p=p^{-N}\). Thus for \(a=p^nm\), the required condition is that \(m\) is a \(p\)-adic unit, not \((a,m)=1\). \oursComputed{} The displayed series equals \(1+2\sum_{j\geq0}4^j=1+2/(1-4)=1/3\) in \(\mathbb Q_2\), since \(|4|_2<1\).}
\GSfield{Reconstruction}{Digits, convergence, and the unit condition are made explicit. The infinite expansion is a completion statement, not a convergent real-number series.}
\GSfield{Fate}{These classical completion facts are checked in \citet[\S7, Theorem 7.23 and the following examples]{milne2020ant}; the source is a modern exposition, not the origin of the result.}
\GSfield{Present status}{As of November 2026: The one-third example is correct. \oursComputed{} For \(S_N=1+2\sum_{j=0}^{N-1}4^j\), exact arithmetic gives \(3S_N-1=2\cdot4^N\), with 2-adic valuation \(2N+1\); this independently verifies the partial sums. The norm\textquotesingle s unit condition is emended.}
\GSfield{Assessment}{literal statement: false as written; publication match: exact; verification: read the result; development: annotated; confidence: high}
\end{GSitem}
\begin{GSitem}{1985-648.3}{E}{The adele ring is a restricted product including the real place}
\GSfield{Kind; source}{definition; \GSpage{P001206}}
\GSfield{The notes}{The notes write adeles as \(\prod_p(\text{p-adic numbers})\).}
\GSfield{Statement}{The adele ring of \(\mathbb Q\) is \(\mathbb A_\mathbb Q=\mathbb R\times\prod'_p\mathbb Q_p\), where \(x_p\in\mathbb Z_p\) for all but finitely many primes. The restricted-product topology is used. Its finite part is \(\mathbb A_f=\prod'_p\mathbb Q_p\); the unrestricted product \(\prod_p\mathbb Q_p\) is a different object. The diagonal rational field is discrete in the full adele ring.}
\GSfield{Reconstruction}{Both the restriction and archimedean factor are missing from the recorded shorthand; neither is optional for the later quotient formula.}
\GSfield{Fate}{\citet[\S2, pp.~10--11]{arthur1981automorphic} gives exactly this adelic setup in a source preceding the talk.}
\GSfield{Present status}{As of November 2026: Standard definition. The notebook formula must be read as shorthand and is false as a literal equality.}
\GSfield{Assessment}{literal statement: false as written; publication match: exact; verification: read the result; development: none; confidence: high}
\end{GSitem}
\begin{GSitem}{1985-648.4}{E}{Principal congruence subgroups and divisibility of levels}
\GSfield{Kind; source}{definition; \GSpage{P001206}}
\GSfield{The notes}{The chain \(SL_2(\mathbb R)\supset SL_2(\mathbb Z)\supset\Gamma_n\) is accompanied by the condition that matrices are the identity modulo \(n\), and \(\Gamma_{n_1}\supset\Gamma_{n_2}\) when \(n_1\mid n_2\).}
\GSfield{Statement}{For \(N\geq1\), \(\Gamma(N)=\ker(SL_2(\mathbb Z)\to SL_2(\mathbb Z/N\mathbb Z))\). If \(N\mid M\), reduction modulo \(N\) factors through reduction modulo \(M\), hence \(\Gamma(M)\subseteq\Gamma(N)\). Levels form a directed set under divisibility, not merely their usual numerical order.}
\GSfield{Reconstruction}{The group and the direction of the inverse system are made explicit.}
\GSfield{Fate}{The definitions and indices are classical; see \citet[\S2, Example 2.23]{milne2017modular}.}
\GSfield{Present status}{As of November 2026: Correct elementary congruence-group statement.}
\GSfield{Assessment}{literal statement: correct; publication match: exact; verification: read the result; development: none; confidence: high}
\end{GSitem}
\begin{GSitem}{1985-648.5}{E}{Adelic quotient as the inverse limit of congruence quotients}
\GSfield{Kind; source}{construction; \GSpage{P001206}}
\GSfield{The notes}{The display identifies \(G(\mathbb A)/G(\mathbb Q)\) with an inverse limit of \(G(\mathbb R)/\Gamma_n\), following the \(SL_2\) example.}
\GSfield{Statement}{For \(G=SL_2\) and principal compact open subgroups \(K(N)\subset SL_2(\widehat{\mathbb Z})\), strong approximation gives \(G(\mathbb Q)\backslash G(\mathbb A)/K(N)\simeq\Gamma(N)\backslash G(\mathbb R)\). Taking the inverse limit over divisible levels recovers \(G(\mathbb Q)\backslash G(\mathbb A)\) as a topological space. Inverting elements gives the right-quotient convention of the notebook. For a general reductive group the finite-level quotient may have several arithmetic components, so the unqualified single-quotient formula need not hold.}
\GSfield{Reconstruction}{The statement is restricted to the explicit \(SL_2\) example; strong approximation is not silently assumed for every \(G\).}
\GSfield{Fate}{\citet[\S2, pp.~12--13]{arthur1981automorphic} explicitly uses strong approximation for \(SL_n\) and relates adelic level quotients to principal congruence quotients.}
\GSfield{Present status}{As of November 2026: Correct in the stated \(SL_2\) setting. A general-group version needs its component set and strong-approximation hypotheses.}
\GSfield{Assessment}{literal statement: correct with missing hypotheses; publication match: partial; verification: read the result; development: none; confidence: high}
\end{GSitem}
\begin{GSitem}{1985-648.6}{E}{Modular curves: cusps are not elliptic points}
\GSfield{Kind; source}{claim; \GSpage{P001206}}
\GSfield{The notes}{The notes call \(\mathfrak H/\Gamma_n\) a punctured sphere, say \(\mathfrak H/SL_2(\mathbb Z)\) is a sphere with two punctures, and correctly give three punctures for \(\Gamma_2\).}
\GSfield{Statement}{As a coarse Riemann surface, \(SL_2(\mathbb Z)\backslash\mathfrak H\simeq\mathbb C\): its compactification is a sphere with one cusp removed; it also has two orbifold points of orders two and three, which are not punctures. \(\Gamma(2)\backslash\mathfrak H\) is a sphere minus three cusps. Principal congruence curves are not all genus zero: the compactification of \(\Gamma(7)\backslash\mathfrak H\) has genus three (already level six has genus one).}
\GSfield{Reconstruction}{The distinction between the coarse curve, its orbifold structure, and removed cusps repairs both inaccurate sphere assertions.}
\GSfield{Fate}{\citet[Proposition 2.21, Theorem 2.22, Example 2.23]{milne2017modular} give the compactification, genus formula, and explicit level table, all read directly.}
\GSfield{Present status}{As of November 2026: The \(\Gamma(2)\) statement is correct; the other two literal descriptions are false. This is a notebook error rather than a newly discovered exception.}
\GSfield{Assessment}{literal statement: false as written; publication match: exact; verification: read the result; development: annotated; confidence: high}
\end{GSitem}
\begin{GSitem}{1985-648.7}{E}{Restricted tensor products of local representations}
\GSfield{Kind; source}{theorem; \GSpage{P001206}}
\GSfield{The notes}{The representation of \(G(\mathbb A)\) is written as a tensor product of its real and all \(p\)-adic factors.}
\GSfield{Statement}{An irreducible admissible automorphic representation factors as \(\pi\simeq\bigotimes'_v\pi_v\), with irreducible admissible local factors (using the appropriate archimedean category). At almost all finite places there is a chosen nonzero \(K_v\)-fixed vector, and the restricted tensor product is the span of tensors equal to those vectors outside a finite set. It is not the unrestricted algebraic tensor product of an infinite collection of spaces.}
\GSfield{Reconstruction}{The prime, almost-everywhere spherical condition, and specified reference vectors are added.}
\GSfield{Fate}{The standard factorization is recorded in \citet[Part I, \S1]{arthurgelbart1991lectures}, with its earlier source cited there, and already in \citet[p.~11]{arthur1981automorphic}.}
\GSfield{Present status}{As of November 2026: Established factorization. Knowing possible local factors is still insufficient to decide which products are automorphic.}
\GSfield{Assessment}{literal statement: correct with missing hypotheses; publication match: exact; verification: read the result; development: none; confidence: high}
\end{GSitem}
\begin{GSitem}{1985-648.8}{E}{Gelfand points to Connes on subfactors and knots}
\GSfield{Kind; source}{remark; \GSpage{P001206}; I. M. Gelfand}
\GSfield{The notes}{Gelfand notes that a Bourbaki lecture by A. Connes relates knot theory to factor isomorphisms.}
\GSfield{Statement}{The identifiable lecture connects finite-index subfactors, Hecke-type relations, braid-group representations and traces, leading to Jones's knot invariant. The notebook supplies no specific assertion that knot equivalence is classified by factor isomorphism.}
\GSfield{Reconstruction}{The literal phrase ``factor isomorphisms'' is not expanded into a nonexistent classification theorem.}
\GSfield{Fate}{The exact match is \citet{connes1986indice}, Bourbaki expos\'e 647, delivered in June 1985 and printed in the 1986 Ast\'erisque volume. The title, dated lecture, and opening mathematical explanation were read.}
\GSfield{Present status}{As of November 2026: A securely identified pre-talk report. Its print year alone would misclassify it as later development.}
\GSfield{Assessment}{literal statement: correct with missing hypotheses; publication match: partial; verification: read the result; development: none; confidence: high}
\end{GSitem}
\begin{GSitem}{1985-648.9}{C}{Which restricted products actually occur automorphically?}
\GSfield{Kind; source}{problem; \GSpage{P001206}}
\GSfield{The notes}{After the tensor-product formula, the page asks which representations occur in \(L^2(G(\mathbb A)/G(\mathbb Q))\), adding that not all do. A preceding sentence says it is enough to know representations of the \(p\)-adic groups.}
\GSfield{Statement}{For a connected semisimple \(G/\mathbb Q\), determine which irreducible admissible restricted products \(\pi=\bigotimes'_v\pi_v\) occur in the discrete automorphic spectrum, and their multiplicities. More generally, identify the inducing automorphic data governing the continuous spectrum. Local admissibility and almost-everywhere sphericality are necessary but do not decide global occurrence; the archimedean factor and global conditions matter.}
\GSfield{Reconstruction}{The discrete multiplicity formulation makes the question precise; it is distinguished from the abstract Eisenstein-series decomposition already known in 1985. No equivalence with every version of the Langlands conjectures is asserted.}
\GSfield{Fate}{\citet[Theorem 3.3]{arthur1981automorphic} supplies the established spectral structure. A major later partial classification is Arthur's orthogonal/symplectic theory: the 2013 monograph \citep{arthur2013endoscopic} was bibliographically checked, and the author explains its scope and the larger classification goal in \citet[\S4]{arthur2018functoriality}.}
\GSfield{Present status}{As of November 2026: A general arithmetic occurrence and multiplicity classification for every connected semisimple group is not supplied by these results. The reviewed primary literature establishes major group-specific classifications, while the unrestricted problem remains a broad program. The supplied earlier index contains related automorphy questions, including GP-1982-43, with different quantifiers; they are not merged.}
\GSfield{Assessment}{literal statement: correct with missing hypotheses; publication match: partial; verification: read the result; development: none; confidence: medium; problem \hyperref[prob:GP-1985-04]{GP-1985-04}}
\GSfield{Worth developing}{Worth pursuing only after fixing a group and a narrowly specified packet or multiplicity question; the unrestricted classification does not offer a plausible bounded contribution in this chapter.}
\end{GSitem}
\begin{GSitem}{1985-648.10}{E}{Spherical representations and unramified principal series}
\GSfield{Kind; source}{definition; \GSpage{P001206-P001207}}
\GSfield{The notes}{The simplest local representations have a vector fixed by a maximal compact subgroup; the notes then use induction from a triangular Borel and characters trivial on the compact part of a torus.}
\GSfield{Statement}{Let \(G/F\) be a split connected reductive group over a nonarchimedean local field, with a fixed split integral model and hyperspecial \(K=G(\mathcal O_F)\). An irreducible smooth representation is \(K\)-spherical if \(\pi^K\ne0\); then \(\dim\pi^K=1\). Every such representation is the unique spherical irreducible subquotient of normalized induction from an unramified character of a split maximal torus. The full induced representation can be reducible.}
\GSfield{Reconstruction}{Hyperspecial \(K\), normalization, and ``subquotient'' are essential. ``A maximal compact'' without specifying the unramified setting is too broad for the displayed parametrization.}
\GSfield{Fate}{\citet[\S6]{satake1963spherical} establishes the spherical framework, and \citet[Part I, \S1]{arthurgelbart1991lectures} explicitly describes the unramified principal-series interpretation.}
\GSfield{Present status}{As of November 2026: Established unramified representation theory, with the stated subgroup and group hypotheses.}
\GSfield{Assessment}{literal statement: correct with missing hypotheses; publication match: exact; verification: read the result; development: none; confidence: high}
\end{GSitem}
\begin{GSitem}{1985-648.11}{E}{Satake parameters as semisimple dual-group conjugacy classes}
\GSfield{Kind; source}{theorem; \GSpage{P001207 (upper)}}
\GSfield{The notes}{Unramified torus characters are identified with points of a dual torus; spherical representations are written as \(T^0(\mathbb C)/W\) and then as conjugacy classes of matrices.}
\GSfield{Statement}{For split \(G/F\) as above, \(X_{\rm nr}(T)=\operatorname{Hom}(T(F)/T(\mathcal O_F),\mathbb C^*)\simeq\widehat T(\mathbb C)\), after fixing a uniformizer. The normalized Satake isomorphism identifies the spherical Hecke algebra with \(\mathbb C[X_*(T)]^W\). Its characters correspond to \(\widehat T(\mathbb C)/W\), equivalently semisimple conjugacy classes in the complex dual group \(\widehat G\). These are not arbitrary matrix conjugacy classes until a representation of \(\widehat G\) is chosen.}
\GSfield{Reconstruction}{The Hecke-algebra step and semisimplicity explain the quotient and remove the ambiguous reference to matrices.}
\GSfield{Fate}{This is Satake's classical theorem, \citet[\S6, Theorem 3]{satake1963spherical}, published in 1963.}
\GSfield{Present status}{As of November 2026: Established classification of the unramified spherical spectrum.}
\GSfield{Assessment}{literal statement: correct with missing hypotheses; publication match: exact; verification: read the result; development: none; confidence: high}
\end{GSitem}
\begin{GSitem}{1985-648.12}{E}{The complex dual group and a chosen representation}
\GSfield{Kind; source}{definition; \GSpage{P001207 (upper)}}
\GSfield{The notes}{The notes introduce \(G^0\), call it the dual group, give \(GL_n\) as self-dual, and choose \(\rho:G^0\to GL_m(\mathbb C)\).}
\GSfield{Statement}{For a split connected reductive group with root datum \((X^*,\Phi,X_*,\Phi^\vee)\), its complex dual group \(\widehat G\) has datum \((X_*,\Phi^\vee,X^*,\Phi)\). In particular \(\widehat{GL_n}=GL_n(\mathbb C)\). A finite-dimensional algebraic representation \(\rho\) converts a semisimple Satake class into a well-defined matrix characteristic polynomial. For nonsplit groups one needs the full \(L\)-group with Galois action; the notes explicitly restrict to split groups here.}
\GSfield{Reconstruction}{The root datum definition is supplied, keeping the stated split hypothesis.}
\GSfield{Fate}{The pre-talk dual-group construction and resulting factors are presented in \citet[\S5, pp.~37--44]{arthur1981automorphic}.}
\GSfield{Present status}{As of November 2026: Standard Langlands formalism; no self-duality of every reductive group is implied by the \(GL_n\) example.}
\GSfield{Assessment}{literal statement: correct with missing hypotheses; publication match: exact; verification: read the result; development: none; confidence: high}
\end{GSitem}
\begin{GSitem}{1985-648.13}{E}{Unramified Euler factors and the partial automorphic L-function}
\GSfield{Kind; source}{definition; \GSpage{P001207 (upper)}}
\GSfield{The notes}{For almost all spherical local factors, the notes write \(L(\pi,\rho,s)=\prod_p\det(1-p^{-s}\rho(x_{\pi_p}))^{-1}\).}
\GSfield{Statement}{For an automorphic representation of a split reductive group over \(\mathbb Q\), choose a finite set \(S\) containing all ramified primes. At \(p\notin S\), let \(t_p\in\widehat G\) be the normalized Satake class and put \(L_p(s,\pi,\rho)=\det(1-p^{-s}\rho(t_p))^{-1}\). The displayed construction defines the partial Euler product \(L^S(s,\pi,\rho)=\prod_{p\notin S}L_p(s,\pi,\rho)\) in its domain of convergence. Bad-prime and archimedean factors are additional data for a completed function; analytic continuation and a functional equation do not follow merely from writing the product.}
\GSfield{Reconstruction}{The superscript \(S\) resolves the inconsistency between ``almost all'' and a product over every prime. No general analytic-continuation conjecture is introduced as a new question stated on this page.}
\GSfield{Fate}{The precise pre-1985 construction is \citet[\S5, pp.~43--44]{arthur1981automorphic}; \citet[\S1]{arthur2018functoriality} retains the partial-product convention and separates it from analytic conjectures.}
\GSfield{Present status}{As of November 2026: The definition is established. Its unrestricted analytic properties remain outside what the notebook asserts; particular triple-product results are indexed at 1985-644.2.}
\GSfield{Assessment}{literal statement: correct with missing hypotheses; publication match: exact; verification: read the result; development: none; confidence: high}
\end{GSitem}
\begin{GSitem}{1985-648.14}{E}{The trivial GL(1) example is the Riemann zeta function}
\GSfield{Kind; source}{remark; \GSpage{P001207 (upper)}}
\GSfield{The notes}{The final example chooses the identity representation \(GL_1\to GL_1\), then the trivial Satake parameter \(x=1\), obtaining \(\prod_p(1-p^{-s})^{-1}\).}
\GSfield{Statement}{For the trivial automorphic character of \(\mathbb A_\mathbb Q^*/\mathbb Q^*\) and the standard one-dimensional representation of \(\widehat{GL_1}\), the finite Euler product is \(\zeta(s)=\prod_p(1-p^{-s})^{-1}\) for \(\operatorname{Re}s>1\). Its usual completed expression adds \(\pi^{-s/2}\Gamma(s/2)\). This example uses a reductive group, extending the opening semisimple shorthand.}
\GSfield{Reconstruction}{The half-plane of convergence and distinction between finite and completed factors are added; the notebook only writes the finite product.}
\GSfield{Fate}{This is the classical Euler product, appearing as the basic \(GL_1\) case in \citet[Part I, \S1]{arthurgelbart1991lectures}.}
\GSfield{Present status}{As of November 2026: Established example; it does not constitute a proof of analytic properties for every higher-dimensional Euler product.}
\GSfield{Assessment}{literal statement: correct; publication match: exact; verification: read the result; development: none; confidence: high}
\end{GSitem}

\GStalk{649}{1985-09-23}{B. L. Feigin}{P001207 (lower)-P001209}{Strings: area, energy, and independent metrics}
The continuation begins halfway down P001207, after the preceding discussion of automorphic \(L\)-functions. Feigin passes from the area of an immersed surface to the Polyakov action on a map and an independent metric, then to the formal quotient by reparametrizations. Several shortcuts conflate the area Hessian with the fixed-metric energy Hessian and harmonicity with conformality. Gelfand's intervention about minimal surfaces is retained without silently identifying the uncertain historical attribution.

\begin{GSitem}{1985-649.1}{E}{Induced metric and the Nambu area}
\GSfield{Kind; source}{definition; \GSpage{P001207 (lower)}}
\GSfield{The notes}{For a map \(x:\mathcal E\to\mathbb R^d\), the notes put \(g_{ij}=(\partial_i x,\partial_j x)\), call this a Riemannian metric, and define \(\mathcal L(x)=\int_{\mathcal E}\sqrt{\det g}\,d\xi_1d\xi_2\). Both Euclidean and Minkowski target metrics are mentioned.}
\GSfield{Statement}{Let \(x:\Sigma\to\mathbb R^d\) be a smooth immersion of an oriented surface, with compact domain or convergent area integral. In a Euclidean target its pullback metric \(h=x^*\langle\ ,\ \rangle\) is positive definite and \(A(x)=\int_\Sigma d\operatorname{vol}_h\) is coordinate independent. In a Lorentzian target this same formula applies to a spacelike immersion. For a timelike world sheet the real area density instead uses \(\sqrt{|\det h|}\).}
\GSfield{Reconstruction}{Immersion, the signature of the induced metric, and integrability are supplied hypotheses. An arbitrary map may have a degenerate pullback, and one negative target direction does not imply a Riemannian world sheet.}
\GSfield{Fate}{This is the classical geometric starting point of \citet[eq.~(3)]{polyakov1981quantum}, published 23 July 1981.}
\GSfield{Present status}{As of November 2026: An established definition with essential signature and nondegeneracy qualifications; no new problem is posed.}
\GSfield{Assessment}{literal statement: false as written; publication match: partial; verification: read the result; development: none; confidence: high}
\end{GSitem}
\begin{GSitem}{1985-649.2}{E}{A symbolic sum over surfaces with marked target points}
\GSfield{Kind; source}{method; \GSpage{P001208}}
\GSfield{The notes}{The functional integral \(\int\exp(i\mathcal L/\hbar)\mathcal D x\) is taken over surfaces passing through prescribed points. Reparametrization invariance is stated.}
\GSfield{Statement}{For a specified surface \(\Sigma\), marked points \(p_a\), and target values \(y_a\), the constraint \(x(p_a)=y_a\) defines a space of parametrized maps. Area is invariant under diffeomorphisms preserving the marked data. The displayed oscillatory integral is a formal expression until a regularization, quotient, measure, and limiting prescription have been supplied.}
\GSfield{Reconstruction}{The distinction between fixed marked domain points and an unparametrized image merely containing the target points is not resolved in the notebook. Neither is the choice between fixed topology and a sum over topologies.}
\GSfield{Fate}{The formal surface-sum program is already explicit in \citet[eqs.~(1)--(3)]{polyakov1981quantum}.}
\GSfield{Present status}{As of November 2026: Identifiable exposition, not a defined probability or oscillatory measure theorem. Its missing choices preclude a single problem-index entry.}
\GSfield{Assessment}{literal statement: not determinable; publication match: thematic only; verification: read the result; development: none; confidence: high}
\end{GSitem}
\begin{GSitem}{1985-649.3}{E}{First variation of area and the minimal-surface equation}
\GSfield{Kind; source}{theorem; \GSpage{P001208}}
\GSfield{The notes}{The first variation is written \(\delta\mathcal L=-\int(\delta x,\Delta x)\), and \(\Delta x=0\) is called the minimal-surface equation.}
\GSfield{Statement}{For a smooth immersed surface in Euclidean space, use its induced metric \(h\) and the nonpositive Laplacian \(\Delta_h=\operatorname{div}_h\nabla_h\). If the variation field \(V\) is compactly supported away from the boundary, then \(\delta A(V)=-\int_\Sigma\langle V,\Delta_hx\rangle\,d\operatorname{vol}_h\). Thus stationarity for all such \(V\) is equivalent to vanishing mean-curvature vector, with \(\Delta_hx\) equal to its trace convention. A closed connected surface has no nonconstant Euclidean harmonic coordinate map.}
\GSfield{Reconstruction}{The measure, boundary condition, induced-metric dependence, and mean-curvature normalization are made explicit. This is not the linear harmonic-map equation for an independently fixed metric.}
\GSfield{Fate}{The same classical equation occurs in \citet[eq.~(5)]{polyakov1981quantum}.}
\GSfield{Present status}{As of November 2026: The geometric first-variation statement is correct with these hypotheses; the closed Euclidean case also explains why nontrivial saddles need additional data.}
\GSfield{Assessment}{literal statement: correct with missing hypotheses; publication match: exact; verification: read the result; development: none; confidence: high}
\end{GSitem}
\begin{GSitem}{1985-649.4}{F}{Gelfand on harmonic maps and the construction of minimal surfaces}
\GSfield{Kind; source}{remark; \GSpage{P001208}; I. M. Gelfand}
\GSfield{The notes}{Gelfand refers to Eells and to a group recorded as \ru{Лашоф}, in connection with constructing minimal surfaces.}
\GSfield{Statement}{A conformal harmonic immersion of a Riemann surface into Euclidean space has minimal image. The note does not specify a theorem, ambient manifold, topology, or construction attached to either name.}
\GSfield{Reconstruction}{Only the general harmonic/conformal implication can be formulated. The manuscript spelling is retained: replacing it by Lawson, Lashof, or a particular joint paper would be conjectural identification.}
\GSfield{Fate}{The named publication has not been identified. The nearby classical statement is documented in \citet[eqs.~(5)--(7)]{polyakov1981quantum}; that is not evidence for the personal attribution.}
\GSfield{Present status}{As of November 2026: Historical attribution unresolved; no orphan research question can be extracted from this sentence.}
\GSfield{Assessment}{literal statement: not determinable; publication match: none found; verification: unresolved lead; development: none; confidence: low}
\end{GSitem}
\begin{GSitem}{1985-649.5}{E}{The proposed second variation of area}
\GSfield{Kind; source}{claim; \GSpage{P001208}}
\GSfield{The notes}{At a critical surface the notes assert \(\delta^2\mathcal L=-\int(\Delta\delta x,\delta x)\), followed by a quadratic Taylor expansion.}
\GSfield{Statement}{For a Euclidean minimal immersion and a compactly supported normal variation \(V\), the area Hessian is \(\int_\Sigma(|\nabla^\perp V|^2-\sum_{a,b}\langle A_{ab},V\rangle^2)\,d\operatorname{vol}_h\), where \(A\) is the second fundamental form and the tangent frame is orthonormal. For a fixed metric \(g\), the Hessian of \(\frac12\int|dx|_g^2\) is instead \(-\int\langle V,\Delta_gV\rangle\,d\operatorname{vol}_g\). A Taylor expansion uses one half of the Hessian and a remainder of order smaller than the quadratic term.}
\GSfield{Reconstruction}{The notebook conflates two functionals. An elementary check already disproves its unrestricted claim: on a planar disk, take a nonzero smooth tangential field vanishing near the boundary and a sufficiently small reparametrization \(x_t(u,v)=(u+t f(u,v),v)\). Its image and area are fixed, whereas the displayed Dirichlet quadratic form is \(\int|\nabla f|^2>0\). This is an editorial countercheck, not a new result attributed to Feigin. Editorial area-versus-energy countercheck \ours.}
\GSfield{Fate}{The distinction between varying the induced metric and an independent metric is explicit in \citet[eqs.~(3),(6)--(7)]{polyakov1981quantum}.}
\GSfield{Present status}{As of November 2026: False literally. The classical Jacobi formula and the fixed-metric energy formula are established; no C/D problem survives the correction.}
\GSfield{Assessment}{literal statement: false as written; publication match: partial; verification: read the result; development: annotated; confidence: high}
\end{GSitem}
\begin{GSitem}{1985-649.6}{E}{Dirichlet energy and independent variation of the metric}
\GSfield{Kind; source}{construction; \GSpage{P001208-P001209}}
\GSfield{The notes}{The area expression is replaced by \(\mathcal L'(x,g)=\int g^{ij}(\partial_i x,\partial_jx)\sqrt g\,d^2\xi\), and on the next page \(x\) and \(g\) are declared independent variables.}
\GSfield{Statement}{On a smooth compact surface, with fixed boundary values if there is boundary, set \(E_g(x)=\frac12\int g^{ij}\partial_i x\cdot\partial_jx\,d\operatorname{vol}_g\). Map variation gives \(\Delta_gx=0\). Metric variation gives \(T_{ij}=\partial_i x\cdot\partial_jx-\frac12g_{ij}|dx|_g^2=0\). For an immersion into a Euclidean target, \(T=0\) means \(x^*\langle\ ,\ \rangle=\lambda g\), \(\lambda>0\). On such pairs \(E_g=A(x)\); the notebook's normalization without \(1/2\) gives \(2A(x)\).}
\GSfield{Reconstruction}{The stress-tensor equation and the factor of two repair the omitted logical step. Weyl invariance in two dimensions leaves the scalar \(\lambda\) undetermined.}
\GSfield{Fate}{This is Polyakov's already published independent-metric reformulation \citep[eqs.~(6)--(7)]{polyakov1981quantum}.}
\GSfield{Present status}{As of November 2026: The classical equivalence is established for nondegenerate critical maps, with the stated boundary conditions and normalization.}
\GSfield{Assessment}{literal statement: correct with missing hypotheses; publication match: exact; verification: read the result; development: none; confidence: high}
\end{GSitem}
\begin{GSitem}{1985-649.7}{E}{Harmonicity does not force conformality}
\GSfield{Kind; source}{claim; \GSpage{P001208}}
\GSfield{The notes}{The text says that at an extremum of the fixed-metric functional the induced metric is conformal to the chosen metric.}
\GSfield{Statement}{Harmonicity \(\Delta_gx=0\) by itself does not imply conformality. On a planar domain with \(g=du^2+dv^2\), \(x(u,v)=(2u,v)\) is a harmonic immersion, but \(x^*g_{\mathbb R^2}=4du^2+dv^2\) is not conformal to \(g\). Conformality follows after also imposing the metric Euler--Lagrange equation \(T=0\).}
\GSfield{Reconstruction}{This counterexample fixes the boundary values to those of \(x\), so it is an actual critical point of the Dirichlet problem. It checks both the hypothesis and the failed conclusion, independently of the area-Hessian countercheck. Editorial harmonic nonconformal counterexample \ours.}
\GSfield{Fate}{The missing second variational equation is present in \citet[eq.~(7)]{polyakov1981quantum}.}
\GSfield{Present status}{As of November 2026: The literal implication is false; its corrected version is classical. No open problem is created by the omission.}
\GSfield{Assessment}{literal statement: false as written; publication match: partial; verification: read the result; development: annotated; confidence: high}
\end{GSitem}
\begin{GSitem}{1985-649.8}{E}{The weak L2 metric on the space of metrics}
\GSfield{Kind; source}{definition; \GSpage{P001209}}
\GSfield{The notes}{The notebook defines \(\|\theta\|_g^2=\int\theta_{ij}\theta_{k\ell}g^{ik}g^{j\ell}\sqrt g\,d\xi_1d\xi_k\) and calls \(\mathcal Dg\) its formal volume form.}
\GSfield{Statement}{For a closed smooth surface \(\Sigma\), the space of positive smooth symmetric 2-tensors has tangent space \(\Gamma(S^2T^*\Sigma)\). Its positive weak Riemannian pairing is \(G_g(\theta,\psi)=\int_\Sigma\operatorname{tr}(g^{-1}\theta g^{-1}\psi)\,d\operatorname{vol}_g\). The coordinate measure is \(d\xi_1d\xi_2\). This infinite-dimensional pairing does not, by itself, define a locally finite Lebesgue-type volume measure.}
\GSfield{Reconstruction}{The last differential is an index error in the notebook. The distinction between a weak metric and an actual measure is essential, not an extra choice of normalization.}
\GSfield{Fate}{The metric-measure heuristic belongs to the earlier gauge-fixed path integral \citep[eqs.~(14),(17)--(20)]{polyakov1981quantum}.}
\GSfield{Present status}{As of November 2026: The weak metric is well defined; its formal volume notation remains formal.}
\GSfield{Assessment}{literal statement: correct with missing hypotheses; publication match: partial; verification: read the result; development: none; confidence: high}
\end{GSitem}
\begin{GSitem}{1985-649.9}{E}{Diffeomorphism redundancy and Faddeev--Popov gauge fixing}
\GSfield{Kind; source}{method; \GSpage{P001209}}
\GSfield{The notes}{The integrand is constant on diffeomorphism orbits, the integral is said to be undefined, and the \ru{трюк Фаддеева--Попова} is proposed.}
\GSfield{Statement}{The action on pairs \((x,g)\) is invariant under \((x,g)\mapsto(x\circ f,f^*g)\). Locally, gauge fixing replaces orbit directions by a Jacobian built from the infinitesimal action, together with residual automorphism and moduli factors. This is a prescription for a regularized functional integral; it is not justified by literally dividing two infinite volumes.}
\GSfield{Reconstruction}{A slice, treatment of stabilizers, and determinant regularization are additions needed to formulate the method. The notebook does not prove existence of a quotient measure.}
\GSfield{Fate}{The concrete string application is in \citet[eqs.~(17)--(23)]{polyakov1981quantum}.}
\GSfield{Present status}{As of November 2026: Established formal gauge-fixing machinery, with rigorous finite-dimensional and elliptic ingredients distinguished from the full path integral.}
\GSfield{Assessment}{literal statement: not determinable; publication match: exact; verification: read the result; development: none; confidence: high}
\end{GSitem}

\GStalk{650}{1985-09-30}{B. L. Feigin}{P001210-P001212}{Spectral determinants, the conformal anomaly, and dimension 26}
The recorder explicitly arrived half an hour late. The surviving notes concern zeta determinants of Dolbeault Laplacians on tensor powers of the tangent bundle, heat regularization, the polynomial \(6j^2+6j+1\), and cancellation of the bosonic string anomaly at target dimension \(26\). The argument is expository: Polyakov's paper appeared in 1981. Real versus complex determinant conventions, zero modes, curvature normalization, and the distinction between a classical Liouville equation and a quantum integral require separate corrections.

\begin{GSitem}{1985-650.1}{E}{Dolbeault Laplacians and primed zeta determinants}
\GSfield{Kind; source}{definition; \GSpage{P001210}}
\GSfield{The notes}{For a complex curve with a compatible metric, \(T\) is its holomorphic tangent bundle, \(T^j\) its tensor power, and \(\Delta_j=\bar\partial^*\bar\partial\). The notes write \(\zeta(s)=\sum\lambda_i^{-s}\) and \(\log\operatorname{Det}\Delta=-\zeta'(0)\).}
\GSfield{Statement}{Let \(\Sigma\) be a closed Riemann surface with a smooth positive metric, let \(j\in\mathbb Z\), and use the induced Hermitian metric on \(T^j\). The positive elliptic operator \(\Delta_j\) has discrete spectrum. Set \(\zeta_j(s)=\sum_{\lambda>0}\lambda^{-s}\), initially for \(\Re s>1\), and \(\det'\Delta_j=\exp[-\zeta'_j(0)]\) by its meromorphic continuation, regular at zero. The prime excludes \(H^0(T^j)\).}
\GSfield{Reconstruction}{Compactness, integer tensor powers, the Hilbert-space realization, and exclusion of zero eigenvalues are implicit in the notes. Without the prime, zero modes make the displayed definition meaningless.}
\GSfield{Fate}{The definition and its use for Cauchy--Riemann operators are in \citet[\S3]{quillen1985determinants}, received June 1984 and published in the first 1985 issue.}
\GSfield{Present status}{As of November 2026: A standard spectral determinant, distinct from an ordinary infinite product.}
\GSfield{Assessment}{literal statement: correct with missing hypotheses; publication match: exact; verification: read the result; development: none; confidence: high}
\end{GSitem}
\begin{GSitem}{1985-650.2}{E}{The formal logarithmic derivative needs regularization}
\GSfield{Kind; source}{method; \GSpage{P001210}}
\GSfield{The notes}{The notebook factors \(\operatorname{Det}(\Delta+\varepsilon\delta\Delta)\) and obtains \(\delta\log\operatorname{Det}\Delta=\operatorname{Tr}(\delta\Delta\,\Delta^{-1})\).}
\GSfield{Statement}{For a smooth family of positive invertible elliptic operators \(A_t\) of fixed order on a closed manifold, the meaningful variation is \(\frac d{dt}\log\det_\zeta A_t=\operatorname{FP}_{s=0}\operatorname{Tr}(\dot A_tA_t^{-s-1})\), where \(\operatorname{FP}\) denotes the constant Laurent coefficient. With constant-dimensional kernels one works on their complements and includes the changing zero-mode metrics when forming a determinant line. In finite dimension this reduces to the ordinary trace formula.}
\GSfield{Reconstruction}{The formal quotient of differential operators is not generally trace class. Zeta determinants are not universally multiplicative, so the factorization is not an independent proof of the variation formula. The finite-part formula follows by differentiating the zeta function where the trace converges and continuing meromorphically. Editorial differentiation of the zeta function \ours; the determinant framework and nonmultiplicativity reference are prior literature.}
\GSfield{Fate}{The zeta-function and determinant-line framework predates the talk \citep[\S3]{quillen1985determinants}. The failure of general multiplicativity is explicitly analyzed in \citet[\S1, especially \S1.1]{kontsevichvishik1994geometry}, circulated in April/June 1994. This later reference verifies the limitation; it does not turn the notebook's formal derivative into a new 1985 discovery.}
\GSfield{Present status}{As of November 2026: The regularized derivative is established; the notebook's unqualified determinant multiplication is false as a general rule.}
\GSfield{Assessment}{literal statement: false as written; publication match: partial; verification: read the result; development: annotated; confidence: high}
\end{GSitem}
\begin{GSitem}{1985-650.3}{E}{Conformal change of the Dolbeault adjoint}
\GSfield{Kind; source}{claim; \GSpage{P001210}}
\GSfield{The notes}{For \(g\mapsto(1+\varepsilon\rho)g\), the page gives \(\Delta+\varepsilon\delta\Delta=(1+\varepsilon\rho)^j\bar\partial^*(1+\varepsilon\rho)^{-j-1}\bar\partial\).}
\GSfield{Statement}{Write \(g'=e^\varphi g\). For sections of the tangent power \(T^j\), the \(L^2\) inner products in Dolbeault degrees zero and one scale by \(e^{(j+1)\varphi}\) and \(e^{j\varphi}\). Therefore \(\bar\partial^*_{g'}=e^{-(j+1)\varphi}\bar\partial^*_g e^{j\varphi}\), and \(\Delta_{j,g'}=e^{-(j+1)\varphi}\bar\partial^*_g e^{j\varphi}\bar\partial\). The displayed notebook formula is instead this expression with \(j\) replaced by \(-j-1\).}
\GSfield{Reconstruction}{This is a direct adjoint calculation using the explicitly declared tangent-bundle convention. The printed expression is the formula for the Serre-dual bundle \(T^{-j-1}\), whose nonzero spectrum is related to that of \(T^j\). Such duality does not make the two differential operators literally equal. Editorial conformal-adjoint computation \oursComputed.}
\GSfield{Fate}{The elliptic-complex metric convention is fixed in \citet[\S3]{quillen1985determinants}; the distinction is checked here directly.}
\GSfield{Present status}{As of November 2026: The operator formula requires the dual-bundle correction or a declared change of convention. Its symmetric determinant anomaly can nevertheless agree.}
\GSfield{Assessment}{literal statement: false as written; publication match: partial; verification: read the result; development: annotated; confidence: high}
\end{GSitem}
\begin{GSitem}{1985-650.4}{E}{Heat finite parts and the one-dimensional regularization ambiguity}
\GSfield{Kind; source}{method; \GSpage{P001210-P001211}}
\GSfield{The notes}{The trace of multiplication by \(\rho\) is regulated by \(\operatorname{Tr}(\rho e^{-t\Delta})=a_{-1}(\rho)t^{-1}+a_0(\rho)+\cdots\). On the codimension-one kernel of \(a_{-1}\) it has a finite limit; extension is written \(\kappa a_{-1}+a_0\).}
\GSfield{Statement}{For a Laplace-type operator on a smooth closed surface and smooth \(\rho\), the heat trace has this expansion. Its finite part is \(a_0(\rho)\). If a linear regularization functional agrees with \(a_0\) on \(\ker a_{-1}\), and \(a_{-1}\ne0\), then elementary linear algebra forces it to equal \(a_0+\kappa a_{-1}\) for a scalar \(\kappa\).}
\GSfield{Reconstruction}{The uniqueness claim concerns a functional at a fixed metric with this agreement condition. It does not classify all metric-dependent, nonlocal, or nonlinear renormalization schemes. Zero modes change the finite part if one replaces the full heat trace by a primed trace. Editorial linear-algebra deduction of the one-dimensional extension ambiguity \ours; the heat expansion itself is classical.}
\GSfield{Fate}{The heat expansion is classical \citep[opening heat-kernel formulas]{mckeansinger1967curvature}; the extension assertion is the linear-algebra argument just stated.}
\GSfield{Present status}{As of November 2026: A valid finite-part prescription with an explicitly limited ambiguity, not an unresolved renormalization theorem.}
\GSfield{Assessment}{literal statement: correct with missing hypotheses; publication match: exact; verification: publisher metadata checked; development: annotated; confidence: high}
\end{GSitem}
\begin{GSitem}{1985-650.5}{E}{The scalar heat coefficient and curvature convention}
\GSfield{Kind; source}{claim; \GSpage{P001211}}
\GSfield{The notes}{The expansion is \((4\pi t)^{-1}[\int\rho+(t/3)\int\rho K+\cdots]\), but \(K\) is called scalar curvature. The subsequent finite-trace line drops \(1/(4\pi)\).}
\GSfield{Statement}{For the nonnegative scalar Laplace--Beltrami operator on a closed Riemannian surface, \(\operatorname{Tr}(\rho e^{-t\Delta})=(4\pi t)^{-1}\int\rho\,dA+(24\pi)^{-1}\int\rho R\,dA+O(t)\), with scalar curvature \(R=2K_G\). Thus the bracket coefficient is \(R/6=K_G/3\).}
\GSfield{Reconstruction}{The displayed coefficient is consistent with Gaussian curvature, not with the accompanying scalar-curvature label. A Dolbeault Laplacian differing by a fixed factor changes the time normalization but not this finite coefficient. A realified complex line doubles both trace coefficients.}
\GSfield{Fate}{The source named in the notes is correctly identifiable as \citet{mckeansinger1967curvature}; the original appeared in 1967, its Russian translation in 1969.}
\GSfield{Present status}{As of November 2026: A normalization error in an old heat-kernel formula. The manuscript, rather than the transcription, contains the conflicting label.}
\GSfield{Assessment}{literal statement: false as written; publication match: partial; verification: publisher metadata checked; development: annotated; confidence: high}
\end{GSitem}
\begin{GSitem}{1985-650.6}{E}{Tensor-spin heat coefficients and Serre symmetry}
\GSfield{Kind; source}{claim; \GSpage{P001211}}
\GSfield{The notes}{For \(T^j\) the page writes a leading coefficient twice the scalar one and a curvature coefficient \(j+2/3\). It then pairs \(j\) with \(-j-1\).}
\GSfield{Statement}{Serre duality on a closed Riemann surface pairs \(H^1(T^j)\) with \(H^0(K\otimes(T^j)^{-1})=H^0(T^{-j-1})\). With compatible metrics, the nonzero singular values of the Dolbeault complex give the corresponding dual spectral pairing. A complex trace, its underlying real trace, and the normalizations of \(\bar\partial^*\bar\partial\) and the scalar Laplacian must be fixed before numerical heat coefficients can be compared.}
\GSfield{Reconstruction}{The dual exponent is unambiguous. The isolated coefficient \(j+2/3\) is not certified as an exact invariant formula because the notebook changes real/complex and operator conventions; the anomaly polynomial is treated separately.}
\GSfield{Fate}{The cohomological interpretation is already part of the determinant-line construction \citep[\S\S1--3]{quillen1985determinants}.}
\GSfield{Present status}{As of November 2026: Duality is established; the shorthand coefficient remains normalization-dependent and cannot be used unqualified.}
\GSfield{Assessment}{literal statement: not determinable; publication match: partial; verification: read the result; development: none; confidence: medium}
\end{GSitem}
\begin{GSitem}{1985-650.7}{E}{The quadratic anomaly coefficient}
\GSfield{Kind; source}{theorem; \GSpage{P001211}}
\GSfield{The notes}{The conformal variation contains a curvature term proportional to \(6j^2+6j+1\), and an area term involving a regularization constant.}
\GSfield{Statement}{For the Dolbeault determinant line of \(T^j=K^{-j}\) on a closed curve, with its induced Hodge and zeta-corrected metric, the local gravitational anomaly relative to that of \(\mathcal O\) has coefficient \(a_j=6j^2+6j+1\). It satisfies \(a_j=a_{-j-1}\), \(a_0=1\), and \(a_1=13\). The determinant without the cohomology norm also has zero-mode contributions; an arbitrary bare determinant is not asserted to obey only the displayed local formula.}
\GSfield{Reconstruction}{The precise object is the determinant-line metric developed explicitly in record 655. An area counterterm belongs to the regularization prescription; it is not part of a universal numerical coefficient.}
\GSfield{Fate}{The \(j=0,1\) string calculation is already in \citet[eqs.~(15)--(25)]{polyakov1981quantum}. The full polynomial, in the cotangent-power convention, is read in \citet[(2.21),(3.22) and \S5]{belavinknizhnik1986complex}; that article cites a 1985 preprint without a verified month. E identifies the expository local-anomaly mechanism, not a claim that this later publication proves its precise public-circulation date.}
\GSfield{Present status}{As of November 2026: The polynomial is standard; it supplies no new unsolved question. The claim is separated from the incorrect operator convention in item 1985-650.3.}
\GSfield{Assessment}{literal statement: correct with missing hypotheses; publication match: partial; verification: read the result; development: none; confidence: high}
\end{GSitem}
\begin{GSitem}{1985-650.8}{E}{Integrating the anomaly into a Liouville action}
\GSfield{Kind; source}{claim; \GSpage{P001211}}
\GSfield{The notes}{The page identifies curvature with \(\partial\bar\partial\log\theta\) for \(ds^2=\theta|dz|^2\), and integrates to an action with a gradient term and \(\kappa e^\varphi\).}
\GSfield{Statement}{With \(ds^2=\theta|dz|^2\), Gaussian curvature is \(K_G=-2\theta^{-1}\partial_z\partial_{\bar z}\log\theta\). If \(g=e^\varphi g_0\) on a closed surface, a global Liouville functional has a kinetic term, a background-curvature term proportional to \(\int R_{g_0}\varphi\,dA_0\), and possibly a cosmological term \(\mu\int e^\varphi dA_0\). Spectral-determinant identities additionally account for zero modes and normalization constants.}
\GSfield{Reconstruction}{The background metric is needed for a globally meaningful statement; a flat coordinate expression cannot be integrated unchanged over an arbitrary compact surface. The curvature formula in the notes omits both the factor of \(\theta^{-1}\) and the convention-dependent sign/factor.}
\GSfield{Fate}{Polyakov gives this action \citep[eqs.~(16), (24)--(25)]{polyakov1981quantum}.}
\GSfield{Present status}{As of November 2026: An established anomaly-to-action construction; the bare global equality recorded here is incomplete and its curvature formula is false.}
\GSfield{Assessment}{literal statement: false as written; publication match: partial; verification: read the result; development: annotated; confidence: high}
\end{GSitem}
\begin{GSitem}{1985-650.9}{F}{Classical Liouville solutions versus the quantum Liouville integral}
\GSfield{Kind; source}{question; \GSpage{P001211-P001212}}
\GSfield{The notes}{Evaluating \(\int e^{\mathcal L(\varphi)}\mathcal D\varphi\) is identified with solving the Liouville equation; the needed integral is left to be computed.}
\GSfield{Statement}{The Euler--Lagrange equation specifies stationary configurations of a Liouville action. It does not determine its full quantum functional integral. One rigorous later setting is the sphere with \(0<\gamma<2\), \(\mu>0\), \(Q=2/\gamma+\gamma/2\), distinct insertion points, and charges \(\alpha_i<Q\), \(\sum_i\alpha_i>2Q\). Gaussian multiplicative chaos then defines nonzero finite renormalized Liouville correlations. These are additional hypotheses, not a reading of the notebook's unspecified integral.}
\GSfield{Reconstruction}{The notebook supplies neither insertion data nor coupling regime, contour, or measure. Consequently no single all-genus or all-central-charge quantum problem can be identified with its final sentence.}
\GSfield{Fate}{For the stated sphere setting, \citet[Thm.~3.2 and Def.~3.6]{davidkupiainenrhodesvargas2016liouville} give a rigorous construction (preprint October 2014; journal 2016). This is a precisely limited later realization, not proof of the asserted equivalence.}
\GSfield{Present status}{As of November 2026: The classical/quantum equivalence is false. A major later construction exists in a specified regime; the original unrestricted request remains unformulated and receives F, rather than an invented C or D problem.}
\GSfield{Assessment}{literal statement: false as written; publication match: thematic only; verification: read the result; development: none; confidence: high}
\end{GSitem}
\begin{GSitem}{1985-650.10}{E}{Critical dimension and mixed determinant powers}
\GSfield{Kind; source}{claim; \GSpage{P001212}}
\GSfield{The notes}{The displayed factor is \([\operatorname{Det}\Delta_1/(\operatorname{Det}\Delta_0)^{d/2}]^{1/2}\), while the following anomaly is proportional to \(13-d/2\), vanishing at \(d=26\).}
\GSfield{Statement}{For \(d\) real scalar matter fields in bosonic string gauge fixing, the nonzero-mode factor is \((\det'_{\mathbb R}P^*P)^{1/2}(\det'\Delta_0)^{-d/2}\), where \(P\) is the trace-free infinitesimal diffeomorphism operator. The ghost contribution has central charge \(-26\), so the total Weyl anomaly is proportional to \(d-26\). Passing to a complex Dolbeault determinant changes the ghost power accordingly.}
\GSfield{Reconstruction}{The outer square root and the scalar exponent in the notebook mix two conventions: without additional definitions the displayed expression does not entail its next coefficient. The conclusion \(d=26\) is correct independently of this shorthand. Vanishing anomaly is not a proof of convergence of the moduli integral.}
\GSfield{Fate}{This critical-dimension cancellation was published in \citet[eqs.~(20)--(25)]{polyakov1981quantum}.}
\GSfield{Present status}{As of November 2026: Known before the seminar. The number 26 is secure, whereas the exact displayed determinant formula needs repair.}
\GSfield{Assessment}{literal statement: false as written; publication match: partial; verification: read the result; development: annotated; confidence: high}
\end{GSitem}
\begin{GSitem}{1985-650.11}{E}{Zero modes turn determinants into cohomology lines}
\GSfield{Kind; source}{construction; \GSpage{P001212}}
\GSfield{The notes}{The last remark invokes exterior powers of \(H^1(T)\) and \(H^1(\mathcal O)\) when zero modes are restored.}
\GSfield{Statement}{For a closed curve define \(\lambda(E)=\det H^0(E)\otimes(\det H^1(E))^{-1}\). The anomaly-free tangent/scalar combination in this convention is \(\lambda(T)\otimes\lambda(\mathcal O)^{-13}\), or its dual if determinant conventions are reversed. At genus \(g\ge2\), \(H^0(T)=0\), but \(H^0(\mathcal O)\) remains one-dimensional; its norm still matters.}
\GSfield{Reconstruction}{The power 13 and the omitted degree-zero factors are restored. The notebook's informal ratio of two exterior powers is not an exact specification of the string measure.}
\GSfield{Fate}{This is the cohomological determinant framework of \citet[\S\S1--3]{quillen1985determinants}; record 655 supplies the seminar's explicit definition.}
\GSfield{Present status}{As of November 2026: An established construction with zero-mode factors essential. The later recurrence should be linked at item level.}
\GSfield{Assessment}{literal statement: correct with missing hypotheses; publication match: partial; verification: read the result; development: none; confidence: high}
\end{GSitem}

\GStalk{651}{1985-10-14}{D. B. Fuchs}{P001213-P001216}{Recapitulation of Feigin: gauge fixing the surface metric}
Fuchs repeats Feigin's account on a closed oriented surface of genus at least two: a formal Gaussian integration over maps, a conformal slice, and the Faddeev--Popov Jacobian. A question about a surface analogue of Wiener measure and Gelfand's physical interpretation of the extra string scale are distinct interventions. The omitted Teichmüller directions and the final identification of the ghost operator with the ordinary Hodge Laplacian are substantive mathematical gaps, rather than new open questions.

\begin{GSitem}{1985-651.1}{E}{The Polyakov action on a closed surface}
\GSfield{Kind; source}{construction; \GSpage{P001213}}
\GSfield{The notes}{The domain is a smooth closed oriented surface, and \(X\) denotes pairs \((x,g)\), with \(x:\Sigma\to\mathbb R^d\) and \(g\) a Riemannian metric. The quadratic action is repeated.}
\GSfield{Statement}{For a closed oriented smooth surface \(\Sigma\), \(g\in\operatorname{Met}(\Sigma)\), and \(x\in C^\infty(\Sigma,\mathbb R^d)\), define \(S(x,g)=\frac12\int_\Sigma|dx|_g^2dA_g\). It is invariant under simultaneous pullback by diffeomorphisms and under \(g\mapsto e^\varphi g\). The notebook omits the inessential factor \(1/2\).}
\GSfield{Reconstruction}{This is the action on parametrized maps with an independent metric. It is not already a measure on embedded or unparametrized surfaces.}
\GSfield{Fate}{A repetition of the construction in \citet[eq.~(6)]{polyakov1981quantum}; earlier seminar occurrences include 1982-569.4.}
\GSfield{Present status}{As of November 2026: Classical expository setup, with its functional integral still requiring interpretation.}
\GSfield{Assessment}{literal statement: correct with missing hypotheses; publication match: exact; verification: read the result; development: none; confidence: high}
\end{GSitem}
\begin{GSitem}{1985-651.2}{E}{Gelfand on a second string scale}
\GSfield{Kind; source}{remark; \GSpage{P001213}; I. M. Gelfand}
\GSfield{The notes}{Gelfand says strings introduce two constants, not only \(\hbar\), and interprets the second as the thickness of a tube swept in place of a particle path.}
\GSfield{Statement}{A string action has a tension \(T\), conventionally expressed through a length-squared parameter \(\alpha'\) after units and \(\hbar\) are chosen. This parameter is independent of the quantum expansion scale. A literal transverse tube radius is not defined by the zero-thickness world-sheet action written on this page.}
\GSfield{Reconstruction}{Tension is a precise modern interpretation of the extra scale; it is not a claim that the notebook defines a thick-string model or proves an equality between thickness and \(\sqrt{\alpha'}\).}
\GSfield{Fate}{The area coefficient is already present in \citet[eq.~(2)]{polyakov1981quantum}; the tube description is Gelfand's recorded physical analogy.}
\GSfield{Present status}{As of November 2026: The dimensional interpretation is identifiable; no independent geometric theorem or parameter determination is stated.}
\GSfield{Assessment}{literal statement: not determinable; publication match: thematic only; verification: read the result; development: none; confidence: medium}
\end{GSitem}
\begin{GSitem}{1985-651.3}{E}{Gaussian integration over the embedding coordinates}
\GSfield{Kind; source}{method; \GSpage{P001214}}
\GSfield{The notes}{Integrating \(e^{-\mathcal L(x,g)}\) over one scalar gives \((\operatorname{Det}\Delta_g)^{-1/2}\), and over \(d\) components gives power \(-d/2\).}
\GSfield{Statement}{For a positive definite real \(N\times N\) matrix \(A\), \[\int_{\mathbb R^N}e^{-x^TAx/2}dx=(2\pi)^{N/2}(\det A)^{-1/2}\]. For the scalar Laplacian on a closed connected surface, constants form its kernel. After fixing the mean, the spectral analogue uses \(\det'\Delta_g\); integrating unrestricted target translations contributes an infinite volume. The associated normalized Gaussian field has covariance \(\Delta_g^{-1}\) on mean-zero test functions and is distribution valued in dimension two.}
\GSfield{Reconstruction}{The finite-dimensional identity does not construct an infinite-dimensional flat measure \(\mathcal Dx\). The choice to remove the constant mode and the zeta normalization are additions.}
\GSfield{Fate}{Polyakov uses this determinant reduction in 1981 \citep[eq.~(14)]{polyakov1981quantum}. A rigorous fixed-background field is described in \citet[\S2]{davidkupiainenrhodesvargas2016liouville}.}
\GSfield{Present status}{As of November 2026: The determinant factor is standard once zero modes and normalization are fixed; smooth random embeddings are not obtained from this Gaussian field.}
\GSfield{Assessment}{literal statement: correct with missing hypotheses; publication match: exact; verification: read the result; development: none; confidence: high}
\end{GSitem}
\begin{GSitem}{1985-651.4}{F}{A surface analogue of Wiener measure}
\GSfield{Kind; source}{question; \GSpage{P001214}}
\GSfield{The notes}{A bracketed question asks whether there is an analogue of Wiener measure for surfaces.}
\GSfield{Statement}{No single probability problem is determined: the notes do not fix parametrized maps versus images, area versus Dirichlet weight, target dimension, topology, embeddedness, or a discretization. On a fixed closed background surface, a mean-zero Gaussian free field is a well-defined probability measure on distributions; it is not supported on smooth immersed surfaces and is not itself the desired reparametrization quotient.}
\GSfield{Reconstruction}{The fixed-background Gaussian field is a comparison, not an identification of the question. Different surface ensembles cannot be treated as interchangeable versions of one conjecture.}
\GSfield{Fate}{The precise field construction in \citet[\S2]{davidkupiainenrhodesvargas2016liouville} and its interacting sphere construction illustrate two distinct choices.}
\GSfield{Present status}{As of November 2026: F: the unspecified ensemble prevents a determinate solved/open assessment. The question is retained as a historical occurrence without a manufactured problem ID.}
\GSfield{Assessment}{literal statement: not determinable; publication match: thematic only; verification: read the result; development: none; confidence: high}
\end{GSitem}
\begin{GSitem}{1985-651.5}{E}{A conformal slice must include Teichmüller parameters}
\GSfield{Kind; source}{claim; \GSpage{P001214}}
\GSfield{The notes}{For genus at least two, \(M_c\) is described using metrics conformal to a standard curvature \(-1\) metric, and \(M_c\times\operatorname{Diff}(\Sigma)\to M\) is said to cover all metrics.}
\GSfield{Statement}{Each conformal class on a closed oriented surface of genus \(g\ge2\) has a unique curvature \(-1\) metric. Locally in Teichmüller space one may choose a family \(\widehat g(t)\), \(t\in\mathcal T_g\), and express a metric as \(f^*(e^\varphi\widehat g(t))\). A single fixed hyperbolic metric and its conformal multiples, even after diffeomorphisms, do not cover inequivalent conformal structures.}
\GSfield{Reconstruction}{The Teichmüller parameter is the omitted variable. The finite mapping-class quotient and finite automorphism stabilizers also prevent a global Cartesian-product identification.}
\GSfield{Fate}{The explicit moduli-space formulation is checked in \citet[\S2]{belavinknizhnik1986complex}; the classical slice is not a discovery of that 1986 paper.}
\GSfield{Present status}{As of November 2026: The literal surjectivity statement is false if ``standard'' denotes one metric. Record 652 partially repairs it by fixing a moduli point explicitly.}
\GSfield{Assessment}{literal statement: false as written; publication match: partial; verification: read the result; development: annotated; confidence: high}
\end{GSitem}
\begin{GSitem}{1985-651.6}{E}{Formal division by the volume of the diffeomorphism group}
\GSfield{Kind; source}{method; \GSpage{P001214}}
\GSfield{The notes}{The metric integral is divided by \(\operatorname{Vol}\operatorname{Diff}(\Sigma)\), and the remaining expression is to be computed on a slice.}
\GSfield{Statement}{In a finite-dimensional free isometric group action, the coarea formula relates orbit integration to a slice Jacobian. In the string calculation the analogous Faddeev--Popov expression is formal: both the metric-space volume and group volume require regularization, and residual stabilizers and moduli must be included.}
\GSfield{Reconstruction}{There is no countably additive infinite-dimensional Haar/Lebesgue volume supplied by the written \(L^2\) metrics. The argument is retained as a computational prescription, not as a measure identity.}
\GSfield{Fate}{The gauge-fixing source is \citet[eqs.~(17)--(23)]{polyakov1981quantum}.}
\GSfield{Present status}{As of November 2026: An established formal reduction; its validity does not follow from cancellation of two infinite constants.}
\GSfield{Assessment}{literal statement: not determinable; publication match: exact; verification: read the result; development: none; confidence: high}
\end{GSitem}
\begin{GSitem}{1985-651.7}{F}{The assertion that the integral collapses to a small set}
\GSfield{Kind; source}{claim; \GSpage{P001214}}
\GSfield{The notes}{The main contribution is said to come from a small set, with everything collapsing to a point.}
\GSfield{Statement}{A localization or saddle-point claim requires a specified family of measures, an asymptotic parameter, an observable, and a topology of convergence. None is given here.}
\GSfield{Reconstruction}{This sentence could refer to a Gaussian saddle, conformal collapse, or a heuristic concentration of random surfaces. The surrounding formulas do not select one.}
\GSfield{Fate}{No exact theorem has been identified. Polyakov's symbolic integral does not supply a proof of the literal concentration statement.}
\GSfield{Present status}{As of November 2026: Unidentifiable claim (F), not evidence that a precise localization problem remained open.}
\GSfield{Assessment}{literal statement: not determinable; publication match: none found; verification: unresolved lead; development: none; confidence: low}
\end{GSitem}
\begin{GSitem}{1985-651.8}{E}{L2 metrics and the infinitesimal diffeomorphism action}
\GSfield{Kind; source}{definition; \GSpage{P001215}}
\GSfield{The notes}{The page gives \(L^2\) norms on metric variations and vector fields, and defines \(D\varepsilon\) by differentiating the change of metric under a diffeomorphism.}
\GSfield{Statement}{For a fixed closed surface \((\Sigma,g)\), use \(\langle h,k\rangle=\int\operatorname{tr}(g^{-1}hg^{-1}k)dA\) on symmetric covariant 2-tensors and \(\langle v,w\rangle=\int g(v,w)dA\). Then \((Dv)_{ab}=(\mathcal L_vg)_{ab}=\nabla_av_b+\nabla_bv_a\). Formal adjoints are defined with these pairings. The trace-free gauge operator is \((Pv)_{ab}=Dv_{ab}-(\operatorname{div}v)g_{ab}\).}
\GSfield{Reconstruction}{Covariant indices and the distinction between \(D\) and its trace-free projection \(P\) are made explicit. Both are needed before interpreting the determinant.}
\GSfield{Fate}{These are the geometric operators behind \citet[eqs.~(17)--(21)]{polyakov1981quantum}.}
\GSfield{Present status}{As of November 2026: Correct geometric constructions, with the projection indispensable for conformal gauge.}
\GSfield{Assessment}{literal statement: correct with missing hypotheses; publication match: exact; verification: read the result; development: none; confidence: high}
\end{GSitem}
\begin{GSitem}{1985-651.9}{E}{The square-root Jacobian and the missing finite-dimensional directions}
\GSfield{Kind; source}{method; \GSpage{P001215}}
\GSfield{The notes}{The orbit Jacobian is \((\det D^*D)^{1/2}\), and symmetric 2-tensors are split into conformal and diffeomorphism parts.}
\GSfield{Statement}{For a closed genus-\(g\ge2\) surface, the elliptic decomposition is \(\Gamma(S^2T^*\Sigma)=\{fg\}\oplus\operatorname{im}P\oplus\ker P^*\) after the appropriate Sobolev completion. The last summand is the real space of holomorphic quadratic differentials, dimension \(6g-6\). The nonzero-mode orbit Jacobian is formally \((\det'P^*P)^{1/2}\); Gram determinants for moduli directions must also be retained.}
\GSfield{Reconstruction}{The notebook's two-summand decomposition omits \(\ker P^*\). At genus at least two there are no nonzero conformal Killing vector fields, but the moduli cokernel is still present.}
\GSfield{Fate}{The complete moduli factors are displayed in \citet[\S2]{belavinknizhnik1986complex}.}
\GSfield{Present status}{As of November 2026: The finite-dimensional missing summand is known and explicit. This is a faithful mathematical error, not an unresolved gauge-fixing question.}
\GSfield{Assessment}{literal statement: false as written; publication match: partial; verification: read the result; development: annotated; confidence: high}
\end{GSitem}
\begin{GSitem}{1985-651.10}{E}{Trace-free tensors as Beltrami differentials}
\GSfield{Kind; source}{construction; \GSpage{P001215}}
\GSfield{The notes}{The symmetric-square representation of rotations is decomposed into a scalar line and a two-dimensional summand; the latter is identified with \(T\otimes\Omega^{0,1}\), and \(D\) is identified with \(\bar\partial\).}
\GSfield{Statement}{On an oriented Riemannian surface the real bundle \(S^2T^*\) splits into its trace and trace-free parts. Raising one index identifies a real trace-free variation with a Beltrami differential in \(T^{1,0}\otimes\Omega^{0,1}\), together with its complex conjugate. In complex coordinates the trace-free infinitesimal action \(P\), rather than the full Lie derivative \(D\), is \(\bar\partial\) on \(T^{1,0}\) up to a fixed real/complex normalization.}
\GSfield{Reconstruction}{The metric identification and conjugate component account for the real ranks. The trace projection is essential; the full Lie derivative includes a scalar divergence term.}
\GSfield{Fate}{This is the Dolbeault description underlying \citet[\S2]{belavinknizhnik1986complex}.}
\GSfield{Present status}{As of November 2026: The representation-theoretic splitting is correct; the operator equality holds only after the stated projection and normalization.}
\GSfield{Assessment}{literal statement: correct with missing hypotheses; publication match: exact; verification: read the result; development: none; confidence: high}
\end{GSitem}
\begin{GSitem}{1985-651.11}{E}{The ghost operator is not the Hodge Laplacian on one-forms}
\GSfield{Kind; source}{claim; \GSpage{P001216}}
\GSfield{The notes}{The final page concludes that \(D^*D\) is the ordinary Laplacian on one-forms, on the ground that exterior differentiation and the Lie derivative are dual.}
\GSfield{Statement}{For the full symmetric gradient \((Dv)_{ab}=\partial_av_b+\partial_bv_a\) in a flat coordinate patch, its positive principal symbol is \(2(|\xi|^2I+\xi\otimes\xi)\). The Hodge Laplacian on one-forms has symbol \(|\xi|^2I\). These cannot agree even after one scalar normalization. For the trace-free operator \(P\), the principal symbols are proportional, but curvature terms still distinguish its Laplacian from the Hodge operator.}
\GSfield{Reconstruction}{The principal-symbol calculation is an editorial countercheck: substituting longitudinal and transverse vectors gives different eigenvalue ratios. It avoids relying on an undeclared sign convention for the curvature term. Editorial flat-symbol computation \oursComputed.}
\GSfield{Fate}{\citet[after eq.~(21)]{polyakov1981quantum} explicitly distinguishes the ghost operator from the ordinary vector Laplacian.}
\GSfield{Present status}{As of November 2026: The literal equality is false and materially affects the anomaly. It is not an unsolved spectral identity.}
\GSfield{Assessment}{literal statement: false as written; publication match: partial; verification: read the result; development: annotated; confidence: high}
\end{GSitem}

\GStalk{653}{1985-10-21}{Yu. S. Ryzhikova}{P001217 (upper)}{Recognizing identities in lattices of subspaces}
Ryzhikova announces an algorithm for identities in lattices of finite-dimensional linear subspaces, based on a computable dimension bound followed by a decision procedure in fixed dimension. This is identifiable with her 1987 Doklady paper. The manuscript heading is \(21/X/85\), whereas the catalogue gives 21 September; the present date follows the manuscript. This contribution precedes Fuchs's record 652 on the same page.

\begin{GSitem}{1985-653.1}{E}{The subspace lattice and the modular identity}
\GSfield{Kind; source}{definition; \GSpage{P001217 (upper)}}
\GSfield{The notes}{Addition and multiplication mean sum and intersection of subspaces. The modular identity is written \(a(ab+c)=ab+ac\), or \(a(b+c)=ab+ac\) when \(b\le a\).}
\GSfield{Statement}{For a vector space \(V\) over a field \(F\), let \(L(V)\) be its lattice of linear subspaces, with \(U\vee W=U+W\), \(U\wedge W=U\cap W\). If \(B\subseteq A\), then \(A\cap(B+C)=B+(A\cap C)\). Substituting \(B=A\cap B'\) gives the first notebook identity.}
\GSfield{Reconstruction}{The operations and order are fixed explicitly. Modularity is weaker than distributivity and does not imply that every modular lattice is a subspace lattice.}
\GSfield{Fate}{Classical background repeated from the seminar's 1973-319.1--2; see \citet[\S1.1]{gelfandponomarev1974modular}.}
\GSfield{Present status}{As of November 2026: Correct expository identity. The separate representability problem GP-1973-30 is not re-posed here.}
\GSfield{Assessment}{literal statement: correct; publication match: exact; verification: read the result; development: none; confidence: high}
\end{GSitem}
\begin{GSitem}{1985-653.2}{A}{A computable dimension bound for a failed lattice identity}
\GSfield{Kind; source}{theorem; \GSpage{P001217 (upper)}}
\GSfield{The notes}{For terms \(p,q\), the proposed test dimension is \[N=1+\min\{\max(d(p),d(q)),\max(k(p),k(q))\}\], with \(d\) counting sums and \(k\) intersections.}
\GSfield{Statement}{Fix a field \(F\). For finite lattice terms \(p,q\) in finitely many variables, using binary sum and intersection, let \(d(p)\) and \(k(p)\) count occurrences of the two operations in the written term. The identity \(p=q\) holds in all finite-dimensional \(F\)-subspace lattices if and only if it holds in \(L(F^N)\), with the stated \(N\). A failure in larger dimension has a witness in dimension at most \(N\), which embeds in dimension \(N\).}
\GSfield{Reconstruction}{The published Lemma 1 gives the join-count bound. Its proof follows a vector witnessing a failed inclusion and retains only the finitely many vectors needed at joins; duality of finite-dimensional subspace lattices gives the meet-count bound. Taking the smaller bound recovers the notebook's symmetric formula. This last combination is explicitly a reconstruction from the lemma, not a verbatim theorem number. Editorial deduction of the meet bound by duality and of the minimum of the two bounds \ours; the join-bound proof is due to Ryzhikova.}
\GSfield{Fate}{\citet[Lemma~1, pp.~1058--1059]{ryzhikova1987decidability} is the exact later paper. It was received 16 May 1986 and published in 1987; no earlier public preprint was verified.}
\GSfield{Present status}{As of November 2026: Established by the speaker soon after the talk (A). This is a genuine quantitative bridge from an unbounded-dimensional question to a finite test.}
\GSfield{Assessment}{literal statement: correct; publication match: exact; verification: read the result; development: annotated; confidence: high; problem \hyperref[prob:GP-1985-20]{GP-1985-20}}
\end{GSitem}
\begin{GSitem}{1985-653.3}{A}{Decidability of identities in linear-subspace lattices}
\GSfield{Kind; source}{theorem; \GSpage{P001217 (upper)}}
\GSfield{The notes}{Ryzhikova announces an algorithm deciding identities for all finite-dimensional vector spaces. After the dimension reduction, the note invokes Tarski for subspaces of a fixed real \(n\)-space.}
\GSfield{Statement}{For each fixed field \(F\), there is an algorithm that, given two finite sum/intersection terms, decides whether their equality holds for every tuple of subspaces of every finite-dimensional \(F\)-vector space. The identity theory depends only on \(\operatorname{char}F\). In characteristic zero one reduces to a fixed-dimensional real subspace lattice, encodes its incidence relations by polynomial conditions, and applies real-closed-field quantifier elimination. In characteristic \(p>0\), reduction to \(\mathbb F_p\) makes the bounded-dimensional test finite.}
\GSfield{Reconstruction}{The notebook's real-case explanation is expanded using the published field-independence lemmas. The theorem concerns identities, not the full first-order theory of arbitrary-dimensional lattices, word problems for all modular lattices, or efficient complexity bounds.}
\GSfield{Fate}{\citet[main theorem and Lemmas~2--3, pp.~1058--1061]{ryzhikova1987decidability} proves precisely this result; the English translation appeared in 1988.}
\GSfield{Present status}{As of November 2026: Solved (A), with the exact quantifiers recovered. No unresolved decision problem is inferred from the absence of a complexity estimate.}
\GSfield{Assessment}{literal statement: correct with missing hypotheses; publication match: exact; verification: read the result; development: none; confidence: high; problem \hyperref[prob:GP-1985-20]{GP-1985-20}}
\end{GSitem}

\GStalk{652}{1985-10-21}{D. B. Fuchs; B. L. Feigin (intervention)}{P001217 (lower)-P001218}{Fixed conformal structure, moduli, and the Liouville action}
After Fuchs's short recap, the notebook explicitly marks Feigin's intervention. Fixing a point of moduli separates the conformal factor from the remaining \(3g-3\) complex parameters. The anomaly calculation and critical dimension recur, and the determinant variation is integrated into a Liouville action. The asserted equivalence of solving its classical equation and evaluating its functional integral is not valid as stated. The date is corrected from the catalogue's September to the manuscript's October.

\begin{GSitem}{1985-652.1}{E}{Fuchs repeats the determinant integral on a metric slice}
\GSfield{Kind; source}{method; \GSpage{P001217 (lower)}}
\GSfield{The notes}{Fuchs starts with \(\int_{M_c}(\operatorname{Det}\Delta)^{-d/2}\nu\), continuing the previous Gaussian reduction.}
\GSfield{Statement}{For a fixed closed oriented surface of genus \(g\ge2\), the formal matter factor is \((\det'\Delta_0)^{-d/2}\). A conformal-gauge slice also carries a ghost Jacobian, Gram factors for moduli, and zero-mode factors. The symbol \(\nu\) alone does not define the resulting measure.}
\GSfield{Reconstruction}{The prime and the distinction between matter and the full measure are supplied. This sentence is Fuchs's recap; the next paragraph is separately attributed to Feigin in the manuscript.}
\GSfield{Fate}{A recurrence of \citet[eqs.~(14),(17)--(23)]{polyakov1981quantum} and of record 651.}
\GSfield{Present status}{As of November 2026: Identifiable expository formula, conditional on a consistent gauge and regularization.}
\GSfield{Assessment}{literal statement: not determinable; publication match: exact; verification: read the result; development: none; confidence: high}
\end{GSitem}
\begin{GSitem}{1985-652.2}{E}{Fixing a conformal structure before integrating the Weyl factor}
\GSfield{Kind; source}{construction; \GSpage{P001217-P001218}; B. L. Feigin}
\GSfield{The notes}{Feigin fixes a moduli parameter \(t\), writes \(g=e^\varphi g_0\) with \(K(g_0)=-1\), and notes that the moduli space has complex dimension \(3g-3\).}
\GSfield{Statement}{For a closed curve of genus \(g\ge2\), choose a marked complex structure \(t\) and its hyperbolic metric \(g_0(t)\). Metrics in that conformal class are precisely \(e^\varphi g_0(t)\), \(\varphi\in C^\infty(\Sigma,\mathbb R)\). The tangent space to Teichmüller space is \(H^1(\Sigma,T)\), of complex dimension \(3g-3\). The unmarked moduli space is an orbifold/stack quotient, not globally a smooth manifold with one coordinate chart.}
\GSfield{Reconstruction}{The genus restriction and the marked/unmarked distinction make the dimension claim precise. Keeping \(t\) fixed repairs the single-conformal-class ambiguity of record 651, but leaves a later integration over \(t\).}
\GSfield{Fate}{The resulting separation is described explicitly in \citet[\S2]{belavinknizhnik1986complex}; the underlying uniformization and deformation theory are classical.}
\GSfield{Present status}{As of November 2026: Established geometry. The moduli directions cannot be discarded by calling the Weyl anomaly zero.}
\GSfield{Assessment}{literal statement: correct with missing hypotheses; publication match: exact; verification: read the result; development: none; confidence: high}
\end{GSitem}
\begin{GSitem}{1985-652.3}{E}{The determinant anomaly with spin j}
\GSfield{Kind; source}{theorem; \GSpage{P001218}; B. L. Feigin}
\GSfield{The notes}{The conformal derivative is written as an area term \((Aj+B)\int\rho\) plus a curvature term \(C(6j^2+6j+1)\int\rho K\); here \(K\) is explicitly Gaussian curvature.}
\GSfield{Statement}{For the determinant of the Dolbeault complex on \(T^j\), completed by its cohomology metrics, the local Weyl anomaly coefficient is \(6j^2+6j+1\). In particular the tangent and scalar coefficients are 13 and 1. Primed determinants and cohomology norms must be treated together; finite area counterterms depend on a common renormalization prescription.}
\GSfield{Reconstruction}{This occurrence corrects the earlier scalar-curvature label but does not specify enough normalization to certify the constants \(A,B,C\). The numerical spin polynomial is distinguished from those unspecified constants.}
\GSfield{Fate}{The earlier Polyakov calculation \citep[eqs.~(15)--(25)]{polyakov1981quantum} is recast in determinant-line language; compare record 655.}
\GSfield{Present status}{As of November 2026: Known anomaly calculation. The general-spin and zero-mode qualifications are needed before using it as an exact metric identity.}
\GSfield{Assessment}{literal statement: correct with missing hypotheses; publication match: partial; verification: read the result; development: none; confidence: high}
\end{GSitem}
\begin{GSitem}{1985-652.4}{E}{What becomes constant at dimension 26}
\GSfield{Kind; source}{claim; \GSpage{P001218}; B. L. Feigin}
\GSfield{The notes}{At \(d=26\) the anomaly cancels, and the notes say that one now integrates a constant.}
\GSfield{Statement}{After consistent matter/ghost normalization and cancellation of any chosen local counterterm, the critical bosonic measure is independent of the Weyl factor at a fixed complex structure. This does not imply constancy as a function on moduli space, finite Weyl-group volume, or convergence at degenerating curves.}
\GSfield{Reconstruction}{``Constant'' can only mean independent of the conformal representative in the calculation at hand. Degeneration behavior remains a different question.}
\GSfield{Fate}{The cancellation is already in \citet[eq.~(23)]{polyakov1981quantum}. \citet[\S\S3--4]{belavinknizhnik1986complex} subsequently analyze nontrivial moduli dependence and boundary singularities.}
\GSfield{Present status}{As of November 2026: The restricted constancy statement is correct; an unrestricted constant integrand is false. No solved anomaly calculation is relabeled as an open string-measure problem.}
\GSfield{Assessment}{literal statement: correct with missing hypotheses; publication match: partial; verification: read the result; development: none; confidence: high}
\end{GSitem}
\begin{GSitem}{1985-652.5}{F}{The remaining integral over moduli}
\GSfield{Kind; source}{question; \GSpage{P001218}; B. L. Feigin}
\GSfield{The notes}{The remaining integral over the \((3g-3)\)-dimensional complex moduli space is said still to require computation.}
\GSfield{Statement}{A definite integral requires a specified genus, inserted states or vacuum normalization, determinant-line trivialization, integration cycle, and a prescription near the Deligne--Mumford boundary. The notebook does not choose these. For the bosonic vacuum measure at genus at least two, a later precise result identifies its holomorphic factor, up to normalization, and its double-pole behavior at degenerations; this does not make the unregularized vacuum integral finite.}
\GSfield{Reconstruction}{The later theorem is a constrained comparison to the stated task of computing moduli dependence. It is not an exact solution of an unspecified numerical integral.}
\GSfield{Fate}{\citet[\S\S3--5]{belavinknizhnik1986complex} provides that geometric analysis. Its bibliography cites a Landau Institute preprint from 1985; no month was verified. The full paper was submitted February 1986 and published August 1986.}
\GSfield{Present status}{As of November 2026: F for the unspecific integral request, with a strong thematic later development. No separate open-problem ID is assigned.}
\GSfield{Assessment}{literal statement: not determinable; publication match: thematic only; verification: read the result; development: none; confidence: high}
\end{GSitem}
\begin{GSitem}{1985-652.6}{E}{The background-curvature term in the Liouville action}
\GSfield{Kind; source}{construction; \GSpage{P001218}; B. L. Feigin}
\GSfield{The notes}{The integrated determinant contains \(\int\varphi\partial\bar\partial\varphi\), \(\int e^\varphi\operatorname{Vol}\), and \(\int\varphi\operatorname{Vol}\), with unspecified constants.}
\GSfield{Statement}{With a fixed hyperbolic background \(g_0\), a global Liouville action has the form \(S_L(\varphi)=a\int|\nabla\varphi|_{g_0}^2dA_0+b\int R_{g_0}\varphi\,dA_0+\mu\int e^\varphi dA_0\), up to a constant and zero-mode conventions. Since \(R_{g_0}\) is constant, the background-curvature term is proportional to the notebook's linear integral.}
\GSfield{Reconstruction}{The sign relating \(\varphi\partial\bar\partial\varphi\) to the positive energy depends on integration conventions; the notes do not fix \(a,b,\mu\). This occurrence restores a term absent from the earlier local formula.}
\GSfield{Fate}{The framework is Polyakov's \citep[eqs.~(16),(24)--(25)]{polyakov1981quantum}. A precise modern action with curvature normalization appears in \citet[eq.~(1.1)]{davidkupiainenrhodesvargas2016liouville}.}
\GSfield{Present status}{As of November 2026: Established action, with its coefficients and zero modes to be fixed before any exact equality is asserted.}
\GSfield{Assessment}{literal statement: correct with missing hypotheses; publication match: partial; verification: read the result; development: none; confidence: high}
\end{GSitem}
\begin{GSitem}{1985-652.7}{F}{Repeating the classical/quantum Liouville equivalence}
\GSfield{Kind; source}{claim; \GSpage{P001218}; B. L. Feigin}
\GSfield{The notes}{Solving the Liouville equation is again called equivalent to computing \(\int e^{\mathcal L_j(\varphi)}\mathcal D\varphi\), and the total action is \(\mathcal L_1-\frac d2\mathcal L_0\).}
\GSfield{Statement}{The combined determinant can formally be written as the exponential of the difference of logarithmic determinants, once conventions are consistent. Its stationary equation does not evaluate the full functional integral. A quantum definition additionally needs the integration law, insertions, zero-mode treatment, and a coupling regime.}
\GSfield{Reconstruction}{The exponential rewriting is an algebraic step, whereas the asserted equivalence with a nonlinear PDE is a separate false implication. Item 1985-650.9 gives a precise later sphere construction and its bounds.}
\GSfield{Fate}{The modern construction in \citet[Thm.~3.2, Def.~3.6]{davidkupiainenrhodesvargas2016liouville} makes the distinction explicit; it is not a universal evaluation by classical solutions.}
\GSfield{Present status}{As of November 2026: The literal equivalence is false. This underdetermined quantum request remains F, without being declared an open problem.}
\GSfield{Assessment}{literal statement: false as written; publication match: thematic only; verification: read the result; development: none; confidence: high}
\end{GSitem}

\GStalk{654}{1985-10-28}{Werner Müller}{P001219-P001221 (upper)}{Characteristic numbers, conic operators, and Hilbert modular cusps}
Müller first reports Cheeger--Gromov's characteristic-number results for complete finite-volume manifolds with bounded curvature and a compact topological core. He then discusses conical heat expansions, eta corrections to the signature, an index question for more general regular singular operators, and locally symmetric examples. The last example matches his published 1984 work on Hilbert modular cusp defects, though its exact handwritten constant has not been checked in that inaccessible full text. The displayed cone metric has the wrong exponent in the manuscript itself.

\begin{GSitem}{1985-654.1}{E}{Gauss--Bonnet for tame complete finite-volume manifolds}
\GSfield{Kind; source}{theorem; \GSpage{P001219}}
\GSfield{The notes}{Under the heading Cheeger/Gromov, \(X\) is complete, \(\operatorname{vol}X<\infty\), \(|K|<1\), and diffeomorphic to the interior of a compact manifold with boundary. The conclusion is \(\int_Xp(\Omega)=\chi(X)\).}
\GSfield{Statement}{Let \(X^{2m}\) be an oriented complete Riemannian manifold of finite volume and uniformly bounded sectional curvature, diffeomorphic to the interior of a compact smooth manifold with boundary. Then its Euler form \(e(TX,\nabla)\), normalized as the Pfaffian in the Chern--Gauss--Bonnet theorem, is integrable and \(\int_Xe(TX,\nabla)=\chi(X)\).}
\GSfield{Reconstruction}{The symbol \(p\) must mean the Euler polynomial here, not an arbitrary characteristic polynomial or a Pontryagin class. The compact-core assumption is retained; finite volume and bounded curvature alone do not justify this Euler-characteristic equality.}
\GSfield{Fate}{\citet[Introduction, pp.~116--117]{cheegergromov1985characteristic} gives the stated tame case and contrasts nontame counterexamples. The paper is dated 1985; no exact circulation day was established.}
\GSfield{Present status}{As of November 2026: Established result explicitly attributed to others in the talk. E denotes reported existing work; within-year publication ordering is not independently known.}
\GSfield{Assessment}{literal statement: correct with missing hypotheses; publication match: exact; verification: read the result; development: none; confidence: high}
\end{GSitem}
\begin{GSitem}{1985-654.2}{E}{Homotopy invariance of the integrated L-form}
\GSfield{Kind; source}{theorem; \GSpage{P001220 (top)}}
\GSfield{The notes}{The second assertion is that \(\int_XL_k(p_1,\ldots,p_k)\) is a homotopy invariant.}
\GSfield{Statement}{For an oriented complete finite-volume \(4k\)-manifold with uniformly bounded sectional curvature, diffeomorphic to the interior of a compact manifold with boundary, its integrated Hirzebruch \(L\)-form depends only on the oriented homotopy type within this class. The differential-form integral converges because the curvature is bounded and volume finite.}
\GSfield{Reconstruction}{The dimension, orientation convention, and all geometric hypotheses carry over from the preceding item. This is not a claim of homotopy invariance for every Pontryagin number on every noncompact manifold.}
\GSfield{Fate}{Exact match to \citet[Thm.~0.6]{cheegergromov1985characteristic}; its accompanying signature/eta relation is eq.~(0.7).}
\GSfield{Present status}{As of November 2026: A known theorem in the stated class, with topology at infinity essential to its formulation.}
\GSfield{Assessment}{literal statement: correct with missing hypotheses; publication match: exact; verification: read the result; development: none; confidence: high}
\end{GSitem}
\begin{GSitem}{1985-654.3}{E}{Cone geometry and the radial regular singular operator}
\GSfield{Kind; source}{construction; \GSpage{P001220}}
\GSfield{The notes}{The notes write \(C(Y)=(0,1]\times Y\), \(ds^2=du^2+u^{-2}g_Y\), and a radial Laplacian \(-\partial_u^2+u^{-2}A(u)\). The same \(A(u)\) is also declared an \(L^2\) section.}
\GSfield{Statement}{For a closed Riemannian link \((Y,g_Y)\), the exact metric cone has metric \(du^2+u^2g_Y\). After a degree-dependent unitary transformation of differential forms to a fixed \(L^2(du)\) space, its Hodge Laplacian is a matrix regular singular operator with radial part \(-\partial_u^2\) and an angular operator multiplied by \(u^{-2}\). For an exact cone the angular coefficient is independent of \(u\); more general conic metrics allow variable coefficients and lower-order terms.}
\GSfield{Reconstruction}{The scan confirms the minus sign in the metric exponent: it is a faithful mathematical error, not an OCR error. The angular operator and the section on which it acts are different objects. The unitary density and degree shifts are missing from the notebook.}
\GSfield{Fate}{The exact cone metric and form decomposition are in \citet[Introduction and \S1]{cheeger1983spectral}, received August 1982, published 1983.}
\GSfield{Present status}{As of November 2026: The cone model is established. The written \(u^{-2}g_Y\) does not describe a cone tip and cannot support the following cone calculation literally.}
\GSfield{Assessment}{literal statement: false as written; publication match: partial; verification: read the result; development: annotated; confidence: high}
\end{GSitem}
\begin{GSitem}{1985-654.4}{E}{Logarithmic heat asymptotics and link spectral data}
\GSfield{Kind; source}{theorem; \GSpage{P001220}}
\GSfield{The notes}{The heat trace has powers of \(t^{1/2}\), a \(c\log t\) term, and a constant \(b\), expressed through \(\zeta_Y(-1)\) and residues at odd integers.}
\GSfield{Statement}{For a compact space with isolated exact conic singularities and a fixed admissible self-adjoint Hodge-Laplacian realization, Cheeger's heat expansion includes ordinary local powers and additional constant and logarithmic contributions. The singular terms are determined by the spectrum of shifted angular operators on the link. In middle degrees, ideal boundary conditions must be specified when uniqueness fails; regular values of a zeta function may need replacement by finite parts at poles.}
\GSfield{Reconstruction}{The notebook does not define whether \(\lambda_j\) are eigenvalues of the link Laplacian, its shifted square root, or the angular matrix operator. Therefore its numerical coefficients \(b_j\), evaluation at \(-1\), and residue arguments are not certified as an exact formula for an unshifted link zeta function. The regular \(t^0\) term must not be counted twice.}
\GSfield{Fate}{\citet[Thms.~4.1 and 4.4, especially eq.~(4.38)]{cheeger1983spectral} supplies the precise expansion and spectral correction.}
\GSfield{Present status}{As of November 2026: Established spectral geometry. The structural statement matches, while the compressed coefficient formula needs the omitted angular and domain conventions.}
\GSfield{Assessment}{literal statement: correct with missing hypotheses; publication match: partial; verification: read the result; development: none; confidence: medium}
\end{GSitem}
\begin{GSitem}{1985-654.5}{E}{The signature operator and its eta correction}
\GSfield{Kind; source}{theorem; \GSpage{P001220}}
\GSfield{The notes}{For \(\dim X=4k\), the signature operator is \(d+\delta:\Lambda_+^*\to\Lambda_-^*\), with index \(\int_XL+\eta_Y(B)\); on the cone it is written \(-\partial_u+u^{-1}B\). The \(L\)-form is said to vanish on the cone.}
\GSfield{Statement}{Let a compact oriented \(4k\)-manifold with boundary be completed by exact cone tips over its closed links. The natural \(L^2\) signature equals the signature of the image of relative middle cohomology in absolute middle cohomology. It is given by the interior \(L\)-integral plus a spectral boundary correction. With the usual outward-boundary odd-signature convention this is written \(\operatorname{sign}M=\int_ML-\eta_{\mathrm{sign}}(\partial M)\). On an exact cone the \(L\)-density vanishes; the remaining tip contribution is spectral.}
\GSfield{Reconstruction}{Changing link orientation or the sign of \(B\) changes the displayed eta sign. For a general Dirac APS index one instead encounters the reduced expression \((\eta+h)/2\). Clifford multiplication, unitary radial rescaling, and degree shifts are needed to interpret the radial formula; the shorthand is not a universal scalar equality.}
\GSfield{Fate}{\citet[\S6, eqs.~(6.1),(6.4),(6.14)--(6.15)]{cheeger1983spectral} gives the exact signature interpretation.}
\GSfield{Present status}{As of November 2026: Known before the talk. Its eta term is meaningful only with the boundary operator, orientation, and realization fixed.}
\GSfield{Assessment}{literal statement: correct with missing hypotheses; publication match: partial; verification: read the result; development: none; confidence: high}
\end{GSitem}
\begin{GSitem}{1985-654.6}{B}{An index formula for general Fuchs-type operators}
\GSfield{Kind; source}{question; \GSpage{P001220-P001221}}
\GSfield{The notes}{For \(D=\partial_u^n+\sum u^{i-n}\partial_u^iA_j\), with differential operators on \(Y\), the notebook asks \ru{индекс --- ?}. It points to the eta term in Cheeger's example.}
\GSfield{Statement}{A well-posed version fixes a compact manifold with isolated conic ends, Hermitian bundles, and an elliptic Fuchs-type operator \(D=u^{-\nu}\sum_{k=0}^\mu A_k(u)(-u\partial_u)^k\), \(\nu>0\), in the cone-elliptic calculus. Its ordinary and rescaled principal symbols must be invertible away from zero. One also specifies a closed domain between the minimal and maximal \(L^2\) domains. These extensions are Fredholm; when the \(A_k\) are constant near each tip, the minimal-domain index is the interior local index contribution plus a generalized eta invariant of the Mellin symbol. Enlarging the domain adds its finite dimension relative to the minimal domain.}
\GSfield{Reconstruction}{Ellipticity, compactness, bundles, density, and domain are supplied requirements. The mixed indices \(i,j\) and missing summation limits in the notebook do not specify a single operator. The reconstructed problem asks for the general index mechanism, rather than a numerical index without data.}
\GSfield{Fate}{A precise later solution is \citet[Prop.~1.3.16, Cor.~2.4.7 and Thm.~2.4.8]{lesch1997fuchs}. The text derives from a June 1993 habilitation, appeared as a July 1996 preprint, and was published in 1997.}
\GSfield{Present status}{As of November 2026: B: the identifiable regular singular index program has an established realization under explicit hypotheses. No exact unproved remainder is asserted for the notebook's underspecified formula.}
\GSfield{Assessment}{literal statement: not determinable; publication match: partial; verification: read the result; development: none; confidence: medium; problem \hyperref[prob:GP-1985-21]{GP-1985-21}}
\end{GSitem}
\begin{GSitem}{1985-654.7}{F}{The endpoint condition does not specify a general cone domain}
\GSfield{Kind; source}{claim; \GSpage{P001221 (top)}}
\GSfield{The notes}{The index is taken on functions satisfying \[\lim_{u\to0}u^{-1/2}\varphi(u)=\lim_{u\to0}u^{1/2}\varphi'(u)\].}
\GSfield{Statement}{An elliptic regular singular operator generally needs conditions on the coefficients of its indicial asymptotics, depending on the angular spectrum. For the scalar repeated-root model, solutions may contain \(a\sqrt u+b\sqrt u\log u\). The two written limits both exist finitely only if \(b=0\), and then are \(a\) and \(a/2\); the stated equality forces \(a=0\). It is therefore not the usual condition merely excluding the logarithmic branch.}
\GSfield{Reconstruction}{This substitution is an editorial check of the formula, not an identification of the intended operator. Since its angular mode and extension are not named, a unique corrected endpoint condition cannot be recovered. Editorial endpoint-asymptotics substitution \oursComputed.}
\GSfield{Fate}{The dependence of domains on indicial data is described in \citet[\S1.3 and Prop.~1.3.16]{lesch1997fuchs}.}
\GSfield{Present status}{As of November 2026: F for the intended extension; the displayed condition cannot serve as a universal conic boundary condition.}
\GSfield{Assessment}{literal statement: not determinable; publication match: thematic only; verification: read the result; development: annotated; confidence: medium}
\end{GSitem}
\begin{GSitem}{1985-654.8}{E}{Finite L2 kernels for invariant operators require ellipticity}
\GSfield{Kind; source}{claim; \GSpage{P001221}}
\GSfield{The notes}{For finite-volume \(X=K\backslash G/\Gamma\), locally defined bundles and a locally invariant operator \(D\), the notes assert finite-dimensional \(\ker\bar D\) and \(\ker\bar D^*\), then define their dimension difference.}
\GSfield{Statement}{For a specified closed densely defined operator \(D:L^2(E)\to L^2(F)\), the \(L^2\) index \(\dim\ker D-\dim\ker D^*\) is defined when both dimensions are finite. Local invariance alone does not imply this: the zero operator on a positive-dimensional finite-volume locally symmetric manifold is invariant and has infinite-dimensional kernel. Relevant theorems concern elliptic homogeneous or generalized Dirac operators with controlled cusp geometry and specified realizations.}
\GSfield{Reconstruction}{The manuscript does not state the needed ellipticity or operator class, and the bundle annotation is uncertain. We do not replace it by an unrestricted theorem about all invariant differential operators. Editorial zero-operator counterexample \ours.}
\GSfield{Fate}{\citet{muller1987cusps} treats rank-one cusp geometry and generalized Dirac operators. Its publisher description and chapter list were checked; the full operator theorem was not read.}
\GSfield{Present status}{As of November 2026: The broad literal assertion is false. The established analytic program is identifiable, but this specific scope is only indirectly matched.}
\GSfield{Assessment}{literal statement: false as written; publication match: partial; verification: publisher metadata checked; development: annotated; confidence: medium}
\end{GSitem}
\begin{GSitem}{1985-654.9}{E}{Hilbert modular surfaces and the one-cusp qualification}
\GSfield{Kind; source}{construction; \GSpage{P001221}}
\GSfield{The notes}{The example is \((\mathbb H\times\mathbb H)/\operatorname{SL}_2(\mathcal O_F)\), \(F=\mathbb Q(\sqrt D)\), followed by ``only 1 cusp''. The cusp lattice and unit group are \(M=\mathcal O_F\), \(V=\mathcal O_F^{*2}\).}
\GSfield{Statement}{Let \(F\) be a real quadratic field and let its two real embeddings define the action of \(\operatorname{SL}_2(\mathcal O_F)\) on \(\mathbb H^2\). This finite-volume quotient has cusp classes indexed by the ideal classes of \(F\); it has one cusp when the ordinary class number is one. The full quotient may have finite isotropy. One must either work with an orbifold index or pass to a torsion-free subgroup, which may introduce more cusps.}
\GSfield{Reconstruction}{The notebook does not impose class number one or remove torsion. The displayed \((M,V)\) describes a chosen standard cusp, not automatically all cusps of every Hilbert modular surface.}
\GSfield{Fate}{The intended application is identifiable with \citet{muller1984signature}, a 1984 paper verified in the author's bibliography. The full paper was inaccessible in this run.}
\GSfield{Present status}{As of November 2026: The geometry is classical; the one-cusp assertion is false without a field/subgroup qualification. Publication match is partial because the exact setup was not recovered from the full source.}
\GSfield{Assessment}{literal statement: false as written; publication match: partial; verification: publisher metadata checked; development: none; confidence: medium}
\end{GSitem}
\begin{GSitem}{1985-654.10}{E}{The cusp signature defect and the Shimizu L-series}
\GSfield{Kind; source}{claim; \GSpage{P001221}}
\GSfield{The notes}{The index is written \(\frac13\int_Xp_1-\frac1{2\pi}L(M,V,1)\operatorname{vol}(\mathbb R^2/M)\). The series is \(\sum\operatorname{sign}N(\mu)/|N(\mu)|^s\), but its summation domain ends with the unexplained \(N_F\).}
\GSfield{Statement}{For a rank-two lattice \(M\) in a real quadratic field and a finite-index group \(V\) of totally positive norm-one units preserving it, define \(L(M,V,s)=\sum_{\mu\in(M\setminus\{0\})/V}\operatorname{sign}N_{F/\mathbb Q}(\mu)|N_{F/\mathbb Q}(\mu)|^{-s}\), initially where it converges absolutely. The orbit quotient is necessary: summing over all nonzero lattice points repeats each norm infinitely often. Relating its continued value at 1 to a cusp signature defect also requires choices of cusp orientation, lattice/covolume normalization, and index realization.}
\GSfield{Reconstruction}{The orbit domain is a mathematical emendation, not a reading of the written \(N_F\). The numerical coefficient \(1/(2\pi)\) in the notebook has not been verified in Müller's original paper, so it is not promoted to a fully checked modern theorem.}
\GSfield{Fate}{The title and scope of \citet{muller1984signature} identify this application; the author's later \citet{muller1987cusps} includes the Hirzebruch conjecture in its final chapter. Only primary metadata was accessible for these exact formulas.}
\GSfield{Present status}{As of November 2026: Established subject and strong publication identification, but exact handwritten normalization remains unverified. This is not classified as an open number-theoretic problem.}
\GSfield{Assessment}{literal statement: correct with missing hypotheses; publication match: partial; verification: publisher metadata checked; development: none; confidence: medium}
\end{GSitem}

\GStalk{655}{1985-10-28}{B. L. Feigin}{P001221 (lower)-P001222}{Determinant lines and anomaly cancellation}
Feigin recasts the previous determinant calculation in the determinant of Dolbeault cohomology. Harmonic representatives give a Hodge metric, and the zeta determinant supplies the Quillen correction. Tensor products whose anomaly coefficients sum to zero carry a metric independent of the chosen conformal representative. The notebook uses the dual of Quillen's own determinant-line convention, explaining why its squared norm is divided, rather than multiplied, by the spectral determinant.

\begin{GSitem}{1985-655.1}{E}{The Dolbeault complex of a tangent tensor power}
\GSfield{Kind; source}{construction; \GSpage{P001221 (lower)}}
\GSfield{The notes}{The sequence is \(0\to H^0(T^j)\to\Gamma(T^j)\xrightarrow{\bar\partial}\Gamma(T^j\otimes\Omega^{0,1})\to H^1(T)\to0\).}
\GSfield{Statement}{On a closed smooth complex curve and for \(j\in\mathbb Z\), the Dolbeault resolution gives \(0\to H^0(T^j)\to C^\infty(T^j)\xrightarrow{\bar\partial}C^\infty(T^j\otimes\Omega^{0,1})\to H^1(T^j)\to0\). The final arrow is the quotient by the image of \(\bar\partial\).}
\GSfield{Reconstruction}{The final cohomology group must retain the exponent \(j\). The scan confirms that it is absent in the notebook, so this is an emendation of the mathematical record.}
\GSfield{Fate}{The complex is the standard setting of \citet[\S1]{quillen1985determinants}.}
\GSfield{Present status}{As of November 2026: Known Dolbeault theory; the last-term typo is substantive when \(j\ne1\).}
\GSfield{Assessment}{literal statement: false as written; publication match: exact; verification: read the result; development: none; confidence: high}
\end{GSitem}
\begin{GSitem}{1985-655.2}{E}{Harmonic representatives give cohomology norms}
\GSfield{Kind; source}{construction; \GSpage{P001222}}
\GSfield{The notes}{Hermitian metrics on both spaces induce metrics on \(H^0(T^j)\) and \(H^1(T^j)\); \(\ker\bar\partial\bar\partial^*\) is identified with the latter.}
\GSfield{Statement}{For a closed Riemann surface with smooth compatible metrics, \(H^0(T^j)=\ker\bar\partial\), and Hodge theory identifies \(H^1(T^j)\) with \(\ker\bar\partial^*\) on \((0,1)\)-forms, equivalently with \(\ker(\bar\partial\bar\partial^*)\). The \(L^2\) inner products on these finite-dimensional harmonic spaces define Hodge norms on their determinant lines.}
\GSfield{Reconstruction}{The compact boundaryless hypothesis makes the harmonic representative unique; no boundary condition is left implicit.}
\GSfield{Fate}{This construction is described in \citet[\S3]{quillen1985determinants}.}
\GSfield{Present status}{As of November 2026: Established and correct with the closed-surface hypothesis.}
\GSfield{Assessment}{literal statement: correct with missing hypotheses; publication match: exact; verification: read the result; development: none; confidence: high}
\end{GSitem}
\begin{GSitem}{1985-655.3}{E}{The determinant of cohomology and its dual convention}
\GSfield{Kind; source}{definition; \GSpage{P001222}}
\GSfield{The notes}{The notes define \(\operatorname{Det}T^j=\Lambda^{\max}H^0(T^j)\otimes(\Lambda^{\max}H^1(T^j))^*\), calling this Quillen's terminology.}
\GSfield{Statement}{Set \(\lambda(T^j)=\det H^0(T^j)\otimes\det H^1(T^j)^{-1}\), with \(\det0=\mathbb C\). This is one complex line. Quillen's operator determinant line for \(D=\bar\partial\) is \(\operatorname{DET}D=(\det\ker D)^{-1}\otimes\det\operatorname{coker}D=\lambda(T^j)^{-1}\). Both conventions are legitimate, but their metrics and sections are dual.}
\GSfield{Reconstruction}{The duality is a convention clarification, not an error in the notebook's definition. In an index-zero family the canonical determinant section is holomorphic in Quillen's convention; one must not transfer that assertion unchanged to its reciprocal.}
\GSfield{Fate}{The convention is explicitly checked against \citet[\S\S1--2]{quillen1985determinants}.}
\GSfield{Present status}{As of November 2026: A standard determinant line; the dual convention is recorded so subsequent signs and powers remain traceable.}
\GSfield{Assessment}{literal statement: correct; publication match: exact; verification: read the result; development: none; confidence: high}
\end{GSitem}
\begin{GSitem}{1985-655.4}{E}{The zeta correction of the Hodge metric}
\GSfield{Kind; source}{construction; \GSpage{P001222}}
\GSfield{The notes}{If \(\varepsilon\) is the Hermitian metric on \(\operatorname{Det}T^j\), the ``Ray--Singer torsion'' metric is \(\varepsilon/\operatorname{Det}\Delta_j\).}
\GSfield{Statement}{On the cohomological convention \(\lambda(T^j)\), the squared Quillen norm is \(\|\cdot\|_Q^2=\|\cdot\|_{L^2}^2/\det'\Delta_j\). Equivalently, on Quillen's dual line it is the Hodge squared norm multiplied by \(\det'\Delta_j\). This definition includes nontrivial kernels through the determinant of cohomology.}
\GSfield{Reconstruction}{``Metric'' is read as the Hermitian pairing or squared norm, not its square root. The prime and the dual convention explain the notebook's division. This is holomorphic analytic torsion of a Dolbeault complex, not an assertion of metric independence for ordinary de Rham torsion in every dimension.}
\GSfield{Fate}{Exact match after dualization to \citet[\S3, metric formula and Proposition]{quillen1985determinants}, published before the talk.}
\GSfield{Present status}{As of November 2026: Established construction. The factor is consistent and should not be corrected merely because a different convention multiplies by the determinant.}
\GSfield{Assessment}{literal statement: correct with missing hypotheses; publication match: exact; verification: read the result; development: none; confidence: high}
\end{GSitem}
\begin{GSitem}{1985-655.5}{E}{Anomaly-free tensor products of determinant lines}
\GSfield{Kind; source}{theorem; \GSpage{P001222}}
\GSfield{The notes}{For integers \(j_\alpha,k_\alpha\), the product metric on \(\bigotimes_\alpha(\operatorname{Det}T^{j_\alpha})^{k_\alpha}\) is unchanged under conformal variation if \(\sum_\alpha(6j_\alpha^2+6j_\alpha+1)k_\alpha=0\).}
\GSfield{Statement}{For a fixed closed complex curve, use induced metrics on the integer powers \(T^{j_\alpha}\) and the canonical zeta-corrected determinant-cohomology metrics, with one consistent normalization. Their local conformal variation is a common Liouville variation multiplied by \(a_{j_\alpha}=6j_\alpha^2+6j_\alpha+1\). Consequently the product metric is independent of the conformal representative whenever \(\sum k_\alpha a_{j_\alpha}=0\). Negative powers mean dual lines. Arbitrary independent finite counterterms are excluded.}
\GSfield{Reconstruction}{The common normalization and zero-mode cohomology factors are supplied. This is a sufficient anomaly-cancellation condition, not a claim that the determinant lines are individually metric independent or that all moduli-space functions are constant.}
\GSfield{Fate}{The ingredients are the pre-talk Quillen metric and Polyakov anomaly \citep{quillen1985determinants,polyakov1981quantum}; \citet[\S\S3,5]{belavinknizhnik1986complex} gives the ensuing string determinant-line formulation. General-spin coefficients there use \(K^n\), so \(n=-j\).}
\GSfield{Present status}{As of November 2026: An established consequence of the local anomaly calculation (E), with the later geometric formulation distinguished from the already known cancellation mechanism.}
\GSfield{Assessment}{literal statement: correct with missing hypotheses; publication match: partial; verification: read the result; development: none; confidence: high}
\end{GSitem}
\begin{GSitem}{1985-655.6}{E}{Descent to the moduli stack and the Faltings comparison}
\GSfield{Kind; source}{construction; \GSpage{P001222}}
\GSfield{The notes}{The one-dimensional spaces form a line bundle on the moduli space of genus-\(g\) curves. The resulting metric depends only on complex structure. A similar construction is tentatively attributed to Faltings.}
\GSfield{Statement}{For a smooth proper family of complex curves, determinant of cohomology forms a line bundle even when individual cohomology dimensions jump. The anomaly-free tensor product is natural under isomorphisms and therefore defines a Hermitian line on the moduli stack; on a fine level cover it is an ordinary line bundle. Descent to the coarse moduli space additionally requires trivial stabilizer action on the fibers.}
\GSfield{Reconstruction}{The notebook's word ``manifold'' suppresses automorphisms. The uncertain name is retained as a comparison, not identified with a specific Faltings theorem or a proof of the exact metric assertion.}
\GSfield{Fate}{The family construction is in \citet[\S\S2--3]{quillen1985determinants}; the string-moduli line is in \citet[\S5]{belavinknizhnik1986complex}.}
\GSfield{Present status}{As of November 2026: Established determinant geometry. The Faltings attribution is not sufficiently specific for a separate publication match.}
\GSfield{Assessment}{literal statement: correct with missing hypotheses; publication match: partial; verification: read the result; development: none; confidence: medium}
\end{GSitem}

\GStalk{656}{1985-11-04}{Yu. S. Ryzhikova}{P001223-P001224}{Perfect elements and cubicles in free modular lattices}
The heading expressly identifies this as an exposition of Gelfand and Ponomarev's work. It defines perfect elements through compatible complements and indecomposable representations, sketches the recursive upper and lower Boolean cubicles, and ends with complete irreducibility in characteristic zero. The exact source is the 1974 Russian Mathematical Surveys paper, with earlier seminar occurrences in 1973. Several indices and the cubicle cardinality need emendation; no completeness conjecture absent from these pages is imported from the source.

\begin{GSitem}{1985-656.1}{E}{Perfect elements and compatible direct-sum decompositions}
\GSfield{Kind; source}{definition; \GSpage{P001223}}
\GSfield{The notes}{An element \(a\in L\) is perfect if every representation \(\rho:L\to L(V)\) admits a complement \(U\) to \(\rho(a)\) splitting every \(\rho(x)\). Equivalently, in every indecomposable representation \(\rho(a)\) is zero or all of \(V\).}
\GSfield{Statement}{For a modular lattice \(L\), quantify over every field \(F\) and every finite-dimensional representation \(\rho:L\to L(V)\). Perfection means there exists \(V=\rho(a)\oplus U\) with \(\rho(x)=(\rho(x)\cap\rho(a))\oplus(\rho(x)\cap U)\) for every \(x\in L\). This is equivalent to \(\rho(a)\in\{0,V\}\) for every indecomposable such representation, by decomposing a finite-dimensional representation into indecomposables.}
\GSfield{Reconstruction}{Finite dimensionality is essential to the elementary decomposition argument and is explicit in the published definition. It should not be silently replaced by arbitrary infinite-dimensional lattice representations.}
\GSfield{Fate}{Exact definition and equivalence in \citet[\S1.2]{gelfandponomarev1974modular}; earlier occurrence 1973-319.7.}
\GSfield{Present status}{As of November 2026: Established expository material (E); the present speaker explicitly credits the 1974 work.}
\GSfield{Assessment}{literal statement: correct with missing hypotheses; publication match: exact; verification: read the result; development: none; confidence: high}
\end{GSitem}
\begin{GSitem}{1985-656.2}{E}{Modular lattice axioms and free generators}
\GSfield{Kind; source}{definition; \GSpage{P001223}}
\GSfield{The notes}{Associativity, commutativity, idempotence, absorption, the induced order, and \(a\le b\Rightarrow b(a+c)=a+bc\) are listed. The free modular lattice \(D^r\) has generators \(e_1,\ldots,e_r\).}
\GSfield{Statement}{A lattice has associative commutative idempotent join and meet satisfying absorption. The additional implication \(a\le b\Rightarrow b\wedge(a\vee c)=a\vee(b\wedge c)\) is modularity. The free modular lattice \(D^r\) has the universal property that any assignment of its \(r\) generators to a modular lattice extends uniquely to a lattice homomorphism.}
\GSfield{Reconstruction}{The notebook's heading ``definition of a lattice'' includes the extra modular axiom. Free modular lattices are not free distributive lattices, and equality in \(D^r\) is stronger than equality in every finite-dimensional linear representation.}
\GSfield{Fate}{\citet[\S\S1.1--1.3]{gelfandponomarev1974modular}; compare occurrences 1973-319.1--4.}
\GSfield{Present status}{As of November 2026: Standard definitions; no representability or word-problem conjecture is posed on this page.}
\GSfield{Assessment}{literal statement: correct with missing hypotheses; publication match: exact; verification: read the result; development: none; confidence: high}
\end{GSitem}
\begin{GSitem}{1985-656.3}{E}{The first upper Boolean cubicle}
\GSfield{Kind; source}{theorem; \GSpage{P001224 (top)}}
\GSfield{The notes}{Put \(h_t(1)=\sum_{j\ne t}e_j\). These generate \(B^+(1)\), a Boolean algebra. The illustrated list starts with \(0\) and then gives products of the \(h_t\).}
\GSfield{Statement}{In the free modular lattice with \(r\ge4\) generators, the sublattice generated by \(h_t(1)=\bigvee_{j\ne t}e_j\) is a Boolean lattice with \(2^r\) elements, and every one of its elements is perfect. Its top is \(\bigvee_je_j\), and its bottom is \(\bigwedge_th_t(1)\). This is a relative Boolean algebra: its bottom need not be the bottom of the ambient free lattice.}
\GSfield{Reconstruction}{The lower bound on \(r\) keeps the statement in the nondegenerate infinite-family setting. In the enumeration the empty meet is the cubicle top, not zero. Under a representation sending every \(e_j\) to the same nonzero space, the cubicle bottom also maps to that space, so identifying it with a universal zero is impossible. Editorial all-generators-equal representation check \ours; the cubicle theorem is due to Gelfand--Ponomarev.}
\GSfield{Fate}{\citet[\S1.3 and \S\S3--4]{gelfandponomarev1974modular}; this repeats 1973-319.9.}
\GSfield{Present status}{As of November 2026: Established Boolean cubicle construction. The notebook's endpoint notation needs correction, not a new theorem.}
\GSfield{Assessment}{literal statement: correct with missing hypotheses; publication match: exact; verification: read the result; development: annotated; confidence: high}
\end{GSitem}
\begin{GSitem}{1985-656.4}{E}{Adjacent-distinct words and the recursive cubicles}
\GSfield{Kind; source}{construction; \GSpage{P001224}}
\GSfield{The notes}{The page defines \(A(r,\ell)\), its subset ending in \(t\), a set \(\Gamma\), and the second-level terms \(e_t(2)\), \(h_t(2)\). Several subscripts conflict.}
\GSfield{Statement}{Let \(A(r,\ell)=\{(i_1,\ldots,i_\ell):1\le i_j\le r,\ i_j\ne i_{j+1}\}\). For \(\alpha\in A(r,\ell)\), set \(\Gamma(\alpha)=\{\beta=(k_1,\ldots,k_{\ell-1})\in A(r,\ell-1):k_j\notin\{i_j,i_{j+1}\}\}\). Define \(e_{(i)}=e_i\) and recursively \(e_\alpha=e_{i_1}\wedge\bigvee_{\beta\in\Gamma(\alpha)}e_\beta\). Then \(e_t(\ell)=\bigvee_{\alpha\in A_t(r,\ell)}e_\alpha\), \(h_t(\ell)=\bigvee_{j\ne t}e_j(\ell)\), and \(B^+(\ell)\) is generated by these \(h_t(\ell)\). In particular \(e_t(2)=\bigvee_{i\ne t}[e_i\wedge\bigvee_{k\notin\{i,t\}}e_k]\).}
\GSfield{Reconstruction}{The full recursion is reconstructed from the identified source; the notes only sketch it. The condition \(\beta\in A(r,\ell-1)\) also requires adjacent \(k\)'s to differ. The corrections are \(i_1\), not \(i_2\), in the inner second-level exclusion and \(i_j,i_{j+1}\) in \(\Gamma\).}
\GSfield{Fate}{Exact construction in \citet[\S1.3, eqs.~(1.1)--(1.3)]{gelfandponomarev1974modular}; earlier occurrence 1973-319.8.}
\GSfield{Present status}{As of November 2026: Established recursive construction; the published formula makes the manuscript's ambiguous indices recoverable.}
\GSfield{Assessment}{literal statement: false as written; publication match: exact; verification: read the result; development: none; confidence: high}
\end{GSitem}
\begin{GSitem}{1985-656.5}{E}{Perfectness, cardinality, dual cubicles, and their order}
\GSfield{Kind; source}{theorem; \GSpage{P001224}}
\GSfield{The notes}{All cubicles are Boolean algebras with \(2^n\) elements, the lower ones are dual, and ``all are ordered''. No completeness conjecture is stated.}
\GSfield{Statement}{For \(r\ge4\), each \(B^+(\ell)\) and its lattice-dual \(B^-(\ell)\) has \(2^r\) perfect elements. If \(\ell<m\), every element of \(B^+(\ell)\) lies above every element of \(B^+(m)\); duality reverses this for lower cubicles. The 1974 source proves comparison of arbitrary lower and upper cubicles after applying every linear representation, while explicitly distinguishing this from the corresponding unproved-at-that-time literal inequality in \(D^r\).}
\GSfield{Reconstruction}{Replace the unexplained \(n\) by the number \(r\) of generators. ``Ordered'' cannot mean that a Boolean cubicle is a total order. The source's separate conjecture that these exhaust all perfect elements is absent from these notes and is not added to this year's problem list.}
\GSfield{Fate}{\citet[\S1.3, pp.~4--5, and main theorem]{gelfandponomarev1974modular} supplies exactly these distinctions.}
\GSfield{Present status}{As of November 2026: The stated perfectness/cardinality theorem is established. The ambiguous short ordering phrase does not determine a further unresolved claim; no C/D classification is warranted.}
\GSfield{Assessment}{literal statement: correct with missing hypotheses; publication match: partial; verification: read the result; development: none; confidence: high}
\end{GSitem}
\begin{GSitem}{1985-656.6}{E}{Complete irreducibility over the rational prime field}
\GSfield{Kind; source}{definition; \GSpage{P001224 (bottom)}}
\GSfield{The notes}{In characteristic zero a representation is completely irreducible if its image is isomorphic to \(L(\mathbb Q^n)\), or projective geometry \(P(n-1,\mathbb Q)\).}
\GSfield{Statement}{Let \(V\) be finite dimensional over a characteristic-zero field \(F\). In Gelfand--Ponomarev's terminology, \(\rho:L\to L(V)\) is completely irreducible when its image lattice is isomorphic to the full lattice of rational linear subspaces of a finite-dimensional \(\mathbb Q\)-space (equivalently, the associated rational projective geometry, including its extremal subspaces). This is stronger than indecomposability and is not the assertion \(\rho(L)=L(V,F)\).}
\GSfield{Reconstruction}{The rational prime field is intentional even when the target field is larger. The notebook gives a definition, not the paper's full classification of the completely irreducible representations.}
\GSfield{Fate}{Exact definition in \citet[\S1.4]{gelfandponomarev1974modular}.}
\GSfield{Present status}{As of November 2026: Established specialized terminology; its intended field must be retained.}
\GSfield{Assessment}{literal statement: correct; publication match: exact; verification: read the result; development: none; confidence: high}
\end{GSitem}

\GStalk{657}{1985-11-04}{B. L. Feigin}{P001225-P001228}{Virasoro modules, Sugawara and coset constructions, and screening integrals}
Feigin links strings to Virasoro representation theory and hypergeometric functions. The notes state the Verma-module reducibility criterion, construct Virasoro actions from affine currents, take the coset difference, and introduce Selberg and Dotsenko--Fateev integrals. These are identifiable established constructions. Important faithful errors include the Kac-polynomial coefficient, incompatible central-charge signs, a missing noncritical-level hypothesis, and incomplete screening-cycle and integral conventions.

\begin{GSitem}{1985-657.1}{E}{Virasoro, the Witt quotient, and a central-sign convention}
\GSfield{Kind; source}{definition; \GSpage{P001225}}
\GSfield{The notes}{The bracket is \([e_i,e_j]=(j-i)e_{i+j}+\delta_{i,-j}(i^3-i)C/12\); the quotient by the center is called the algebra of vector fields on the circle. Positive-index modes are creators.}
\GSfield{Statement}{Over \(\mathbb C\), use \([L_m,L_n]=(m-n)L_{m+n}+\delta_{m,-n}(m^3-m)c/12\), with \(c\) central. The quotient is the Witt algebra of Laurent-polynomial vector fields, with \(L_m=-z^{m+1}\partial_z\). The notebook's bracket is isomorphic by \(e_i=L_{-i}\), \(C=-c\). Its algebraic direct sum is not the Fr\'echet Lie algebra of all smooth vector fields on the circle.}
\GSfield{Reconstruction}{This change of indices also changes the sign of the central generator. The Kac polynomial and positive Sugawara central charge on the next pages use the usual \(c\); they cannot be compared with the first bracket without this conversion. The central extension itself is legitimate. Editorial reversal-of-modes and central-sign computation \oursComputed.}
\GSfield{Fate}{The exact conventions and Verma-module setup can be checked in \citet[\S3.1]{feiginfuchs1982invariant}, received September 1981 and published in 1982.}
\GSfield{Present status}{As of November 2026: Established definition; a convention mismatch is recorded without treating the existence of the algebra as an open problem.}
\GSfield{Assessment}{literal statement: correct with missing hypotheses; publication match: exact; verification: read the result; development: annotated; confidence: high}
\end{GSitem}
\begin{GSitem}{1985-657.2}{E}{Verma modules and their grading}
\GSfield{Kind; source}{construction; \GSpage{P001225}}
\GSfield{The notes}{A vacuum satisfies \(e_0v=hv\), \(Cv=cv\), and \(e_{-i}v=0\) for \(i>0\). Dixmier and Bernstein--Gelfand--Gelfand are mentioned; the category name reads \ru{Оре}.}
\GSfield{Statement}{In the usual convention, let \(M(h,c)=U(\mathrm{Vir})\otimes_{U(\mathfrak b)}\mathbb C_{h,c}\), where \(\mathfrak b\) is spanned by \(L_n\) for \(n\ge0\) and the central generator, \(L_nv=0\) for \(n>0\), \(L_0v=hv\), and the center acts by \(c\). For arbitrary \(h,c\in\mathbb C\), the PBW basis consists of \(L_{-\lambda_1}\cdots L_{-\lambda_s}v\) with \(\lambda_1\ge\cdots\ge\lambda_s>0\); grade \(N\) has dimension the partition number \(p(N)\).}
\GSfield{Reconstruction}{The central eigenvalue is converted as in the preceding item. No integrability or unitarity assumption belongs to the general Verma module. The uncertain category word is not silently expanded into a particular title or paper.}
\GSfield{Fate}{The construction and grading occur in \citet[\S3.1]{feiginfuchs1982invariant}. The notebook's short bibliography does not identify the two BGG papers uniquely.}
\GSfield{Present status}{As of November 2026: Known highest-weight representation theory. The historical reading uncertainty is narrower than the fully identifiable mathematical construction.}
\GSfield{Assessment}{literal statement: correct with missing hypotheses; publication match: exact; verification: read the result; development: none; confidence: high}
\end{GSitem}
\begin{GSitem}{1985-657.3}{E}{The reducibility criterion and the miscopied Kac polynomial}
\GSfield{Kind; source}{question; \GSpage{P001225}}
\GSfield{The notes}{The question is when \(M_{h,c}\) is reducible. The displayed \(\Phi_{rs}\) has final coefficient \(1/2\), followed by an if-and-only-if statement attributed to Feigin and Fuchs.}
\GSfield{Statement}{For the usual Virasoro bracket, put \(c(t)=13-6(t+t^{-1})\) and \(h_{r,s}(t)=((rt-s)^2-(t-1)^2)/(4t)\), where \(t\in\mathbb C^\times\) and \(r,s\ge1\) are integers. Then \(M(h,c)\) is reducible exactly when \(c=c(t)\), \(h=h_{r,s}(t)\) for some such data. Equivalently some \(\Phi_{rs}(h,c)\) vanishes, with \(\Phi_{rs}=\prod_{j\in\{r,s\}}[h+(c-13)(j^2-1)/24+(rs-1)/2]+(r^2-s^2)^2/16\); the product uses both entries, even if \(r=s\).}
\GSfield{Reconstruction}{The manuscript visibly has \(1/2\), whereas the identified paper has \(1/16\). The difference is not a harmless scalar multiple of the polynomial. For example, at \(c=1,r=1,s=2\) the correct factor is \((h-1/4)^2\). The central-sign conversion from the first item is also required. Editorial substitution \(c=1,r=1,s=2\) \oursComputed; the corrected coefficient itself is read in Feigin--Fuchs.}
\GSfield{Fate}{\citet[\S3.2, Thm.~3.1 and its corollary, printed pp.~58--59]{feiginfuchs1982invariant} gives the determinant and reducibility result. This repeats problem GP-1982-03, occurrence 1982-566.7; publication preceded this talk by three years.}
\GSfield{Present status}{As of November 2026: Solved before the seminar. The theorem is E; its displayed polynomial is false as written.}
\GSfield{Assessment}{literal statement: false as written; publication match: exact; verification: read the result; development: annotated; confidence: high; problem \hyperref[prob:GP-1982-03]{GP-1982-03}}
\end{GSitem}
\begin{GSitem}{1985-657.4}{E}{Singular vectors and the first singular grade}
\GSfield{Kind; source}{theorem; \GSpage{P001225-P001226}}
\GSfield{The notes}{A reducible Verma module contains a non-vacuum vector killed by all annihilation modes, and its level is asserted to be \(rs\).}
\GSfield{Statement}{Every nonzero proper submodule of a Virasoro Verma module contains a homogeneous vector of positive grade killed by all \(L_n\), \(n>0\): choose its least occupied grade. If \(M(h,c)\) is reducible, its first such grade is \(\min\{rs:\Phi_{rs}(h,c)=0\}\). At a point on just one reducibility component this is the corresponding product \(rs\). Intersections require specifying which pair or which singular vector is meant.}
\GSfield{Reconstruction}{The elementary least-grade argument proves existence but does not by itself compute the grade. The determinant formula supplies the minimum. It is not legitimate to select an arbitrary pair among several vanishing factors and call its product the first level. Editorial least-grade proof of singular-vector existence \ours; computation of the first grade uses the published determinant theorem.}
\GSfield{Fate}{\citet[\S3.2, corollary to Thm.~3.1]{feiginfuchs1982invariant} explicitly gives the existence criterion for a singular vector of degree at most a prescribed integer.}
\GSfield{Present status}{As of November 2026: Established before 1985, with the collision qualification made explicit. It belongs to the same earlier indexed problem.}
\GSfield{Assessment}{literal statement: correct with missing hypotheses; publication match: exact; verification: read the result; development: annotated; confidence: high; problem \hyperref[prob:GP-1982-03]{GP-1982-03}}
\end{GSitem}
\begin{GSitem}{1985-657.5}{E}{Laurent vector fields, infinite matrices, and Fock realizations}
\GSfield{Kind; source}{method; \GSpage{P001226}}
\GSfield{The notes}{The first proof mechanism maps the Witt quotient into \(gl_\infty\), \(O_\infty\), or \(sp_\infty\), then uses suitable representations.}
\GSfield{Statement}{Laurent-polynomial vector fields act by Lie derivative on tensor densities. On Laurent half-densities the pairing \(\operatorname{Res}(fg\,dz)\) is symmetric, invariant, and nondegenerate. On \(\mathbb C[z,z^{-1}]/\mathbb C\), the pairing \(\operatorname{Res}(f\,dg)\) is alternating, invariant, and nondegenerate. In Laurent bases these actions have finitely many nonzero diagonals, not finitely many nonzero entries. Passing to a polarized wedge or spin representation requires its central extension and a specified completion.}
\GSfield{Reconstruction}{These explicit models explain the three classical-algebra labels without falsely embedding Witt in the finitary matrix algebra. They are an editorial reconstruction of a one-line construction, not three complete theorems asserted in the notes. Editorial residue-pairing reconstruction of the indicated classical-algebra models \ours; the Fock constructions are prior literature.}
\GSfield{Fate}{The relevant semi-infinite wedge mechanism is used in \citet[\S4]{feiginfuchs1982invariant}. Kac--Peterson's 1981 spin and wedge paper is an earlier primary reference identified in that paper; only its author-maintained metadata was checked here \citep{kacpeterson1981spin}.}
\GSfield{Present status}{As of November 2026: Established representation-theoretic method. The notebook does not specify enough completion data to identify every indicated model exactly.}
\GSfield{Assessment}{literal statement: correct with missing hypotheses; publication match: partial; verification: read the result; development: annotated; confidence: medium}
\end{GSitem}
\begin{GSitem}{1985-657.6}{E}{The affine current algebra and Gelfand's intervention}
\GSfield{Kind; source}{definition; \GSpage{P001226}; B. L. Feigin; I. M. Gelfand (intervention)}
\GSfield{The notes}{A simple finite-dimensional \(\mathfrak g\) is replaced by a one-dimensional extension of its current algebra. Gelfand explains currents as matrix functions on the circle. A Killing form and dual root/Cartan bases are chosen.}
\GSfield{Statement}{Let \(\mathfrak g\) be a finite-dimensional complex simple Lie algebra, with invariant form \(B\) normalized so long roots have squared length 2. Its untwisted affine algebra without the degree derivation is \(\widehat{\mathfrak g}=\mathfrak g[z,z^{-1}]\oplus\mathbb C K\), with \([xz^m,yz^n]=[x,y]z^{m+n}+m\delta_{m,-n}B(x,y)K\). A level-\(k\) module has \(K=k\,\mathrm{id}\). The basis used in the quadratic construction is dual with respect to this same form.}
\GSfield{Reconstruction}{The algebraic loop algebra suffices for the formulas. Smooth loops require a topological completion. The normalized form is generally a scalar multiple of the raw Killing form; changing it changes the numerical level and the denominator involving the dual Coxeter number.}
\GSfield{Fate}{The current-algebra setup appears in \citet[p.~132, eq.~(10)]{kent1985unitary}; the paper is a contemporary account of the already circulating Sugawara and coset constructions.}
\GSfield{Present status}{As of November 2026: Known construction. Gelfand's explanatory intervention is preserved without treating matrix-valued smooth functions and algebraic Laurent currents as identical categories.}
\GSfield{Assessment}{literal statement: correct with missing hypotheses; publication match: exact; verification: read the result; development: none; confidence: high}
\end{GSitem}
\begin{GSitem}{1985-657.7}{E}{Sugawara operators, normal ordering, and the noncritical level}
\GSfield{Kind; source}{construction; \GSpage{P001226-P001227}}
\GSfield{The notes}{The quadratic sums \(S(m)\) are called elements of the enveloping algebra; one inner sum has \(r\in\Delta\). The notes give their commutator and divide by \(2(g+k)\).}
\GSfield{Statement}{On a restricted level-\(k\) \(\widehat{\mathfrak g}\)-module, each vector is killed by \(xz^m\) for sufficiently large positive \(m\). If \(k\ne-h^\vee\), the normally ordered operators \(L_n=[2(k+h^\vee)]^{-1}\sum_a\sum_{m\in\mathbb Z}:J^a_mJ_{a,n-m}:\) are well defined on each vector and satisfy the Virasoro relations with \(c=k\dim\mathfrak g/(k+h^\vee)\), as well as \([L_n,J^a_m]=-mJ^a_{n+m}\). Normal ordering places annihilation modes to the right; the zero-mode convention is fixed consistently with the quadratic Casimir.}
\GSfield{Reconstruction}{The root sum is over \(\alpha\in\Delta\), but both mode sums are over \(\mathbb Z\). The infinite quadratic expression lies in a suitable completion, or defines operators on a restricted module, not in the ordinary finite-sum enveloping algebra. At the critical level the displayed division is undefined. Indices and the central sign are converted to the usual convention. Editorial quadratic-scaling and central-sign check \oursComputed; the Sugawara theorem itself is from the cited source.}
\GSfield{Fate}{\citet[p.~133, eqs.~(15)--(17)]{kent1985unitary} gives the Sugawara formula and charge. Multiplication by \(2(k+h^\vee)\) recovers the noncentral scaling of the quadratic commutator. Reversing mode indices also reverses the central sign; the notebook's displayed positive central term and later standard positive charge are not one consistent convention.}
\GSfield{Present status}{As of November 2026: Established construction. The missing denominator hypothesis, summation domain, completion, and form normalization are essential mathematical qualifications.}
\GSfield{Assessment}{literal statement: false as written; publication match: exact; verification: read the result; development: annotated; confidence: high}
\end{GSitem}
\begin{GSitem}{1985-657.8}{E}{The affine sl2 example and the spin/Dynkin-label mismatch}
\GSfield{Kind; source}{construction; \GSpage{P001227}}
\GSfield{The notes}{For \(\widehat{sl}_2\), the notes put \(k=k_1+k_2\), \(c=3k/(k+2)\), and \(h=k_1(k_1+1)/(k+2)\).}
\GSfield{Statement}{For a highest-weight \(\widehat{sl}_2\)-module of level \(k\ne-2\), let \(\ell\) denote the finite \(sl_2\) highest weight in Dynkin-label normalization. Then the Sugawara highest weight is \(h=\ell(\ell+2)/(4(k+2))\) and \(c=3k/(k+2)\). If \(j=\ell/2\) is the spin, this becomes \(h=j(j+1)/(k+2)\). Affine Dynkin labels satisfy \(\lambda_0+\lambda_1=k\), with \(\ell=\lambda_1\).}
\GSfield{Reconstruction}{The notebook uses \(k_1\) both as an affine label contributing with coefficient one to \(k\) and as a spin. Those are different normalizations. Neither integrality nor unitarity is needed for this formal highest-weight formula; integrable unitary modules further have \(k\in\mathbb Z_{\ge0}\), \(0\le\ell\le k\). Editorial spin-to-Dynkin-label conversion \oursComputed.}
\GSfield{Fate}{The \(SU(2)\) charge and highest-weight formulas are stated in \citet[p.~133]{kent1985unitary}, immediately following the Sugawara construction.}
\GSfield{Present status}{As of November 2026: A standard example with a genuine normalization mismatch in its paired formulas.}
\GSfield{Assessment}{literal statement: false as written; publication match: exact; verification: read the result; development: annotated; confidence: high}
\end{GSitem}
\begin{GSitem}{1985-657.9}{E}{The coset difference commutes with the subalgebra currents}
\GSfield{Kind; source}{theorem; \GSpage{P001227}}
\GSfield{The notes}{For \(\mathfrak g\supset\mathfrak g_1\), subtract the two Sugawara actions. The difference is Virasoro, commutes with \(\widehat{\mathfrak g}_i\), and has the difference of the central charges.}
\GSfield{Statement}{Let \(\mathfrak h\subset\mathfrak g\) be complex simple Lie algebras, with embedding index \(I\) in the normalized invariant forms. On a restricted level-\(k\) \(\widehat{\mathfrak g}\)-module, the induced \(\widehat{\mathfrak h}\)-level is \(Ik\). Assume \(k+h^\vee_{\mathfrak g}\ne0\) and \(Ik+h^\vee_{\mathfrak h}\ne0\). Then \(L_n^{\mathrm{coset}}=L_n^{\mathfrak g}-L_n^{\mathfrak h}\) commutes with all \(\widehat{\mathfrak h}\) currents and satisfies Virasoro with \(c_{\mathrm{coset}}=k\dim\mathfrak g/(k+h^\vee_{\mathfrak g})-Ik\dim\mathfrak h/(Ik+h^\vee_{\mathfrak h})\).}
\GSfield{Reconstruction}{The subscript is \(1\), not the running mode index \(i\). The embedding index and both noncritical conditions are supplied. The charge here is a scalar on the fixed-level module; subtracting abstract central generators without fixing levels is not this statement.}
\GSfield{Fate}{\citet[pp.~133--134, eqs.~(18)--(20)]{kent1985unitary} states the construction, commutation, and charge difference, as an account of the Goddard--Kent--Olive work.}
\GSfield{Present status}{As of November 2026: Established by 1985. The seminar is an exposition of the coset mechanism, not an unattributed new open problem.}
\GSfield{Assessment}{literal statement: correct with missing hypotheses; publication match: exact; verification: read the result; development: none; confidence: high}
\end{GSitem}
\begin{GSitem}{1985-657.10}{E}{The Selberg integral: exponent and chamber normalization}
\GSfield{Kind; source}{theorem; \GSpage{P001227}}
\GSfield{The notes}{The first integral has a Vandermonde power \(\rho\), unspecified product ranges, and limits 0 and 1. Its gamma product starts with \(\prod_{i=1}^m\Gamma(i\rho)/\Gamma(\rho)\). It is called a special case of Fateev--Dotsenko.}
\GSfield{Statement}{For \(m\ge1\), real \(\alpha,\beta>-1\) and \(\rho>0\), the integral over \(0<t_1<\cdots<t_m<1\) of \(\prod_{i<j}(t_j-t_i)^{2\rho}\prod_i t_i^\alpha(1-t_i)^\beta\,dt_1\cdots dt_m\) equals exactly the gamma product printed in (17). The symmetric integral over \([0,1]^m\), with \(|t_i-t_j|^{2\rho}\), is \(m!\) times that product. Meromorphic continuation is a later statement about the resulting functions of the exponents, not unconditional convergence of the original integral.}
\GSfield{Reconstruction}{These are sufficient convergence hypotheses. The Vandermonde exponent is doubled, and an ordered chamber is specified to match the printed right side; alternatively keep the cube and multiply that side by \(m!\). No branch ambiguity remains under the displayed real hypotheses. Editorial chamber-to-cube normalization check \oursComputed; the Selberg evaluation is classical.}
\GSfield{Fate}{This is Selberg's classical integral, announced in 1941 and proved in 1944; see the authors' derivation and historical account in \citet[eqs.~(1.1)--(1.2) and pp.~1--2]{forresterwarnaar2008selberg}. Dotsenko--Fateev's use is a later conformal-field-theory application, not its origin.}
\GSfield{Present status}{As of November 2026: Established well before the seminar; the exact notebook formula fails without correcting its exponent and integration convention.}
\GSfield{Assessment}{literal statement: false as written; publication match: exact; verification: read the result; development: annotated; confidence: high}
\end{GSitem}
\begin{GSitem}{1985-657.11}{E}{Two screening charges and the missing integration cycle}
\GSfield{Kind; source}{construction; \GSpage{P001228}}
\GSfield{The notes}{The displayed integral has \(n\) charges \(\alpha_+\), \(m\) charges \(\alpha_-\), \(\alpha_+\alpha_-=-2\), same-species powers \(\alpha_\pm^2\), mixed power \(-2\), and powers of \(t_i\) and \(z-t_i\). One product ends at \(m+1\) and the differentials end at \(dt_n\).}
\GSfield{Statement}{In a chiral free-boson normalization with charge weight \(\Delta(q)=q(q-2q_0)/2\), choose screening charges solving \(\Delta(\alpha_\pm)=1\); then \(\alpha_++\alpha_-=2q_0\) and \(\alpha_+\alpha_-=-2\). For nonnegative integers \(n,m\), impose sphere charge neutrality \(\sum_aq_a+n\alpha_++m\alpha_-=2q_0\), including the charge at infinity. Wick contractions produce pair factors with exponents equal to charge products. Integrating the screening variables over compatible complex contours or twisted cycles gives chiral blocks, when those integrals converge or have a specified regularization. The mixed range is \(1\le u\le n<n+1\le v\le n+m\), and all \(n+m\) differentials are required.}
\GSfield{Reconstruction}{The charge convention is specified explicitly; Dotsenko--Fateev also use a rescaled convention. For four insertions fixed at \(0,z,1,\infty\), powers of \(1-t_i\), generally different species-dependent external-charge powers, branch choices, and cycles must also be supplied. A common real cube with both species present is not an absolutely convergent integral across the mixed double poles. The missing external charges cannot be uniquely recovered from this page. Editorial mixed-double-pole divergence check \ours; the contour construction is due to Dotsenko--Fateev.}
\GSfield{Fate}{\citet[\S3, especially eqs.~(3.30)--(3.36)]{dotsenkofateev1984conformal} constructs screening integrals and discusses their contours; the paper was received 17 April 1984 and published in 1984.}
\GSfield{Present status}{As of November 2026: An identifiable established method; the literal integral is incomplete and its ordinary-real-integral interpretation is false.}
\GSfield{Assessment}{literal statement: false as written; publication match: partial; verification: read the result; development: annotated; confidence: high}
\end{GSitem}
\begin{GSitem}{1985-657.12}{E}{Four-point blocks depend on a cross-ratio}
\GSfield{Kind; source}{claim; \GSpage{P001228}}
\GSfield{The notes}{The integral is first called a two-point correlation function and then a four-point function depending on the cross-ratio \(z\).}
\GSfield{Statement}{For four distinct points on \(\mathbb P^1\), projective transformations fix three to \(0,1,\infty\) and leave one complex cross-ratio \(z\). Chiral four-point conformal blocks may therefore have one nontrivial complex argument. A physical single-valued nonchiral correlation is an appropriate sesquilinear combination of such blocks and their antiholomorphic partners; a chiral contour integral alone is not automatically that correlation.}
\GSfield{Reconstruction}{The lower four-point statement resolves the internal two-/four-point mismatch. A two-point function on the sphere has no independent cross-ratio. The distinction between a block and a full correlation is added because the notebook uses the latter term indiscriminately.}
\GSfield{Fate}{\citet[\S3, eqs.~(3.39)--(3.42), and \S4]{dotsenkofateev1984conformal} explicitly fixes the four points and combines the solutions.}
\GSfield{Present status}{As of November 2026: Known conformal covariance and block construction; the first label is a faithful mathematical slip.}
\GSfield{Assessment}{literal statement: false as written; publication match: exact; verification: read the result; development: none; confidence: high}
\end{GSitem}
\begin{GSitem}{1985-657.13}{F}{Analytic continuation in the number of screenings}
\GSfield{Kind; source}{claim; \GSpage{P001228 (bottom)}}
\GSfield{The notes}{The ``real'' correlation functions are said to be analytic continuations of the integral formulas to complex \(n,m\).}
\GSfield{Statement}{A sequence of integrals defined only for \((n,m)\in\mathbb Z_{\ge0}^2\) does not by itself determine a unique holomorphic or meromorphic function of two complex variables. If one continuation \(F(n,m)\) exists, then \(F(n,m)+\sin(\pi n)H(n,m)\) has the same values on every integer pair for any holomorphic \(H\). A well-posed continuation claim therefore needs a target class, singularity and growth conditions, functional equations, or an independently defined correlation function.}
\GSfield{Reconstruction}{No such selection principle is stated. This is not the familiar analytic continuation of the Selberg integral in its continuous exponents: there the integral is first defined on an open parameter domain. The displayed ambiguity is an elementary editorial check. Editorial nonuniqueness example using \(\sin(\pi n)\) \ours.}
\GSfield{Fate}{The source screening construction is \citet[\S3]{dotsenkofateev1984conformal}. For comparison, \citet[pp.~1--2]{forresterwarnaar2008selberg} carefully invokes a growth theorem in a continuation argument; neither source establishes the notebook's unrestricted two-integer-variable assertion.}
\GSfield{Present status}{As of November 2026: Insufficiently formulated (F), not an open problem licensed by the vague claim. No universal continuation theorem is asserted here.}
\GSfield{Assessment}{literal statement: not determinable; publication match: thematic only; verification: read the result; development: annotated; confidence: high}
\end{GSitem}

\GStalk{658}{1985-11-11}{A. M. Vershik}{P001229-P001231}{Configuration spaces, polytope universality and lattices}
The talk moves from Grassmannian configuration strata and isotopy of classes defined by determinant signs to realization spaces of polytopes. Mn\"ev universality was already circulating and published; the statement about simplicial four-polytopes retains an informal hedge and is not established by either general four-dimensional universality or variable-dimensional neighborly universality. The final part concerns geometric lattices, the Fano obstruction, lattice representations and the Pudl\'ak--T\r{u}ma theorem. Gelfand adds a historical intervention. The manuscript begins on P001229 and ends on the upper part of P001231; the supplied catalogue range overlaps adjacent talks.

\begin{GSitem}{1985-658.1}{E}{Topology of Grassmannian configuration strata}
\GSfield{Kind; source}{problem; \GSpage{P001229}}
\GSfield{The notes}{For \(E\in G(n,m)\), the notes set \(\tau_E=(P_Ee_1,\ldots,P_Ee_n)\) and ask for the topology of the stratum with prescribed combinatorial type.}
\GSfield{Statement}{Fix labels and a rank-\(m\) incidence or oriented-matroid type \(\mathcal M\) on \(n\) elements. Study the locus of \(m\)-planes whose coordinate projections realize \(\mathcal M\). The oriented version fixes signs of maximal minors after an orientation is chosen; the unoriented version fixes ranks of subsets. These are different stratifications.}
\GSfield{Reconstruction}{The notebook does not choose between incidence data, an oriented matroid, and a face lattice. The statement makes these alternatives explicit rather than identifying their spaces.}
\GSfield{Fate}{Configuration varieties and their normalized forms were already treated in \citet[definitions and Theorems~1--3]{mnev1985varieties}, published in 1985 from a manuscript received in July 1984. The broader question is therefore expository here, although its fixed-dimensional polytope variant below is not settled by that paper.}
\GSfield{Present status}{As of November 2026, universality rules out a general simple classification of these strata. Exact topology requires specifying the combinatorial data and normalization.}
\GSfield{Assessment}{literal statement: correct with missing hypotheses; publication match: partial; verification: read the result; development: none; confidence: high; problem \hyperref[prob:GP-1985-40]{GP-1985-40}}
\end{GSitem}
\begin{GSitem}{1985-658.2}{E}{Every polytope type as an orthogonal projection of a simplex}
\GSfield{Kind; source}{construction; \GSpage{P001229}}
\GSfield{The notes}{The exercise asks to obtain combinatorially every \(m\)-dimensional convex polytope with \(n\) vertices by projecting \(S=\operatorname{conv}(e_1,\ldots,e_n)\subset\mathbb R^n\) to an \(m\)-plane.}
\GSfield{Statement}{If \(Q=\operatorname{conv}(v_1,\ldots,v_n)\subset\mathbb R^m\) is full-dimensional, let \(A:\mathbb R^n\to\mathbb R^m\) have columns \(v_i\), and put \(E=(\ker A)^\perp\). Then \(A|_E\) is invertible and \(A=(A|_E)P_E\). Thus \(Q\) is linearly isomorphic to \(P_E S\), so their labelled face lattices agree.}
\GSfield{Reconstruction}{Full-dimensionality ensures that \(A\) has rank \(m\). The proof supplies the omitted linear isomorphism; the projected vertices need not equal the given numerical coordinates. \ours\ The factorization through \((\ker A)^\perp\) is the elementary proof supplied here; no universality assertion is being proved by it.}
\GSfield{Fate}{This is an elementary linear-algebra verification of the notebook exercise; it does not require a new universality theorem. The configuration/Grassmannian framework is contemporary with \citet[p.~1312]{mnev1985varieties}.}
\GSfield{Present status}{As of November 2026, this remains a standard result under the stated hypotheses.}
\GSfield{Assessment}{literal statement: correct; publication match: exact; verification: read the result; development: annotated; confidence: high}
\end{GSitem}
\begin{GSitem}{1985-658.3}{E}{Isotopy of matrices with prescribed determinant signs}
\GSfield{Kind; source}{question; \GSpage{P001229-P001230}}
\GSfield{The notes}{Two \(3\times n\) matrices \(x,x'\) have all minors nonzero with the same signs. Can one join them through matrices in the same class?}
\GSfield{Statement}{For a realizable uniform rank-three chirotope \(\chi\), is the space \(\{A\in\mathbb R^{3\times n}:\operatorname{sign}\det A_I=\chi(I)\text{ for every }I\in\binom{[n]}3\}\) path connected? With an oriented basis fixed, this is the corresponding normalized realization-space question.}
\GSfield{Reconstruction}{Maximal \(3\times3\) minors are intended in the oriented-matroid interpretation. Literally imposing signs of minors of every size is an additional condition not identified by these references.}
\GSfield{Fate}{The general negative answer follows from rank-three universality already announced by \citet[Theorem~3]{mnev1985varieties}; the later complete uniform formulation is recorded in \citet[Theorem~1(iii)]{adiprasitopadrol2017universality}.}
\GSfield{Present status}{As of November 2026, the general isotopy assertion is false; realization spaces may be disconnected and have much more complicated topology.}
\GSfield{Assessment}{literal statement: correct with missing hypotheses; publication match: exact; verification: read the result; development: none; confidence: high; problem \hyperref[prob:GP-1985-41]{GP-1985-41}}
\end{GSitem}
\begin{GSitem}{1985-658.4}{E}{The small-configuration bound}
\GSfield{Kind; source}{claim; \GSpage{P001230}}
\GSfield{The notes}{The first answer recorded is \(n\leq7\Rightarrow\) a path exists.}
\GSfield{Statement}{For every realizable rank-three oriented matroid on at most seven labelled elements, its normalized realization space is path connected. The uniform version gives the positive answer to the preceding question.}
\GSfield{Fate}{\citet[introduction, p.~287]{tsukamoto2013disconnected} records the stronger established bound of fewer than nine elements, citing the classical small-configuration literature. This verifies the notebook inequality but does not identify its precise historical source.}
\GSfield{Present status}{As of November 2026, the stated bound is valid and not optimal as a general small-size assertion.}
\GSfield{Assessment}{literal statement: correct with missing hypotheses; publication match: exact; verification: indirect; development: none; confidence: medium; problem \hyperref[prob:GP-1985-41]{GP-1985-41}}
\end{GSitem}
\begin{GSitem}{1985-658.5}{F}{The nineteen-column counterexample attributed to Mnëv}
\GSfield{Kind; source}{claim; \GSpage{P001230}}
\GSfield{The notes}{The notes state \(n=19\Rightarrow\) no path, with N.~E.~Mn\"ev named.}
\GSfield{Statement}{There should exist two realizations of a single uniform rank-three oriented matroid on nineteen labelled elements lying in distinct path components of its realization space.}
\GSfield{Reconstruction}{The formulation assumes that ``all minors'' means maximal minors. A later smaller nonuniform example does not verify this exact claim.}
\GSfield{Fate}{The qualitative counterexample is compatible with Mn\"ev universality, but the exact nineteen-element construction and its uniformity were not located in the opened primary texts. \citet[Theorem~1.3]{tsukamoto2013disconnected} constructs thirteen-element disconnected rank-three spaces with prescribed collinearities; these are not substitutes for the notebook's nonzero-minor assertion.}
\GSfield{Present status}{As of November 2026, the general phenomenon is established; the notebook's exact numerical example remains an unresolved identification in this companion.}
\GSfield{Assessment}{literal statement: not determinable; publication match: partial; verification: unresolved lead; development: none; confidence: low; problem \hyperref[prob:GP-1985-41]{GP-1985-41}}
\end{GSitem}
\begin{GSitem}{1985-658.6}{E}{Universality for real configuration spaces}
\GSfield{Kind; source}{theorem; \GSpage{P001230}}
\GSfield{The notes}{The topology can be that of any real algebraic variety.}
\GSfield{Statement}{Every primary basic semialgebraic set defined over \(\mathbb Z\) is stably equivalent to a realization space of a rank-three oriented matroid. For an open such set, the oriented matroid may be chosen uniform. In particular these spaces realize arbitrary finite simplicial-complex homotopy types; equality or homeomorphism to an arbitrary variety is not asserted.}
\GSfield{Reconstruction}{Stable equivalence, integral defining equations and the uniform/open restriction replace the notebook's informal universal quantifier. These qualifications carry arithmetic information beyond homotopy type.}
\GSfield{Fate}{\citet[Theorem~3]{mnev1985varieties} announced the open-set version before the seminar. The standard stable-equivalence formulation and the open/uniform qualification are stated in \citet[Theorem~1]{adiprasitopadrol2017universality}; the latter first circulated in February 2014, appeared online in May 2016, and was printed in 2017.}
\GSfield{Present status}{As of November 2026, this remains a standard result under the stated hypotheses.}
\GSfield{Assessment}{literal statement: correct with missing hypotheses; publication match: exact; verification: read the result; development: none; confidence: high; problem \hyperref[prob:GP-1985-40]{GP-1985-40}}
\end{GSitem}
\begin{GSitem}{1985-658.7}{F}{A complex analogue of convexity}
\GSfield{Kind; source}{question; \GSpage{P001230}}
\GSfield{The notes}{Over \(\mathbb C\), the analogue is said to be computing \(\pi_1\); the notes then ask what the convex analogue over \(\mathbb C\) is.}
\GSfield{Statement}{Two prompts are recorded: determine fundamental groups of suitably defined complex configuration strata, and identify an appropriate complex counterpart of real convex-polytopal realization spaces. No category of complex convex objects or equivalence relation is specified.}
\GSfield{Fate}{The real universality statements in \citet{mnev1985varieties} do not select a unique complex problem. Searches for complex configuration varieties and realization spaces did not identify an exact result matching both prompts.}
\GSfield{Present status}{As of November 2026, no precise open problem is extracted: choosing a complex analogue would itself be an editorial addition.}
\GSfield{Assessment}{literal statement: not determinable; publication match: none found; verification: unresolved lead; development: none; confidence: low}
\end{GSitem}
\begin{GSitem}{1985-658.8}{E}{Steinitz characterization of three-dimensional polytope graphs}
\GSfield{Kind; source}{theorem; \GSpage{P001230}}
\GSfield{The notes}{Steinitz's theorem is cited as characterizing the one-dimensional skeleton of a convex polytope in \(\mathbb R^3\).}
\GSfield{Statement}{A finite simple graph is the graph of a convex three-dimensional polytope if and only if it is planar and three-vertex-connected.}
\GSfield{Fate}{\citet[section~2]{richtergebertziegler1995realization} restates the classical theorem and dates the original result to Steinitz, 1922. The theorem predates this seminar by decades.}
\GSfield{Present status}{As of November 2026, this remains a standard result under the stated hypotheses.}
\GSfield{Assessment}{literal statement: correct; publication match: exact; verification: read the result; development: none; confidence: high}
\end{GSitem}
\begin{GSitem}{1985-658.9}{E}{Contractibility of normalized three-polytope realization spaces}
\GSfield{Kind; source}{theorem; \GSpage{P001230}}
\GSfield{The notes}{A ``corollary'' says that the space of three-polytopes of a fixed combinatorial type is contractible.}
\GSfield{Statement}{For a realizable labelled three-polytope type, its realization space after fixing an affine basis is contractible; a standard normalized realization-space model is an open ball. The unnormalized collection of all vertex coordinates is not generally contractible.}
\GSfield{Reconstruction}{An affine basis and orientation component must be specified. This does not follow merely from the graph characterization without the stronger realization-space form of Steinitz's theorem.}
\GSfield{Fate}{\citet[section~2, Steinitz theorem]{richtergebertziegler1995realization} gives the normalized open-ball assertion. \citet[p.~1312]{mnev1985varieties} explicitly describes orientation components in an unnormalized model, showing why the normalization cannot be omitted.}
\GSfield{Present status}{As of November 2026, this remains a standard result under the stated hypotheses.}
\GSfield{Assessment}{literal statement: correct with missing hypotheses; publication match: exact; verification: read the result; development: none; confidence: high}
\end{GSitem}
\begin{GSitem}{1985-658.10}{C}{Universality for simplicial four-polytope realization spaces}
\GSfield{Kind; source}{claim; \GSpage{P001230}}
\GSfield{The notes}{In \(\mathbb R^4\), even simplicial polytopes are said to have \ru{более или менее любой} (``more or less any'') topology for a fixed combinatorial type.}
\GSfield{Statement}{A precise homotopy interpretation would ask: for every nonempty finite simplicial complex \(K\), is there a labelled simplicial four-polytope whose appropriately normalized realization space is homotopy equivalent to \(K\)? The notebook neither specifies this equivalence nor supplies exact quantifiers. A stronger literature conjecture asks for stable equivalence to every nonempty open primary basic semialgebraic set, and even requires the four-polytope to be neighborly.}
\GSfield{Reconstruction}{The displayed homotopy question is an editorial sharpening of the hedge, not a verbatim conjecture attributed to Vershik. Stable-equivalence and homotopy versions must remain separate. Adiprasito--Padrol use a homogenized cone-realization model modulo \(\mathrm{GL}_{d+1}(\mathbb R)\), whereas the bounded proof here fixes an affine frame and a stacking history; their normalizations are not silently identified.}
\GSfield{Fate}{General four-polytope universality was proved in work of 1994 and published by \citet[Main Theorem]{richtergebertziegler1995realization} in 1995. \citet[Theorem~2]{adiprasitopadrol2017universality}, circulated in 2014, published online in 2016 and printed in 2017, proves simplicial/neighborly universality with dimension growing as \(2n-4\); their Conjecture~4 retains dimension four and additionally requires neighborliness. These are distinct results. Adiprasito--Padrol explicitly judge the earlier simplicial proof incomplete and supply its missing argument; the 1985 announcement is not itself a fully checked proof of that part.}
\GSfield{Present status}{As of November 2026, no resolution of fixed-dimensional simplicial universality was found in this bounded search. The latest opened source explicitly recording the constant-dimension simplicial problem as open is the 2024 compendium, entries CG68--CG69. The chapter proves only a bounded contractibility result for fixed sequences of facet stackings.}
\GSfield{Assessment}{literal statement: not determinable; publication match: partial; verification: read the result; development: developed in this chapter; confidence: high; problem \hyperref[prob:GP-1985-42]{GP-1985-42}}
\GSfield{Worth developing}{Worth developing in bounded families: explicit insertion fibers distinguish tractable history-compatible spaces from the unresolved universal quantifier.}
\end{GSitem}
\begin{GSitem}{1985-658.11}{F}{Fields of polytopes in tangent spaces}
\GSfield{Kind; source}{problem; \GSpage{P001230}}
\GSfield{The notes}{The stated task is to construct a field of polytopes, one in every tangent space of a manifold; a classification problem is mentioned.}
\GSfield{Statement}{For a smooth manifold \(M\), one could prescribe a convex-polytope type \(P\) and ask for a continuous section of the bundle of copies of \(P\) inside \(TM\), then classify such sections up to a specified equivalence. The notebook does not prescribe regularity, a type, or equivalence.}
\GSfield{Reconstruction}{The associated-bundle formulation is a possible modern interpretation only.}
\GSfield{Fate}{\citet[bibliography, item~6]{mnev1985varieties} points to Vershik--Chernyakov's 1982 work on fields of cones as a nearby motivation. This is a lead, not a verified solution of the notebook's classification programme.}
\GSfield{Present status}{As of November 2026, the recorded programme is too broad to assign a definite unresolved theorem or a new problem-index identifier.}
\GSfield{Assessment}{literal statement: not determinable; publication match: thematic only; verification: indirect; development: none; confidence: low}
\end{GSitem}
\begin{GSitem}{1985-658.12}{E}{Geometric lattices and finite matroids}
\GSfield{Kind; source}{definition; \GSpage{P001230}}
\GSfield{The notes}{A geometric lattice is described as semimodular and ``finite'', with a parenthesis saying that every element is a finite join of atoms.}
\GSfield{Statement}{A finite lattice is geometric when it is atomistic and upper semimodular. Atomistic means every element is a join of atoms; finiteness and atomisticity are separate requirements. The lattice of flats of a finite simple matroid is geometric, and conversely every finite geometric lattice arises this way.}
\GSfield{Fate}{This is classical matroid theory, not a new definition in 1985. \citet[sections~1--2]{oxleybriefly} provides a modern author exposition of the underlying finite-matroid definitions; the notebook's conflation is corrected explicitly.}
\GSfield{Present status}{As of November 2026, this remains a standard result under the stated hypotheses.}
\GSfield{Assessment}{literal statement: false as written; publication match: partial; verification: indirect; development: none; confidence: high}
\end{GSitem}
\begin{GSitem}{1985-658.13}{E}{Representability of a geometric lattice by a vector configuration}
\GSfield{Kind; source}{question; \GSpage{P001230}}
\GSfield{The notes}{After ``configuration --- geometric lattice'', the converse question asks when a geometric lattice is given by a configuration.}
\GSfield{Statement}{Given a finite geometric lattice \(L\) and a specified field \(F\), does there exist a finite vector configuration over \(F\) whose lattice of flats is isomorphic to \(L\)? For \(F=\mathbb R\), prescribed ranks can be expressed by vanishing minors and nonvanishing minors, producing a finite semialgebraic feasibility problem.}
\GSfield{Reconstruction}{The field is supplied explicitly. This matroid representation problem is not equivalent to arbitrary join-preserving embeddings of the same lattice into a subspace lattice.}
\GSfield{Fate}{The field dependence is fundamental in classical representability; see \citet[section~4, Theorem~4.1]{oxleybriefly}. The semialgebraic encoding places real representability in the setting of \citet{mnev1985varieties}; it is not a simple combinatorial criterion.}
\GSfield{Present status}{As of November 2026, individual finite real instances are decidable by real quantifier elimination, while universality explains their potential complexity.}
\GSfield{Assessment}{literal statement: correct with missing hypotheses; publication match: partial; verification: read the result; development: none; confidence: high}
\end{GSitem}
\begin{GSitem}{1985-658.14}{E}{The Fano configuration and characteristic two}
\GSfield{Kind; source}{construction; \GSpage{P001230}}
\GSfield{The notes}{The seven-point diagram is called the Fano configuration and declared unrealizable over \(\mathbb R\) or \(\mathbb Z\). The scan shows only six of the seven Fano lines.}
\GSfield{Statement}{The Fano matroid \(F_7\) is representable over a field \(F\) exactly when \(\operatorname{char}F=2\). A representation in characteristic two uses the seven nonzero columns of \(\mathbb F_2^3\). No real matrix has exactly its rank incidences; consequently no integer matrix does either.}
\GSfield{Reconstruction}{The named Fano matroid supplies the missing seventh line, but no absent stroke is restored to the manuscript diagram. The characteristic claim is valid for the named object, not for every seven-point drawing with those six lines.}
\GSfield{Fate}{\citet[Theorem~4.1]{oxleybriefly} states the exact characteristic criterion. The result is classical and predates the talk.}
\GSfield{Present status}{As of November 2026, this remains a standard result under the stated hypotheses.}
\GSfield{Assessment}{literal statement: correct with missing hypotheses; publication match: exact; verification: read the result; development: none; confidence: high}
\end{GSitem}
\begin{GSitem}{1985-658.15}{F}{Faithful join representations and an unresolved subspace example}
\GSfield{Kind; source}{definition; \GSpage{P001230-P001231}}
\GSfield{The notes}{A faithful representation \(a:\Gamma\to\mathcal L(F^n)\) is required to preserve strict order and suprema. The following question-marked fragment mentions \(m=4\), \(E,F\), and \(\mathbb R^6\).}
\GSfield{Statement}{An injective join-semilattice map satisfies \(a(x\vee y)=a(x)+a(y)\) and preserves strict order. It need not preserve meets unless that is imposed separately. The fragment involving four-dimensional subspaces of \(\mathbb R^6\) does not determine a particular lattice or a theorem.}
\GSfield{Fate}{\citet[section~1]{gelfandponomarev1974modular} is a genuine pre-1985 source on representations of modular lattices, but the stated join-only definition and the incomplete dimensional fragment do not identify one of its theorems.}
\GSfield{Present status}{As of November 2026, the definition is interpretable; the intended four-dimensional example and claimed representability remain unidentified.}
\GSfield{Assessment}{literal statement: not determinable; publication match: thematic only; verification: publisher metadata checked; development: none; confidence: low}
\end{GSitem}
\begin{GSitem}{1985-658.16}{E}{Representability and systems of algebraic inequalities}
\GSfield{Kind; source}{claim; \GSpage{P001231}}
\GSfield{The notes}{The notes identify the realization task with solving a system of algebraic inequalities.}
\GSfield{Statement}{For a fixed finite real vector-configuration type, rank constraints can be written as polynomial equalities and strict nonvanishing conditions on minors. Conversely, real configuration realization problems are universal for suitable primary basic semialgebraic sets up to stable equivalence. An arbitrary faithful join embedding of a lattice is a different representation notion and is not covered without further specifications.}
\GSfield{Reconstruction}{The forward encoding and the converse universality statement are separated. Strict inequalities, equations and the equivalence relation cannot all be suppressed.}
\GSfield{Fate}{The universality direction is already present in \citet[Theorem~3]{mnev1985varieties}; \citet[Theorem~1]{adiprasitopadrol2017universality} gives the modern precise formulation.}
\GSfield{Present status}{As of November 2026, this remains a standard result under the stated hypotheses.}
\GSfield{Assessment}{literal statement: correct with missing hypotheses; publication match: partial; verification: read the result; development: none; confidence: high; problem \hyperref[prob:GP-1985-40]{GP-1985-40}}
\end{GSitem}
\begin{GSitem}{1985-658.17}{E}{Gelfand on lattices and their representation theory}
\GSfield{Kind; source}{remark; \GSpage{P001231}; I. M. Gelfand}
\GSfield{The notes}{Gelfand names Dedekind, Kurosh, Ore, von Neumann, and work with Ponomarev in discussing structures/lattices.}
\GSfield{Statement}{The intervention situates the representation of lattices by subspaces within an older algebraic programme. It contains no further explicitly stated classification theorem.}
\GSfield{Fate}{The concrete Gelfand--Ponomarev connection is \citet{gelfandponomarev1974modular}, published in 1974. The list of names alone does not warrant importing additional results.}
\GSfield{Present status}{As of November 2026, this is a historical attribution, not an independent open problem.}
\GSfield{Assessment}{literal statement: correct; publication match: thematic only; verification: publisher metadata checked; development: none; confidence: high}
\end{GSitem}
\begin{GSitem}{1985-658.18}{E}{Every finite lattice embeds in a finite partition lattice}
\GSfield{Kind; source}{theorem; \GSpage{P001231}}
\GSfield{The notes}{A theorem attributed to Pudl\'ak and T\r{u}ma embeds every finite lattice in a lattice of equivalence relations; infimum is intersection and supremum is described using union.}
\GSfield{Statement}{For every finite lattice \(L\), there exists a finite set \(X\) and an injective lattice homomorphism \(L\hookrightarrow\operatorname{Eq}(X)\). Meets in \(\operatorname{Eq}(X)\) are intersections; joins are the equivalence relations generated by unions, not the raw unions.}
\GSfield{Reconstruction}{The lattice operations are made explicit. This embedding theorem is not the finite algebra congruence-lattice representation problem.}
\GSfield{Fate}{The preliminary announcement is \citet{pudlaktuma1977every}. \citet[Conclusions]{pudlaktuma1980finite} gives a bounded embedding into \(\operatorname{Eq}(n)\) for sufficiently large \(n\); the full article appeared in December 1980 after receipt in 1977.}
\GSfield{Present status}{As of November 2026, this remains a standard result under the stated hypotheses.}
\GSfield{Assessment}{literal statement: correct with missing hypotheses; publication match: exact; verification: read the result; development: none; confidence: high}
\end{GSitem}

\GStalk{659}{1985-11-11}{B. L. Feigin}{P001231-P001235}{Dual amplitudes, free bosons and vertex operators}
The first string lecture introduces elastic-scattering invariants and the Veneziano amplitude, then reconstructs amplitudes through oscillator Fock spaces, Virasoro modes, normal-ordered exponential vertices and contraction identities. The notebook asserts uniqueness under conditions that do not even fix an overall normalization. Its five-point evaluation remark lacks the exponents and a specified meaning of evaluation, so no open problem is manufactured. These are mainly established constructions; zero modes, metric signs and plane versus cylinder conventions account for several essential corrections. The Feigin heading occurs midway down P001231.

\begin{GSitem}{1985-659.1}{E}{Elastic scattering and Mandelstam invariants}
\GSfield{Kind; source}{definition; \GSpage{P001231}}
\GSfield{The notes}{For equal masses, \(p_1+q_1=p_2+q_2\), all four squares equal \(m^2\), and \(s=(p_1+q_1)^2\), \(t=(p_1-p_2)^2\), \(u=(p_1-q_2)^2\), with \(s+t+u=4m^2\). The notes also say total momentum is zero.}
\GSfield{Statement}{For four on-shell Minkowski momenta with the stated conservation law, the displayed invariants obey \(s+t+u=4m^2\) by bilinear expansion. A zero total spatial momentum may specify the centre-of-mass frame; a zero sum of all incoming four-momenta requires reversing the outgoing momenta.}
\GSfield{Reconstruction}{The frame convention is supplied to avoid interpreting positive-energy incoming four-momenta as having zero total energy. \ours\ Bilinear expansion, using \(p_2+q_2=p_1+q_1\) and the four mass-shell equations, verifies \(s+t+u=4m^2\).}
\GSfield{Fate}{This is standard kinematic input to the pre-1985 dual-resonance construction, as in \citet{veneziano1968construction}. The invariant identity is verified directly from the notebook equations.}
\GSfield{Present status}{As of November 2026, this remains a standard result under the stated hypotheses.}
\GSfield{Assessment}{literal statement: correct with missing hypotheses; publication match: exact; verification: publisher metadata checked; development: annotated; confidence: high}
\end{GSitem}
\begin{GSitem}{1985-659.2}{E}{The listed duality conditions}
\GSfield{Kind; source}{definition; \GSpage{P001231-P001232}}
\GSfield{The notes}{The amplitude is analytic, symmetric in \(s,t\), has only simple \(s\)-poles at positions independent of \(t\), and has polynomial residues \(R_n(t)\) in \(\sum_{n\geq0}R_n(t)/(n-\alpha(s))\).}
\GSfield{Statement}{The intended class consists of meromorphic crossing-symmetric functions whose pole divisors are \(\alpha(s)=n\) and \(\alpha(t)=n\), with polynomial residues in the crossed variable when \(\alpha\) is affine. A partial-fraction series additionally needs a convergence prescription or subtraction terms.}
\GSfield{Reconstruction}{``Analytic'' means holomorphic off the poles, not entire. Growth and normalization conditions are absent.}
\GSfield{Fate}{The conditions belong to the historical dual-amplitude programme of \citet{veneziano1968construction}. Only publisher metadata was accessible for that original paper; the notebook conditions themselves are analyzed here without attributing a completeness theorem to it.}
\GSfield{Present status}{As of November 2026, this remains a standard result under the stated hypotheses.}
\GSfield{Assessment}{literal statement: correct with missing hypotheses; publication match: exact; verification: publisher metadata checked; development: none; confidence: high}
\end{GSitem}
\begin{GSitem}{1985-659.3}{E}{The Veneziano beta amplitude and the unsupported uniqueness assertion}
\GSfield{Kind; source}{claim; \GSpage{P001232}}
\GSfield{The notes}{The five conditions are said to determine uniquely \[A(s,t)=B(-\alpha(s),-\alpha(t))\].}
\GSfield{Statement}{The beta function gives one meromorphic crossing-symmetric solution. For affine \(\alpha\), its residue at \(\alpha(s)=n\) is polynomial in \(\alpha(t)\). The listed properties alone do not fix even its overall nonzero scalar factor; unless a fully convergent pole expansion without an entire part is imposed, adding a crossing-symmetric polynomial also preserves the listed local pole properties.}
\GSfield{Reconstruction}{\ours\ If \(A\) satisfies the listed homogeneous conditions, so does \(cA\) for every nonzero constant \(c\); this supplies the stated countercheck of uniqueness without normalization.}
\GSfield{Fate}{The amplitude was published by \citet{veneziano1968construction} in September 1968, after receipt on 29 July. No uniqueness theorem under exactly the notebook hypotheses was verified; elementary counterexamples show why extra assumptions are necessary.}
\GSfield{Present status}{As of November 2026, the beta-function construction is classical; uniqueness under only the recorded conditions is false.}
\GSfield{Assessment}{literal statement: false as written; publication match: partial; verification: publisher metadata checked; development: annotated; confidence: high}
\end{GSitem}
\begin{GSitem}{1985-659.4}{E}{Linear Regge trajectories as physical input}
\GSfield{Kind; source}{remark; \GSpage{P001232}}
\GSfield{The notes}{The linearity of \(\alpha(s)\) is attributed to analysis of experimental data.}
\GSfield{Statement}{The model takes \(\alpha(s)=\alpha_0+\alpha's\), with intercept and slope supplied as parameters. This is a phenomenological assumption in the historical amplitude, not a mathematical consequence of the five conditions.}
\GSfield{Fate}{The original title of \citet{veneziano1968construction} explicitly specifies linearly rising trajectories. The notebook supplies no experiment or numerical fit to audit.}
\GSfield{Present status}{As of November 2026, this item is a historical model assumption; no claim of exact linearity for all observed hadrons is inferred.}
\GSfield{Assessment}{literal statement: correct with missing hypotheses; publication match: partial; verification: publisher metadata checked; development: none; confidence: high}
\end{GSitem}
\begin{GSitem}{1985-659.5}{E}{Euler integral for the four-point amplitude}
\GSfield{Kind; source}{construction; \GSpage{P001232}}
\GSfield{The notes}{The beta amplitude is rewritten as a one-variable integral with exponents involving pairwise momentum products; the accompanying two-in/two-out diagram is followed by \(\sum_i p_i=0\).}
\GSfield{Statement}{For \(\operatorname{Re}\alpha(s)<0\) and \(\operatorname{Re}\alpha(t)<0\), \[B(-\alpha(s),-\alpha(t))=\int_0^1 x^{-\alpha(s)-1}(1-x)^{-\alpha(t)-1}\,dx.\] Other values are defined by meromorphic continuation. Expressing the exponents through momenta additionally fixes the metric, mass shell, slope and all-incoming convention.}
\GSfield{Reconstruction}{The unambiguous beta-integral formula replaces an identification of physical parameters that the notes do not fully normalize. Outgoing arrows alone do not implement all-incoming signed momenta.}
\GSfield{Fate}{This is Euler's beta integral in the amplitude of \citet{veneziano1968construction}; no post-1985 mathematical development is needed.}
\GSfield{Present status}{As of November 2026, this remains a standard result under the stated hypotheses.}
\GSfield{Assessment}{literal statement: correct with missing hypotheses; publication match: exact; verification: publisher metadata checked; development: none; confidence: high}
\end{GSitem}
\begin{GSitem}{1985-659.6}{F}{Five-point Euler integrals and the meaning of closed evaluation}
\GSfield{Kind; source}{claim; \GSpage{P001232}}
\GSfield{The notes}{A double integral with factors \(x,y,1-x,1-y,1-xy\) and omitted exponents is said not to have been calculated in general.}
\GSfield{Statement}{For parameters \(a,b,c,d>0\) and real \(e<b+d\), the model integral \[\int_0^1\!\int_0^1x^{a-1}(1-x)^{b-1}y^{c-1}(1-y)^{d-1}(1-xy)^{-e}\,dx\,dy\] equals \[B(a,b)B(c,d)\,{}_3F_2\!\left(\begin{matrix}a,c,e\\a+b,c+d\end{matrix};1\right).\] This follows by the binomial expansion and beta integration, with continuation where justified. A special-function representation is different from a gamma-product evaluation.}
\GSfield{Reconstruction}{The five-parameter integral and its evaluation are an editorial calculation, not a formula supplied by the notebook. No exact physical five-point amplitude or gamma-product conjecture is claimed. \ours\ Expanding \((1-xy)^{-e}\) and integrating each power gives the displayed beta factors and the \({}_3F_2\) series. \oursComputed\ The terminating specialization \(a=b=c=d=1,e=-2\) gives \(11/18\) independently by polynomial integration and by the hypergeometric expression; code and output are supplied. This checks an editorial model, not missing notebook exponents.}
\GSfield{Fate}{The notebook omits the exponents and the intended class of answers. The displayed elementary reconstruction establishes that ``not calculated'' cannot be promoted to a precise general open problem. The four-point comparator is \citet{veneziano1968construction}.}
\GSfield{Present status}{As of November 2026, the notebook's intended evaluation question is not identifiable.}
\GSfield{Assessment}{literal statement: not determinable; publication match: none found; verification: unresolved lead; development: annotated; confidence: medium}
\end{GSitem}
\begin{GSitem}{1985-659.7}{E}{The Heisenberg algebra and infinitely many string modes}
\GSfield{Kind; source}{construction; \GSpage{P001232-P001233}}
\GSfield{The notes}{Generators \(a_m^\mu,x^\mu\) satisfy \([a_m^\mu,a_n^\nu]=mg^{\mu\nu}\delta_{m,-n}\), \(a_0=p\), \([x^\mu,p^\nu]=ig^{\mu\nu}\), and \([x,a_m]=0\) for \(m\ne0\). The text first says finitely many oscillators and immediately says finitely many do not suffice.}
\GSfield{Statement}{Let \(g\) be a nondegenerate form on a \(d\)-dimensional space, indexed by \(0\leq\mu<d\). Use one oscillator of each positive integer frequency in each component, together with zero modes. The algebraic Fock space consists of finite products of negative-frequency generators acting on a momentum vacuum; infinitely many generators exist although each vector uses finitely many.}
\GSfield{Reconstruction}{The index endpoint is corrected from \(d\) to \(d-1\). Finitely many species and countably many frequencies resolve the contradictory prose.}
\GSfield{Fate}{The same Heisenberg representation is written in \citet[introduction, equations~(1)--(2)]{frenkel1985representations}, tracing the construction to the early dual-resonance literature.}
\GSfield{Present status}{As of November 2026, this remains a standard result under the stated hypotheses.}
\GSfield{Assessment}{literal statement: false as written; publication match: partial; verification: read the result; development: none; confidence: high}
\end{GSitem}
\begin{GSitem}{1985-659.8}{E}{Chiral coordinate fields and their distributional commutators}
\GSfield{Kind; source}{construction; \GSpage{P001233}}
\GSfield{The notes}{The fields are \(P(\theta)=\sum_m a_me^{-im\theta}\) and \(Q(\theta)=x+p\theta+i\sum_{n\ne0}a_n e^{-in\theta}/n\). Their commutators involve \(\delta\), \(\delta'\), and a function called Heaviside.}
\GSfield{Statement}{With the corrected common index \(n\) in \(Q\), \(\partial_\theta Q=P\). In the distributional sense, \([Q^\mu(\theta),P^\nu(\theta')]=2\pi i g^{\mu\nu}\delta_{2\pi}(\theta-\theta')\) and \([P^\mu(\theta),P^\nu(\theta')]=2\pi i g^{\mu\nu}\delta'_{2\pi}(\theta-\theta')\). On a lifted interval with \(0<|\theta-\theta'|<2\pi\), \([Q^\mu(\theta),Q^\nu(\theta')]=-i\pi g^{\mu\nu}\operatorname{sgn}(\theta-\theta')\), extended by the zero-mode monodromy.}
\GSfield{Reconstruction}{The plain Heaviside function cannot be the antisymmetric commutator kernel. The local sign function, a lift of the angular coordinate, and distributional interpretation are supplied. \ours\ On \(0<|u|<2\pi\), the Fourier identity \(\sum_{n\ne0}e^{-inu}/n=-i(\pi\operatorname{sgn}u-u)\), together with the zero-mode commutator \(-iu\), yields the stated odd kernel. Differentiating the distributional series gives the two current commutators.}
\GSfield{Fate}{These identities follow directly from the displayed oscillator brackets; the field formalism is the one used in \citet[introduction, equations~(1)--(4)]{frenkel1985representations}.}
\GSfield{Present status}{As of November 2026, this remains a standard result under the stated hypotheses.}
\GSfield{Assessment}{literal statement: false as written; publication match: partial; verification: read the result; development: annotated; confidence: high}
\end{GSitem}
\begin{GSitem}{1985-659.9}{E}{Sugawara operators for free bosons}
\GSfield{Kind; source}{theorem; \GSpage{P001233-P001234}}
\GSfield{The notes}{\(L_m=(4\pi)^{-1}\int e^{im\theta}:P(\theta)^2:\,d\theta\), equivalently \(2L_m=\sum_{r+s=m}:a_ra_s:\), is asserted to satisfy Virasoro with central charge \(d\); this is an exercise. The formula for \(L_0\) is given.}
\GSfield{Statement}{On algebraic Fock space, \[L_m=\tfrac12\sum_{r\in\mathbb Z}g_{\mu\nu}:a_r^\mu a_{m-r}^\nu:,\qquad L_0=\tfrac12(p,p)+\sum_{n>0}(a_{-n},a_n),\] and \[[L_m,L_n]=(m-n)L_{m+n}+\tfrac d{12}(m^3-m)\delta_{m,-n}I.\] Normal ordering puts annihilation operators on the right; each mode acts by a finite sum on a fixed finite-energy vector.}
\GSfield{Fate}{The construction and the double-contraction calculation are given in \citet[section~2 and Theorem~3.3]{frenkel1985representations}. The signature of \(g\) does not change the central charge \(d\).}
\GSfield{Present status}{As of November 2026, this remains a standard result under the stated hypotheses.}
\GSfield{Assessment}{literal statement: correct with missing hypotheses; publication match: exact; verification: read the result; development: none; confidence: high}
\end{GSitem}
\begin{GSitem}{1985-659.10}{E}{Normal-ordered exponential vertex operators}
\GSfield{Kind; source}{construction; \GSpage{P001234}}
\GSfield{The notes}{\(V(k,\theta)=\exp(i\langle k,Q(\theta)\rangle)\) is factorized at \(\theta=0\) into creation and annihilation exponentials with a momentum-shift factor \(e^{ikx}\).}
\GSfield{Statement}{The exponential must be normal ordered. At \(z=1\), one may write \[Y_k(1)=\exp\!\left(\sum_{n>0}(k,a_{-n})/n\right)D_k\exp\!\left(-\sum_{n>0}(k,a_n)/n\right).\] The complete field is \(Y_k(z)=E_-(k,z)D_k z^{(k,a_0)}E_+(k,z)\), with an additional lattice cocycle when mutual locality is required.}
\GSfield{Reconstruction}{Normal ordering and the zero-mode power for variable \(z\) are supplied; an unregularized exponential of the field has divergent self-contraction.}
\GSfield{Fate}{Compare \citet[section~2.1, equation~(2.3)]{frenkel1985representations} and the original affine construction of \citet{frenkelkac1980basic} (published in 1980).}
\GSfield{Present status}{As of November 2026, this remains a standard result under the stated hypotheses.}
\GSfield{Assessment}{literal statement: correct with missing hypotheses; publication match: exact; verification: read the result; development: none; confidence: high}
\end{GSitem}
\begin{GSitem}{1985-659.11}{E}{Vertex covariance and conformal weight}
\GSfield{Kind; source}{theorem; \GSpage{P001234}}
\GSfield{The notes}{The notes give \(V(k,\theta)=e^{i\theta L_0}V(k)e^{-i\theta L_0}\) and \([L_n-L_0,V(k)]=n(k,k)V(k)/2\).}
\GSfield{Statement}{For the plane primary field \(Y_k(z)\) of weight \(h=(k,k)/2\), \[[L_n,Y_k(z)]=\bigl(z^{n+1}\partial_z+(n+1)hz^n\bigr)Y_k(z).\] Evaluation at \(z=1\) gives the stated difference commutator. Angular translation and the cylinder field require a consistent choice of \(z=e^{\pm i\theta}\) and the factor \(z^h\).}
\GSfield{Reconstruction}{The field of conformal weight \(h\), its plane/cylinder convention, and the zero-mode factor are essential. The note's two angular signs should not be combined without this convention choice.}
\GSfield{Fate}{This is the standard free-boson covariance relation; \citet[section~2, vertex and Virasoro relations]{frenkel1985representations} provides the pre-1985 operator construction.}
\GSfield{Present status}{As of November 2026, this remains a standard result under the stated hypotheses.}
\GSfield{Assessment}{literal statement: correct with missing hypotheses; publication match: partial; verification: read the result; development: none; confidence: high}
\end{GSitem}
\begin{GSitem}{1985-659.12}{E}{Factorizing amplitudes through vertices and a propagator}
\GSfield{Kind; source}{method; \GSpage{P001232,P001234}}
\GSfield{The notes}{The proposed vacuum expectation becomes a chain \[\langle0|V(p_2)D(s_2)V(p_3)\cdots D(s_{n-2})V(p_{n-1})|0\rangle\], with \(D(s)=\int_0^1x^{L_0-\alpha(s)-1}(1-x)^{\alpha(0)-1}\,dx\).}
\GSfield{Statement}{On an eigenvector of the operator denoted by \(L_0\) with eigenvalue \(\lambda\), the displayed propagator is \(B(\lambda-\alpha(s),\alpha(0))\), initially for positive real parts of both arguments, and then meromorphically continued. In a full string amplitude one must specify whether this operator is the oscillator number operator or the full zero-mode-inclusive \(L_0\), the incoming/outgoing vacuum momenta, and vertex normalizations.}
\GSfield{Reconstruction}{The spectral meaning of the integral is made explicit. The notebook reuses \(L_0\) in a way that may double-count the intermediate momentum if read literally. \ours\ The spectral formula follows by applying the displayed operator integral to an \(L_0\)-eigenvector and then recognizing Euler's beta integral.}
\GSfield{Fate}{The method belongs to the established dual-resonance operator formalism discussed by \citet[introduction]{frenkel1985representations}. No exact equality of the notebook's incompletely normalized chain with every physical amplitude is asserted.}
\GSfield{Present status}{As of November 2026, this remains a standard result under the stated hypotheses.}
\GSfield{Assessment}{literal statement: correct with missing hypotheses; publication match: partial; verification: read the result; development: annotated; confidence: high}
\end{GSitem}
\begin{GSitem}{1985-659.13}{E}{Momentum Fock spaces and mode shifts}
\GSfield{Kind; source}{construction; \GSpage{P001234-P001235}}
\GSfield{The notes}{Vacua satisfy \(a_n|0,p\rangle=0\) for \(n>0\) and \(a_0|0,p\rangle=p|0,p\rangle\). Vertex modes map \(H(p)\to H(p+k)\) and satisfy \([a_m^\mu,V_n(k)]=k_\mu V_{n+m}(k)\).}
\GSfield{Statement}{Let \(H(p)=S(\bigoplus_{n>0}\mathbb C^d t^{-n})|p\rangle\). Then \([a_0^\mu,Y_k(z)]=k^\mu Y_k(z)\) and \([a_m^\mu,Y_k(z)]=k^\mu z^mY_k(z)\). With the convention \(Y_k(z)=\sum_nY_{k,n}z^{-n}\), the mode relation is \([a_m^\mu,Y_{k,n}]=k^\mu Y_{k,n+m}\). Integral mode labels require the appropriate integral momentum pairings; otherwise shifted powers occur.}
\GSfield{Reconstruction}{An algebraic sum over chosen momenta is sufficient for formal fields; a Hilbert-space direct integral over continuous momenta is a different completion.}
\GSfield{Fate}{\citet[section~2.1, equations~(2.2)--(2.4)]{frenkel1985representations} gives this construction. The upper/lower target index must be raised or lowered consistently with \(g\).}
\GSfield{Present status}{As of November 2026, this remains a standard result under the stated hypotheses.}
\GSfield{Assessment}{literal statement: correct with missing hypotheses; publication match: exact; verification: read the result; development: none; confidence: high}
\end{GSitem}
\begin{GSitem}{1985-659.14}{E}{The two-vertex contraction and binomial generating series}
\GSfield{Kind; source}{theorem; \GSpage{P001235}}
\GSfield{The notes}{An exercise asserts \(\sum_n A_nz^n=(1-z)^{(k_1,k_2)}\) for two vertex modes. The final momentum in the bra is written \(p+k_1+k\), without the second subscript.}
\GSfield{Statement}{For the oscillator parts of two normal-ordered vertices, in the radial region \(|z|<1\), contraction gives \[\exp\!\left(-(k_1,k_2)\sum_{n>0}\frac{z^n}{n}\right)=(1-z)^{(k_1,k_2)}.\] Thus the normalized oscillator coefficients are \(A_n=(-1)^n\binom{(k_1,k_2)}n\). The full matrix element also contains the zero-mode monomial and has final momentum \(p+k_1+k_2\).}
\GSfield{Reconstruction}{``Normalized oscillator coefficients'' separates the stated contraction from omitted zero-mode powers and the possible shift of Laurent indices. \ours\ The oscillator commutator of the two exponential exponents is the scalar \(-(k_1,k_2)\sum_{n>0}z^n/n\); exponentiation and \(-\log(1-z)=\sum_{n>0}z^n/n\) give the generating series.}
\GSfield{Fate}{The calculation is the two-field instance of the vertex product formula in \citet[equation~(2.5)]{frenkel1985representations}.}
\GSfield{Present status}{As of November 2026, this remains a standard result under the stated hypotheses.}
\GSfield{Assessment}{literal statement: correct with missing hypotheses; publication match: exact; verification: read the result; development: annotated; confidence: high}
\end{GSitem}

\GStalk{660}{1985-11-18}{B. L. Feigin}{P001236-P001238}{Physical quotients, the no-ghost theorem and lattice vertices}
Feigin defines the constrained physical space and spurious quotient, states the critical no-ghost result, and connects screening residues with Kac--Moody algebras and transverse DDF oscillators. Gelfand identifies the Monster/string connection. The strongest needed qualifications are the nonzero-momentum exception, the quotient by null states, the distinction between ghosts and tachyons, and the momentum normalization in DDF brackets. The last part on P001238 ends immediately before Zelevinsky begins; it does not include that page in its entirety.

\begin{GSitem}{1985-660.1}{E}{Virasoro weights and the physical constraint space}
\GSfield{Kind; source}{definition; \GSpage{P001236}}
\GSfield{The notes}{The continuation introduces weight vectors, highest weight, and \(\widetilde{\mathcal H}=\{v:L_0v=v,\ L_iv=0\ (i>0)\}\).}
\GSfield{Statement}{For a Virasoro module with diagonalizable \(L_0\), the relation \([L_0,L_n]=-nL_n\) makes \(L_n\) lower the \(L_0\)-eigenvalue by \(n\). The physical constraint space at intercept one is \(P^1=\ker(L_0-1)\cap\bigcap_{n>0}\ker L_n\). It is a subspace of highest-weight vectors, not a Virasoro submodule in general.}
\GSfield{Fate}{This is the standard constraint space in \citet[equations~(8) and (4.1)]{frenkel1985representations}.}
\GSfield{Present status}{As of November 2026, this remains a standard result under the stated hypotheses.}
\GSfield{Assessment}{literal statement: correct with missing hypotheses; publication match: exact; verification: read the result; development: none; confidence: high}
\end{GSitem}
\begin{GSitem}{1985-660.2}{E}{Spurious states and the physical quotient}
\GSfield{Kind; source}{definition; \GSpage{P001236}}
\GSfield{The notes}{The set of spurions is written as those physical \(v\) expressible as one \(L_iv'\), \(i<0\), and the physical space is its quotient.}
\GSfield{Statement}{In an algebraic free-boson Fock module with its invariant Hermitian form, put \(S=\sum_{n>0}L_{-n}H\) and \(N=P^1\cap S\). The linear quotient \(P^1/N\) is well defined. In each finite-dimensional energy space, adjointness \(L_n^*=L_{-n}\) identifies \(N\) with the radical of the restricted form on \(P^1\). A union of individual images need not be a subspace.}
\GSfield{Reconstruction}{The sum of images and the invariant form are supplied. The equivalence with the radical uses finite-dimensionality of each fixed-momentum energy level and its nondegenerate ambient form. \ours\ At a fixed finite-dimensional energy level, adjointness gives \(P^1=S^\perp\); nondegeneracy of the ambient form gives \((P^1)^\perp=S\). Thus its radical is \(P^1\cap S\), verifying the corrected quotient algebraically.}
\GSfield{Fate}{The transverse/spurious decomposition occurs in \citet[section~4, Lemmas~4.2--4.3]{frenkel1985representations}; \citet[section~5, Theorem~5.1]{borcherds1992monstrous} explicitly uses the quotient by the null space in a later application.}
\GSfield{Present status}{As of November 2026, this remains a standard result under the stated hypotheses.}
\GSfield{Assessment}{literal statement: false as written; publication match: exact; verification: read the result; development: annotated; confidence: high}
\end{GSitem}
\begin{GSitem}{1985-660.3}{E}{The weight of a preimage under L minus one}
\GSfield{Kind; source}{claim; \GSpage{P001236}}
\GSfield{The notes}{Equation (5) says that if \(L_{-1}v'=v\) has weight one then \(v'\) has weight two.}
\GSfield{Statement}{If \(v'\) is a weight vector, \(L_{-1}v'\ne0\), and \(L_0L_{-1}v'=L_{-1}v'\), then \(L_0v'=0\). Indeed, \([L_0,L_{-1}]=L_{-1}\), so \(L_{-1}\) raises the weight by one.}
\GSfield{Reconstruction}{\ours\ For a weight-\(h\) vector, the bracket gives \(L_0L_{-1}v'=(h+1)L_{-1}v'\). A nonzero weight-one image forces \(h=0\).}
\GSfield{Fate}{The correction is forced by the Virasoro bracket, for example \citet[equation~(2.9)]{frenkel1985representations}. The scan clearly contains both the word ``two'' and \(2v'\); this is a faithful mathematical error, not a translation correction.}
\GSfield{Present status}{As of November 2026, this remains a standard result under the stated hypotheses.}
\GSfield{Assessment}{literal statement: false as written; publication match: exact; verification: read the result; development: annotated; confidence: high}
\end{GSitem}
\begin{GSitem}{1985-660.4}{E}{Half constraints followed by a quotient}
\GSfield{Kind; source}{method; \GSpage{P001236}}
\GSfield{The notes}{Genuine invariants under diffeomorphisms are obstructed by the nontrivial centre. The proposed procedure takes invariants for one half and quotients by the action of the other half.}
\GSfield{Statement}{The construction \(P^1/(P^1\cap\sum_{n>0}L_{-n}H)\) imposes the positive-mode constraints and removes the spurious descendants. It is not the subspace annihilated by every Virasoro operator: with nonzero central charge, such an ordinary invariant subspace is zero.}
\GSfield{Reconstruction}{The quotation marks around ``invariants'' are retained in substance; no unstated equivalence with an arbitrary cohomology theory is attributed to the speaker.}
\GSfield{Fate}{The physical quotient is the established pre-1985 construction summarized in \citet[section~4]{frenkel1985representations}. A full cohomological or BRST reformulation is not specified in this notebook passage.}
\GSfield{Present status}{As of November 2026, this remains a standard result under the stated hypotheses.}
\GSfield{Assessment}{literal statement: correct with missing hypotheses; publication match: partial; verification: read the result; development: none; confidence: high}
\end{GSitem}
\begin{GSitem}{1985-660.5}{E}{The critical no-ghost theorem}
\GSfield{Kind; source}{theorem; \GSpage{P001236}}
\GSfield{The notes}{The two goals are a description by \(d-2\) oscillators and strict positivity on the physical quotient; both are said to hold for \(d=26\). The parenthesis calls this absence of tachyons.}
\GSfield{Statement}{For the bosonic Fock space with one negative and twenty-five positive directions, intercept one, and nonzero real momentum \(p\), the form on \(P^1(p)\) is positive semidefinite. Its quotient by the radical is positive definite and has the transverse oscillator multiplicities of twenty-four bosons at level \(N=1-(p,p)/2\), when \(N\) is a nonnegative integer. Strict positivity applies to nonzero quotient vectors. Zero momentum must be excluded: its level-one physical space contains the full indefinite target space.}
\GSfield{Reconstruction}{The momentum hypothesis, quotient by null vectors and level condition are supplied. Positivity removes negative-norm ghosts; it does not remove the bosonic ground-state tachyon. \ours\ The zero-momentum exception can be checked directly: every level-one vector \(a_{-1}(v)|0\rangle\) obeys all positive-mode constraints and inherits the target-space form, which is indefinite. This is a check of a known exception, not a proof of no-ghost.}
\GSfield{Fate}{No-ghost was established in the early 1970s. The exact zero-momentum exception and the critical transverse proof can be checked in \citet[Theorem~4.1 and Lemmas~4.2--4.3]{frenkel1985representations}; a later general form is \citet[Theorem~5.1]{borcherds1992monstrous}.}
\GSfield{Present status}{As of November 2026, the theorem is classical. The notebook assertion becomes correct with these hypotheses; its identification with absence of tachyons is false.}
\GSfield{Assessment}{literal statement: false as written; publication match: exact; verification: read the result; development: annotated; confidence: high}
\end{GSitem}
\begin{GSitem}{1985-660.6}{E}{The creation coefficient and zero modes in the continued vertex formula}
\GSfield{Kind; source}{construction; \GSpage{P001237}}
\GSfield{The notes}{Equation (6) writes the creation exponential without \(1/n\), followed by \(D_p\) and the correctly divided annihilation exponential. Equation (7) defines the momentum shift.}
\GSfield{Statement}{The oscillator factors must both have frequency denominators: \(E_-(p,z)=\exp(\sum_{n>0}(p,a_{-n})z^n/n)\) and \(E_+(p,z)=\exp(-\sum_{n>0}(p,a_n)z^{-n}/n)\). The plane field is \(Y_p(z)=E_-(p,z)D_pz^{(p,a_0)}E_+(p,z)\). Its cylinder-normalized counterpart is \(X_p(z)=z^{(p,p)/2}Y_p(z)\), up to consistently ordered zero-mode factors.}
\GSfield{Reconstruction}{The scan confirms the missing creation denominator. Plane and cylinder fields are distinguished so that the subsequent \(dz/z\) residue is meaningful. \ours\ In a single mode, represent \(a_{-n}\) by multiplication by \(u\) and \(a_n\) by \(n\partial_u\). The creation exponential then has the required commutator precisely with coefficient \(1/n\). \oursComputed\ This differential-operator identity was checked for frequencies one through six on \(u^3+u+1\); omitting the denominator multiplies the result by \(n\).}
\GSfield{Fate}{Compare \citet[equations~(2.2)--(2.3)]{frenkel1985representations}. In that source the exponential of the sum of zero-mode operators incorporates a scalar power through their commutator; replacing it by a bare translation loses this factor.}
\GSfield{Present status}{As of November 2026, this remains a standard result under the stated hypotheses.}
\GSfield{Assessment}{literal statement: false as written; publication match: exact; verification: read the result; development: annotated; confidence: high}
\end{GSitem}
\begin{GSitem}{1985-660.7}{E}{Multipoint oscillator contractions}
\GSfield{Kind; source}{theorem; \GSpage{P001237}}
\GSfield{The notes}{Equation (8) gives a vacuum matrix element as \(\prod_{j>1}(1-z_j/z_i)^{(p_i,p_j)}\).}
\GSfield{Statement}{For corrected, oscillator-stripped vertices in the radial region \(|z_1|>\cdots>|z_r|\), their normalized contraction is \[\prod_{1\leq i<j\leq r}(1-z_j/z_i)^{(p_i,p_j)}.\] A full matrix element of \(Y_{p_i}\), between initial momentum \(k\) and final momentum \(k+\sum_i p_i\), also contains the zero-mode monomial \(\prod_i z_i^{(p_i,k+\sum_{j>i}p_j)}\), as well as any chosen cocycle factors.}
\GSfield{Reconstruction}{The product index is changed from \(j>1\) to \(i<j\); the extra powers assume the displayed order of plane-field zero modes. The oscillator identity is not the full correlation function without this qualification. \ours\ Acting from right to left, the momentum immediately before \(Y_{p_i}\) is \(k+\sum_{j>i}p_j\), which gives the displayed zero-mode monomial. This verifies that the stripped contraction and full matrix element differ.}
\GSfield{Fate}{The repeated two-field contraction follows from \citet[equation~(2.5)]{frenkel1985representations}.}
\GSfield{Present status}{As of November 2026, this remains a standard result under the stated hypotheses.}
\GSfield{Assessment}{literal statement: false as written; publication match: exact; verification: read the result; development: annotated; confidence: high}
\end{GSitem}
\begin{GSitem}{1985-660.8}{E}{Norm-two screening operators commute with Virasoro}
\GSfield{Kind; source}{theorem; \GSpage{P001237}}
\GSfield{The notes}{The integrals \(\int X(p,z)\,dz/z\) are to commute with Virasoro when \((p,p)=2\).}
\GSfield{Statement}{If \(Y_p(z)\) is a primary field of weight one with integral Laurent powers on the chosen momentum sectors, then \(Q_p=\operatorname{Res}_{z=0}Y_p(z)\,dz\) commutes with every \(L_n\), because its commutator is the residue of \(\partial_z(z^{n+1}Y_p)\). For \(X_p=zY_p\), this is equivalently \(\operatorname{Res}X_p(z)\,dz/z\). Free-boson weight one means \((p,p)=2\).}
\GSfield{Reconstruction}{The notebook operator lacks the necessary zero-mode normalization; its formula cannot be used literally in this residue assertion. Residue notation replaces an unspecified contour and factor \(2\pi i\). \ours\ The residue of the derivative of any integral Laurent series is zero, proving the displayed commutation assertion from the primary-field transformation law.}
\GSfield{Fate}{The commuting zero modes are part of \citet[equations~(2.10)--(2.12)]{frenkel1985representations}.}
\GSfield{Present status}{As of November 2026, this remains a standard result under the stated hypotheses.}
\GSfield{Assessment}{literal statement: correct with missing hypotheses; publication match: exact; verification: read the result; development: annotated; confidence: high}
\end{GSitem}
\begin{GSitem}{1985-660.9}{E}{A symmetric Cartan matrix from Lorentzian vectors}
\GSfield{Kind; source}{construction; \GSpage{P001237}}
\GSfield{The notes}{Vectors \(p_i\in\mathbb R^{26}\) have square two and integral products \(A_{ij}\), followed by the unrestricted inequality \(A_{ij}<0\).}
\GSfield{Statement}{The intended symmetric generalized Cartan conditions are \(A_{ii}=2\), \(A_{ij}\in\mathbb Z_{\leq0}\) for \(i\ne j\), and \(A_{ij}=0\iff A_{ji}=0\). A realization as a Gram matrix in signature \((25,1)\) is an additional restriction: its positive and negative inertia indices are at most twenty-five and one, respectively, and its rank is at most twenty-six.}
\GSfield{Reconstruction}{The negative sign applies off the diagonal and allows zero. The inertia restrictions are elementary consequences of a Gram realization. \ours\ The Gram matrix is the pullback of the ambient form under the column map \(\mathbb R^n\to\mathbb R^{26}\); a positive or negative definite subspace maps injectively into one of the corresponding sign in the target. This proves the stated necessary inertia bounds.}
\GSfield{Fate}{The exact Cartan conditions appear in \citet[section~1.1]{frenkel1985representations}. No assertion that every symmetric Cartan matrix fits this fixed Lorentzian space is warranted.}
\GSfield{Present status}{As of November 2026, this remains a standard result under the stated hypotheses.}
\GSfield{Assessment}{literal statement: false as written; publication match: exact; verification: read the result; development: annotated; confidence: high}
\end{GSitem}
\begin{GSitem}{1985-660.10}{E}{Kac--Moody Chevalley operators from screening residues}
\GSfield{Kind; source}{theorem; \GSpage{P001237}}
\GSfield{The notes}{The norm-two vertex residues \(e_i,f_i\), together with \(h_i|k\rangle=(k,p_i)|k\rangle\), are asserted to represent the algebra with Cartan matrix \(A\). Macdonald at Bourbaki is recommended.}
\GSfield{Statement}{For a symmetric nonaffine Kac--Moody algebra and its standard lattice/Heisenberg vertex construction, a compatible bimultiplicative sign cocycle gives screening operators for the simple roots and their negatives satisfying the Chevalley--Serre relations. The sign of \(f_i\) is chosen to match either \([e_i,f_i]=h_i\) or the opposite Chevalley convention. In the affine case one uses the separate basic current-algebra construction, rather than assuming the same nondegenerate root-space model.}
\GSfield{Reconstruction}{The cocycle, zero modes and representation-space hypotheses are indispensable additions. The notebook outline is not an unrestricted construction from any list of Lorentzian vectors.}
\GSfield{Fate}{\citet[Theorem~2.1]{frenkel1985representations} proves the nonaffine statement and explicitly separates its affine Theorem~2.2. The suggested reading is identifiable as \citet{macdonald1981affine}, Bourbaki expos\'e~577, June 1981.}
\GSfield{Present status}{As of November 2026, this remains a standard result under the stated hypotheses.}
\GSfield{Assessment}{literal statement: correct with missing hypotheses; publication match: partial; verification: read the result; development: none; confidence: high}
\end{GSitem}
\begin{GSitem}{1985-660.11}{E}{Gelfand on the Monster, strings, and moonshine}
\GSfield{Kind; source}{remark; \GSpage{P001237}; I. M. Gelfand}
\GSfield{The notes}{Gelfand connects the group of approximately \(10^{54}\) elements with strings, says ``the largest finite simple group'', and recalls work of Frenkel, Lepowsky and an unnamed third author.}
\GSfield{Statement}{The intended group is the Monster, the largest sporadic finite simple group; finite simple groups as a whole have unbounded orders. The third author is Arne Meurman. The mathematical connection is the construction of a graded Monster representation with normalized modular \(J\)-character, and the later proof of the moonshine conjectures by vertex-algebra and no-ghost methods.}
\GSfield{Fate}{Frenkel--Lepowsky--Meurman announced their module in May 1984 \citep{frenkellepowskymeurman1984natural}, before this intervention. \citet[Theorem~1.1 and sections~5--6]{borcherds1992monstrous} subsequently proved the moonshine assertions. That later theorem is a development of the connection, not proof of a precise conjecture stated here.}
\GSfield{Present status}{As of November 2026, this is an identifiable pre-1985 mathematical connection with major later development. The notebook records no exact moonshine conjecture to classify separately.}
\GSfield{Assessment}{literal statement: false as written; publication match: partial; verification: read the result; development: none; confidence: high}
\end{GSitem}
\begin{GSitem}{1985-660.12}{E}{Null directions and transverse DDF operators}
\GSfield{Kind; source}{construction; \GSpage{P001238}}
\GSfield{The notes}{For a null vector \(p\), the vectors \(v^1,\ldots,v^{d-2},p\) form a basis of \(p^\perp\). The residues \(U_n(v)=\int v(z)X(np,z)\,dz/z\) are asserted to be oscillators.}
\GSfield{Statement}{Let \(c\ne0\) be null in a real space of signature \((d-1,1)\). The quotient \(c^\perp/\mathbb Rc\) is positive definite of dimension \(d-2\). Choose orthonormal transverse representatives. Correctly normal-ordered DDF residues then satisfy \([A_m^i,A_n^j]=m\delta_{ij}\delta_{m,-n}(c,a_0)\); on a momentum tower with \((c,a_0)=1\), this is the standard Heisenberg relation. Sign changes in the exponential and mode labels must be made together.}
\GSfield{Reconstruction}{Nonzero \(c\), orthonormality in the quotient and normalization of its pairing with the momentum are absent from the notes and supplied here. Mere linear independence gives the Gram matrix in the commutator, not \(\delta_{ij}\).}
\GSfield{Fate}{\citet[introduction, equations~(4)--(5), and section~4]{frenkel1985representations} gives the residues and their central factor. The original Del Giudice--Di Vecchia--Fubini construction dates to 1972, as identified in that bibliography.}
\GSfield{Present status}{As of November 2026, this remains a standard result under the stated hypotheses.}
\GSfield{Assessment}{literal statement: correct with missing hypotheses; publication match: exact; verification: read the result; development: none; confidence: high}
\end{GSitem}
\begin{GSitem}{1985-660.13}{E}{Completeness of transverse oscillators at the critical dimension}
\GSfield{Kind; source}{theorem; \GSpage{P001238}}
\GSfield{The notes}{After the Heisenberg assertion, the notes say that the operators generate the whole space.}
\GSfield{Statement}{At \(d=26\), fixed nonzero physical momentum and level, DDF transverse states provide representatives for every class in the physical quotient by null/spurious states. A chosen DDF vacuum generates its compatible momentum tower, not every momentum sector from one vector. At zero momentum the usual positive-definite completeness statement has an exceptional sector.}
\GSfield{Reconstruction}{The words ``whole space'' are restricted to the physical quotient and the momentum sectors on which the chosen DDF operators act.}
\GSfield{Fate}{\citet[section~4, Lemmas~4.2--4.3 and equations~(4.8)--(4.9)]{frenkel1985representations} supplies the critical-dimension transverse description.}
\GSfield{Present status}{As of November 2026, this remains a standard result under the stated hypotheses.}
\GSfield{Assessment}{literal statement: correct with missing hypotheses; publication match: exact; verification: read the result; development: none; confidence: high}
\end{GSitem}

\GStalk{661}{1985-11-18}{A. V. Zelevinsky}{P001238-P001240}{Schubert cells, divided differences and Schubert polynomials}
This is explicitly an exposition of Lascoux--Sch\"utzenberger, following the 1973 Bernstein--Gelfand--Gelfand paper. It compares the Borel presentation with the Schubert basis and constructs polynomial representatives by divided differences. The missing top-class factor \(1/n!\) and the divided-difference sign are genuine mathematical inconsistencies in the manuscript. Earlier related occurrences include 1971-239.1--2 (Schubert homology), 1979-483.1 (reduced words), 1979-484.3 (Schubert closures), and 1982-575.1 (Grassmannian cohomology); no earlier problem identifier matches a new open question here.

\begin{GSitem}{1985-661.1}{E}{Borel presentation of the cohomology of the full flag variety}
\GSfield{Kind; source}{theorem; \GSpage{P001238}}
\GSfield{The notes}{For complete flags in \(V^n\), the first description is \(H^*(X,\mathbb Q)=\mathbb Q[x_1,\ldots,x_n]/J\), where \(J\) is generated by symmetric polynomials with no constant term.}
\GSfield{Statement}{For \(X=\operatorname{Fl}(\mathbb C^n)\), let \(x_i\) be consistently signed first Chern classes of the tautological line quotients. Then \(H^*(X;\mathbb Q)\cong\mathbb Q[x_1,\ldots,x_n]/(e_1,\ldots,e_n)\), with \(\deg x_i=2\). The ideal equals that generated by all positive-degree symmetric polynomials.}
\GSfield{Reconstruction}{The field \(\mathbb C\), cohomological grading and signs of Chern generators are specified.}
\GSfield{Fate}{The Borel presentation is recalled by \citet[section~1]{bernsteingelfandgelfand1973schubert}. The notebook explicitly presents the lecture as exposition of BGG and Lascoux--Sch\"utzenberger.}
\GSfield{Present status}{As of November 2026, this remains a standard result under the stated hypotheses.}
\GSfield{Assessment}{literal statement: correct with missing hypotheses; publication match: exact; verification: read the result; development: none; confidence: high}
\end{GSitem}
\begin{GSitem}{1985-661.2}{E}{Bruhat cells and the homology/cohomology bases}
\GSfield{Kind; source}{construction; \GSpage{P001238-P001239}}
\GSfield{The notes}{The orbit of the permuted standard flag under transformations preserving the standard flag is \(C_w\), with complex dimension \(\ell(w)\). Cells are said to form the homology basis; dual classes have degree \(2\ell(w)\).}
\GSfield{Statement}{The Borel orbits \(C_w=BwB/B\) are affine spaces \(\mathbb C^{\ell(w)}\). Their closures \(\overline{C_w}\) have fundamental classes forming a basis of ordinary homology. The Kronecker-dual cohomology basis has degrees \(2\ell(w)\). With the standard opposite-Schubert convention, these are the Schubert classes represented by the polynomials below.}
\GSfield{Reconstruction}{The duality here is evaluation on homology classes, not an unqualified identification of a cell with its Poincar\'e-dual class of complementary degree.}
\GSfield{Fate}{\citet[section~2, in particular Theorem~2.11]{bernsteingelfandgelfand1973schubert} treats the cells and their closures. The passage from cells to fundamental classes of closures is essential.}
\GSfield{Present status}{As of November 2026, this remains a standard result under the stated hypotheses.}
\GSfield{Assessment}{literal statement: correct with missing hypotheses; publication match: exact; verification: read the result; development: none; confidence: high}
\end{GSitem}
\begin{GSitem}{1985-661.3}{E}{Schubert polynomial representatives}
\GSfield{Kind; source}{construction; \GSpage{P001239}}
\GSfield{The notes}{BGG classes modulo \(J\) have been lifted to actual polynomials of degree \(\ell(w)\), with \(P_{\mathrm{id}}=1\) and \(P_{w_0}=x_1^{n-1}x_2^{n-2}\cdots x_{n-1}\).}
\GSfield{Statement}{The Lascoux--Sch\"utzenberger Schubert polynomial \(\mathfrak S_w\in\mathbb Z[x_1,x_2,\ldots]\) is homogeneous of ordinary polynomial degree \(\ell(w)\); for \(w\in S_n\), the longest-element polynomial is \(x^\delta=x_1^{n-1}x_2^{n-2}\cdots x_{n-1}\). Its class in the Borel quotient represents the corresponding Schubert class.}
\GSfield{Reconstruction}{The canonical choice involves stability and divided differences, not merely an arbitrary lifting of a residue class.}
\GSfield{Fate}{\citet[Definition~1.3]{lascouxschutzenberger1982schubert} defines these polynomials; the note appeared on 29 March 1982. Their 1983 longer account also predates this seminar \citep{lascouxschutzenberger1983symmetry}.}
\GSfield{Present status}{As of November 2026, this remains a standard result under the stated hypotheses.}
\GSfield{Assessment}{literal statement: correct with missing hypotheses; publication match: exact; verification: read the result; development: none; confidence: high}
\end{GSitem}
\begin{GSitem}{1985-661.4}{E}{Normalization of the top cohomology class}
\GSfield{Kind; source}{claim; \GSpage{P001239}}
\GSfield{The notes}{The Vandermonde \(\prod_{i<j}(x_i-x_j)\) is called the same class modulo \(J\) as \(x^\delta\). An alternative is \(\rho^{m(n-1)/2}/(m(n-1)/2)!\).}
\GSfield{Statement}{Write \(N=n(n-1)/2\) and \(\rho=(n-1)x_1+(n-2)x_2+\cdots+x_{n-1}\). With compatible positive-root conventions, \[x^\delta\equiv\frac1{n!}\prod_{i<j}(x_i-x_j)\equiv\frac{\rho^N}{N!}\pmod J.\] The unnormalized Vandermonde represents \(n!\) times the top generator. Already for \(n=2\), \(x_1-x_2\equiv2x_1\).}
\GSfield{Reconstruction}{\ours\ The \(n=2\) identity \(x_2\equiv-x_1\pmod J\) already detects the missing factor. \oursComputed\ Exact Gr\"obner reductions for \(n=2,3,4\) separately checked the Vandermonde and \(\rho\)-power normalizations. The general identities are the cited BGG results.}
\GSfield{Fate}{\citet[Theorem~3.15 and Corollary~3.16]{bernsteingelfandgelfand1973schubert} gives the normalization and the \(\rho\)-power formula. The scan confirms both the absent \(1/n!\) and the letter \(m\) in the exponent.}
\GSfield{Present status}{As of November 2026, this remains a standard result under the stated hypotheses.}
\GSfield{Assessment}{literal statement: false as written; publication match: exact; verification: read the result; development: annotated; confidence: high}
\end{GSitem}
\begin{GSitem}{1985-661.5}{E}{The sign convention for divided differences}
\GSfield{Kind; source}{construction; \GSpage{P001240}}
\GSfield{The notes}{The numerator is \(f-s_if\), but the denominator is written \(x_{i+1}-x_i\).}
\GSfield{Statement}{The convention compatible with \(\mathfrak S_{w_0}=x^\delta\), \(\mathfrak S_{\mathrm{id}}=1\), and the usual positive Schubert polynomials is \[\partial_i f=\frac{f-s_if}{x_i-x_{i+1}}.\] This is again a polynomial. The opposite denominator changes the sign of each operator and hence contributes \((-1)^{\ell(w)}\) to a reduced product.}
\GSfield{Reconstruction}{Opposite-sign divided differences are mathematically legitimate if every accompanying formula is adjusted. The recorded top and identity classes make the unadjusted collection inconsistent. \ours\ The countercheck \(\partial_1x_1=1\) determines the compatible denominator sign. \oursComputed\ Exact polynomial calculations also checked square-zero and the braid identity on a nontrivial input and the longest-word top-to-one recurrence for \(n=2,3,4\). These finite audits do not replace the cited general theorem.}
\GSfield{Fate}{Compare \citet[equations~(1.1)--(1.2)]{lascouxschutzenberger1982schubert} and \citet[Lemma~3.3]{bernsteingelfandgelfand1973schubert}. The smallest check is \(\partial_1x_1=1\); the notebook convention gives \(-1\).}
\GSfield{Present status}{As of November 2026, this remains a standard result under the stated hypotheses.}
\GSfield{Assessment}{literal statement: false as written; publication match: exact; verification: read the result; development: annotated; confidence: high}
\end{GSitem}
\begin{GSitem}{1985-661.6}{E}{Independence of reduced expressions}
\GSfield{Kind; source}{theorem; \GSpage{P001240}}
\GSfield{The notes}{An operator \(A_w\) is obtained from a reduced product of adjacent transpositions and asserted independent of the decomposition. The last displayed transposition has a misplaced subscript.}
\GSfield{Statement}{The operators \(\partial_i\) satisfy \(\partial_i^2=0\), \(\partial_i\partial_j=\partial_j\partial_i\) for \(|i-j|>1\), and \(\partial_i\partial_{i+1}\partial_i=\partial_{i+1}\partial_i\partial_{i+1}\). Thus \(\partial_w\) is independent of a reduced expression \(w=s_{i_1}\cdots s_{i_{\ell(w)}}\). A nonreduced product of this length pattern vanishes when the corresponding word is not reduced.}
\GSfield{Reconstruction}{The factors are adjacent simple reflections, not arbitrary inversions. The last factor is \((i_{\ell(w)},i_{\ell(w)}+1)\).}
\GSfield{Fate}{The exact reduced/nonreduced statement is \citet[Theorem~3.4]{bernsteingelfandgelfand1973schubert}.}
\GSfield{Present status}{As of November 2026, this remains a standard result under the stated hypotheses.}
\GSfield{Assessment}{literal statement: correct with missing hypotheses; publication match: exact; verification: read the result; development: none; confidence: high}
\end{GSitem}
\begin{GSitem}{1985-661.7}{E}{Generating all Schubert polynomials from the longest element}
\GSfield{Kind; source}{theorem; \GSpage{P001239-P001240}}
\GSfield{The notes}{The displayed formula is \(P_w=A_{w_0w^{-1}}P_{w_0}\).}
\GSfield{Statement}{Under the usual left composition convention \(\partial_u=\partial_{i_1}\cdots\partial_{i_r}\) for \(u=s_{i_1}\cdots s_{i_r}\), one has \[\mathfrak S_w=\partial_{w^{-1}w_0}x^\delta.\] Equivalently, apply \(\partial_i\) whenever replacing the current permutation \(v\) by \(vs_i\) lowers its length. An opposite action convention reverses the product indexing.}
\GSfield{Reconstruction}{The notebook does not specify left versus right operator action. Its \(w_0w^{-1}\) is therefore treated as a convention ambiguity rather than an isolated transcription error.}
\GSfield{Fate}{\citet[Definition~1.3]{lascouxschutzenberger1982schubert} gives the construction; \citet[section~3]{bernsteingelfandgelfand1973schubert} provides its cohomological precursor.}
\GSfield{Present status}{As of November 2026, this remains a standard result under the stated hypotheses.}
\GSfield{Assessment}{literal statement: correct with missing hypotheses; publication match: exact; verification: read the result; development: none; confidence: high}
\end{GSitem}
\begin{GSitem}{1985-661.8}{E}{Symmetry at a non-descent}
\GSfield{Kind; source}{theorem; \GSpage{P001240}}
\GSfield{The notes}{\(P_w\) is symmetric in \(x_i,x_{i+1}\) if \((i,i+1)\) is not an inversion of \(w\).}
\GSfield{Statement}{For the usual Schubert polynomial, \(s_i\mathfrak S_w=\mathfrak S_w\) if and only if \(w(i)<w(i+1)\), equivalently \(\ell(ws_i)>\ell(w)\). Indeed \(\partial_i\mathfrak S_w=0\) in this case, whereas for a right descent it equals the nonzero polynomial \(\mathfrak S_{ws_i}\).}
\GSfield{Reconstruction}{The notebook records the valid sufficient direction. ``Inversion'' is made precise as an adjacent right descent, not a value-based left descent.}
\GSfield{Fate}{The exact equivalence is stated immediately after \citet[Definition~1.3]{lascouxschutzenberger1982schubert}.}
\GSfield{Present status}{As of November 2026, this remains a standard result under the stated hypotheses.}
\GSfield{Assessment}{literal statement: correct with missing hypotheses; publication match: exact; verification: read the result; development: none; confidence: high}
\end{GSitem}

\GStalk{662}{1985-11-25}{B. L. Feigin}{P001241-P001244}{Normal ordering, even lattices and affine current representations}
The continuation revisits Virasoro and no-ghost, then constructs root Lie algebras and affine current representations from even lattices and sign cocycles, concluding with the DDF recipe. Gelfand asks how normal ordering depends on the underlying geometry. Krichever--Novikov supplies a nearby later construction, but the intervention does not specify a unique theorem. The module assertion depending only on momentum norm fails at zero versus nonzero null momentum; the sign cocycle itself forces the evenness temporarily waived in the notes. The Leech-lattice aside must be kept distinct from the norm-two-root construction because that lattice is rootless. The next dated talk starts on P001245 and is outside this record.

\begin{GSitem}{1985-662.1}{E}{Normal ordering and the missing Sugawara half}
\GSfield{Kind; source}{construction; \GSpage{P001241}}
\GSfield{The notes}{The lecture restarts with \(L_i=:\sum_{m+n=i}g_{\mu\nu}a_m^\mu a_n^\nu:\), followed by the standard Virasoro bracket. Normal ordering is said to make the expression meaningful and only add a constant.}
\GSfield{Statement}{The correctly normalized free-boson modes are \(L_i=\frac12\sum_{m+n=i}g_{\mu\nu}:a_m^\mu a_n^\nu:\). Moving annihilation modes to the right makes each mode locally finite on algebraic Fock vectors. Reordering a finite truncation changes it by a scalar, but the infinite unrenormalized scalar need not converge; it is not an ordinary finite constant that may simply be added.}
\GSfield{Reconstruction}{\ours\ At zero momentum, \(L_{-2}|0\rangle=\frac12\sum_{\mu,\nu}g_{\mu\nu}a_{-1}^\mu a_{-1}^\nu|0\rangle\); the corresponding annihilation operator gives \(L_2L_{-2}|0\rangle=(d/2)|0\rangle\). Omitting both halves gives \(2d\), inconsistent with the listed central term. The same central coefficient is obtained with positive-sign species because the metric contraction is \(g_{\mu\nu}g^{\mu\nu}=d\) in either signature. \oursComputed\ Exact oscillator-polynomial computations give these values for \(d=1,2,26\).}
\GSfield{Fate}{\citet[equation~(2.10) and the calculation following Theorem~3.3]{frenkel1985representations} fixes the normalization. The absent half is confirmed by the scan and contradicts the correct \(2L_m\) formula on P001233.}
\GSfield{Present status}{As of November 2026, this remains a standard result under the stated hypotheses.}
\GSfield{Assessment}{literal statement: false as written; publication match: exact; verification: read the result; development: annotated; confidence: high}
\end{GSitem}
\begin{GSitem}{1985-662.2}{F}{Gelfand asks for normal ordering on other underlying spaces}
\GSfield{Kind; source}{question; \GSpage{P001241}; I. M. Gelfand}
\GSfield{The notes}{Gelfand asks what is needed for the operation \(:\ :\), and what happens for a lattice, torus or sphere.}
\GSfield{Statement}{The question asks which choices of geometry, modes, polarization and vacuum permit normal ordering and an analogue of the resulting centrally extended algebra. It does not say whether ``torus'' and ``sphere'' are world sheets, spatial manifolds or target spaces, nor which algebra should replace Virasoro.}
\GSfield{Reconstruction}{The two-point Riemann-surface formulation is a literature comparison, not a reconstructed exact question attributed to Gelfand.}
\GSfield{Fate}{A precise nearby development is the higher-genus operator-field and normal-ordering formalism of Krichever--Novikov, begun in 1986 and described in \citet[section~2, Theorem~2.1]{krichevernovikov1989riemann}. That theorem establishes a Virasoro-type central extension for its own two-point Riemann-surface setup; it does not answer every interpretation of this intervention.}
\GSfield{Present status}{As of November 2026, this remains an identifiable research direction but not a sufficiently specified original open problem. No C/D designation is warranted.}
\GSfield{Assessment}{literal statement: not determinable; publication match: thematic only; verification: read the result; development: none; confidence: medium}
\end{GSitem}
\begin{GSitem}{1985-662.3}{E}{Coordinate and current Fourier series}
\GSfield{Kind; source}{construction; \GSpage{P001241}}
\GSfield{The notes}{The notes list \(p^\mu(t)=\sum a_m^\mu e^{-imt}\), a coordinate series with coefficients \(a_m^\mu/m\), and \([x^\mu,a_0^\nu]=g^{\mu\nu}\), without the earlier factors of \(i\) or the linear zero-mode term.}
\GSfield{Statement}{In the Hermitian convention of the previous lecture, \([x^\mu,a_0^\nu]=ig^{\mu\nu}\) and \(Q^\mu(t)=x^\mu+a_0^\mu t+i\sum_{m\ne0}a_m^\mu e^{-imt}/m\), so \(\partial_tQ=P\). One can instead rescale generators in a complex algebraic convention, but then all factors of \(i\), zero modes and derivative identities must be adjusted together.}
\GSfield{Reconstruction}{The formula restores the earlier P001233 convention rather than silently combining the two displays.}
\GSfield{Fate}{The coordinate field is standard in \citet[introduction, equations~(1)--(3)]{frenkel1985representations}. The P001241 series is treated as abbreviated conventions, not as an independently valid derivative identity.}
\GSfield{Present status}{As of November 2026, this remains a standard result under the stated hypotheses.}
\GSfield{Assessment}{literal statement: correct with missing hypotheses; publication match: partial; verification: read the result; development: annotated; confidence: high}
\end{GSitem}
\begin{GSitem}{1985-662.4}{E}{Gelfand identifies the classical vector fields and their central extension}
\GSfield{Kind; source}{remark; \GSpage{P001241}; I. M. Gelfand}
\GSfield{The notes}{Gelfand writes \(L_n=e^{int}d/dt\) in the classical theory and compares its one-dimensional quantum extension with \([d/dx,x]=1\).}
\GSfield{Statement}{The complex vector fields \(D_n=i e^{int}d/dt\) satisfy \([D_m,D_n]=(m-n)D_{m+n}\). Their standard nontrivial one-dimensional central extension has cocycle \((m^3-m)\delta_{m,-n}/12\), giving Virasoro. The displayed unscaled fields satisfy the same relation with an extra factor of \(i\). The Heisenberg comparison is structural, not an identification of the two cocycles.}
\GSfield{Reconstruction}{A normalization of the Fourier vector fields is supplied to match the notebook bracket.}
\GSfield{Fate}{\citet[equations~(2.8)--(2.9)]{frenkel1985representations} states the vector-field/central-extension relation before this talk.}
\GSfield{Present status}{As of November 2026, this remains a standard result under the stated hypotheses.}
\GSfield{Assessment}{literal statement: correct with missing hypotheses; publication match: exact; verification: read the result; development: none; confidence: high}
\end{GSitem}
\begin{GSitem}{1985-662.5}{E}{When the Virasoro Fock module depends only on momentum norm}
\GSfield{Kind; source}{claim; \GSpage{P001241}}
\GSfield{The notes}{The notes assert \(\mathcal H_p\cong\mathcal H_{p'}\) whenever \((p,p)=(p',p')\), and label the module only by this scalar.}
\GSfield{Statement}{An isometry of the target form carrying \(p\) to \(p'\) induces a Virasoro-module isomorphism. In a real Lorentzian space, equal nonzero norms give such an isometry, as do two nonzero null vectors. Zero and a nonzero null vector are an exception despite equal norm: on level one, \(L_1(a_{-1}(v)|p\rangle)=(v,p)|p\rangle\), a rank-zero map for \(p=0\) and rank-one map for \(p\ne0\). Hence these graded Virasoro modules are not isomorphic.}
\GSfield{Reconstruction}{The isometry-orbit argument gives a sufficient statement. The level-one calculation isolates a literal counterexample in the Lorentzian setting actually used by these lectures. \ours\ The level-one rank calculation supplies a literal counterexample to the norm-only assertion. \oursComputed\ A two-dimensional Lorentzian block with momenta \(0\) and \((1,1)\) has equal norm zero but the two ranks are zero and one; the block embeds in the twenty-six-dimensional target.}
\GSfield{Fate}{The countercheck is immediate from the standard relations in \citet[equations~(2.2), (2.10), (2.14)]{frenkel1985representations}. It is a correction of an expository claim, not a new unresolved representation problem.}
\GSfield{Present status}{As of November 2026, the norm-only assertion is false without excluding the zero/nonzero-null distinction.}
\GSfield{Assessment}{literal statement: false as written; publication match: exact; verification: read the result; development: annotated; confidence: high}
\end{GSitem}
\begin{GSitem}{1985-662.6}{E}{Summing momentum sectors before imposing the constraints}
\GSfield{Kind; source}{construction; \GSpage{P001241-P001242}}
\GSfield{The notes}{The module \(\bigoplus_{(p,p)\in\mathbb R}\mathcal H_{(p,p)}\) is used to construct physical states, with the comment that the construction depends only on the Virasoro action.}
\GSfield{Statement}{One may form an algebraic direct sum of chosen Virasoro modules and impose \(P^1\) and its spurious quotient componentwise. For the subsequent vertex action, however, the summands must retain actual momenta in a set closed under the required shifts; one representative for each numerical norm does not specify those maps. Positivity also uses the invariant form, not the abstract Virasoro action alone.}
\GSfield{Reconstruction}{The statement separates a legitimate abstract module sum from the extra momentum and form data required later.}
\GSfield{Fate}{The momentum-labelled lattice sum \(S(\mathfrak h^-)\otimes\mathbb C[\Gamma]\) is explicit in \citet[section~2.1, equation~(2.2)]{frenkel1985representations}.}
\GSfield{Present status}{As of November 2026, this remains a standard result under the stated hypotheses.}
\GSfield{Assessment}{literal statement: correct with missing hypotheses; publication match: partial; verification: read the result; development: none; confidence: high}
\end{GSitem}
\begin{GSitem}{1985-662.7}{E}{The repeated no-ghost statement uses the quotient, not all constrained states}
\GSfield{Kind; source}{theorem; \GSpage{P001242}}
\GSfield{The notes}{The theorem says that twenty-four commuting oscillators generate \(\widetilde{\mathcal H}\) from a vacuum, and that its inherited form is positive definite, only when \(d=26\).}
\GSfield{Statement}{In the critical Lorentzian bosonic theory at nonzero momentum, constrained physical states have a positive semidefinite form; the quotient by its radical is positive definite and is represented by transverse DDF states. Spurious null vectors remain before quotienting. A single chosen DDF vacuum generates a normalized momentum tower, not the unrestricted direct sum of all momenta.}
\GSfield{Reconstruction}{The inserted tilde in the manuscript is preserved in notes but corrected in the mathematical statement. Zero momentum is again excluded.}
\GSfield{Fate}{The exact qualifications are given in \citet[Theorem~4.1 and section~4]{frenkel1985representations}; compare \citet[Theorem~5.1]{borcherds1992monstrous}. The positive-semidefinite assertion also has subcritical variants, so ``only 26'' should refer to the precise twenty-four-oscillator critical completeness statement.}
\GSfield{Present status}{As of November 2026, this remains a standard result under the stated hypotheses.}
\GSfield{Assessment}{literal statement: false as written; publication match: exact; verification: read the result; development: annotated; confidence: high}
\end{GSitem}
\begin{GSitem}{1985-662.8}{E}{The corrected oscillator vertex formula still omits zero-mode powers}
\GSfield{Kind; source}{construction; \GSpage{P001242}}
\GSfield{The notes}{Equation (9) now includes \(1/n\) in both exponentials and defines \(D_p|k\rangle=|k+p\rangle\).}
\GSfield{Statement}{This is a valid oscillator-stripped momentum-shifting series. To obtain the plane primary field one must also include \(z^{(p,a_0)}\); the cylinder-normalized field includes \(z^{(p,p)/2}\) as well. Integral pairings are needed for unshifted integral Laurent modes on a lattice sum.}
\GSfield{Fate}{The complete expression and mode convention are \citet[equation~(2.3)]{frenkel1985representations}. This passage repairs the denominator missing on P001237, but does not itself restore the zero modes.}
\GSfield{Present status}{As of November 2026, this remains a standard result under the stated hypotheses.}
\GSfield{Assessment}{literal statement: correct with missing hypotheses; publication match: exact; verification: read the result; development: none; confidence: high}
\end{GSitem}
\begin{GSitem}{1985-662.9}{E}{The Laurent-mode sign in the Heisenberg commutator}
\GSfield{Kind; source}{claim; \GSpage{P001242}}
\GSfield{The notes}{The expansion is \(X(p,z)=\sum_iX_i(p)z^i\), while equation (12) asserts \([a_n^\mu,X_i(p)]=p^\mu X_{i+n}(p)\).}
\GSfield{Statement}{The field relation is \([a_n^\mu,X(p,z)]=p^\mu z^nX(p,z)\). With \(X=\sum_iX_i z^i\), coefficient comparison gives \([a_n^\mu,X_i]=p^\mu X_{i-n}\). The recorded plus sign is correct only for the alternative expansion \(X=\sum_iX_i z^{-i}\).}
\GSfield{Reconstruction}{\ours\ The coefficient of \(z^i\) in \(z^n\sum_jX_jz^j\) is \(X_{i-n}\). \oursComputed\ Twenty-five exact Laurent-coefficient comparisons confirm this indexing; the general correction is the coefficient identity itself.}
\GSfield{Fate}{The negative-exponent convention occurs in \citet[equations~(2.3)--(2.4)]{frenkel1985representations}. Both conflicting signs in the notebook are clear in the scan.}
\GSfield{Present status}{As of November 2026, this remains a standard result under the stated hypotheses.}
\GSfield{Assessment}{literal statement: false as written; publication match: exact; verification: read the result; development: annotated; confidence: high}
\end{GSitem}
\begin{GSitem}{1985-662.10}{E}{Integral, even, and unimodular lattices}
\GSfield{Kind; source}{definition; \GSpage{P001242}}
\GSfield{The notes}{The lattice \(Q\subset\mathbb R^d\) is integral and unimodular when its Gram entries are integers and their determinant is one. Evenness is \((\gamma,\gamma)\in2\mathbb Z\).}
\GSfield{Statement}{A full-rank lattice in a nondegenerate real quadratic space is integral if \((Q,Q)\subset\mathbb Z\), even if every norm is even, and unimodular if \(Q=Q^*\), equivalently the Gram determinant has absolute value one. In the positive-definite case used for the finite root system below, this determinant is \(+1\).}
\GSfield{Reconstruction}{The rank and signature assumptions are stated. Indefinite unimodular Gram determinants can be \(-1\).}
\GSfield{Fate}{These standard definitions are recalled in \citet[section~1]{brouwerleech}. They precede the seminar; the reference is an undated modern exposition, not a priority claim.}
\GSfield{Present status}{As of November 2026, this remains a standard result under the stated hypotheses.}
\GSfield{Assessment}{literal statement: correct with missing hypotheses; publication match: exact; verification: read the result; development: none; confidence: high}
\end{GSitem}
\begin{GSitem}{1985-662.11}{E}{Norm-two roots of an integral positive-definite lattice}
\GSfield{Kind; source}{definition; \GSpage{P001243}}
\GSfield{The notes}{A root is a lattice vector of squared length two; later \(\Delta\) denotes the set of all such vectors and positivity is imposed.}
\GSfield{Statement}{For a positive-definite integral lattice \(Q\), \(\Delta=\{\alpha\in Q:(\alpha,\alpha)=2\}\) is finite. The reflections \(s_\alpha(v)=v-(v,\alpha)\alpha\) preserve \(Q\) and \(\Delta\). Thus \(\Delta\) is a reduced simply laced crystallographic root system in its span, possibly empty or reducible; its components have types \(A,D,E\).}
\GSfield{Reconstruction}{Positivity is essential to finiteness. Roots need not span the entire lattice space.}
\GSfield{Fate}{The finite simply laced root-system construction is the setting of \citet[section~1.1, equation~(1.3)]{frenkel1985representations}.}
\GSfield{Present status}{As of November 2026, this remains a standard result under the stated hypotheses.}
\GSfield{Assessment}{literal statement: correct with missing hypotheses; publication match: exact; verification: read the result; development: none; confidence: high}
\end{GSitem}
\begin{GSitem}{1985-662.12}{E}{The rank restriction for even unimodular positive lattices}
\GSfield{Kind; source}{theorem; \GSpage{P001243}}
\GSfield{The notes}{Immediately after defining roots, the notes say ``such'' objects exist only in dimensions divisible by eight.}
\GSfield{Statement}{A positive-definite even unimodular lattice has rank divisible by eight. Conversely, for every positive multiple \(8r\), the orthogonal sum \(E_8^{\oplus r}\) provides one. This statement concerns lattices with all three hypotheses, not norm-two roots; for example the rank-one root lattice \(A_1\) already has roots.}
\GSfield{Reconstruction}{The antecedent of ``such'' is resolved using the preceding integral/even/unimodular assumptions. In indefinite signature the appropriate restriction is on the signature difference.}
\GSfield{Fate}{\citet[Theorem~1.2]{brouwerleech} states the divisibility result and its theta-series context. This is a classical theorem, not a post-1985 development.}
\GSfield{Present status}{As of November 2026, this remains a standard result under the stated hypotheses.}
\GSfield{Assessment}{literal statement: correct with missing hypotheses; publication match: exact; verification: read the result; development: none; confidence: high}
\end{GSitem}
\begin{GSitem}{1985-662.13}{E}{The eight-dimensional lattice exercise}
\GSfield{Kind; source}{construction; \GSpage{P001243}}
\GSfield{The notes}{The exercise is to construct such a lattice in \(\mathbb R^8\).}
\GSfield{Statement}{Let \(D_8=\{x\in\mathbb Z^8:\sum_i x_i\equiv0\pmod2\}\) and \(h=(1/2,\ldots,1/2)\). Then \(E_8=D_8\cup(D_8+h)\) is a lattice: \(2h\in D_8\). It is integral and even since \((h,h)=2\) and \((x,h)=\frac12\sum x_i\in\mathbb Z\) on \(D_8\). Its covolume is \(\operatorname{covol}(D_8)/2=1\), hence it is unimodular.}
\GSfield{Reconstruction}{The coordinates and the elementary verification are supplied by the companion. \ours\ The parity, integrality and index-two covolume verification constitute the supplied proof of the standard \(E_8\) exercise.}
\GSfield{Fate}{This explicit standard construction solves the notebook exercise; \citet[section~1]{brouwerleech} supplies the surrounding even-unimodular lattice context. No novelty is claimed for the construction.}
\GSfield{Present status}{As of November 2026, this remains a standard result under the stated hypotheses.}
\GSfield{Assessment}{literal statement: correct; publication match: exact; verification: read the result; development: annotated; confidence: high}
\end{GSitem}
\begin{GSitem}{1985-662.14}{E}{ADE root lattices are even}
\GSfield{Kind; source}{claim; \GSpage{P001243}}
\GSfield{The notes}{An aside says that root lattices of the \(A,D,E\) series satisfy condition (14), namely evenness.}
\GSfield{Statement}{With every root normalized to squared length two, an integral linear combination of simple roots has even norm, because diagonal terms are even and cross terms occur twice. Thus simply laced root lattices are even. They need not be unimodular; evenness, not condition (13), is the recorded assertion.}
\GSfield{Reconstruction}{\ours\ For \(v=\sum n_i\alpha_i\), the norm is \(2\sum_i n_i^2+2\sum_{i<j}n_in_j(\alpha_i,\alpha_j)\), an even integer.}
\GSfield{Fate}{The normalization and root-lattice setting are in \citet[section~1.1]{frenkel1985representations}. The elementary expansion proves the exact notebook assertion.}
\GSfield{Present status}{As of November 2026, this remains a standard result under the stated hypotheses.}
\GSfield{Assessment}{literal statement: correct; publication match: exact; verification: read the result; development: annotated; confidence: high}
\end{GSitem}
\begin{GSitem}{1985-662.15}{E}{The sign cocycle forces evenness}
\GSfield{Kind; source}{construction; \GSpage{P001243}}
\GSfield{The notes}{The lattice is now said not necessarily to be even, yet \(\varepsilon(\alpha,\beta)\varepsilon(\beta,\alpha)=(-1)^{(\alpha,\beta)}\) is required for every pair, with \(\varepsilon(\alpha,\alpha)=-1\) on roots.}
\GSfield{Statement}{For an even integral free abelian lattice, choose a bimultiplicative normalized cocycle \(\varepsilon:Q\times Q\to\{\pm1\}\) with commutator \((-1)^{(\alpha,\beta)}\) and diagonal \(\varepsilon(\alpha,\alpha)=(-1)^{(\alpha,\alpha)/2}\). Such a choice is made on an ordered lattice basis and extended multiplicatively. The displayed commutator identity itself forces evenness on setting \(\alpha=\beta\), since the left side equals one. An arbitrary function with the two displayed identities need not satisfy the cocycle identity.}
\GSfield{Reconstruction}{Evenness and bimultiplicativity are restored. Odd-lattice superalgebra conventions require a different commutator sign and are not silently substituted. \ours\ Substitution \(\beta=\alpha\) gives \(1=(-1)^{(\alpha,\alpha)}\), directly contradicting the allowance of odd norms.}
\GSfield{Fate}{The word ``bilinear'' and the existence argument on a basis are explicit in \citet[equation~(1.2) and following paragraph]{frenkel1985representations}.}
\GSfield{Present status}{As of November 2026, this remains a standard result under the stated hypotheses.}
\GSfield{Assessment}{literal statement: false as written; publication match: exact; verification: read the result; development: annotated; confidence: high}
\end{GSitem}
\begin{GSitem}{1985-662.16}{E}{A finite Lie algebra from the lattice roots}
\GSfield{Kind; source}{construction; \GSpage{P001243}}
\GSfield{The notes}{The space \(\mathfrak g=\mathfrak h\oplus\bigoplus_{\alpha\in\Delta}\mathbb Cx_\alpha\) is given root brackets using \(\varepsilon\), including \([x_\alpha,x_{-\alpha}]=-\alpha\), and called an important construction of simple algebras.}
\GSfield{Statement}{For a positive-definite even integral lattice and the cocycle just specified, set \([h,h']=0\), \([h,x_\alpha]=(h,\alpha)x_\alpha\), use \(\varepsilon(\alpha,\beta)x_{\alpha+\beta}\) when \(\alpha+\beta\in\Delta\), zero when the sum is neither a root nor zero, and \([x_\alpha,x_{-\alpha}]=-\alpha\). This gives a reductive complex Lie algebra. It is semisimple when roots span \(\mathfrak h\), and simple when they also form an irreducible root system. The orthogonal complement of their span is central.}
\GSfield{Reconstruction}{The Cartan bracket and the spanning/irreducibility hypotheses are supplied. For a rootless lattice the displayed algebra is abelian, not simple. \ours\ If \(\Delta=\varnothing\), the displayed vector space is just the abelian Cartan space, which checks why simplicity cannot be claimed under the notebook assumptions alone.}
\GSfield{Fate}{\citet[equation~(1.3)]{frenkel1985representations} gives the same sign convention for the irreducible ADE case.}
\GSfield{Present status}{As of November 2026, this remains a standard result under the stated hypotheses.}
\GSfield{Assessment}{literal statement: correct with missing hypotheses; publication match: exact; verification: read the result; development: annotated; confidence: high}
\end{GSitem}
\begin{GSitem}{1985-662.17}{E}{Twisted lattice vertices and the affine basic representation}
\GSfield{Kind; source}{construction; \GSpage{P001243-P001244}}
\GSfield{The notes}{The factor \(\varepsilon_\gamma|p\rangle=\varepsilon(\gamma,p)|p\rangle\) modifies the vertex field, and its modes together with Heisenberg modes are said to give a projective current-algebra representation.}
\GSfield{Statement}{For an ADE root lattice \(Q\), the space \(S(\bigoplus_{n>0}(Q\otimes\mathbb C)t^{-n})\otimes\mathbb C[Q]\), with the corrected zero-mode factors and lattice cocycle, is the level-one basic module of the untwisted affine algebra of the corresponding finite simple Lie algebra. More general momentum lattices between the root and weight lattices decompose into level-one modules by root-lattice cosets; an arbitrary even unimodular lattice is not automatically a single ADE basic module.}
\GSfield{Reconstruction}{The finite ADE algebra, root lattice, level and completion conventions are specified. The result is linear for the central extension and projective for the unextended loop algebra.}
\GSfield{Fate}{The original construction is \citet{frenkelkac1980basic}, published in 1980. The exact representation statement is also \citet[Theorem~2.2]{frenkel1985representations}.}
\GSfield{Present status}{As of November 2026, this remains a standard result under the stated hypotheses.}
\GSfield{Assessment}{literal statement: correct with missing hypotheses; publication match: exact; verification: read the result; development: none; confidence: high}
\end{GSitem}
\begin{GSitem}{1985-662.18}{E}{The special rank-twenty-four lattice and sporadic groups}
\GSfield{Kind; source}{remark; \GSpage{P001244}}
\GSfield{The notes}{A special lattice in \(\mathbb R^{24}\) is said to occur here and in finite simple groups; the number of sporadic simple groups is twenty-six.}
\GSfield{Statement}{The intended lattice is the Leech lattice: up to isometry it is the unique positive-definite even unimodular rank-twenty-four lattice with no norm-two vectors. Its symmetry groups and the associated vertex constructions connect to sporadic groups and moonshine. Under the standard convention excluding the Tits group, there are twenty-six sporadic simple groups.}
\GSfield{Reconstruction}{The name Leech is inferred from the rank and group-theoretic context. Its absence of roots means that the preceding norm-two-current construction alone yields no nonabelian root currents for this lattice.}
\GSfield{Fate}{\citet[Theorem~1.1]{brouwerleech} states the precise lattice characterization. The maintained Atlas lists the twenty-six groups \citep{wilsonparkerbray2003sporadic}; the pre-seminar Monster-module link is identified by \citet{frenkellepowskymeurman1984natural}.}
\GSfield{Present status}{As of November 2026, this remains a standard result under the stated hypotheses.}
\GSfield{Assessment}{literal statement: correct with missing hypotheses; publication match: exact; verification: read the result; development: none; confidence: high}
\end{GSitem}
\begin{GSitem}{1985-662.19}{E}{The current-algebra brackets and the omitted opposite-root case}
\GSfield{Kind; source}{theorem; \GSpage{P001244}}
\GSfield{The notes}{The brackets vanish for \((\gamma_1,\gamma_2)>0\), give a signed \(\gamma_1+\gamma_2\) mode for inner product \(-1\), and other relations are left unwritten.}
\GSfield{Statement}{With consistent negative-exponent mode labels, lattice root currents satisfy the affine relations \[[x\otimes t^m,y\otimes t^n]=[x,y]\otimes t^{m+n}+m\delta_{m,-n}\kappa(x,y)I.\] For roots, nonnegative inner product gives commuting modes, inner product \(-1\) gives the signed sum-root mode, and opposite roots supply the Cartan mode and the central term. In the convention \([x_\alpha,x_{-\alpha}]=-\alpha\), one takes \(\kappa(x_\alpha,x_{-\alpha})=-1\).}
\GSfield{Reconstruction}{The notebook gives valid partial cases but does not display the central term responsible for projectivity. Positive-exponent mode labels reverse indices and must be handled consistently.}
\GSfield{Fate}{All three cases occur in \citet[equations~(1.4), (2.4), Theorem~2.2]{frenkel1985representations}.}
\GSfield{Present status}{As of November 2026, this remains a standard result under the stated hypotheses.}
\GSfield{Assessment}{literal statement: correct with missing hypotheses; publication match: exact; verification: read the result; development: none; confidence: high}
\end{GSitem}
\begin{GSitem}{1985-662.20}{E}{The final DDF residue recipe}
\GSfield{Kind; source}{construction; \GSpage{P001244}}
\GSfield{The notes}{Equations (26)--(30) use a null \(c\), a current with polarization \(\alpha\perp c\), and twenty-four such polarizations forming with \(c\) a basis of \(c^\perp\).}
\GSfield{Statement}{Take \(c\ne0\) null in signature \((25,1)\), choose an orthonormal basis of \(c^\perp/\mathbb Rc\), and use the normal-ordered residues of the corresponding currents times \(Y_{nc}\). On a sector with the appropriate normalized value of \((c,a_0)\), these are twenty-four Heisenberg oscillators commuting with Virasoro; their classes generate the nonzero-momentum physical quotient as in record~660.}
\GSfield{Reconstruction}{The omitted momentum normalization, normal ordering and orthonormal quotient basis are restored. The displayed arbitrary basis alone does not imply standard delta-normalized brackets.}
\GSfield{Fate}{This repeats the construction in \citet[equations~(4)--(5) and (4.8)--(4.9)]{frenkel1985representations}.}
\GSfield{Present status}{As of November 2026, this remains a standard result under the stated hypotheses.}
\GSfield{Assessment}{literal statement: correct with missing hypotheses; publication match: exact; verification: read the result; development: none; confidence: high}
\end{GSitem}

\GStalk{663}{1985-12-16}{S. P. Novikov}{P001245-P001248 (upper)}{Two-dimensional Schrödinger operators at one energy}
The talk moves from the failure of a naive all-energy analogy with one-dimensional inverse scattering to Manakov triples and the two-point Baker--Akhiezer construction. It explains potential and reality reductions by involutions of the spectral curve, Prym theta functions, and the two-dimensional analogue of KdV. The main constructions were published in 1976 and 1984. Several derivatives in the notebook are internally inconsistent; these are mathematical emendations of faithfully transcribed formulas. The final rapidly decaying assertion is less precise than the periodic construction. The actual heading is on P001245; P001244 still belongs to Feigin.

\begin{GSitem}{1985-663.1}{E}{The blanket exclusion of ordinary Lax pairs}
\GSfield{Kind; source}{claim; \GSpage{P001245}}
\GSfield{The notes}{For \(L=-\partial\bar\partial+U=-\Delta+U\), the notes say \ru{не бывает пар Лакса} and infer that \(L_t=[A,L]\) gives only Euclidean motions.}
\GSfield{Statement}{This universal assertion is false without restrictions on \(U\) and \(A\). For a smooth solution of \(u_t=6uu_x-u_{xxx}\), take \(L=-\partial_x^2-\partial_y^2+u(x,t)\) and \(A=-4\partial_x^3+6u\partial_x+3u_x\). Then \(L_t=[A,L]\), including nontravelling KdV solutions. What motivates the talk is instead the richer genuinely two-dimensional fixed-energy equation \(L_t=[A,L]+BL\).}
\GSfield{Reconstruction}{\ours\ The counterexample and its differential-operator check are ours, not a reported seminar intervention. The equality \(\partial\bar\partial=\Delta\) also requires the notebook normalization; with Wirtinger derivatives it equals \(\Delta/4\).}
\GSfield{Fate}{The one-dimensional Lax mechanism already existed; the two-dimensional fixed-energy replacement is in \citet[pp.~245--246]{manakov1976inverse}. This is not evidence for an unrestricted no-Lax theorem.}
\GSfield{Present status}{As of November 2026: the unrestricted wording is refuted by the displayed reduction; the intended genuinely two-dimensional obstruction is not specified sufficiently to recover a theorem.}
\GSfield{Assessment}{literal statement: false as written; publication match: partial; verification: read the result; development: annotated; confidence: high}
\end{GSitem}
\begin{GSitem}{1985-663.2}{E}{Counting the variables of scattering data}
\GSfield{Kind; source}{remark; \GSpage{P001245}}
\GSfield{The notes}{The scattering amplitude \(S(k_1,k_2)\), with \(|k_1|=|k_2|\), has \(2n-1\) variables; at fixed energy and \(n=2\) it has two.}
\GSfield{Statement}{For real positive energy \(E\), incoming and outgoing momenta belong to \(S_E^{n-1}\times S_E^{n-1}\), of dimension \(2n-2\). Allowing \(E>0\) adds one parameter. Thus the fixed-energy data in two dimensions and the potential on \(\mathbb R^2\) both depend on two real variables. For an actual scattering amplitude one must additionally impose sufficient decay and regularity, e.g. a smooth compactly supported real potential.}
\GSfield{Reconstruction}{The dimension count is exact. It is not a proof of injectivity, surjectivity, or compatibility of inverse-scattering data; the notes supply none of these analytic hypotheses.}
\GSfield{Fate}{The fixed-energy viewpoint precedes the talk in \citet[\S1]{dubrovinkrichevernovikov1976schrodinger}; later analytic inverse-scattering results are discussed by \citet{novikovhenkin1987dbar}.}
\GSfield{Present status}{As of November 2026: established in the stated scope; the qualifications above are essential.}
\GSfield{Assessment}{literal statement: correct with missing hypotheses; publication match: partial; verification: read the result; development: annotated; confidence: high}
\end{GSitem}
\begin{GSitem}{1985-663.3}{E}{The Manakov L-A-B triple preserves a kernel}
\GSfield{Kind; source}{method; \GSpage{P001245}}
\GSfield{The notes}{The notebook replaces \(L_t=[A,L]\) by \(L_t=[A,L]+BL\), or equivalently \(C_1L+LC_2\), and imposes \(L\psi=0\).}
\GSfield{Statement}{Let \(L,A,B\) be time-dependent differential operators on a common class of smooth functions. If \(L_t=[A,L]+BL\) and \(\psi_t=A\psi\), then \(\partial_t(L\psi)=(A+B)L\psi\). Consequently the zero-energy equation is formally compatible with the evolution. At energy \(E_0\), replace every occurrence of \(L\) in the right ideal by \(L-E_0\). This is a formal compatibility assertion, not a claim that the entire spectrum is isospectral.}
\GSfield{Reconstruction}{\ours\ The equation for \(\psi_t\) and the distinction between one level and the whole spectrum are supplied to make the recorded method precise.}
\GSfield{Fate}{Published before the talk in \citet{manakov1976inverse}; \citet[\S4]{dubrovinkrichevernovikov1976schrodinger} gives operator identities modulo the Schrödinger operator.}
\GSfield{Present status}{As of November 2026: established in the stated scope; the qualifications above are essential.}
\GSfield{Assessment}{literal statement: correct; publication match: exact; verification: read the result; development: annotated; confidence: high}
\end{GSitem}
\begin{GSitem}{1985-663.4}{E}{Two-point Baker--Akhiezer data and uniqueness}
\GSfield{Kind; source}{construction; \GSpage{P001245-P001246}}
\GSfield{The notes}{A genus-\(g\) curve \(\Gamma\), two marked points \(\infty_1,\infty_2\), local parameters \(w_i=k_i^{-1}\), and a degree-\(g\) pole divisor determine \(\psi\) with exponential singularities \(c_1e^{ik_1z}\), \(c_2e^{ik_2\bar z}\). The normalization is \(c_1=1\).}
\GSfield{Statement}{Let \(\Gamma\) be a smooth compact connected complex curve; let the distinct marked points and the effective nonspecial degree-\(g\) divisor \(D\) be disjoint. Prescribe the two indicated exponential germs and at most the poles \(D\). For parameter values where the corresponding evaluation problem is nonspecial, the normalized Baker--Akhiezer function exists uniquely. Its dependence on the flow parameters is meromorphic; exceptional theta-divisor values must be excluded from a local uniqueness assertion.}
\GSfield{Reconstruction}{The nonspecialness condition and the exclusion of exceptional parameter values are reconstructed. Distinct poles alone do not supply them. Earlier occurrence 1975-374.4 concerns the one-point construction, not an identical two-point inverse problem.}
\GSfield{Fate}{\citet[\S3, Lemma~1]{dubrovinkrichevernovikov1976schrodinger} supplies the two-point construction (received 19 February 1976).}
\GSfield{Present status}{As of November 2026: established in the stated scope; the qualifications above are essential.}
\GSfield{Assessment}{literal statement: correct with missing hypotheses; publication match: exact; verification: read the result; development: annotated; confidence: high}
\end{GSitem}
\begin{GSitem}{1985-663.5}{E}{Gauge normalization and recovery of the operator}
\GSfield{Kind; source}{construction; \GSpage{P001246}}
\GSfield{The notes}{Equations (9)--(10) read \(L=\partial\bar\partial+A\partial+B\), \(A=\bar\partial\log c_2\), following \(c_1=1\).}
\GSfield{Statement}{Fix \(\partial=\partial_z\), \(\bar\partial=\partial_{\bar z}\) and the displayed exponential asymptotics. In the gauge \(\psi\sim e^{ik_1z}(1+\xi_1/k_1+\cdots)\), \(\psi\sim c_2e^{ik_2\bar z}(1+\cdots)\), the compatible operator with leading term \(\partial\bar\partial\) is \(L=\partial\bar\partial+a\bar\partial+b\), with \(a=-\partial\log c_2\) and \(b=-i\bar\partial\xi_1\). Multiplication \(\psi\mapsto c\psi\), \(c\ne0\), conjugates the operator by \(c\).}
\GSfield{Reconstruction}{\ours\ The signs and derivative placements are rederived by cancelling the leading Laurent terms at each puncture, rather than copied from a potentially differently normalized publication. The source has both the unbarred derivative after \(A\) and the positive coefficient as transcribed.}
\GSfield{Fate}{The construction is the pre-existing one of \citet[\S3, Lemma~2]{dubrovinkrichevernovikov1976schrodinger}; the notebook formula is not consistent with its own normalization.}
\GSfield{Present status}{As of November 2026: established in the stated scope; the qualifications above are essential.}
\GSfield{Assessment}{literal statement: false as written; publication match: partial; verification: read the result; development: annotated; confidence: high}
\end{GSitem}
\begin{GSitem}{1985-663.6}{E}{Quasiperiodicity does not guarantee a smooth potential}
\GSfield{Kind; source}{claim; \GSpage{P001246}}
\GSfield{The notes}{\(A,B\) are called quasiperiodic for a Riemann surface in general position.}
\GSfield{Statement}{In the algebro-geometric construction the coefficients are restrictions of meromorphic theta expressions to an affine plane \(Uz+V\bar z+Z\) in a Jacobian. Periodic or quasiperiodic dependence follows when the relevant real flow lies on a real compact torus. Smooth bounded real coefficients additionally require that the real flow avoid the denominator theta divisor. General complex spectral data need not give a globally smooth or bounded real potential.}
\GSfield{Reconstruction}{The distinction between meromorphic theta dependence and real quasiperiodicity is essential; neither smoothness nor reality follows from the notebook phrase alone.}
\GSfield{Fate}{\citet[\S3, Lemma~3 and Theorem~1]{dubrovinkrichevernovikov1976schrodinger} separates real-oval regularity from the theta formulas.}
\GSfield{Present status}{As of November 2026: established in the stated scope; the qualifications above are essential.}
\GSfield{Assessment}{literal statement: correct with missing hypotheses; publication match: partial; verification: read the result; development: annotated; confidence: high}
\end{GSitem}
\begin{GSitem}{1985-663.7}{E}{Zero magnetic flux in a periodic vector-potential gauge}
\GSfield{Kind; source}{claim; \GSpage{P001246}}
\GSfield{The notes}{The magnetic field is required to have zero flux through a period cell; the condition is called \ru{паразитное}.}
\GSfield{Statement}{If a smooth magnetic potential \(A_1dx+A_2dy\) is globally periodic on a two-torus, then its curvature \(B=\partial_xA_2-\partial_yA_1\) has integral zero over a fundamental cell, by cancellation of opposite boundary integrals. This does not exclude nonzero quantized flux for a connection on a nontrivial line bundle, which has no globally periodic vector potential of this form.}
\GSfield{Reconstruction}{\ours\ The bundle qualification is reconstructed to avoid interpreting a gauge restriction as a prohibition on magnetic flux in general.}
\GSfield{Fate}{This is an explicit starting assumption of \citet[\S1]{dubrovinkrichevernovikov1976schrodinger}, not an unresolved mathematical obstruction.}
\GSfield{Present status}{As of November 2026: established in the stated scope; the qualifications above are essential.}
\GSfield{Assessment}{literal statement: correct with missing hypotheses; publication match: exact; verification: read the result; development: annotated; confidence: high}
\end{GSitem}
\begin{GSitem}{1985-663.8}{E}{The fixed-energy complex Fermi curve}
\GSfield{Kind; source}{definition; \GSpage{P001246}}
\GSfield{The notes}{Bloch functions form a complex set; its fixed-energy level is called the Fermi surface. A two-dimensional operator is finite-gap when that surface has finite genus.}
\GSfield{Statement}{For a doubly periodic two-dimensional elliptic Schrödinger operator, consider Bloch solutions \(\psi(x+\ell)=\chi(\ell)\psi(x)\) with complex characters \(\chi:\Lambda\to\mathbb C^*\), at a fixed energy \(E_0\). The resulting analytic spectral locus is generically a complex curve, possibly singular. In the rank-one algebro-geometric class its normalization compactifies to a finite-genus curve by adding the prescribed asymptotic points. This is a fixed-energy definition; it does not mean finitely many gaps in the full two-dimensional spectrum.}
\GSfield{Reconstruction}{Normalization, compactification and rank-one hypotheses are made explicit. A literal finite-genus assertion about an arbitrary singular analytic set is insufficient.}
\GSfield{Fate}{\citet[\S2, Definitions~1--2]{dubrovinkrichevernovikov1976schrodinger} distinguishes algebraic and weakly algebraic operators.}
\GSfield{Present status}{As of November 2026: established in the stated scope; the qualifications above are essential.}
\GSfield{Assessment}{literal statement: correct with missing hypotheses; publication match: exact; verification: read the result; development: annotated; confidence: high}
\end{GSitem}
\begin{GSitem}{1985-663.9}{E}{The holomorphic involution removes the magnetic term}
\GSfield{Kind; source}{theorem; \GSpage{P001246-P001247}}
\GSfield{The notes}{For \(\sigma^2=1\), \(\sigma(\infty_i)=\infty_i\), \(\sigma(w_i)=-w_i\), impose \(D+\sigma D\sim K+\infty_1+\infty_2\). The notes say the first-order term can then be removed.}
\GSfield{Statement}{In the smooth two-point, rank-one Baker--Akhiezer setup, assume a holomorphic involution with exactly the two marked fixed points, odd local parameters, and an admissible nonspecial divisor satisfying the displayed linear equivalence. Equivalently, the required meromorphic differential has simple poles of opposite residues at the marked points and zero divisor \(D+\sigma D\). Its residue identity forces the second leading coefficient to be constant after normalization; the constructed operator is purely potential, on the locus where the coefficients are regular.}
\GSfield{Reconstruction}{The precise fixed-point and admissibility conditions are supplied. These are sufficient spectral-data conditions; necessity is not asserted for arbitrary operators or arbitrary redundant spectral data.}
\GSfield{Fate}{Already published in \citet[Theorem~1, smooth \(k=0\) case]{veselovnovikov1984potential}, received 5 March 1984.}
\GSfield{Present status}{As of November 2026: established in the stated scope; the qualifications above are essential.}
\GSfield{Assessment}{literal statement: correct with missing hypotheses; publication match: partial; verification: read the result; development: annotated; confidence: high}
\end{GSitem}
\begin{GSitem}{1985-663.10}{E}{Reality requires the compatible antiholomorphic involution}
\GSfield{Kind; source}{theorem; \GSpage{P001247}}
\GSfield{The notes}{The commuting involutions \(\sigma,\tau\) satisfy \(\tau\infty_1=\infty_2\), \(\tau D=D\) and \(\tau w_1=\pm\overline{w_2}\). Conditions (17)--(19) are called necessary and sufficient for reality and potentiality.}
\GSfield{Statement}{With the preceding admissible potential spectral data, let \(\tau\) be an antiholomorphic involution commuting with \(\sigma\), exchanging the punctures, preserving \(D\), and satisfying the sign convention \(\tau^*w_1=-\overline{w_2}\) for the chosen oscillatory exponentials. The reconstructed potential is real wherever nonsingular. Smoothness on all of \(\mathbb R^2\) requires additional admissibility/real-locus conditions. A necessity claim must specify minimal spectral data and the involutions induced by the operator; the notebook does not do so.}
\GSfield{Reconstruction}{The sign is not a free choice once the exponential and operator conventions are fixed. The universally quantified necessary-and-sufficient wording is therefore not retained.}
\GSfield{Fate}{The sufficient statement is exactly the smooth specialization of \citet[Theorem~2]{veselovnovikov1984potential}; the notebook is broader than that checked theorem.}
\GSfield{Present status}{As of November 2026: established in the stated scope; the qualifications above are essential.}
\GSfield{Assessment}{literal statement: correct with missing hypotheses; publication match: partial; verification: read the result; development: annotated; confidence: medium}
\end{GSitem}
\begin{GSitem}{1985-663.11}{E}{Prym theta functions and potential spectral data}
\GSfield{Kind; source}{construction; \GSpage{P001247}}
\GSfield{The notes}{The eigenfunction is expressed through Prym theta functions, described as more numerous than Jacobian theta functions.}
\GSfield{Statement}{For the involution \(\sigma\) in the potential reduction, the connected kernel of the norm map \(\operatorname{Jac}(\Gamma)\to\operatorname{Jac}(\Gamma/\sigma)\) is the Prym variety. When the double cover is ramified at exactly two points, the admissible divisor classes form a translate of this Prym, whose dimension is half the genus of \(\Gamma\). The potential and Baker--Akhiezer function admit theta formulas on that variety, away from their polar loci.}
\GSfield{Reconstruction}{The informal comparison with Jacobians is not a numerical count. The Prym description of the divisor family is the recoverable assertion.}
\GSfield{Fate}{\citet[Theorem~1(2)]{veselovnovikov1984potential} and \citet{veselovnovikov1984explicit} published the reduction and explicit formulas in 1984.}
\GSfield{Present status}{As of November 2026: established in the stated scope; the qualifications above are essential.}
\GSfield{Assessment}{literal statement: correct with missing hypotheses; publication match: exact; verification: read the result; development: annotated; confidence: high}
\end{GSitem}
\begin{GSitem}{1985-663.12}{F}{An underspecified higher-dimensional finite-gap criterion}
\GSfield{Kind; source}{remark; \GSpage{P001247}}
\GSfield{The notes}{Under \ru{Многомерный случай}, finite-gapness is described by the existence of a nontrivial algebraic curve among all Bloch functions at arbitrary energies.}
\GSfield{Statement}{No sufficient definition is recoverable from this sentence alone. It omits the parameter space, compactification, rank, puncture asymptotics and the sense in which the curve determines the operator. Merely containing some algebraic curve is weaker than reconstructing an operator from complete spectral data on it.}
\GSfield{Reconstruction}{The omission is retained; no classification theorem or open inverse problem is supplied on the speaker\textquotesingle s behalf.}
\GSfield{Fate}{\citet[final paragraph]{dubrovinkrichevernovikov1976schrodinger} states the need for spectral data on a curve that uniquely determine the operator. This is only a thematic comparator for the notebook shorthand.}
\GSfield{Present status}{As of November 2026: this notebook occurrence remains too underspecified for an exact theorem or a definite research problem.}
\GSfield{Assessment}{literal statement: not determinable; publication match: thematic only; verification: read the result; development: annotated; confidence: low}
\end{GSitem}
\begin{GSitem}{1985-663.13}{E}{The Novikov--Veselov evolution and exchanged derivatives}
\GSfield{Kind; source}{theorem; \GSpage{P001247-P001248}}
\GSfield{The notes}{With \(L=\partial\bar\partial+V\), \(A=\partial^3+u\partial\), the notes give \(\partial u=3\bar\partial V\), \(B=\partial\bar u+\bar\partial u\), and \(V_t=(\partial^3+\partial u)V+\mathrm{c.c.}\).}
\GSfield{Statement}{For smooth real \(V\) and complex \(u\) satisfying \(\bar\partial u=3\partial V\), set \(P=\partial^3+\bar\partial^3+u\partial+\bar u\bar\partial\) and \(B=\partial u+\bar\partial\bar u\). Then \([P,L]+BL\) is multiplication by \(\partial^3V+\bar\partial^3V+\partial(uV)+\bar\partial(\bar uV)\). Thus \(L_t=[P,L]+BL\) is equivalent to this evolution of \(V\), with the auxiliary equation and boundary/normalization conditions for \(u\).}
\GSfield{Reconstruction}{\ours\ The derivative placements are corrected by a full commutator expansion. In the final expression \(\partial u\) must mean composition \(\partial\circ M_u\), not only multiplication by \(\partial u\). The symbolic check and independent polynomial substitutions are supplied in \texttt{code/check\_mathematics.py}.}
\GSfield{Fate}{The hierarchy was published before the seminar in \citet{veselovnovikov1984explicit}. The calculation here verifies the convention in this chapter, not a new evolution equation.}
\GSfield{Present status}{As of November 2026: established in the stated scope; the qualifications above are essential.}
\GSfield{Assessment}{literal statement: false as written; publication match: partial; verification: read the result; development: annotated; confidence: high}
\end{GSitem}
\begin{GSitem}{1985-663.14}{B}{Rational fixed-energy soliton potentials}
\GSfield{Kind; source}{claim; \GSpage{P001248 (upper)}}
\GSfield{The notes}{The talk ends by mentioning exact rational soliton solutions.}
\GSfield{Statement}{The identifiable topic is rational or determinant-form potentials for the two-dimensional Schrödinger problem at one fixed energy, with associated Novikov--Veselov evolutions. It is not the assertion that every one-dimensional multisoliton construction extends to arbitrary two-dimensional scattering data. No particular rational formula or parameter conditions are recorded here.}
\GSfield{Reconstruction}{A specific family is not reconstructed because the note supplies none. The bibliographic identification is to a construction, not to a missing individual solution.}
\GSfield{Fate}{\citet{grinevichnovikov1985analogs} appeared in the 1985 fourth issue; \citet{grinevich1986rational} explicitly identifies rational Novikov--Veselov solitons with reflectionless fixed-energy potentials. R.~G.~Novikov here is distinct from the seminar speaker S.~P.~Novikov.}
\GSfield{Present status}{As of November 2026: established explicit families; the particular solution alluded to in the notebook cannot be identified. Exact pre-talk appearance of the 1985 issue has not been established, so the chronology is not reduced to a year alone.}
\GSfield{Assessment}{literal statement: correct with missing hypotheses; publication match: partial; verification: publisher metadata checked; development: annotated; confidence: medium}
\end{GSitem}
\begin{GSitem}{1985-663.15}{F}{The rapidly decaying inverse problem is said to be essentially solved}
\GSfield{Kind; source}{claim; \GSpage{P001248 (upper)}}
\GSfield{The notes}{\ru{Быстроубыв. задача тоже по существу решена} is followed by R.~Novikov, P.~Grinevich, Manakov.}
\GSfield{Statement}{The sentence does not specify energy, decay rate, reality, exceptional points, smallness, admissible data or uniqueness class. Therefore it does not support a theorem solving the unrestricted two-dimensional inverse-scattering problem. Fixed-energy results below the ground state and results using generalized analytic functions are distinct precisely formulated successors.}
\GSfield{Reconstruction}{No unstated analytic hypothesis is attributed to the 1985 talk. Later work is used only to delimit the announcement.}
\GSfield{Fate}{Compare \citet{novikovhenkin1987dbar} (received 23 October 1986) and \citet{grinevichnovikov1988negative}. They do not retroactively establish all possible meanings of the notebook sentence.}
\GSfield{Present status}{As of November 2026: the recorded global claim has no recoverable quantifiers; the checked literature treats specified energy and potential classes. It is not classified as an undeveloped problem.}
\GSfield{Assessment}{literal statement: not determinable; publication match: thematic only; verification: publisher metadata checked; development: annotated; confidence: low}
\end{GSitem}

\GStalk{664}{1985-12-16}{V. G. Knizhnik}{P001248 (middle)-P001252}{Strings, moduli and quadratic differentials}
The second talk of 16 December constructs the critical bosonic-string measure in conformal gauge, explains the role of quadratic and Beltrami differentials, and works through the torus determinant before sketching higher-genus theta formulas. Novikov intervenes about the Schottky problem. The notes combine classical geometry, Polyakov\textquotesingle s 1981 formalism, and material soon published by Belavin and Knizhnik. They do not provide a rigorous measure on all embedded surfaces. The heading says \ru{Д. Книжник. (Дима)}; the catalogue identification as V.~G.~Knizhnik is retained with the discrepancy explicit.

\begin{GSitem}{1985-664.1}{E}{The Polyakov partition function is a formal gauge integral}
\GSfield{Kind; source}{construction; \GSpage{P001248}}
\GSfield{The notes}{The sum \(Z_g=\sum_{\text{surfaces of genus }g}e^{-A}\) is rewritten as \(\int\mathcal Dg\,\mathcal DX\,e^{-\int\sqrt g g^{ab}\partial_aX\cdot\partial_bX}\).}
\GSfield{Statement}{For a fixed closed oriented smooth surface \(\Sigma\), a Riemannian metric \(g\), and a smooth map \(X:\Sigma\to\mathbb R^D\), the Polyakov action is the Dirichlet energy, with a chosen positive normalization. It is invariant under pullback by diffeomorphisms and under \(g\mapsto e^{2\sigma}g\) in dimension two. Eliminating \(g\) classically at an immersion gives the area action, up to normalization. The functional integral is a perturbative prescription modulo those symmetries, not a defined Lebesgue measure on embeddings.}
\GSfield{Reconstruction}{Maps rather than only embeddings, the gauge quotient and the distinction between a classical equivalence and a quantum integration prescription are supplied.}
\GSfield{Fate}{Pre-existing in \citet[eqs.~(3)--(12)]{polyakov1981quantum}. Earlier occurrence 1982-569.9 already treated conformal gauge and residual moduli.}
\GSfield{Present status}{As of November 2026: established in the stated scope; the qualifications above are essential.}
\GSfield{Assessment}{literal statement: correct with missing hypotheses; publication match: partial; verification: read the result; development: annotated; confidence: high}
\end{GSitem}
\begin{GSitem}{1985-664.2}{E}{Critical dimension and conformal gauge leave moduli}
\GSfield{Kind; source}{theorem; \GSpage{P001248}}
\GSfield{The notes}{After coordinate gauge fixing \(g_{ab}=\rho\delta_{ab}\), the notes set \(D=26\) and use the additional conformal symmetry to reduce the integral to \(M_g\).}
\GSfield{Statement}{Locally every oriented two-dimensional metric has isothermal coordinates. Globally one chooses a reference conformal structure and retains Teichm\"uller parameters, then divides by the mapping class group and residual automorphisms. In the flat-target bosonic model the matter central charge is \(D\), and that of the reparametrization ghosts is \(-26\); cancellation occurs at \(D=26\). This cancels the local Weyl anomaly, not divergences at the boundary of moduli space.}
\GSfield{Reconstruction}{The genus-dependent automorphisms and the local/global distinction are missing from the shorthand gauge statement.}
\GSfield{Fate}{\citet[eqs.~(17)--(26)]{polyakov1981quantum}; the higher-genus implementation appears in \citet[\S2]{belavinknizhnik1986complex}.}
\GSfield{Present status}{As of November 2026: established in the stated scope; the qualifications above are essential.}
\GSfield{Assessment}{literal statement: correct with missing hypotheses; publication match: exact; verification: read the result; development: annotated; confidence: high}
\end{GSitem}
\begin{GSitem}{1985-664.3}{E}{Quadratic differentials and their Gram matrix}
\GSfield{Kind; source}{construction; \GSpage{P001248-P001249}}
\GSfield{The notes}{Holomorphic quadratic differentials \(f^i(dz)^2\) enter the denominator \(\det(f^i,f^j)\); the written integrand in (5) is \(\rho^{-1}\overline{f^i}f^i\).}
\GSfield{Statement}{For a smooth compact curve \(C\) with metric \(\rho|dz|^2\), the Hermitian Gram matrix of a basis \(f^1,\ldots,f^N\) of \(H^0(C,K_C^2)\) is \(G_{ij}=\int_C\rho^{-1}\overline{f^i}f^j\,d^2z\), using a consistent positive area convention. For \(g\ge2\), \(N=3g-3\). The determinant transforms by \(|\det A|^2\) under a basis change and cancels the corresponding coordinate-volume factor.}
\GSfield{Reconstruction}{The repeated \(i\) on the right of (5) is a genuine notebook error; without changing the second index to \(j\), the alleged Gram matrix has identical columns.}
\GSfield{Fate}{See \citet[eqs.~(2.13)--(2.19)]{belavinknizhnik1986complex} for its use in deformation theory; the pairing predates the talk.}
\GSfield{Present status}{As of November 2026: established in the stated scope; the qualifications above are essential.}
\GSfield{Assessment}{literal statement: false as written; publication match: partial; verification: read the result; development: annotated; confidence: high}
\end{GSitem}
\begin{GSitem}{1985-664.4}{E}{Products of Abelian differentials and the hyperelliptic exception}
\GSfield{Kind; source}{theorem; \GSpage{P001249}}
\GSfield{The notes}{In genus two, \(\omega_1^2,\omega_1\omega_2,\omega_2^2\) form the quadratic differentials. They are said always to be obtained this way, with a surplus for \(g\ge4\).}
\GSfield{Statement}{For a smooth compact connected complex curve of genus \(g\ge2\), the multiplication map \(\operatorname{Sym}^2H^0(K_C)\to H^0(K_C^2)\) is surjective for \(g=2\), and for nonhyperelliptic \(g\ge3\). Its kernel then has dimension \((g-2)(g-3)/2\). For hyperelliptic \(g\ge3\), its image instead has dimension \(2g-1\), so the cokernel has dimension \(g-2\). Thus the unrestricted word “always” is false.}
\GSfield{Reconstruction}{\ours\ The omitted nonhyperellipticity hypothesis is decisive. On \(y^2=f(x)\), \(\deg f=2g+2\) with simple roots, the products span \(x^k(dx)^2/y^2\), \(0\le k\le2g-2\); the additional holomorphic forms \(x^j(dx)^2/y\), \(0\le j\le g-3\), prove the defect. The local calculation is included in the mathematical-check appendix.}
\GSfield{Fate}{This is Max Noether\textquotesingle s pre-existing theorem, explicitly restated in \citet[\S3.1, before Proposition~3.2]{matonevolpato2013determinantal}. That later source is a checked formulation, not the discovery date.}
\GSfield{Present status}{As of November 2026: established in the stated scope; the qualifications above are essential.}
\GSfield{Assessment}{literal statement: false as written; publication match: partial; verification: read the result; development: annotated; confidence: high}
\end{GSitem}
\begin{GSitem}{1985-664.5}{B}{Novikov identifies the Schottky recognition problem}
\GSfield{Kind; source}{question; \GSpage{P001249}; S. P. Novikov}
\GSfield{The notes}{In the margin the Schottky problem is queried; S.~P.~Novikov says that writing the relations is equivalent to it.}
\GSfield{Statement}{Given a symmetric matrix \(\tau\) with positive definite imaginary part and the principally polarized abelian variety \(A_\tau\), characterize when it is the Jacobian of a smooth compact curve. In the KP characterization one requires \(A_\tau\) indecomposable and vectors \(U\ne0,V,W\) such that \(u=-2\partial_x^2\log\theta(Ux+Vy+Wt+Z\mid\tau)\) satisfies \(3u_{yy}=(4u_t+6uu_x-u_{xxx})_x\), as a meromorphic identity for the phase parameter \(Z\).}
\GSfield{Reconstruction}{The recognition question is clearly identified; the KP criterion is later precision, not a formula in this occurrence. Quadratic relations on an already given canonical curve do not by themselves provide equations cutting out the Jacobian locus. The all-abelian-varieties question 1982-592.6 is related but not the same problem.}
\GSfield{Fate}{\citet{shiota1986characterization} published the KP characterization in 1986; its bibliography records separate Kyoto notes circulated in 1983. The original proof was not accessible here. The precise theorem and a later refinement were read in \citet[Introduction and Theorem~1.1]{krichever2006linear}.}
\GSfield{Present status}{As of November 2026: the recognition problem has established KP and auxiliary-linear-equation characterizations. This does not identify an unspecified finite list of the notebook\textquotesingle s relations with a full set of projective Schottky equations.}
\GSfield{Assessment}{literal statement: correct with missing hypotheses; publication match: partial; verification: read the result; development: annotated; confidence: high; problem \hyperref[prob:GP-1985-60]{GP-1985-60}}
\end{GSitem}
\begin{GSitem}{1985-664.6}{E}{Beltrami coordinates are dual to quadratic differentials}
\GSfield{Kind; source}{construction; \GSpage{P001249}}
\GSfield{The notes}{Equation (6) writes \(\int\eta_i f^i dz\wedge d\bar z=\delta_i^j\), followed by \(\rho|dz+y^i\eta_i d\bar z|^2\).}
\GSfield{Statement}{The tangent space to Teichm\"uller space is \(H^1(C,T_C)\), dual by Serre duality to \(H^0(C,K_C^2)\). Choose Beltrami representatives \(\mu_i\) with \(\int_C\mu_i f^j\,d^2z=\delta_i^j\). For small parameters \(y^i\) so that \(\|\sum_i y^i\mu_i\|_\infty<1\), the metric \(\rho|dz+\sum_i y^i\mu_i d\bar z|^2\) is positive and gives a deformed conformal structure. Different representatives modulo \(\bar\partial\)-exact fields give the same infinitesimal moduli tangent.}
\GSfield{Reconstruction}{The superscript in (6) must be \(j\), and finite coordinates require smallness and a local slice; linear representatives need not give a single global coordinate system.}
\GSfield{Fate}{\citet[eqs.~(2.13)--(2.17)]{belavinknizhnik1986complex} gives the explicit duality and local construction.}
\GSfield{Present status}{As of November 2026: established in the stated scope; the qualifications above are essential.}
\GSfield{Assessment}{literal statement: false as written; publication match: partial; verification: read the result; development: annotated; confidence: high}
\end{GSitem}
\begin{GSitem}{1985-664.7}{E}{Elliptic moduli and the modular fundamental domain}
\GSfield{Kind; source}{definition; \GSpage{P001249}}
\GSfield{The notes}{A torus is \(\mathbb C/(\mathbb Z+\tau\mathbb Z)\); changes of cycles act by \(SL_2(\mathbb Z)\), generated by \(\tau\mapsto\tau+1\), \(\tau\mapsto-1/\tau\). The drawn domain is called a hyperbolic triangle of angles \(0,60,60\) degrees.}
\GSfield{Statement}{For pointed complex elliptic curves, \(\tau\in\mathbb H\), and the moduli orbifold is \([\mathbb H/SL_2(\mathbb Z)]\). A closed standard domain is \(\{\tau:|\Re\tau|\le1/2,\ |\tau|\ge1\}\), with identified boundary arcs. Its two finite corner angles are \(\pi/3\), and the cusp has angle zero. The right vertical boundary is \(\Re\tau=+1/2\). The effective action is by \(PSL_2(\mathbb Z)\), while the stack retains the central automorphism.}
\GSfield{Reconstruction}{The word triangle refers to its hyperbolic geometry, not a Euclidean straight-edged triangle. Identification of a genus-one curve with its Jacobian requires a chosen origin.}
\GSfield{Fate}{This is classical elliptic-curve geometry, also used in \citet[\S6, genus-one discussion]{belavinknizhnik1986complex}. Earlier 1968-150.8 explains why a universal family does not exist over the coarse moduli space without extra structure.}
\GSfield{Present status}{As of November 2026: established in the stated scope; the qualifications above are essential.}
\GSfield{Assessment}{literal statement: correct with missing hypotheses; publication match: exact; verification: read the result; development: annotated; confidence: high}
\end{GSitem}
\begin{GSitem}{1985-664.8}{E}{A torus Beltrami deformation and a wedge-sign error}
\GSfield{Kind; source}{construction; \GSpage{P001249-P001250}}
\GSfield{The notes}{With \(y=\Im\tau\), the notes take \(u(z)=z+\eta\bar z/y\), \(\tau\mapsto u(\tau)/u(1)\), then identify \(d\bar\eta\wedge d\eta/y\) with \(d\tau\wedge d\bar\tau/(4y)\).}
\GSfield{Statement}{For \(|\eta|<y\), the deformed period ratio is \(\tau_\eta=(\tau+\eta\bar\tau/y)/(1+\eta/y)\). At \(\eta=0\), \(d\tau=-2i\,d\eta\), hence \(d\bar\eta\wedge d\eta=-\tfrac14d\tau\wedge d\bar\tau\). The claimed identity has the wrong sign with its written wedge orders. These are tangent-coordinate identities at the base point, not globally constant differentials.}
\GSfield{Reconstruction}{\ours\ The computation is ours and is independently checked by a real Jacobian in \texttt{code/check\_mathematics.py}. The letter \(\eta\) replaces the notebook\textquotesingle s \(y^1\) to separate it from \(\Im\tau\).}
\GSfield{Fate}{The underlying Beltrami parametrization is the classical one used in \citet[eq.~(2.14)]{belavinknizhnik1986complex}; the printed wedge identity has no exact match because its sign is wrong.}
\GSfield{Present status}{As of November 2026: established in the stated scope; the qualifications above are essential.}
\GSfield{Assessment}{literal statement: false as written; publication match: partial; verification: read the result; development: annotated; confidence: high}
\end{GSitem}
\begin{GSitem}{1985-664.9}{E}{Gaussian modes, zero modes and the power minus twelve}
\GSfield{Kind; source}{method; \GSpage{P001250}}
\GSfield{The notes}{The 26 scalar fields contribute \((\det^{\prime}\Delta)^{-13}\); the denominator contributes \((\det^{\prime}\Delta)^{-1}\). The zero-mode ratio is called the universe volume divided by the torus area, yielding \((\det^{\prime}\Delta)^{-12}\).}
\GSfield{Statement}{On the orthogonal complement of constants, a normalized real Gaussian with quadratic operator \(\Delta>0\) contributes \((\det^{\prime}\Delta)^{-1/2}\). Twenty-six independent scalars give exponent \(-13\); the appropriate torus ghost determinant contributes exponent \(+1\). Scalar translation zero modes and conformal Killing vector zero modes must be factored separately. Their norms, the reference area and the moduli-coordinate volume must be kept in the same convention.}
\GSfield{Reconstruction}{The elementary nonzero-mode exponents are exact. The symbolic infinite-dimensional Gaussian is defined spectrally, and the notebook\textquotesingle s isolated zero-mode factor is not meaningful without its normalization of area and field norms.}
\GSfield{Fate}{\citet[\S2, especially eq.~(2.19)]{belavinknizhnik1986complex} displays the missing Gram factors explicitly. This implements the already existing Polyakov construction.}
\GSfield{Present status}{As of November 2026: established in the stated scope; the qualifications above are essential.}
\GSfield{Assessment}{literal statement: correct with missing hypotheses; publication match: partial; verification: read the result; development: annotated; confidence: high}
\end{GSitem}
\begin{GSitem}{1985-664.10}{E}{The torus determinant depends on the area convention}
\GSfield{Kind; source}{theorem; \GSpage{P001250-P001251}}
\GSfield{The notes}{The answer is \((\det^{\prime}\Delta)^{-1/2}=y^{-1/2}|\eta(\tau)|^{-2}\), with the Dedekind product.}
\GSfield{Statement}{For \(C_\tau=\mathbb C/(\mathbb Z+\tau\mathbb Z)\), \(y=\Im\tau>0\), and the positive scalar Laplacian of the unit-area metric \(|dz|^2/y\), \(\det^{\prime}\Delta=C\,y|\eta(\tau)|^4\), where \(C\) depends only on the Laplacian convention. For the metric \(|dz|^2\), of area \(y\), the answer is instead \(C\,y^2|\eta(\tau)|^4\). Here \(\eta(\tau)=q^{1/24}\prod_{n\ge1}(1-q^n)\), \(q=e^{2\pi i\tau}\).}
\GSfield{Reconstruction}{\oursComputed\ The formula is correct after unit-area normalization, or when the determinant is divided by area. The notebook\textquotesingle s reference to “the usual” Laplacian does not specify this. Scaling follows from \(\zeta_\Delta(0)=-1\) on a flat torus with its zero mode removed.}
\GSfield{Fate}{The classical determinant evaluation goes back to \citet{raysinger1973complex}; the exact flat-metric convention was read in \citet[\S5.0.1, eq.~(5.1)]{kokotov2007polyhedral}. Original 1973 metadata, not its full proof, was checked.}
\GSfield{Present status}{As of November 2026: established in the stated scope; the qualifications above are essential.}
\GSfield{Assessment}{literal statement: correct with missing hypotheses; publication match: partial; verification: read the result; development: annotated; confidence: high}
\end{GSitem}
\begin{GSitem}{1985-664.11}{E}{The modular invariant one-loop bosonic measure}
\GSfield{Kind; source}{construction; \GSpage{P001250-P001251}}
\GSfield{The notes}{The notebook combines the invariant area element with \((\det^{\prime}\Delta)^{-12}\).}
\GSfield{Statement}{With the preceding unit-area convention and per unit target volume, the resulting density is \(C\,d^2\tau\,y^{-14}|\eta(\tau)|^{-48}=C\,d^2\tau\,y^{-14}|\Delta(\tau)|^{-2}\), with \(\Delta=\eta^{24}\). Under \(\tau\mapsto(a\tau+b)/(c\tau+d)\), the weights of \(d^2\tau\), \(y^{-14}\), and \(|\eta|^{-48}\) are respectively \(-4,28,-24\), whose sum is zero. Orbifold stabilizers and overall constants remain separate conventions.}
\GSfield{Reconstruction}{\oursComputed\ The exponent \(-14\), not \(-13\), includes the moduli measure. This is a density, not yet a convergent value of the vacuum amplitude.}
\GSfield{Fate}{The standard one-loop expression is reproduced by \citet[\S6]{belavinknizhnik1986complex}. It reappears correctly in the later Belavin intervention on P001256, despite the conflicting P001253 formula.}
\GSfield{Present status}{As of November 2026: established in the stated scope; the qualifications above are essential.}
\GSfield{Assessment}{literal statement: correct with missing hypotheses; publication match: exact; verification: read the result; development: annotated; confidence: high}
\end{GSitem}
\begin{GSitem}{1985-664.12}{E}{Metric variation inserts the stress tensor}
\GSfield{Kind; source}{method; \GSpage{P001251}}
\GSfield{The notes}{The variation of \((\det^{\prime}\Delta)^{-1/2}\) is computed by varying the Gaussian. The stress tensor is written with repeated free indices: \(T_{ab}=\partial_a\psi\partial_b\psi-\tfrac12g_{ab}\partial^a\psi\partial^b\psi\).}
\GSfield{Statement}{For \(S[\psi,g]=\int\sqrt g\,g^{ab}\partial_a\psi\partial_b\psi\,d^2\xi\), variation in the inverse metric gives \(\delta S=\int\sqrt g\,T_{ab}\delta g^{ab}\,d^2\xi\), where \(T_{ab}=\partial_a\psi\partial_b\psi-\tfrac12g_{ab}g^{cd}\partial_c\psi\partial_d\psi\). The regulated logarithmic Gaussian variation is the negative expectation of the corresponding action variation, plus the measure and zero-mode terms of the chosen regularization.}
\GSfield{Reconstruction}{\ours\ The density factor and contracted indices are corrected. The notebook\textquotesingle s tentative formulas for \(\delta g^{++},\delta g^{--}\) have insufficient metric conventions to fix their numerical constants.}
\GSfield{Fate}{The stress-tensor variational method is pre-existing in \citet{polyakov1981quantum}; \citet[\S3]{belavinknizhnik1986complex} implements a compatible determinant regularization.}
\GSfield{Present status}{As of November 2026: established in the stated scope; the qualifications above are essential.}
\GSfield{Assessment}{literal statement: false as written; publication match: partial; verification: read the result; development: annotated; confidence: high}
\end{GSitem}
\begin{GSitem}{1985-664.13}{E}{Point splitting, normal ordering and the projective anomaly}
\GSfield{Kind; source}{definition; \GSpage{P001251}}
\GSfield{The notes}{The coincident product \(\partial\psi(z)\partial\psi(z)\) is replaced by a contour-defined normal-ordered product; \(T\) is called an almost holomorphic differential.}
\GSfield{Statement}{In a fixed local coordinate, subtract the universal double-pole singularity from the two-point function before taking the coincident limit. The contour prescription extracts its regular Laurent coefficient. For a chiral scalar with standard normalization, the resulting stress tensor has central charge one and transforms with a Schwarzian term. A difference of projective connections, or a stress tensor after subtracting a chosen reference projective connection, is a genuine quadratic differential.}
\GSfield{Reconstruction}{An operator, its expectation and a tensor density are different objects, although the notebook uses the same \(T\) for them. The anomalous transformation cannot be omitted on a higher-genus surface.}
\GSfield{Fate}{The local conformal-field-theory framework is in \citet{belavinpolyakovzamolodchikov1984infinite}; occurrence 1976-417.18 already distinguishes Schwarzians from quadratic differentials.}
\GSfield{Present status}{As of November 2026: established in the stated scope; the qualifications above are essential.}
\GSfield{Assessment}{literal statement: correct with missing hypotheses; publication match: partial; verification: read the result; development: annotated; confidence: high}
\end{GSitem}
\begin{GSitem}{1985-664.14}{E}{Period normalization and theta characteristics}
\GSfield{Kind; source}{definition; \GSpage{P001252}}
\GSfield{The notes}{The Abel map \(z(Q)=\int_{Q_0}^Q\omega\) and periods \(\int_{a_i}\omega_j=\delta_i^j\), \(\int_{b_i}\omega_j=\tau_{ij}\) are followed by a corrected second theta series and parity \(4\delta\cdot\varepsilon\pmod2\).}
\GSfield{Statement}{For a symplectic homology marking, \(\tau=\tau^T\) and \(\Im\tau>0\). The Abel map is valued in \(\mathbb C^g/(\mathbb Z^g+\tau\mathbb Z^g)\). For half-characteristics \(\delta,\varepsilon\in\tfrac12\mathbb Z^g/\mathbb Z^g\), define \(\theta[\delta,\varepsilon](z|\tau)=\sum_{n\in\mathbb Z^g}\exp\{\pi i(n+\delta)^T\tau(n+\delta)+2\pi i(n+\delta)^T(z+\varepsilon)\}\). Its parity is \((-1)^{4\delta^T\varepsilon}\).}
\GSfield{Reconstruction}{The second displayed series is the coherent one. The first crossed-over exponent is not independently a definition and must not be repaired by silently combining fragments.}
\GSfield{Fate}{Classical material in \citet{fay1973theta}; the full formulation was read in \citet[\S4.1]{matonevolpato2013determinantal}.}
\GSfield{Present status}{As of November 2026: established in the stated scope; the qualifications above are essential.}
\GSfield{Assessment}{literal statement: correct with missing hypotheses; publication match: exact; verification: read the result; development: annotated; confidence: high}
\end{GSitem}
\begin{GSitem}{1985-664.15}{E}{An odd characteristic gives a half-differential only when nonsingular}
\GSfield{Kind; source}{construction; \GSpage{P001252}}
\GSfield{The notes}{For an odd characteristic the differential \(\theta^*=\sum_k\theta_{,k}(0)\omega_k\) is said to admit a square root because all zeros have order two.}
\GSfield{Statement}{Choose a nonsingular odd characteristic \(\nu\), meaning \(\nabla\theta[\nu](0)\ne0\). The nonzero holomorphic differential \(\alpha_\nu=\sum_k\partial_k\theta[\nu](0)\omega_k\) has an even zero divisor and is the square of a section \(h_\nu\) of the corresponding theta-characteristic line bundle \(L_\nu\), \(L_\nu^{\otimes2}\simeq K_C\). Zero multiplicities need only be even, not all exactly two. This is a square root in a line bundle, not a globally single-valued scalar square root.}
\GSfield{Reconstruction}{The nonsingularity condition and the bundle interpretation are essential. An odd theta characteristic with vanishing gradient gives the zero differential, so the later quotients are then undefined.}
\GSfield{Fate}{\citet[\S4.1, immediately before eq.~(4.1)]{matonevolpato2013determinantal} states the nonsingular hypothesis and the half-differential construction, referring to the pre-existing theory in \citet{fay1973theta}.}
\GSfield{Present status}{As of November 2026: established in the stated scope; the qualifications above are essential.}
\GSfield{Assessment}{literal statement: false as written; publication match: partial; verification: read the result; development: annotated; confidence: high}
\end{GSitem}
\begin{GSitem}{1985-664.16}{F}{The incomplete higher-genus stress-tensor formula}
\GSfield{Kind; source}{claim; \GSpage{P001252}}
\GSfield{The notes}{Equation (28) combines \(\tfrac12[\theta^{*\prime\prime}/\theta^*-\tfrac32(\theta^{*\prime}/\theta^*)^2]\), a third theta derivative divided by \(\theta^{\prime}\), and \(-\tfrac\pi2\omega_i(\Im\tau)^{-1}_{ik}\omega_k\), then calls the sum a quadratic differential.}
\GSfield{Statement}{If \(\theta^*=a(z)\) denotes the coefficient of a one-form and \(F^{\prime}=a\), the bracket is the Schwarzian \(\{F,z\}=a^{\prime\prime}/a-\tfrac32(a^{\prime}/a)^2\). The other displayed terms, interpreted using \(a\) in the denominator, transform tensorially; they cannot cancel the Schwarzian inhomogeneity. Hence the literal sum is a projective-connection-type expression, not intrinsically a quadratic differential. Its exact normalization is not recoverable without the definition of \(\theta^{\prime}\), the scalar Green function and the normal-ordering convention.}
\GSfield{Reconstruction}{We identify the mechanism but do not fabricate a corrected coefficient formula. The prime-form/Green-function construction is a precise route to one after these choices are fixed.}
\GSfield{Fate}{\citet{knizhnik1986analytic} subsequently developed analytic fields on Riemann surfaces; publisher metadata was checked, not a formula-by-formula match. The pre-existing theta framework is \citet{fay1973theta}.}
\GSfield{Present status}{As of November 2026: the standard higher-genus stress tensor is established, but the exact handwritten expression cannot be verified in its unspecified conventions. This is an identification limit, not an open problem.}
\GSfield{Assessment}{literal statement: false as written; publication match: thematic only; verification: publisher metadata checked; development: annotated; confidence: medium}
\end{GSitem}

\GStalk{665}{}{}{P001253 (catalogue boundary)}{The catalogue heading: Year 1986}
This is the catalogue heading \ru{Год 1986}, not a separate mathematical talk. The next supplied leaf begins Beilinson\textquotesingle s dated talk on 17 February 1986. No independent mathematical item is assigned to the heading, and its page reference denotes the boundary only. The academic-year item prefix remains 1985.


\GStalk{666}{1986-02-17}{A. A. Beilinson}{P001253-P001257}{Quillen determinants, anomaly cancellation and the Mumford form}
Beilinson recasts the string determinant in the language of determinant lines and their Quillen metrics. A finite-dimensional comparison exposes a dual-line inconsistency in the written norm formula. The anomaly polynomial leads to the combination \(\lambda_2\lambda_1^{-13}\); Belavin intervenes with the holomorphic measure, low-genus theta formulas and divergences. The final uniqueness argument needs an overall normalization. P001258 and P001259 begin separately headed Kontsevich and Morozov talks (catalogue records 667 and 668), and are treated as adjacent evidence in a separate supplement, not assigned to Beilinson or counted as additional records in this chapter.

\begin{GSitem}{1985-666.1}{E}{Genus expansion and the erroneous one-loop q-density}
\GSfield{Kind; source}{construction; \GSpage{P001253}}
\GSfield{The notes}{The partition function is \(Z=\sum_{g\ge0}\int_{M_g}\mu_g\). The genus-one density is written \(|dq/q|^2\Delta(q)^{-2}\log|q|^{-13}\), attributed to Kronecker (1890) and Shapiro (1972).}
\GSfield{Statement}{The genus expansion is a formal perturbative series. For the critical uncompactified bosonic string per unit target volume, the consistent one-loop density is \(C|dq/q|^2|\Delta(q)|^{-2}(-\log|q|)^{-14}\), where \(q=e^{2\pi i\tau}\), \(\Delta(q)=q\prod_{n\ge1}(1-q^n)^{24}\), and constants absorb powers of \(2\pi\). The modular fundamental domain, boundary identifications and automorphism weights remain part of the integral.}
\GSfield{Reconstruction}{The missing modulus of \(\Delta\) and exponent \(-13\) are not compatible with the later formula on P001256. Kronecker\textquotesingle s limit formula is relevant to the determinant calculation, but the specific dates and Shapiro attribution in this sentence were not independently verified and are not treated as an exact publication match.}
\GSfield{Fate}{The corrected measure is given in \citet[\S6]{belavinknizhnik1986complex} and \citet[\S1.6.4]{bost1987determinants}; it is an old one-loop result, not a new measure computed in this talk.}
\GSfield{Present status}{As of November 2026: established in the stated scope; the qualifications above are essential.}
\GSfield{Assessment}{literal statement: false as written; publication match: partial; verification: read the result; development: annotated; confidence: high}
\end{GSitem}
\begin{GSitem}{1985-666.2}{E}{Dolbeault cohomology and the determinant of cohomology}
\GSfield{Kind; source}{definition; \GSpage{P001253}}
\GSfield{The notes}{For a holomorphic line bundle \(L\) on \(X\), \(H^0=\ker\bar\partial\), \(H^1=\operatorname{coker}\bar\partial=\ker\bar\partial^+\), and \(d(X,L)=\det H^0/\det H^1\).}
\GSfield{Statement}{Let \(X\) be a compact connected Riemann surface, with smooth Hermitian metrics on \(X\) and a holomorphic line bundle \(L\). Elliptic Hodge theory identifies the finite-dimensional Dolbeault cohomology with \(\ker\bar\partial_L\) and \(\ker\bar\partial_L^*\). Define \(\lambda(L)=\det H^0(X,L)\otimes\det H^1(X,L)^{-1}\). The empty top exterior power is \(\mathbb C\). The metric on the oriented surface determines its complex structure; it does not alone give a holomorphic structure on an arbitrary smooth line bundle.}
\GSfield{Reconstruction}{The identifications use harmonic representatives and a specified holomorphic structure on \(L\). We fix the symbol \(\lambda\) for the notebook\textquotesingle s convention, which is dual to Quillen\textquotesingle s original determinant line.}
\GSfield{Fate}{The construction predates the talk; \citet[\S\S1--2]{quillen1985determinants} uses the dual convention. \citet[\S2.2]{bost1987determinants} explicitly warns about the two conventions.}
\GSfield{Present status}{As of November 2026: established in the stated scope; the qualifications above are essential.}
\GSfield{Assessment}{literal statement: correct with missing hypotheses; publication match: exact; verification: read the result; development: annotated; confidence: high}
\end{GSitem}
\begin{GSitem}{1985-666.3}{E}{Line bundles of fixed degree form a Picard torsor}
\GSfield{Kind; source}{definition; \GSpage{P001253}}
\GSfield{The notes}{Holomorphic line bundles are described by an integral Chern number and \(g\) continuous parameters, a point on the Jacobian.}
\GSfield{Statement}{For a fixed compact genus-\(g\) complex curve, isomorphism classes of degree-\(d\) holomorphic line bundles form \(\operatorname{Pic}^d(X)\), a complex torus torsor of dimension \(g\) under \(\operatorname{Jac}(X)=\operatorname{Pic}^0(X)\). An identification with the Jacobian requires a reference degree-\(d\) bundle; no canonical such identification exists for arbitrary \(d\).}
\GSfield{Reconstruction}{The parameters are complex, and the torsor qualification is needed. Changing the curve itself would introduce further moduli.}
\GSfield{Fate}{Classical Picard theory precedes the talk. The family determinant/Picard setup is used in \citet[\S2]{quillen1985determinants}; earlier occurrence 1975-385.7 concerns the Picard group and its tangent space.}
\GSfield{Present status}{As of November 2026: established in the stated scope; the qualifications above are essential.}
\GSfield{Assessment}{literal statement: correct with missing hypotheses; publication match: partial; verification: read the result; development: annotated; confidence: high}
\end{GSitem}
\begin{GSitem}{1985-666.4}{E}{Zeta regularization removes zero eigenvalues before differentiating}
\GSfield{Kind; source}{definition; \GSpage{P001254}}
\GSfield{The notes}{For \(\Delta=\bar\partial^*\bar\partial\), \(\zeta_{\Delta^{\prime}}(s)=\sum_{\lambda\ne0}\lambda^{-s}\), \(\Re s>1\), and \(\det\Delta^{\prime}=e^{-\zeta^{\prime}_{\Delta^{\prime}}(0)}\).}
\GSfield{Statement}{For a nonnegative elliptic Laplacian of order two on a smooth closed real two-dimensional manifold, sum the positive eigenvalues with multiplicities. The series converges for \(\Re s>1\) and extends meromorphically to \(\mathbb C\), regularly at zero. Its zeta determinant is \(\det^{\prime}\Delta=\exp[-\zeta_\Delta^{\prime}(0)]>0\). The zero eigenspace is kept separately in the determinant line, not assigned an artificial determinant factor.}
\GSfield{Reconstruction}{Compactness, self-adjointness, smooth coefficients, multiplicities and analytic continuation are made explicit. The prime on \(\Delta^{\prime}\) means restriction to the positive spectrum, whereas the prime on \(\zeta^{\prime}\) means differentiation.}
\GSfield{Fate}{\citet[\S3]{quillen1985determinants}; the heat-kernel justification is spelled out in \citet[\S1.5]{bost1987determinants}.}
\GSfield{Present status}{As of November 2026: established in the stated scope; the qualifications above are essential.}
\GSfield{Assessment}{literal statement: correct; publication match: exact; verification: read the result; development: annotated; confidence: high}
\end{GSitem}
\begin{GSitem}{1985-666.5}{E}{The Quillen norm is inverted on the notebook determinant line}
\GSfield{Kind; source}{construction; \GSpage{P001254}}
\GSfield{The notes}{The Hodge norm on \(d=\det H^0/\det H^1\) is multiplied by \(\sqrt{\det^{\prime}\Delta}\) in equation (7).}
\GSfield{Statement}{With \(\lambda=\det H^0\otimes(\det H^1)^{-1}\), the Quillen norm is \(\|\cdot\|_{Q,\lambda}=\|\cdot\|_{L^2,\lambda}(\det^{\prime}\Delta)^{-1/2}\). Equivalently, on \(\lambda^{-1}\), which is Quillen\textquotesingle s determinant-line convention, \(\|\cdot\|_Q=\|\cdot\|_{L^2}(\det^{\prime}\Delta)^{1/2}\). Squared norms use the corresponding first powers. The corrected norm is smooth in smooth elliptic families, even across jumps of kernel dimension.}
\GSfield{Reconstruction}{\ours\ The source clearly writes multiplication and clearly defines \(d=\lambda\). Its formula is therefore a dual-convention error, not a transcription error. The scalar map \(T=a:\mathbb C\to\mathbb C\) gives an independent check: under the cohomology-to-complex isomorphism, \(1\) maps to \(a^{-1}e/f\) and has norm \(|a|^{-1}\).}
\GSfield{Fate}{\citet[\S\S1,3]{quillen1985determinants}, received 18 June 1984 and published in the first 1985 issue, is the precise pre-existing source. Compare \citet[Theorem~2.1]{bost1987determinants}.}
\GSfield{Present status}{As of November 2026: established in the stated scope; the qualifications above are essential.}
\GSfield{Assessment}{literal statement: false as written; publication match: partial; verification: read the result; development: annotated; confidence: high}
\end{GSitem}
\begin{GSitem}{1985-666.6}{E}{Finite-dimensional determinant comparison and operator variance}
\GSfield{Kind; source}{construction; \GSpage{P001254-P001255}}
\GSfield{The notes}{For \(T:E_1\to E_2\), the cohomological determinant is canonically \(\det E_1/\det E_2\). In the invertible one-dimensional case \(\det T\) is then placed in \(E_1/E_2\).}
\GSfield{Statement}{The determinant-of-complex isomorphism \(\det\ker T\otimes(\det\operatorname{coker}T)^{-1}\simeq\det E_1\otimes(\det E_2)^{-1}\) is canonical up to the standard consistent sign convention. Under it the induced norm is the cohomology norm divided by the product of the nonzero singular values of \(T\). In contrast, \(\det T\) itself belongs to \(\operatorname{Hom}(\det E_1,\det E_2)\), the dual line. For invertible \(T\), the canonical acyclic generator maps to \((\det T)^{-1}\).}
\GSfield{Reconstruction}{\ours\ The reciprocal explains both the misplaced determinant and the preceding Quillen-norm discrepancy. A complete singular-value argument and exact diagonal examples are included in the check appendix.}
\GSfield{Fate}{The finite-dimensional model is already used in \citet[\S1]{quillen1985determinants}. It motivates the analytic construction but does not erase the dependence of the infinite-dimensional metric on the operator.}
\GSfield{Present status}{As of November 2026: established in the stated scope; the qualifications above are essential.}
\GSfield{Assessment}{literal statement: false as written; publication match: partial; verification: read the result; development: annotated; confidence: high}
\end{GSitem}
\begin{GSitem}{1985-666.7}{E}{Operator independence holds only in the fixed finite-dimensional model}
\GSfield{Kind; source}{claim; \GSpage{P001255}}
\GSfield{The notes}{The correct metric is said not to change when \(\bar\partial\) changes. Defining a functorial infinite-dimensional determinant independent of \(\bar\partial\) is called an unattainable ideal.}
\GSfield{Statement}{With fixed finite-dimensional Hermitian spaces, the metric on \(\det E_1\otimes(\det E_2)^{-1}\) is independent of \(T\); the isomorphism from cohomology depends on \(T\). In a family of elliptic operators, the Quillen determinant is a line bundle with generally nonzero curvature, not a fixed one-dimensional space with a flat operator-independent norm. A functorial determinant of every bare infinite-dimensional vector space, under all linear isomorphisms, is not defined by the notes.}
\GSfield{Reconstruction}{The finite-dimensional statement is valid; its unrestricted infinite-dimensional extrapolation is not. The aspirational sentence omits a category and regularization, so it is not promoted to a new problem.}
\GSfield{Fate}{\citet[\S4]{quillen1985determinants} explicitly computes nonzero curvature in operator space. \citet[\S\S2.1--2.5]{bost1987determinants} explains the family construction.}
\GSfield{Present status}{As of November 2026: established in the stated scope; the qualifications above are essential.}
\GSfield{Assessment}{literal statement: correct with missing hypotheses; publication match: partial; verification: read the result; development: annotated; confidence: high}
\end{GSitem}
\begin{GSitem}{1985-666.8}{E}{Metrics on powers of the canonical bundle}
\GSfield{Kind; source}{definition; \GSpage{P001255}}
\GSfield{The notes}{For \(L=\omega^j\), sections are \(f(z)(dz)^j\), and the metric \(\rho dzd\bar z\) gives \(|(dz)^j|^2=\rho^{-j}\), \(j\in\mathbb Z\).}
\GSfield{Statement}{For every integer \(j\), use the dual and tensor-power metric on \(K_X^j\) induced from a smooth positive metric on \(T_X\). In the convention \(\|dz\|^2=\rho^{-1}\), its local norm is exactly \(\|f(dz)^j\|^2=|f|^2\rho^{-j}\). Its Dolbeault determinant line depends on the complex curve, while its Hodge and Quillen metrics also use the chosen fiber metric.}
\GSfield{Reconstruction}{The index labels tensor powers of \(K\), not real differential-form degree. In particular negative \(j\) is meaningful.}
\GSfield{Fate}{This pre-existing tensor construction is the setup of \citet[\S\S2.5--2.6]{bost1987determinants}.}
\GSfield{Present status}{As of November 2026: established in the stated scope; the qualifications above are essential.}
\GSfield{Assessment}{literal statement: correct; publication match: exact; verification: read the result; development: annotated; confidence: high}
\end{GSitem}
\begin{GSitem}{1985-666.9}{E}{The anomaly polynomial and its normalization}
\GSfield{Kind; source}{theorem; \GSpage{P001255}}
\GSfield{The notes}{Polyakov\textquotesingle s theorem is written \(\delta\log|\cdot|_j/\delta\rho=(6j^2-6j+1)R(\rho)\); the polynomial is called the second Bernoulli polynomial.}
\GSfield{Statement}{For the Quillen metrics of \(\lambda_j=\lambda(K_X^j)\) with all bundle metrics induced from the same surface metric, conformal variation has the universal coefficient \(c_j=6j^2-6j+1\). In an explicit relative formulation, for two surface metrics \(\rho_0,\rho_1\), the ratio of Quillen norms on each fixed \(\lambda_j\) is the \(c_j\)-th power of the corresponding ratio on \(\lambda_0\). With the usual convention \(B_2(t)=t^2-t+1/6\), one has \(c_j=6B_2(j)\).}
\GSfield{Reconstruction}{The bare functional derivative omits whether \(\rho\), \(\log\rho\), norms or squared norms are varied, and suppresses area and numerical constants. The relative formulation fixes the invariant content without assigning a false absolute coefficient.}
\GSfield{Fate}{The conformal anomaly is pre-existing; see \citet[eq.~(2.21)]{belavinknizhnik1986complex} and the precise metric comparison in \citet[Proposition~2.3]{bost1987determinants}.}
\GSfield{Present status}{As of November 2026: established in the stated scope; the qualifications above are essential.}
\GSfield{Assessment}{literal statement: correct with missing hypotheses; publication match: partial; verification: read the result; development: annotated; confidence: high}
\end{GSitem}
\begin{GSitem}{1985-666.10}{E}{Cancellation makes the tensor-product norm metric independent}
\GSfield{Kind; source}{theorem; \GSpage{P001255}}
\GSfield{The notes}{For integers \(b_i,j_i\) with \(\sum_i b_i c_{j_i}=0\), the Quillen metric on \(\bigotimes_i d(X,\omega^{j_i})^{b_i}\) depends only on complex structure.}
\GSfield{Statement}{Fix a compact complex curve, use the compatible induced metrics and determinant convention in the previous item, and allow negative tensor powers as duals. The product of relative Quillen-norm changes is \(R_0^{\sum_i b_ic_{j_i}}=1\). Thus the tensor-product norm is independent of the auxiliary smooth surface metric. It may still vary when the complex structure varies.}
\GSfield{Reconstruction}{\ours\ The conclusion follows from the conformal anomaly and must not be confused with an automatic trivialization of the resulting holomorphic line over moduli.}
\GSfield{Fate}{\citet[\S2, eq.~(2.21)]{belavinknizhnik1986complex}; \citet[\S3.1]{bost1987determinants}. The linear cancellation argument is already implicit in the notes.}
\GSfield{Present status}{As of November 2026: established in the stated scope; the qualifications above are essential.}
\GSfield{Assessment}{literal statement: correct; publication match: exact; verification: read the result; development: annotated; confidence: high}
\end{GSitem}
\begin{GSitem}{1985-666.11}{E}{Why thirteen and twenty-six occur}
\GSfield{Kind; source}{remark; \GSpage{P001255}}
\GSfield{The notes}{The important example is \((j_1,b_1)=(2,1)\), \((j_2,b_2)=(1,-13)\); the notes infer no conformal anomaly in dimension 26.}
\GSfield{Statement}{Since \(c_1=1\) and \(c_2=13\), the product \(\lambda_2\otimes\lambda_1^{-13}\) has vanishing anomaly. Each real free scalar contributes one-half of a scalar determinant power, so 26 scalar coordinates supply the exponent 13 that cancels the reparametrization ghosts. This concerns the flat-target critical bosonic model with the stated field content.}
\GSfield{Reconstruction}{\oursComputed\ The dimension is not a general theorem about arbitrary two-dimensional quantum theories. The elementary arithmetic is independently checked in the supplied code.}
\GSfield{Fate}{\citet{polyakov1981quantum}; the geometric relation to the Todd coefficient is explained in \citet[\S2.6]{bost1987determinants}.}
\GSfield{Present status}{As of November 2026: established in the stated scope; the qualifications above are essential.}
\GSfield{Assessment}{literal statement: correct; publication match: exact; verification: read the result; development: annotated; confidence: high}
\end{GSitem}
\begin{GSitem}{1985-666.12}{E}{Determinant lines form a holomorphic bundle even across cohomology jumps}
\GSfield{Kind; source}{construction; \GSpage{P001256}}
\GSfield{The notes}{For a holomorphic family \(X_M\to M\), the spaces \(d(X_t,\omega^j)\) are fibers of a holomorphic line bundle \(d_j\).}
\GSfield{Statement}{Let \(\pi:X\to S\) be a proper holomorphic submersion with compact connected one-dimensional fibers and let \(K_{X/S}\) be its relative canonical bundle. The determinant of the perfect derived direct image \(\lambda_j=\det R\pi_*K_{X/S}^j\) is a holomorphic line bundle, with the stated fiber description and base-change compatibility. The individual cohomology dimensions may jump in general determinant constructions; local finite complexes, not exterior powers of nonexistent cohomology vector bundles, define the line across jumps.}
\GSfield{Reconstruction}{Properness, smoothness and the relative canonical bundle are supplied. The moduli stack or a local/level cover is the correct base for a universal family; the coarse moduli space need not carry one.}
\GSfield{Fate}{Pre-existing determinant functor theory is summarized precisely in \citet[\S2.2]{bost1987determinants}; Quillen\textquotesingle s analytic construction is \citet{quillen1985determinants}.}
\GSfield{Present status}{As of November 2026: established in the stated scope; the qualifications above are essential.}
\GSfield{Assessment}{literal statement: correct with missing hypotheses; publication match: exact; verification: read the result; development: annotated; confidence: high}
\end{GSitem}
\begin{GSitem}{1985-666.13}{B}{Belavin--Knizhnik curvature cancellation}
\GSfield{Kind; source}{theorem; \GSpage{P001256}}
\GSfield{The notes}{Belavin and Knizhnik\textquotesingle s theorem says the canonical metric on \(\bigotimes_i d_{j_i}^{b_i}\) has zero curvature when \(\sum_i b_ic_{j_i}=0\).}
\GSfield{Statement}{For the preceding smooth holomorphic family with compatible Hermitian metrics, the family-index curvature formula in the cohomological determinant convention is \(c_1(\lambda_j,h_Q)=\tfrac{6j^2-6j+1}{12}\,\pi_*[c_1(K_{X/S},h)^2]\). Consequently every indicated anomaly-free tensor product has zero Chern curvature. This is flatness of the Chern connection; global trivial holonomy is a further issue.}
\GSfield{Reconstruction}{\oursComputed\ The normalization uses \(c_1(L,h)=-\tfrac{i}{2\pi}\partial\bar\partial\log h(s,s)\) in a holomorphic frame. The coefficient is independently reproduced by the degree-two expansion of \(e^{jx}(1-x/2+x^2/12)\).}
\GSfield{Fate}{The result was current in the seminar and published in \citet[\S3, eqs.~(3.3)--(3.5)]{belavinknizhnik1986complex}; the later family-index formulation is \citet[Theorem~2.2 and \S3.1]{bost1987determinants}. Class B reflects the named authors of the result, while Beilinson\textquotesingle s role in its geometric formulation is explicitly acknowledged in their paper.}
\GSfield{Present status}{As of November 2026: established in the stated scope; the qualifications above are essential.}
\GSfield{Assessment}{literal statement: correct with missing hypotheses; publication match: exact; verification: read the result; development: annotated; confidence: high}
\end{GSitem}
\begin{GSitem}{1985-666.14}{E}{Zero curvature alone cannot normalize a metric}
\GSfield{Kind; source}{claim; \GSpage{P001256}}
\GSfield{The notes}{Theorem (19) is said to determine the canonical metric uniquely.}
\GSfield{Statement}{If \(h\) is a Hermitian metric of zero curvature, so is \(Ch\) for every \(C>0\). Locally the ratio of any two flat metrics is \(e^{2\Re f}\) for a holomorphic function \(f\). Global constancy of this ratio requires a separate argument on the chosen base; even then a positive scalar remains. Flatness of a line bundle does not itself produce a nowhere-vanishing global holomorphic section.}
\GSfield{Reconstruction}{\ours\ The counterexample by constant rescaling is independent of every analytic convention. Later P001257 narrows the ambiguity but still suppresses the scalar.}
\GSfield{Fate}{\citet[Theorem~3.3]{bost1987determinants} gives the precise statement with a genus-dependent constant, consistent with \citet[\S\S4--5]{belavinknizhnik1986complex}.}
\GSfield{Present status}{As of November 2026: established in the stated scope; the qualifications above are essential.}
\GSfield{Assessment}{literal statement: false as written; publication match: partial; verification: read the result; development: annotated; confidence: high}
\end{GSitem}
\begin{GSitem}{1985-666.15}{A}{The Mumford form is a section, not a global scalar function}
\GSfield{Kind; source}{construction; \GSpage{P001256}}
\GSfield{The notes}{Belavin writes the scalar/ghost determinant ratio as the squared modulus of a holomorphic function on moduli, calling it almost unique.}
\GSfield{Statement}{For \(g\ge2\), the Mumford isomorphism \(\lambda_2\simeq\lambda_1^{13}\) on the stack of smooth curves supplies a nowhere-zero section \(\mu_g\) of \(\lambda_2\lambda_1^{-13}\), unique up to the chosen normalization. Via \(H^0(K^2)\simeq T^*\mathcal M_g\), it is a top differential valued in \(\lambda_1^{-13}\). In a homology marking with period matrix \(\Omega\), the bosonic measure is locally \(C_g|\mu_g|^2/(\det\Im\Omega)^{13}\), interpreted with the Hodge frame and its transformation law.}
\GSfield{Reconstruction}{The local coefficient is holomorphic only in a holomorphic frame. Gram determinants and zero-mode factors are suppressed by the notebook\textquotesingle s schematic ratio. The Mumford isomorphism is stronger global information than curvature zero.}
\GSfield{Fate}{\citet[\S5, theorem after eq.~(5.3)]{belavinknizhnik1986complex} gives the global formulation and credits Beilinson--Drinfeld. Beilinson subsequently coauthored \citet{beilinsonmanin1986mumford} (September 1986); its publisher abstract confirms an explicit Mumford-form formula, but that paper\textquotesingle s full proof was not accessible. Class A records this direct subsequent publication by the speaker.}
\GSfield{Present status}{As of November 2026: established in the stated scope; the qualifications above are essential.}
\GSfield{Assessment}{literal statement: correct with missing hypotheses; publication match: partial; verification: read the result; development: annotated; confidence: high}
\end{GSitem}
\begin{GSitem}{1985-666.16}{B}{The displayed local determinant Hessian suppresses zero modes}
\GSfield{Kind; source}{claim; \GSpage{P001256}; A. A. Belavin}
\GSfield{The notes}{Equations (21)--(21\textquotesingle) write both complex-structure and metric variations of \(\log\det\Delta_j\) as \(c_j\) times universal curvature integrals, with \(f=\rho\eta\).}
\GSfield{Statement}{The invariant statement is the Quillen curvature and conformal-anomaly formula for determinant lines with their cohomology metrics. Raw \(\log\det^{\prime}\Delta_j\) has additional variations from the Gram determinants of harmonic sections and from the area/constant mode. Thus a bare determinant Hessian is not, in general, just \(c_j\) times the displayed local expression. The precise combination in Belavin--Knizhnik (2.21) includes the zero-mode normalizations.}
\GSfield{Reconstruction}{The source does not state whether its determinant notation already contains those factors; we treat this as an ambiguity instead of fixing an arbitrary local regularization.}
\GSfield{Fate}{\citet[eqs.~(2.19)--(2.21), (3.3)]{belavinknizhnik1986complex} is the precise comparator; \citet[Theorem~2.2]{bost1987determinants} gives the coordinate-free version.}
\GSfield{Present status}{As of November 2026: established in the stated scope; the qualifications above are essential.}
\GSfield{Assessment}{literal statement: correct with missing hypotheses; publication match: partial; verification: read the result; development: annotated; confidence: high}
\end{GSitem}
\begin{GSitem}{1985-666.17}{E}{The bosonic moduli integral diverges at the cusp}
\GSfield{Kind; source}{claim; \GSpage{P001256}; A. A. Beilinson and A. A. Belavin}
\GSfield{The notes}{Beilinson says the integral diverges and must be renormalized; Belavin writes the genus-one integral and repeats that it diverges.}
\GSfield{Statement}{For the uncompactified critical bosonic vacuum density, \(|\eta(x+iy)|^{-48}=e^{4\pi y}(1+O(e^{-2\pi y}))\) uniformly for \(|x|\le1/2\) as \(y\to\infty\). Therefore the positive cusp integral \(\int_Y^\infty y^{-14}e^{4\pi y}\,dy\) diverges. This is an infrared degeneration associated with the tachyon, even though the local conformal anomaly vanishes. The mere word renormalization does not specify a canonical finite value.}
\GSfield{Reconstruction}{\ours\ The elementary comparison proves the recorded divergence in the displayed genus-one case; it does not construct a subtraction scheme at all genera.}
\GSfield{Fate}{\citet[\S4]{belavinknizhnik1986complex} studies boundary asymptotics; \citet[\S\S2.3,6.4]{witten2012revisited} distinguishes tachyon and tadpole infrared problems from ultraviolet behavior.}
\GSfield{Present status}{As of November 2026: established in the stated scope; the qualifications above are essential.}
\GSfield{Assessment}{literal statement: correct; publication match: exact; verification: read the result; development: annotated; confidence: high}
\end{GSitem}
\begin{GSitem}{1985-666.18}{F}{The unqualified hope that everything converges for superstrings}
\GSfield{Kind; source}{conjecture; \GSpage{P001256}; A. A. Belavin}
\GSfield{The notes}{\ru{Но есть гипотеза, что в суперструне всё сходится}.}
\GSfield{Statement}{The sentence supplies no string model, vacuum, external states, momenta, genus, integration cycle or distinction between a fixed-genus integral and the all-genus series. It therefore does not define one precise convergence conjecture. Tachyon-free supersymmetric vacua have substantially better behavior, but physical on-shell poles and possible infrared singularities are not excluded by the word superstring.}
\GSfield{Reconstruction}{We do not silently replace this with a universal all-orders theorem or a new narrowly chosen open problem. Conditions eliminating massless tadpoles and handling degeneration contours are additional input.}
\GSfield{Fate}{\citet[\S\S6.4,7--8]{witten2012revisited}, first posted September 2012 and revised through November 2023, develops the integration-cycle and infrared analysis. This is a thematic successor, not a proof that every possible meaning of the notebook sentence converges.}
\GSfield{Present status}{As of November 2026: the recorded slogan has no definite quantifiers, and should not be indexed as an unresolved mathematical problem. The checked modern treatment explicitly retains physical infrared/on-shell qualifications.}
\GSfield{Assessment}{literal statement: not determinable; publication match: thematic only; verification: read the result; development: annotated; confidence: low}
\end{GSitem}
\begin{GSitem}{1985-666.19}{B}{The genus-two measure uses ten even theta constants}
\GSfield{Kind; source}{construction; \GSpage{P001256}; A. A. Belavin}
\GSfield{The notes}{The notes give \(Z_2=\int d^2\tau\,(\Im\tau)^{-3}(\Im\tau)^{-10}|F(\tau)|^{-2}\), \(F=\prod_m\theta_m^2\), naming Knizhnik, Morozov and Perelomov.}
\GSfield{Statement}{On the marked genus-two Siegel space, write \(d^6\Omega=\prod_{1\le i\le j\le2}d\Re\Omega_{ij}\,d\Im\Omega_{ij}\). Up to an overall constant, the bosonic density is \(d^6\Omega\,(\det\Im\Omega)^{-13}|\chi_{10}(\Omega)|^{-2}\), where \(\chi_{10}\) is a nonzero scalar multiple of \(\prod_{m\ \mathrm{even}}\theta_m(0|\Omega)^2\), over the ten even characteristics. Remove the decomposable locus to describe smooth genus-two curves and quotient by the symplectic action with orbifold factors.}
\GSfield{Reconstruction}{The handwritten powers must be powers of a determinant, and its \(d^2\tau\) is shorthand for three complex variables. Including odd characteristics would make the product identically zero. The overall constant and the normalization of \(\chi_{10}\) are not fixed by the notes.}
\GSfield{Fate}{\citet[eqs.~(6.9),(6.14)--(6.17)]{belavinknizhnik1986complex} was read. The named low-genus work appeared as \citet{belavinetal1986loops}; the institutional author bibliography was checked. Belavin is an omitted coauthor in the notebook attribution.}
\GSfield{Present status}{As of November 2026: established in the stated scope; the qualifications above are essential.}
\GSfield{Assessment}{literal statement: correct with missing hypotheses; publication match: partial; verification: read the result; development: annotated; confidence: high}
\end{GSitem}
\begin{GSitem}{1985-666.20}{B}{A genus-three formula exists, but general moduli are not absent}
\GSfield{Kind; source}{remark; \GSpage{P001256}; A. A. Belavin}
\GSfield{The notes}{\(Z_3\) is said also to exist; for general genus there is no good description of moduli space.}
\GSfield{Statement}{The recoverable assertion is that low-genus period-matrix formulas are particularly explicit. In genus three the product \(\chi_{18}\) of the 36 even theta constants enters a bosonic density; its square-root interpretation requires attention to the Torelli marking and hyperelliptic divisor. For arbitrary genus, moduli stacks, compactifications and deformation theory already exist; the difficulty is obtaining comparably simple global coordinates or theta expressions, not existence of moduli.}
\GSfield{Reconstruction}{No exact genus-three formula was written in the notebook, so none is attributed to this occurrence. The phrase “no good description” is retained as a judgment about explicit calculational convenience.}
\GSfield{Fate}{\citet[eqs.~(6.18)--(6.20)]{belavinknizhnik1986complex} supplies the later formula; \citet[\S3.4]{bost1987determinants} discusses its geometric scope.}
\GSfield{Present status}{As of November 2026: established in the stated scope; the qualifications above are essential.}
\GSfield{Assessment}{literal statement: correct with missing hypotheses; publication match: partial; verification: read the result; development: annotated; confidence: medium}
\end{GSitem}
\begin{GSitem}{1985-666.21}{E}{Global constancy on moduli and the remaining scalar ambiguity}
\GSfield{Kind; source}{theorem; \GSpage{P001257}}
\GSfield{The notes}{For \(g>2\), every holomorphic function on moduli is declared constant; flat metrics differ by a holomorphic modulus square, so the metric and section are said to be unique.}
\GSfield{Statement}{On the full complex moduli space \(\mathcal M_g\), \(g\ge3\), global holomorphic functions are constant. More directly for the metric argument, a real pluriharmonic function is constant because any two moduli points can be joined by a complete curve inside moduli. Consequently two globally defined flat metrics on the same holomorphic line differ by a positive constant. A chosen nonzero holomorphic section is likewise unique only up to a nonzero scalar when the quotient is globally holomorphic. Local charts and Teichm\"uller space do have nonconstant holomorphic functions.}
\GSfield{Reconstruction}{The complete-curve argument supplies the missing global hypothesis and avoids falsely assuming that every pluriharmonic function globally has a single-valued holomorphic conjugate. The scalar is still not determined by flatness.}
\GSfield{Fate}{\citet[proof of Theorem~3.3, pp.~141--142]{bost1987determinants} explicitly gives this argument for \(g\ge3\); \citet[Theorem~3.4]{favalepirolatorelli2023forms} proves the corresponding modern statement on the smooth locus for \(g\ge4\). The latter range is not incorrectly substituted for the full-moduli genus-three statement.}
\GSfield{Present status}{As of November 2026: established in the stated scope; the qualifications above are essential.}
\GSfield{Assessment}{literal statement: correct with missing hypotheses; publication match: partial; verification: read the result; development: annotated; confidence: high}
\end{GSitem}

\section{Developments}\label{dev:1985}
\subsection{Simplicial four-polytopes: an explicit contractible family}
\label{sec:gy1985-stacking}

This development belongs to item~1985-658.10 and problem~GP-1985-42.
The notebook on P001230 says only that even simplicial polytopes in
\(\mathbb R^4\) can have \ru{более или менее любой} topology of their
fixed-type realization spaces. The qualification means ``more or less any'';
it does not specify a homeomorphism, homotopy equivalence, stable equivalence,
or the normalization of realizations. An editorial homotopy formulation is:
does every nonempty finite simplicial-complex homotopy type occur as a
normalized realization space of a simplicial four-polytope? The stronger
precise conjecture of \citet[Conjecture~4]{adiprasitopadrol2017universality}
asks for stable equivalence to every nonempty open primary basic
semialgebraic set defined over \(\mathbb Z\), even requiring the four-polytope
to be neighborly. This latter formulation is a literature conjecture, not
an exact quotation from the seminar. Adiprasito--Padrol use a homogenized
realization model: vector configurations realizing the cone over the
polytope, modulo \(\mathrm{GL}_{d+1}(\mathbb R)\). The fixed affine-frame
and prescribed-history model proved contractible below is specified
separately; no identification of these normalization conventions is
assumed.

\paragraph{The established results and their limits.}
Mn\"ev's 1985 announcement already gives configuration universality and a
simplicial-polytope statement with dimension unrestricted
\citep[Theorems~3--4]{mnev1985varieties}. Adiprasito--Padrol explicitly regard
the earlier arguments for the simplicial part as incomplete and supply
the missing proof; the 1985 announcement must therefore not be treated
as equivalent to an independently checked complete proof of that part. The proof of general four-polytope
universality reported in \citet[Main Theorem]{richtergebertziegler1995realization}
concerns work of 1994 and appeared in 1995; it does not impose simpliciality.
\citet[Theorem~2]{adiprasitopadrol2017universality} proves the neighborly
simplicial result in dimension \(2n-4\) with \(2n\) vertices. That paper first
circulated in February 2014, appeared online in May 2016, and was printed
in 2017. Its Conjecture~4 preserves the difficult fixed dimension.
Their discussion after Conjecture~4 also cites established disconnected
realization spaces of simplicial four-polytopes, so nontrivial topology
in dimension four is already known; universality is the remaining
quantifier. The same paper records trivial realization spaces for stacked polytopes
in the discussion after Conjecture~4; its Theorem~3 proves a ball result
for specified sewing families. The proof below supplies an explicit
continuous insertion and contraction for a narrowly specified stacking
space. It is not a new universality theorem or a new discovery of
contractibility for stacked polytopes.

As of November 2026, no resolution of fixed-dimensional simplicial
universality was found in the bounded search for this chapter. The latest
opened source explicitly retaining the constant-dimension simplicial
question is the 2024 compendium, entries CG68--CG69
\citep{schaefercardinalmiltzow2024existential}. This is dated evidence, not a
claim to have excluded every subsequent publication. A later 2026 search
lead concerning noncontractible realization spaces is recorded as
unverified in the research log and is not used in the argument.

\paragraph{The space to which the proof applies.}
Fix \(d\geq2\), and fix the labelled simplex
\(P_0=\operatorname{conv}(0,e_1,\ldots,e_d)\subset\mathbb R^d\).
A \emph{prescribed stacking history} \(\mathcal H\) of length \(r\) specifies,
at step \(k\), a labelled simplex facet \(F_k\) of the combinatorial polytope
produced by the previous steps. The next labelled point must lie strictly
beyond \(F_k\) and strictly beneath every other current facet.
Let \(\mathcal R_0\) be the one-point space consisting of the fixed simplex.
Recursively, \(\mathcal R_k\) consists of the coordinates of the first \(k\)
inserted points satisfying every one of these conditions at its own step.
These are \emph{history-compatible} realizations. Coordinates of the initial
simplex remain fixed, so no affine group factor is present. We do not
identify \(\mathcal R_r\) with the full realization space of the final
combinatorial type without a separate deletion argument.

For a configuration \(q\in\mathcal R_{k-1}\), write \(P(q)\) for its convex
hull. For each facet \(G\), choose an inward affine supporting form
\[
  \ell_{G,q}(x)=a_G(q)\cdot x+b_G(q),\qquad
  \|a_G(q)\|=1,
\]
which vanishes on \(G\) and is positive on the interior of \(P(q)\).
These choices are continuous in \(q\). Indeed, the labelled vertices of a
simplex facet give a nonzero affine cofactor vector. Divide its spatial
part by its norm and choose its sign to be positive at the average of all
vertices of \(P(q)\). That average is in the interior; the sign test never
vanishes, so it is continuous, including if the parameter space were
initially disconnected. This construction avoids an unproved global
choice of a facet normal.

\begin{lemma}[The stacking insertion region \ours]
For a full-dimensional simplicial polytope \(P\) and a facet \(F\), set
\[
 U_F(P)=\{x\in\mathbb R^d:\ell_F(x)<0,
                  \ \ell_G(x)>0\text{ for every facet }G\ne F\}.
\]
This is a nonempty open convex set. For every \(x\in U_F(P)\),
\(\operatorname{conv}(P\cup\{x\})\) has the stacking combinatorial type:
it retains all facets except \(F\), and replaces \(F\) by the simplices
\(\operatorname{conv}(R\cup\{x\})\), one for each ridge \(R\) of \(F\).
\end{lemma}
\begin{proof}
The inequalities define an intersection of finitely many open half-spaces,
hence an open convex set. Let \(b\) be the barycenter of the vertices of
\(F\), and let \(n_F=-a_F\) be the outward unit normal. Since \(b\) lies in
the relative interior of \(F\), \(\ell_G(b)>0\) for every other facet.
For sufficiently small \(\varepsilon>0\), the point
\(b+\varepsilon n_F\) therefore satisfies all the required inequalities;
nonemptiness is established quantitatively in the next proposition.

For the combinatorial assertion, a segment from \(x\) to any point of \(P\)
meets the plane \(\ell_F=0\) at a point of \(F\), unless it is wholly on the
\(P\) side. At the crossing, all other facet inequalities hold because they
hold at both endpoints. Consequently
\[
 \operatorname{conv}(P\cup\{x\})
       =P\cup\operatorname{conv}(F\cup\{x\}).
\]
Every old facet \(G\ne F\) remains a supporting facet, since
\(\ell_G(x)>0\). If a ridge \(R\subset F\) is shared by \(F\) and \(G\),
then
\[
 h_R(y)=\ell_G(x)\ell_F(y)-\ell_F(x)\ell_G(y)
\]
is nonnegative on \(P\), vanishes precisely on \(R\) there, and also
vanishes at \(x\). It exposes exactly
\(\operatorname{conv}(R\cup\{x\})\), which is a simplex facet. The boundary
of the union displayed above consists of the old boundary without the
relative interior of \(F\), and these lateral faces of the added pyramid.
Thus the listed facets exhaust the new boundary and give the stated type.
\end{proof}

\begin{proposition}[A continuous section and a fiberwise contraction \ours]
The map
\[
 \pi_k:\mathcal R_k\longrightarrow\mathcal R_{k-1}
\]
forgetting the last vertex is surjective and has a continuous section.
Moreover, \(\mathcal R_k\) strongly deformation retracts onto the graph of
that section while every previous vertex is held fixed.
\end{proposition}
\begin{proof}
Let \(F=F_k\). Its facet barycenter \(b(q)\) and outward unit normal
\(n_F(q)=-a_F(q)\) are continuous. Define
\[
 \varepsilon(q)=\frac12\min\left(
       1,\ \min_{G\ne F}
       \frac{\ell_{G,q}(b(q))}{1+|a_G(q)\cdot n_F(q)|}
                         \right)>0,
 \qquad s(q)=b(q)+\varepsilon(q)n_F(q).
\]
The minimum is over the fixed finite labelled facet set at this step, so
\(\varepsilon\) and \(s\) are continuous. For the selected facet,
\(\ell_{F,q}(s(q))=-\varepsilon(q)<0\). For every other facet, write
\(u=|a_G(q)\cdot n_F(q)|\). Then
\[
 \ell_{G,q}(s(q))
 \geq \ell_{G,q}(b(q))-\varepsilon(q)u
 \geq \ell_{G,q}(b(q))
             \left(1-\frac{u}{2(1+u)}\right)>0.
\]
Thus \(s(q)\in U_F(P(q))\), proving both surjectivity and the section
\(q\mapsto(q,s(q))\).

The fiber over \(q\) is exactly \(U_F(P(q))\). Its convexity proves that
\[
 H_k((q,x),t)=\bigl(q,(1-t)x+t\,s(q)\bigr),\qquad 0\leq t\leq1,
\]
remains in \(\mathcal R_k\). This map is jointly continuous, equals the
identity at \(t=0\), has image in the section at \(t=1\), and fixes every
point of the section for every \(t\). These are precisely the strong
deformation-retraction conditions. No claim that contractible fibers
alone imply a homotopy equivalence is needed.
\end{proof}

\begin{theorem}[Contractibility for a prescribed stacking history \ours]
For every finite prescribed facet-stacking history on the fixed simplex,
the normalized history-compatible realization space \(\mathcal R_r\) is
nonempty and contractible. In particular this holds in dimension four.
\end{theorem}
\begin{proof}
The space \(\mathcal R_0\) is a point. At each step the section proves
nonemptiness, and the preceding proposition retracts \(\mathcal R_k\)
onto a copy of \(\mathcal R_{k-1}\). If a contraction of
\(\mathcal R_{k-1}\) is already given, first use \(H_k\) to move to the
section and then apply the contraction there through the continuous
section map. Concatenating these two homotopies contracts
\(\mathcal R_k\). Induction gives a finite, continuous contraction of
\(\mathcal R_r\). The conclusion concerns exactly the recursively defined
space: no assertion that every realization of the final type admits
this history has entered the proof.
\end{proof}

\paragraph{A second check of the section estimate.}\ours\
The argument can be rederived without unit normals, which is useful for
exact arithmetic. Take inward forms with nonzero gradients \(a_G\), put
\[
 \eta=\min\left(1,\ \min_{G\ne F}
          \frac{\ell_G(b)}{2(1+|a_G\cdot a_F|)}\right),
 \qquad x=b-\eta a_F.
\]
Then \(\ell_F(x)=-\eta\|a_F\|^2<0\), while
\(\ell_G(x)\geq\ell_G(b)-\eta|a_G\cdot a_F|>0\).
This independently checks the sign and the strict-slack bound. If the
vertices and inward forms are rational, \(b,\eta,x\) are rational; no
square root is required. These identities check the proof for every
admissible step, independently of the finite examples computed elsewhere
in this chapter. Numerical or exact examples do not establish equality
with the full realization space.

\begin{proposition}[Rational points in simplicial realization spaces \ours]
Every nonempty realization space of a labelled simplicial polytope
contains a rational realization. The same holds after normalizing a
specified affinely independent vertex frame to a rational simplex, on
any nonempty such chart.
\end{proposition}
\begin{proof}
Start with a labelled simplicial realization. Each facet has exactly
\(d\) affinely independent vertices, and its supporting affine form is
strictly positive on every other vertex. These are finitely many open
conditions on the vertex coordinates: affine independence and the signs
of the appropriate determinants. They continue to hold in a sufficiently
small neighbourhood. Every listed simplex facet therefore persists.

For completeness, this persistence rules out extra facets. Every ridge
of a listed simplex facet belongs to exactly two listed facets in the
original simplicial boundary complex. In the perturbed convex hull those
two facets still meet along this ridge, and a ridge of a convex polytope
belongs to exactly two facets. Hence no facet adjacent across a ridge to
a listed facet can be unlisted. The facet-adjacency graph of a convex
polytope is connected, so every facet is listed. Each old labelled vertex
remains a vertex because it is a vertex of a persisting simplex facet.
Thus the combinatorial type is constant throughout that neighbourhood.

Rational points are dense in the coordinate space, so choose all vertex
coordinates rational within this neighbourhood. For a normalized chart,
first apply the unique affine map sending the chosen vertex frame to the
specified rational frame. Hold those frame coordinates fixed and perturb
the remaining coordinates rationally within the same finite strict
inequalities. This gives the normalized rational realization as well.
\end{proof}

\begin{corollary}[The arithmetic limit of a simplicial universality claim \ours]
Stable equivalence over \(\mathbb Q\), generated by rational equivalences
and stable projections with nonempty relatively open polyhedral fibers
defined over the base, cannot identify a nonempty simplicial-polytope
realization space with every primary basic semialgebraic set over
\(\mathbb Z\) without an openness or comparable arithmetic restriction.
\end{corollary}
\begin{proof}
Existence of a rational point is preserved by rational equivalences
whose maps and inverses are defined on the spaces. It is also preserved
in both directions by a stable projection. One direction follows by
projection. For the other, over a rational base point the fiber is a
nonempty relatively open polyhedron with rational linear equations and
strict rational inequalities. Its affine hull has a rational point and
a rational direction basis; density in that affine hull yields a rational
point satisfying all strict inequalities.

The primary basic set
\(S=\{t\in\mathbb R:t^2-2=0,\ t>0\}\) is nonempty and has no rational
point. The proposition supplies a rational point for every nonempty
simplicial realization space, so such a space cannot be stably equivalent
to \(S\) in this sense. By contrast, a nonempty open subset of Euclidean
space always has rational points. This explains why the established and
conjectural simplicial universality statements restrict to open primary
sets. It imposes no corresponding obstruction on homotopy universality:
this particular \(S\) already has the homotopy type of a point.
\end{proof}

\paragraph{The remaining obstruction and a bounded next step.}
A fixed stacking history admits a global section and convex insertion
fibers, so it cannot encode arbitrary topology. The unresolved universal
claim needs mechanisms beyond this restricted construction. In the
known approaches, recording complicated order types inside auxiliary
configurations does not by itself ensure that a fixed simplicial
four-dimensional face lattice forces those order types. The discussion
after Conjecture~4 in \citet{adiprasitopadrol2017universality} identifies
precisely this distinction using stacked vertex figures and contractions.
The present proof does not address that enforcement problem.

A further bounded run could take one explicitly labelled small final
stacked four-polytope and determine, by exact inequalities, whether every
realization in the chosen normalization permits a specified last-vertex
deletion. A proof would identify its full realization space with the
history-compatible space for that case; a counterexample would delimit
the identification. This is a concrete extension of the certificate,
not a proposal to settle fixed-dimensional universality by checking a
finite list of examples.

\subsection{Independent mathematical checks of editorial emendations}
\label{checks:1985}

This subsection checks corrections to established mathematics; it is not a second
research development selected from class C or D. The original notebook formulas
remain in the occurrence index. The labels \ours\ and \oursComputed\ mean that
the proof or computation was supplied for this companion, not that the underlying
mathematics is new. Exact code and output are in
\texttt{code/check\_mathematics.py} and its JSON output. A coefficient-algebra
implementation and direct application to functions are separate checks.

\begin{proposition}[The corrected Novikov--Veselov commutator]\ours\
Let \(D=\partial_z\), \(E=\partial_{\bar z}\), and let the smooth coefficients
\(V,u,v\) satisfy \(Eu=3DV\), \(Dv=3EV\). Set
\[
 L=DE+V,\qquad P=D^3+E^3+uD+vE,\qquad B=Du+Ev.
\]
Here a coefficient without a derivative to its right denotes multiplication.
Then
\[
 [P,L]+BL=M_{D^3V+E^3V+D(uV)+E(vV)}.
\]
For real \(V\), one may put \(v=\bar u\). Consequently the right hand side
defines the stated fixed-energy evolution, with the boundary or normalization
conditions on \(u\) still to be specified.
\end{proposition}
\begin{proof}
Leibniz's rule gives the following coefficients, before imposing the auxiliary
equations:
\[
\begin{array}{c|l}
\text{derivative} & \text{coefficient in }[P,L]+BL\\\hline
D^2&3DV-Eu\\
E^2&3EV-Dv\\
DE&B-Du-Ev\\
D&3D^2V-DEu\\
E&3E^2V-DEv\\
1&D^3V+E^3V+uDV+vEV+BV.
\end{array}
\]
All positive-order coefficients vanish; the scalar coefficient is the one
claimed. This calculation also fixes the meaning of the notebook expression
\((D u)V\): it must be \(D(uV)\), not \((Du)V\) alone. Independently, the code
applies the two operators successively to two different test polynomials, for
each of two polynomial potentials, constructing \(u,v\) by integration of the
auxiliary equations. Every residual is exactly zero. This second check tests the
implementation; the displayed algebra proves the identity for all smooth
coefficients. The hierarchy itself is the prior result of
\citet[\S II, equations (4)--(6)]{veselovnovikov1984explicit}.
\end{proof}

\begin{proposition}[An explicit counterexample to the unrestricted no-Lax wording]\ours\
There is a smooth real potential \(U(x,t)\), independent of \(y\), for which
\[
 L=-\partial_x^2-\partial_y^2+U,\qquad
 A=-4\partial_x^3+6U\partial_x+3U_x
\]
satisfy \(L_t=[A,L]\), but the potentials do not evolve by Euclidean motions.
The example is not a genuinely two-dimensional decaying scattering potential;
that extra hypothesis is precisely absent from the universal notebook claim.
\end{proposition}
\begin{proof}
Put
\[
 r=e^{2x-8t},\quad s=e^{4x-64t},\quad
 T=1+r+s+\tfrac19rs,\quad U=-2\partial_x^2\log T.
\]
Since \(T>0\) for real \(x,t\), the formula is smooth and real for all such
parameters. To verify the equation without relying on a remembered soliton
formula, use the four exponential terms of \(T\), with pairs
\((k,\omega)=(0,0),(2,-8),(4,-64),(6,-72)\) and coefficients
\(1,1,1,1/9\). In the bilinear expression
\((D_x^4+D_xD_t)T\mathbin{\cdot}T\), each ordered pair contributes
\((k_i-k_j)^4+(k_i-k_j)(\omega_i-\omega_j)\).
Every unordered pair gives zero except \((0,3)\) and \((1,2)\). They have the
same exponential \(rs\) and weighted coefficients \(864/9=96\) and \(-96\),
respectively, so they cancel. Writing \(q=\log T\), the bilinear identity is
\(q_{xt}+q_{xxxx}+6q_{xx}^2=0\); differentiation gives
\(U_t=6UU_x-U_{xxx}\).

A direct Leibniz expansion gives
\([A,L]=M_{6UU_x-U_{xxx}}\), since \(A\) commutes with
\(\partial_y^2\). Thus the Lax equation holds. To check that this is not a
time-dependent translation, compare \(U_t/U_x\) at \(t=0\), first with
\((r,s)=(1,1)\), then with \((r,s)=(2,4)\). The exact ratios are $-92/3$ and $-89524/4861$, respectively, and both denominators $U_x$ are nonzero. These values were also obtained independently from moment and cumulant formulas for the four exponential terms. A translation would give a
ratio depending on \(t\) alone. A continuous Euclidean motion of a nonconstant
potential independent of \(y\), remaining independent of \(y\), cannot rotate
its direction of variation, so it reduces to this excluded translation case.
The code independently verifies the rational PDE identity in \(r,s\) and the
formal commutator by two implementations. This counterexample does not assert
that the intended genuinely two-dimensional obstruction has been identified.
\end{proof}

\begin{lemma}[Leading Baker--Akhiezer coefficients]\ours\
Use the notebook's exponential conventions and suppose locally that
\[
 \psi=e^{ik_1z}(1+\xi_1/k_1+O(k_1^{-2})),\qquad
 \psi=c_2e^{ik_2\bar z}(1+O(k_2^{-1})),\quad c_2\ne0.
\]
For an operator \(L=DE+aE+b\) annihilating \(\psi\), cancellation of the
leading terms requires
\(a=-D\log c_2\), \(b=-iE\xi_1\).
\end{lemma}
\begin{proof}
After dividing by \(e^{ik_1z}\), the constant term of \(DE\psi\) is
\(iE\xi_1\); the term \(aE\psi\) begins at order \(k_1^{-1}\).
At the other puncture the coefficient of \(ik_2e^{ik_2\bar z}\) is
\(Dc_2+ac_2\). These two coefficients give the formulas. The recorded
operator \(DE+A D+B\) instead has an uncancelled \(ik_1 A\) term at the first
puncture. Symbolic differentiation of two truncated Laurent expansions gives
the same cancellations independently. Existence and uniqueness from spectral
data are supplied by the cited theorem of
\citet[\S3, Lemmas 1--2]{dubrovinkrichevernovikov1976schrodinger}; this lemma
only determines the necessary coefficients in the convention used here.
\end{proof}

\begin{proposition}[The hyperelliptic defect in products of differentials]\ours\
Let \(X\) be a smooth compact connected complex hyperelliptic curve of genus \(g\ge2\).
The multiplication map
\(\operatorname{Sym}^2H^0(X,K_X)\to H^0(X,K_X^2)\) has image dimension
\(2g-1\) and cokernel dimension \(g-2\).
\end{proposition}
\begin{proof}
Choose a coordinate on the quotient so that infinity is unbranched, and write
\(y^2=f(x)\) with \(f\) squarefree of degree \(2g+2\). The one-forms
\(x^i dx/y\), \(0\le i\le g-1\), are a basis, by local regularity and
Riemann--Roch. Their products span the independent forms
\[
 x^k\frac{dx^2}{y^2},\quad 0\le k\le 2g-2.
\]
The complementary forms are
\[
 x^j\frac{dx^2}{y},\quad 0\le j\le g-3,
\]
with the second list empty when \(g=2\). At a branch point use
\(x-a=t^2\), \(y=t\) times a nonvanishing analytic factor: the two basic
quadratic differentials have orders zero and one, respectively. At either
point at infinity put \(s=1/x\). Their orders are \(2g-2-k\) and
\(g-3-j\), so all listed forms are holomorphic. The two families have opposite
parity under the hyperelliptic involution; inside each family their rational
monomials are independent. There are \(3g-3=h^0(K_X^2)\) forms in total,
so they form a basis. This proves both dimensions for every \(g\ge2\).
Independently the code checks the orders, product exponents and dimensions for
\(g=2,\ldots,9\), explicitly including the exceptional zero-cokernel case
\(g=2\). The general nonhyperelliptic surjectivity statement is Max Noether's
theorem, recalled in \citet[\S3.1]{matonevolpato2013determinantal}.
\end{proof}

\begin{lemma}[The sign of the torus Beltrami wedge]\ours\
Write \(y=\Im\tau>0\), and use the deformation
\(\tau_\eta=(\tau+\eta\bar\tau/y)/(1+\eta/y)\). At \(\eta=0\),
\[
 d\tau=-2i\,d\eta,\qquad
 d\bar\eta\wedge d\eta=-\tfrac14d\tau\wedge d\bar\tau.
\]
\end{lemma}
\begin{proof}
Differentiate the quotient with the base \(\tau\) fixed; the derivative is
\((\bar\tau-\tau)/y=-2i\). Conjugation and antisymmetry give the wedge
identity. For an independent real-coordinate check, \(\eta=p+iq\) gives
\((\delta\Re\tau,\delta\Im\tau)=(2q,-2p)\), with Jacobian determinant four.
The coordinate identities \(d\bar\eta\wedge d\eta=2i\,dp\wedge dq\) and
\(d\tau\wedge d\bar\tau=-2i\,d\Re\tau\wedge d\Im\tau\) give exactly
the same minus sign. These are tangent identities at the base point.
\end{proof}

\begin{proposition}[The finite-dimensional determinant metric fixes the reciprocal]\ours\
For a linear map \(T:E_1\to E_2\) of finite-dimensional Hermitian spaces, put
\(\lambda_T=\det\ker T\otimes(\det\operatorname{coker}T)^{-1}\).
Under the determinant-of-complex isomorphism
\(\lambda_T\simeq\det E_1\otimes(\det E_2)^{-1}\), the norm induced from the
fixed spaces is the cohomology norm divided by the product of the nonzero
singular values of \(T\).
\end{proposition}
\begin{proof}
Take orthonormal bases of \(\ker T\), its orthogonal complement, the image,
and its orthogonal complement, with \(Te_i=\sigma_i f_i\), \(\sigma_i>0\).
Complete the kernel basis by the \(e_i\), and the cokernel basis by the
\(Te_i\), in a fixed determinant sign convention. The corresponding element
of \(\det E_1\otimes(\det E_2)^{-1}\) contains the scalar
\(\prod_i\sigma_i^{-1}\). The remaining bases are orthonormal, which proves
the formula; signs have modulus one. If \(T\) is invertible, the canonical
acyclic generator therefore maps to \((\det T)^{-1}\), whereas \(\det T\)
belongs to the dual line. Empty products give the rank-zero case.

The positive eigenvalues of \(T^*T\) are \(\sigma_i^2\), so the squared norm
factor is \((\det{}'T^*T)^{-1}\). Two unrelated rational diagonal examples
and the empty example are checked in the code. Replacing the product by a
zeta determinant yields the direction of the spectral factor on the
notebook's line; the analytic existence and smoothness theorem remains
\citet[\S3]{quillen1985determinants}, who uses its dual.
\end{proof}

\begin{example}[Anomaly coefficients and modular-weight check]\oursComputed\
For \(x=c_1(K_{X/S})\), the degree-two term of
\[
 e^{jx}\left(1-\frac{x}{2}+\frac{x^2}{12}\right)
\]
is \((6j^2-6j+1)x^2/12\). This is also obtained by multiplying the three
terms by hand: \(j^2/2-j/2+1/12\). Thus \(c_1=1\), \(c_2=13\), and
\(c_2-13c_1=0\). This is only the coefficient calculation; the analytic
family-index theorem is cited in the occurrence index.

For a modular substitution, the genus-one factors
\(d^2\tau\), \((\Im\tau)^{-14}\), and \(|\eta(\tau)|^{-48}\) have
absolute weights \(-4,28,-24\), whose sum is zero. The corresponding
\(q\)-density has \((-\log|q|)^{-14}\). Separately, the torus formula
\(C y^2|\eta(\tau)|^4\) for \(|dz|^2\), read in
\citet[\S5.0.1, equation (5.1)]{kokotov2007polyhedral}, becomes
\(C y|\eta(\tau)|^4\) in the unit-area metric: its Laplacian is multiplied
by \(y\), and \(\zeta_\Delta(0)=-1\). The latter follows from the flat
two-dimensional heat expansion after subtracting the one-dimensional
constant eigenspace. Area conventions cannot be changed in only one factor
of the measure.
\end{example}

\paragraph{Scope of the second checks.}
The Novikov--Veselov and KdV identities were rederived in coefficient form and
by direct application. The KdV example also has a separate bilinear proof.
The determinant reciprocal was checked by both a one-dimensional acyclic map
and the singular-value argument. The Beltrami sign was checked in complex and
real coordinates. The hyperelliptic dimension result was checked by local
orders and Riemann--Roch, with finite examples used only to catch boundary
mistakes. The stacking development has an independent exact facet-enumeration
check and a separate proof audit recorded with its source files. No finite
computation is presented as a proof of an unbounded universality assertion.

\subsection{Adjacent-page supplement: Kontsevich and Morozov}
\label{supp:1985-adjacent}
These two pages are preserved as supplementary evidence only. Their records
667--668 lie outside the requested core records 641--666; they are excluded
from the year's core talk and occurrence counts and receive ADJ-prefixed
identifiers. The manuscript heading on P001258 reads \(24/\mathrm{II}/85\),
whereas catalogue record 667 gives 24 February 1987. P001259 reads
24 February 1986, whereas record 668 also gives 24 February 1987.
Neither conflict is silently repaired. Kontsevich's matching paper was
submitted 3 December 1986 and published in 1987; Morozov explicitly reports
Belavin--Knizhnik and Beilinson rather than claiming new results. The technical
assessment below is independent of resolving these dates.

\subsubsection{ADJ-667.1: A curve with a point and a formal parameter is not unmarked Teichmueller space}
\noindent\textbf{Source.} P001258; M. L. Kontsevich.\par
\noindent\textbf{Notebook.} Kontsevich's page calls \(M_g\) Teichm\"uller space and writes \(\widetilde M_g=(C,p,x)\), where \(x\) is a coordinate at \(p\).\par

\noindent\textbf{Precise statement.} For \(g\ge2\), let \(\widehat{\mathcal M}_{g,1}\) parameterize smooth connected projective complex curves \(C\) of genus \(g\), a point \(p\), and a formal parameter \(x\in\widehat{\mathfrak m}_p\) with nonzero first derivative. It is the inverse limit of the finite-jet moduli spaces. Forgetting \(x\) gives pointed moduli \(\mathcal M_{g,1}\); forgetting both \(p,x\) gives \(\mathcal M_g\). A topological marking would instead place these constructions over Teichm\"uller space.\par

\noindent\textbf{Reconstruction and limits.} The triple records no topological marking. We use a moduli stack or a local smooth cover rather than silently identify it with Teichm\"uller space. A formal parameter need not be a convergent analytic coordinate on a fixed disk.\par

\noindent\textbf{Literature fate.} \citet[\S1]{kontsevich1987virasoro} explicitly constructs these finite-jet moduli and their inverse limit. The paper was submitted 3 December 1986 and published in issue 2 of 1987.\par

\noindent\textbf{Present status.} As of November 2026: The objects are identifiable standard deformation-theoretic data (E). The notebook's two uses of the word Teichm\"uller require this distinction.\par

\noindent\textbf{Assessment.} Class E; literal correctness: correct with missing hypotheses; publication match: exact; verification: read the result; development: none; confidence: high.\par

\subsubsection{ADJ-667.2: Formal local vector fields and the curve-dependent global subalgebra}
\noindent\textbf{Source.} P001258; M. L. Kontsevich.\par
\noindent\textbf{Notebook.} The page defines \(\mathfrak g=\mathbb C((x))\partial_x\), \(\mathfrak k=\mathbb C[[x]]\partial_x\), and \(\gamma\) as Laurent expansions of vector fields on \(C\) with a pole only at \(p\).\par

\noindent\textbf{Precise statement.} Set \(\gamma_{C,p,x}=H^0(C,T_C(*p))\subset\mathbb C((x))\partial_x\) by Laurent expansion. Its elements are algebraic vector fields on \(C\setminus\{p\}\), meromorphic on \(C\) with finite pole order at \(p\). The local subalgebra \(\mathfrak k\) consists of fields regular on the formal disk, while \(\mathfrak k_0=x\mathbb C[[x]]\partial_x\) consists of fields fixing its marked origin.\par

\noindent\textbf{Reconstruction and limits.} The space \(\gamma\) depends on the curve and parameter; \(\mathfrak k\) does not. The marginal phrase \ru{всю C!!} cannot mean holomorphic extension through \(p\): for \(g\ge2\), \(H^0(C,T_C)=0\), whereas meromorphic fields with allowed pole order are essential here. Editorial check of the marginal whole-curve interpretation \ours.\par

\noindent\textbf{Literature fate.} These subalgebras and the global punctured-curve condition are defined in \citet[\S1]{kontsevich1987virasoro}.\par

\noindent\textbf{Present status.} As of November 2026: The finite-pole global subalgebra is exactly identified (E). The ambiguous marginal instruction is retained as provenance rather than imposed as a false extra condition.\par

\noindent\textbf{Assessment.} Class E; literal correctness: correct with missing hypotheses; publication match: exact; verification: read the result; development: annotated; confidence: high.\par

\subsubsection{ADJ-667.3: Virasoro uniformization: the infinitesimal action and stabilizer}
\noindent\textbf{Source.} P001258; M. L. Kontsevich.\par
\noindent\textbf{Notebook.} Theorem1 says \(\mathfrak g\) acts on \(\widetilde M_g\), with stabilizer \(\gamma\), so it is locally the homogeneous space \(\mathfrak g/\gamma\).\par

\noindent\textbf{Precise statement.} For the moduli of smooth genus-\(g\ge2\) pointed curves with a formal parameter, the Lie algebra \(\operatorname{Der}\mathbb C((x))\) acts infinitesimally. At \((C,p,x)\), evaluation of that action is surjective and its kernel is \(H^0(C,T_C(*p))\). Thus \(T_{(C,p,x)}\widehat{\mathcal M}_{g,1}\cong\mathbb C((x))\partial_x/\gamma_{C,p,x}\), in the formal inverse-limit sense.\par

\noindent\textbf{Reconstruction and limits.} This is an infinitesimal uniformization theorem. It does not assert the existence of an ordinary global Lie group, a single fixed subgroup for all curves, or a global homogeneous-space presentation. The completed formal setting and genus hypothesis are supplied.\par

\noindent\textbf{Literature fate.} Exact match to \citet[\S1, proposition and preceding tangent formula]{kontsevich1987virasoro}. The paper credits independent results of Beilinson and Schechtman. Publication in 1987 follows either the manuscript's 1985 heading or the catalogue's February 1987 date; the latter conflict is not resolved by the mathematics.\par

\noindent\textbf{Present status.} As of November 2026: Published soon after by the speaker (A), with the precise infinitesimal meaning verified in the primary paper.\par

\noindent\textbf{Assessment.} Class A; literal correctness: correct with missing hypotheses; publication match: exact; verification: read the result; development: none; confidence: high.\par

\subsubsection{ADJ-667.4: A meromorphic lift of a one-parameter deformation}
\noindent\textbf{Source.} P001258; M. L. Kontsevich.\par
\noindent\textbf{Notebook.} For a one-parameter family the notebook chooses a connection with poles only at the section \(p\), locally \(\nabla_t=\partial_t+F(x,t)\partial_x\).\par

\noindent\textbf{Precise statement.} Locally on a smooth one-dimensional base, a family of smooth projective curves with section admits a lift of the base tangent vector to a relative meromorphic vector field after allowing a sufficiently high finite pole order along that section. In a formal trivialization it has the written form. Two such lifts differ by a relative vertical meromorphic field; at a fixed fiber their difference belongs to \(\gamma_{C,p,x}\).\par

\noindent\textbf{Reconstruction and limits.} Allowing a sufficiently large pole divisor kills the relevant first cohomology by Serre vanishing, giving local existence. This is a noncanonical local lift, not a globally chosen holomorphic connection. Its ambiguity explains the quotient by \(\gamma\). The notebook neither fixes a pole order nor proves uniqueness. Editorial local-lift argument from vanishing of the first cohomology after a sufficiently large pole divisor \ours.\par

\noindent\textbf{Literature fate.} This is the deformation-theoretic interpretation of \citet[\S1]{kontsevich1987virasoro}; the article states the equivalent Kodaira--Spencer and Laurent-field description rather than this exact local connection formula.\par

\noindent\textbf{Present status.} As of November 2026: The construction is explained within the established infinitesimal uniformization. It is E as deformation-theoretic machinery, with a partial publication match rather than a claim of a separately stated theorem.\par

\noindent\textbf{Assessment.} Class E; literal correctness: correct with missing hypotheses; publication match: partial; verification: read the result; development: annotated; confidence: high.\par

\subsubsection{ADJ-667.5: The erroneous g/k tangent quotient and the two forgetting maps}
\noindent\textbf{Source.} P001258; M. L. Kontsevich.\par
\noindent\textbf{Notebook.} One line says \(T\widetilde M_g=\mathfrak g/\mathfrak k\), but the following display divides instead by fields extending to \(C\setminus\{p\}\).\par

\noindent\textbf{Precise statement.} The tangent space with full formal parameter is \(\mathfrak g/\gamma\). The tangent space of unpointed moduli is \(\mathfrak g/(\gamma+\mathfrak k)\cong H^1(C,T_C)\), of complex dimension \(3g-3\). The tangent space of pointed moduli is \(\mathfrak g/(\gamma+\mathfrak k_0)\cong H^1(C,T_C(-p))\), of dimension \(3g-2\). The intermediate quotient \(\mathfrak g/\mathfrak k\) is merely the space of formal principal parts; it is not any of these complete tangent spaces.\par

\noindent\textbf{Reconstruction and limits.} These formulas follow from the two-open-set formal gluing calculation: changes of the trivialization on the punctured curve contribute \(\gamma\), while those on the disk contribute \(\mathfrak k\) or \(\mathfrak k_0\) depending on whether its center is remembered. They also distinguish \(\operatorname{Der}\mathbb C[[x]]\) from the Lie algebra \(x\mathbb C[[x]]\partial_x\) of formal coordinate changes fixing the point. Editorial formal-gluing derivation of the two forgetting-map tangent quotients \ours; the full-parameter uniformization theorem is due to Kontsevich.\par

\noindent\textbf{Literature fate.} \citet[\S1]{kontsevich1987virasoro} gives the full-parameter tangent formula and notes that regular formal fields act transitively along the fibers of forgetting the point and parameter. The two finite-dimensional cohomology quotients are the explicit editorial reconstruction.\par

\noindent\textbf{Present status.} As of November 2026: The first written quotient is false (E correction); the next display contains the intended curve-dependent stabilizer. No open deformation problem remains.\par

\noindent\textbf{Assessment.} Class E; literal correctness: false as written; publication match: exact; verification: read the result; development: annotated; confidence: high.\par

\subsubsection{ADJ-667.6: The central extension lifts to determinant-of-cohomology lines}
\noindent\textbf{Source.} P001258; M. L. Kontsevich.\par
\noindent\textbf{Notebook.} The final lines say that the Virasoro central extension acts in the determinant bundle \(d_k\) pulled back from \(M_g\) to \(\widetilde M_g\).\par

\noindent\textbf{Precise statement.} For an integer \(k\), let \(\lambda_k(C)=\det H^0(C,K_C^k)\otimes\det H^1(C,K_C^k)^{-1}\). The infinitesimal Laurent-vector-field action lifts to a Virasoro central extension acting by first-order differential operators on the pullback of \(\lambda_k\) to curves with a formal parameter. Relative to the determinant extension for \(k=0\) (equivalently \(k=1\)), its cocycle class is multiplied by \(6k^2-6k+1\). The numerical central charge changes with central-generator and dual-line conventions.\par

\noindent\textbf{Reconstruction and limits.} The notebook does not define \(d_k\); identifying it with \(\lambda_k\) is the exact source-based reconstruction. A Lie algebra action on sections with nontrivial central action is not an ordinary action of the center-free algebra on the base. No particular scalar central charge is assigned without declaring the cocycle normalization.\par

\noindent\textbf{Literature fate.} \citet[\S\S2--3]{kontsevich1987virasoro} constructs the determinant extension on the Laurent-space Grassmannian and pulls it back to the curve moduli. Section3 supplies the coefficient and the corresponding cancellation corollary.\par

\noindent\textbf{Present status.} As of November 2026: Published by the speaker in 1987 (A). The determinant-line action is verified directly; no unrelated later Kontsevich-model claim is attached to the page.\par

\noindent\textbf{Assessment.} Class A; literal correctness: correct with missing hypotheses; publication match: exact; verification: read the result; development: none; confidence: high.\par

\subsubsection{ADJ-668.1: The critical bosonic measure and the exponent13}
\noindent\textbf{Source.} P001259; Morozov.\par
\noindent\textbf{Notebook.} Morozov, under ``Consequences of the Belavin--Knizhnik theorem'', summarizes Beilinson's talk and writes a sum over genera with measure \(|d\mu_p|^2/(\det\operatorname{Im}T)^{13}\), multiplied by a vertex factor \(K\{V\}\).\par

\noindent\textbf{Precise statement.} For a smooth compact Riemann surface of fixed genus \(p\ge2\), the gauge-fixed critical bosonic vacuum measure can be expressed locally as \(|\mu_p|^2/(\det\operatorname{Im}\tau)^{13}\), where \(\tau\) is the period matrix in a symplectic marking and \(\mu_p\) is the corresponding holomorphic Mumford form, a section of the determinant-line combination \(\lambda_2\otimes\lambda_1^{-13}\). Insertions require the associated normalized matter/ghost correlation factor and their integration data.\par

\noindent\textbf{Reconstruction and limits.} This is a fixed-genus determinant identity with specified zero modes. The infinite genus sum and integration near degenerating curves are not proved convergent. Genus 0 and 1 have automorphisms and additional special zero-mode conventions, so they are not covered by silently repeating the same formula. The scan confirms the exponent13.\par

\noindent\textbf{Literature fate.} \citet[\S\S2--5, especially (2.19), (3.4), and (5.1)]{belavinknizhnik1986complex} proves the determinant-line/holomorphy statement and discusses degeneration poles. The paper cites its own Landau Institute 1985 preprint; the full article was submitted 18 February 1986 and published August 1986.\par

\noindent\textbf{Present status.} As of November 2026: Exposition of already circulating work (E), explicitly credited in the notebook. The recorded theorem does not justify a finite total string partition function.\par

\noindent\textbf{Assessment.} Class E; literal correctness: correct with missing hypotheses; publication match: exact; verification: read the result; development: none; confidence: high.\par

\subsubsection{ADJ-668.2: World sheets, topology, and replacing external circles by punctures}
\noindent\textbf{Source.} P001259; Morozov.\par
\noindent\textbf{Notebook.} Two sketches show a propagating loop and a surface capped at two points. The text says that interactions are accounted for by different topologies, and that it is technically convenient to begin and end at points.\par

\noindent\textbf{Precise statement.} In the perturbative world-sheet description of closed strings, the genus labels the loop order, while punctures and their assigned states describe external insertions. Turning an external cylindrical end into a puncture requires a local coordinate and a state/vertex prescription. This is a formal sewing description of amplitudes, not a topological theorem that arbitrary prescribed boundary-loop data equal a point delta function.\par

\noindent\textbf{Reconstruction and limits.} The caption abbreviation is retained without a forced expansion. The diagrams specify the physical construction but not a measure, sewing theorem, or convergence statement. The point insertion model below is one chosen Fourier representation.\par

\noindent\textbf{Literature fate.} The earlier surface-sum and insertion formulas are in \citet[eqs.~(1)--(3),(26)--(31)]{polyakov1981quantum}; the fixed-genus consequences appear in \citet[\S6]{belavinknizhnik1986complex}.\par

\noindent\textbf{Present status.} As of November 2026: Known formal perturbative interpretation (E); no new precise topological or analytic question is posed.\par

\noindent\textbf{Assessment.} Class E; literal correctness: not determinable; publication match: thematic only; verification: read the result; development: none; confidence: high.\par

\subsubsection{ADJ-668.3: Fourier-transforming target point constraints into exponential vertices}
\noindent\textbf{Source.} P001259; Morozov.\par
\noindent\textbf{Notebook.} The page first calls \(V_\alpha\) delta functions, then writes \(V_\alpha(p)=\int e^{ip_\alpha\cdot X(\xi)}d^2\xi\).\par

\noindent\textbf{Precise statement.} For a fixed smooth map, Fourier transformation of \(\delta^{(d)}(X(\xi)-y)\) in the target point \(y\), with convention \(\int e^{ip\cdot y}\delta(X-y)dy\), gives \(e^{ip\cdot X(\xi)}\). In the quantum free-boson theory this exponential needs normal ordering. An integrated closed-string vertex must also be a coordinate-invariant world-sheet density; in conformal coordinates the usual on-shell scalar vertex has conformal weights \((1,1)\) before integration against \(d^2z\).\par

\noindent\textbf{Reconstruction and limits.} The elementary Fourier identity is exact for a smooth map. Its promotion to an integrated physical vertex requires a normalization of the action, a momentum/mass-shell condition, regularization, and the ghost prescription. None is specified by the bare coordinate integral. Editorial Fourier identity in the stated sign convention \oursComputed; the quantum vertex construction is prior literature.\par

\noindent\textbf{Literature fate.} \citet[eqs.~(26)--(31)]{polyakov1981quantum} gives the vertex construction; \citet[\S6, (6.21)--(6.26)]{belavinknizhnik1986complex} computes the resulting tachyon factors with momentum conservation and mass-shell conditions.\par

\noindent\textbf{Present status.} As of November 2026: Established construction (E), with the distinction between a classical Fourier identity and a renormalized quantum vertex retained.\par

\noindent\textbf{Assessment.} Class E; literal correctness: correct with missing hypotheses; publication match: exact; verification: read the result; development: annotated; confidence: high.\par

\subsubsection{ADJ-668.4: A target delta derivative is not a world-sheet derivative vertex}
\noindent\textbf{Source.} P001259; Morozov.\par
\noindent\textbf{Notebook.} The sentence about derivatives of delta functions is followed by \(\int\partial X^\mu e^{ip\cdot X}\,d^2\xi\).\par

\noindent\textbf{Precise statement.} Fourier transformation of a target-coordinate derivative \(\partial_{y^\mu}\delta(X-y)\) gives a factor \(-ip_\mu e^{ip\cdot X}\) under the preceding convention. It does not give \(\partial X^\mu\). World-sheet derivatives of the free field produce oscillator or descendant insertions. For a closed-string tensor vertex a typical full density includes both \(\partial X^\mu\) and \(\bar\partial X^\nu\), appropriate polarization and on-shell conditions; the isolated chiral factor has an uncontracted world-sheet index.\par

\noindent\textbf{Reconstruction and limits.} The notebook does not say which delta variable is differentiated. If it means the same target point whose Fourier transform was just taken, the indicated implication is false. If it means a different state construction, that is missing data, not a justification for silently adding the antiholomorphic factor. Editorial target-delta derivative calculation \oursComputed.\par

\noindent\textbf{Literature fate.} The distinction is an elementary Fourier calculation, checked against the exponential-insertion setup in \citet[eqs.~(26)--(31)]{polyakov1981quantum}.\par

\noindent\textbf{Present status.} As of November 2026: E with a faithful mathematical gap; no general new vertex-operator theorem is inferred from the incomplete line.\par

\noindent\textbf{Assessment.} Class E; literal correctness: false as written; publication match: partial; verification: read the result; development: annotated; confidence: high.\par

\subsubsection{ADJ-668.5: Gaussian integration with sources, determinant powers, and the zero mode}
\noindent\textbf{Source.} P001259; Morozov.\par
\noindent\textbf{Notebook.} Integrating over \(X\) is said to give \((\det\Delta_0)^{-13}\) times products of exponentials \(p_\alpha p_\beta G_{\mathrm{reg}}(\xi_\alpha,\xi_\beta)\). The preceding action has a positive exponent.\par

\noindent\textbf{Precise statement.} Fix a closed connected surface and a nonnegative scalar Laplacian \(\Delta\). For 26 real mean-zero fields with action \(S=\frac12\langle X,\Delta X\rangle\), a finite-mode Gaussian integral with sources \(\sum_a p_a\delta_{z_a}\) equals a normalization times \((\det{}'\Delta)^{-13}\exp[-\frac12\sum_{a,b}p_a\cdot p_b\,G(z_a,z_b)]\), before diagonal regularization, where \(G=\Delta^{-1}\) on mean-zero functions. Integrating the constant target mode separately gives \((2\pi)^{26}\delta^{(26)}(\sum_a p_a)\). The infinite-dimensional determinant requires zeta regularization.\par

\noindent\textbf{Reconstruction and limits.} The sign and factor 1/2 are stated for this explicit action and covariance convention; another convention for \(G\) changes their appearance. A positive Euclidean exponential of positive energy is divergent. The prime, momentum-conservation delta, and contractions of target indices are missing from the notebook's shorthand. Derivative vertices require additional differentiated Green-function contractions. Editorial finite-mode Gaussian completion of the square and zero-mode check \oursComputed; the string-theory application is the cited published calculation.\par

\noindent\textbf{Literature fate.} \citet[\S6, (6.21)--(6.24)]{belavinknizhnik1986complex} performs this Gaussian source calculation and explicitly uses momentum conservation. The finite-mode completion-of-the-square identity provides the normalization check.\par

\noindent\textbf{Present status.} As of November 2026: Established Gaussian calculation (E). The result read is a regularized fixed-surface building block, not a rigorous value for the total genus sum.\par

\noindent\textbf{Assessment.} Class E; literal correctness: false as written; publication match: exact; verification: read the result; development: annotated; confidence: high.\par

\subsubsection{ADJ-668.6: The Green-function diagonal and Wick ordering}
\noindent\textbf{Source.} P001259; Morozov.\par
\noindent\textbf{Notebook.} The final line defines \(G_{\mathrm{reg}}\) as a Green function smoothed when the two points coincide.\par

\noindent\textbf{Precise statement.} The two-dimensional scalar covariance has a logarithmic diagonal singularity. For a mollified mean-zero Gaussian field \(X_\varepsilon\), define a point exponential at finite cutoff by \(:e^{ip\cdot X_\varepsilon(z)}:=\exp(ip\cdot X_\varepsilon(z)+\frac12|p|^2\mathbb E[X_\varepsilon^1(z)^2])\). At pairwise distinct points, its normalized correlation has limit \(\exp[-\sum_{a<b}p_a\cdot p_bG(z_a,z_b)]\). Integrated collisions, local metric factors, and alternative diagonal finite parts require additional prescriptions.\par

\noindent\textbf{Reconstruction and limits.} Here the components have identical scalar covariance and are independent. This formula identifies one precise diagonal subtraction. Merely replacing the diagonal by an arbitrary smooth value does not select the physical normalization or prove integrability when insertion points collide. Editorial Wick-subtraction and off-diagonal contraction calculation \oursComputed; normal ordering itself is established machinery.\par

\noindent\textbf{Literature fate.} \citet[\S6, between (6.24) and (6.26)]{belavinknizhnik1986complex} explicitly regularizes the diagonal and then imposes the mass-shell condition; \citet[\S3]{dotsenkofateev1984conformal} supplies the corresponding normal-ordered free-boson contraction framework.\par

\noindent\textbf{Present status.} As of November 2026: Known renormalized Gaussian building block (E). No unqualified convergence theorem is extracted from the notebook's word ``smoothed''.\par

\noindent\textbf{Assessment.} Class E; literal correctness: correct with missing hypotheses; publication match: partial; verification: read the result; development: annotated; confidence: high.\par

\section{Problem index}\label{problems:1985}
Earlier-year occurrence identifiers are retained as text; current-year identifiers link to their entries.

\subsection{GP-1985-01: Uniform degree bounds for zeros of Abelian integrals}\label{prob:GP-1985-01}
For real polynomials \(H,P,Q\) of bounded degrees, bound the total isolated zeros, counted with multiplicity, of \(\int_{\gamma(t)}(P\,dx+Q\,dy)\) on the regular compact oval branches of \(H=t\), uniformly in the coefficients. Identically zero branches are excluded from the isolated-zero count.
\par\noindent\textit{Occurrences:} \hyperref[item:1985-642.1]{1985-642.1}.
\par\noindent\textit{Status (2026-11):} As of November 2026, a finite explicit uniform bound is established by \citet[Theorem 2]{binyamininovikovyakovenko2010zeros} (preprint 2008; publication 2010). This addresses the precise first-period question, not Hilbert's full sixteenth problem. No exact earlier occurrence with this problem's quantifiers was found in the supplied index, which ends in 1982.

\subsection{GP-1985-02: Local Langlands correspondence for GSp(4)}\label{prob:GP-1985-02}
For every characteristic-zero nonarchimedean local field \(F\), construct the finite \(L\)-packet classification of irreducible smooth \(GSp_4(F)\)-representations by admissible Weil--Deligne parameters in \(GSp_4(\mathbb C)\), with the specified natural compatibilities.
\par\noindent\textit{Occurrences:} \hyperref[item:1985-644.6]{1985-644.6}.
\par\noindent\textit{Status (2026-11):} As of November 2026, solved in this scope by \citet[Main Theorem]{gantakeda2011local}, first circulated in 2007 and published in 2011. Vigneras's 1986 work is an important predecessor, not evidence that the complete exhaustion had already been proved in 1985.

\subsection{GP-1985-03: Central section of a four-cube perpendicular to a main diagonal}\label{prob:GP-1985-03}
Identify, up to Euclidean similarity, \([-1,1]^4\cap\{x_1+x_2+x_3+x_4=0\}\).
\par\noindent\textit{Occurrences:} \hyperref[item:1985-647.1]{1985-647.1}.
\par\noindent\textit{Status (2026-11):} As of November 2026: an elementary solved exercise. The section is a regular octahedron, proved in 1985-647.1 and independently checked by exact enumeration in this chapter. The index retains this explicit posed question while the E classification records that no research novelty is claimed.

\subsection{GP-1985-04: Arithmetic occurrence and multiplicity in the automorphic spectrum}\label{prob:GP-1985-04}
For a fixed connected semisimple \(G/\mathbb Q\), determine which irreducible admissible restricted products of local representations occur in the discrete spectrum of \(L^2(G(\mathbb Q)\backslash G(\mathbb A))\), and their multiplicities; the unrestricted seminar question ranges over such groups. Distinguish this arithmetic classification from the already established decomposition by Eisenstein series.
\par\noindent\textit{Occurrences:} \hyperref[item:1985-648.9]{1985-648.9}.
\par\noindent\textit{Status (2026-11):} As of November 2026: partly developed, with major classifications for specific families of groups, exemplified by Arthur's orthogonal/symplectic theory. The cited results do not give the unrestricted classification. The problem is too broad for a bounded new contribution without choosing a group and a particular packet or multiplicity. GP-1982-43 concerns a different automorphy question and is linked contextually, not merged.

\subsection{GP-1985-20: Deciding identities in all finite-dimensional subspace lattices}\label{prob:GP-1985-20}
For a fixed field \(F\), is there an algorithm deciding whether two finite terms in sum and intersection have equal values on every tuple of subspaces in every finite-dimensional \(F\)-vector space? The concrete proposed reduction is to dimension \(1+\min\{\max(d(p),d(q)),\max(k(p),k(q))\}\), where \(d\) and \(k\) count joins and meets in the written terms.
\par\noindent\textit{Occurrences:} \hyperref[item:1985-653.2]{1985-653.2}, \hyperref[item:1985-653.3]{1985-653.3}.
\par\noindent\textit{Status (2026-11):} Solved by the speaker: \citet[Lemma~1, Lemmas~2--3 and main theorem]{ryzhikova1987decidability}; received May 1986, Russian publication 1987, English translation 1988. As of November 2026, the exact identity-decision problem is settled; no complexity or general first-order decidability claim is added.

\subsection{GP-1985-21: Index formulas for elliptic Fuchs-type operators with conic ends}\label{prob:GP-1985-21}
For a cone-elliptic regular singular differential operator on a compact manifold with isolated conic ends, and a specified closed \(L^2\) domain between its minimal and maximal domains, express its index by interior data, Mellin-symbol spectral data, and the chosen indicial-domain contribution. This is a reconstruction of the notebook's general ``index?'' question; its printed operator and domain are incomplete.
\par\noindent\textit{Occurrences:} \hyperref[item:1985-654.6]{1985-654.6}.
\par\noindent\textit{Status (2026-11):} Developed by others. \citet[Prop.~1.3.16 and Cor.~2.4.7]{lesch1997fuchs} proves Fredholmness and an index formula in the constant-near-tip setting; Thm.~2.4.8 specializes it for first-order operators. Habilitation submitted 1993; preprint 1996; book 1997. As of November 2026 this is an established formulation, not a claim that every unqualified regular singular expression has a canonical index.

\subsection{GP-1982-03: Reducibility and first singular vectors of Virasoro Verma modules}\label{prob:GP-1982-03}
For \(h,c\in\mathbb C\), determine when the Virasoro Verma module \(M(h,c)\) is reducible and determine its first positive grade containing a singular vector. The central-generator convention must be fixed before comparing formulas.
\par\noindent\textit{Occurrences:} 1982-566.7, \hyperref[item:1985-657.3]{1985-657.3}, \hyperref[item:1985-657.4]{1985-657.4}.
\par\noindent\textit{Status (2026-11):} Solved before the present talk. \citet[\S3.2, Thm.~3.1 and corollary]{feiginfuchs1982invariant} gives the determinant factors and the first-grade criterion. Received September 1981, published 1982. As of November 2026, this recurrence remains E; the 1985 notebook's \(1/2\) coefficient must be corrected to \(1/16\).

\subsection{GP-1985-40: Topology of real configuration realization spaces}\label{prob:GP-1985-40}
For a prescribed finite rank configuration or oriented-matroid type, determine the topology of its normalized real realization space; in particular, determine what homotopy types can occur. Incidence types and uniform chirotopes are distinguished. The exact modern universality theorem concerns stable equivalence to primary basic semialgebraic sets, with the uniform case restricted to open sets.
\par\noindent\textit{Occurrences:} \hyperref[item:1985-658.1]{1985-658.1}, \hyperref[item:1985-658.6]{1985-658.6}, \hyperref[item:1985-658.16]{1985-658.16}.
\par\noindent\textit{Status (2026-11):} As of November 2026, general rank-three universality is established and was already announced before the talk by \citet[Theorem~3]{mnev1985varieties}. The modern formulation is \citet[Theorem~1]{adiprasitopadrol2017universality}. No claim of literal homeomorphism to every arbitrary algebraic variety is intended.

\subsection{GP-1985-41: Isotopy within a fixed determinant-sign class}\label{prob:GP-1985-41}
For a uniform realizable rank-three oriented matroid on \(n\) labelled elements, is its normalized real realization space path connected? Equivalently, can matrices with the same nonzero maximal-minor signs be joined through the same sign class after fixing an oriented basis?
\par\noindent\textit{Occurrences:} \hyperref[item:1985-658.3]{1985-658.3}, \hyperref[item:1985-658.4]{1985-658.4}, \hyperref[item:1985-658.5]{1985-658.5}.
\par\noindent\textit{Status (2026-11):} As of November 2026, the general answer is negative by universality, already available before this talk. The small-size assertion \(n\leq7\) is supported by a stronger known bound cited in \citet[introduction]{tsukamoto2013disconnected}. The exact uniform nineteen-element example attributed to Mn\"ev in the notebook remains an unresolved source identification; this does not make the general mathematical question open.

\subsection{GP-1985-42: Universality in dimension four under simpliciality}\label{prob:GP-1985-42}
The notebook informally suggests ``more or less any topology'' for fixed-type simplicial four-polytope realization spaces. An editorial homotopy sharpening asks whether every nonempty finite simplicial-complex homotopy type occurs. The stronger precise modern conjecture asks whether every nonempty open primary basic semialgebraic set defined over \(\mathbb Z\) is stably equivalent to the homogenized realization space, modulo \(\mathrm{GL}_5(\mathbb R)\), of a neighborly (hence simplicial) four-polytope.
\par\noindent\textit{Occurrences:} \hyperref[item:1985-658.10]{1985-658.10}.
\par\noindent\textit{Status (2026-11):} As of November 2026, no resolution of either fixed-four-dimensional universality formulation was found in the bounded search. General four-polytopes are universal by \citet[Main Theorem]{richtergebertziegler1995realization}; simplicial neighborly universality in \citet[Theorem~2]{adiprasitopadrol2017universality} allows growing dimension, and their Conjecture~4 retains dimension four and additionally requires neighborliness. The latest opened explicit open-status evidence is the 2024 compendium CG68--CG69; this is not an exhaustive proof of absence of a later result. The chapter establishes only contractibility for normalized history-compatible spaces of prescribed facet-stackings.

\subsection{GP-1985-60: Recognizing Jacobians among principally polarized abelian varieties}\label{prob:GP-1985-60}
Characterize the period matrices of smooth compact connected complex curves among symmetric matrices with positive definite imaginary part, modulo the symplectic group. Require indecomposability in the KP formulation; distinguish this recognition problem from finding a chosen explicit ideal of the Jacobian locus.
\par\noindent\textit{Occurrences:} \hyperref[item:1985-664.5]{1985-664.5}.
\par\noindent\textit{Status (2026-11):} As of November 2026: solved in the KP-recognition sense by Shiota, with earlier circulation documented separately, and refined by Krichever\textquotesingle s auxiliary linear equation. The exact unnamed relations in the seminar margin cannot be reconstructed. Earlier items 1974-336.5 and 1982-592.6 are contextual links, not asserted occurrences of this precise question.

\section{Corrections to the transcription and translation}\label{corr:1985}
The evidence distinguishes a transcription correction from a mathematical emendation of wording faithfully preserved in the notebook.
\subsection*{P001199 (en; translation): Opening Baker-function phrase}
\noindent\textit{Transcription reads.} English: Baker function of an arbitrary set\par
\noindent\textit{Critical reading.} Baker function on an arbitrary manifold (or variety)\par
\noindent\textit{Evidence.} Russian abbreviation mn-ziya is the standard manuscript abbreviation for manifold; the scalar exponential construction needs a complex analytic space. The scan was inspected.\par
\subsection*{P001199 (both; ambiguity): Superdimension notation}
\noindent\textit{Transcription reads.} \(\mathbb C^{10/16}\), lines of dimension \(1/8\)\par
\noindent\textit{Critical reading.} Interpret as \(\mathbb C^{10|16}\) and \(1|8\), retaining the diplomatic slash\par
\noindent\textit{Evidence.} The slash is visible in the scan; even and odd dimensions in the immediately following discussion determine the meaning.\par
\subsection*{P001199 (both; faithful error): Spinor representation}
\noindent\textit{Transcription reads.} Spinor representation of \(SO(10)\)\par
\noindent\textit{Critical reading.} Linear half-spin representation of \(\operatorname{Spin}(10)\), projective on \(SO(10)\)\par
\noindent\textit{Evidence.} The notebook faithfully writes SO; the double-cover qualification follows from the Clifford construction.\par
\subsection*{P001200 (both; faithful error): Marginal Clifford relation}
\noindent\textit{Transcription reads.} \(xy+yx=(x,y)I\)\par
\noindent\textit{Critical reading.} \(xy+yx=2(x,y)I\) with the same bilinear-form normalization\par
\noindent\textit{Evidence.} The main generator relation visibly has \(2\delta_{ik}\). The margin omits 2; alternatively it would have to redefine the form.\par
\subsection*{P001200 (both; faithful error): Inner automorphism inference}
\noindent\textit{Transcription reads.} Innerness gives a representation of the orthogonal group\par
\noindent\textit{Critical reading.} Innerness gives a projective representation; the spin/pin cover gives the linear representation\par
\noindent\textit{Evidence.} The scan contains no scalar-lifting argument. Implementing matrices are determined only up to scalars.\par
\subsection*{P001200 (both; ambiguity): Projective diagram}
\noindent\textit{Transcription reads.} Projective superscript appears as 12 (overwritten); \((\mathbb C^{12}\setminus0)/\mathbb C^*\)\par
\noindent\textit{Critical reading.} If the vector space exponent is 12, its projectivization is \(\mathbb P^{11}\); retain the superscript ambiguity\par
\noindent\textit{Evidence.} Direct inspection cannot securely distinguish all retraced strokes; the dimension calculation is definite. No identification of the remaining arrows is asserted.\par
\subsection*{P001201 (both; provenance): Theorem attribution}
\noindent\textit{Transcription reads.} German Petrov-Tankin (diplomatic personal-name reading)\par
\noindent\textit{Critical reading.} Retain the reading; bibliographic identification of the theorem: G. S. Petrov, 1986\par
\noindent\textit{Evidence.} The handwritten personal wording was inspected. The exact Hamiltonian, degree convention and bound match the opened original theorem by G. S. Petrov. The marginal student biography is not verified.\par
\subsection*{P001201 (both; faithful error): Linear-independence comment}
\noindent\textit{Transcription reads.} A combination of \(N\) independent functions has \(N-1\) zeros\par
\noindent\textit{Critical reading.} One can force zeros at prescribed \(N-1\) points; an upper bound requires a Chebyshev hypothesis\par
\noindent\textit{Evidence.} No inequality sign appears in the scan. The functions \(1\text{ and }\sin(t)\text{ on }(0,10\pi)\) disprove the unqualified exact/upper-bound reading.\par
\subsection*{P001201 (both; ambiguity): Period-kernel shorthand}
\noindent\textit{Transcription reads.} Multiply by \(H-c\) to obtain zero integral\par
\noindent\textit{Critical reading.} This is a fixed-fiber identity; in the moving family use exact forms and multiples of \(dH\)\par
\noindent\textit{Evidence.} The text uses a fixed c. The statement is correct there but cannot be applied unchanged as the fiber parameter varies.\par
\subsection*{P001202 (both; faithful error): Pochhammer product}
\noindent\textit{Transcription reads.} \((1;q^2)_\infty(-zq^{-1};q^2)_\infty(-z^{-1}q^{-1};q^2)_\infty\)\par
\noindent\textit{Critical reading.} \((q^2;q^2)_\infty(-zq;q^2)_\infty(-z^{-1}q;q^2)_\infty\)\par
\noindent\textit{Evidence.} All three recorded arguments are visible. Expanding the immediately supplied zero-based definition verifies the correction; \((1;q^2)_\infty\) is zero.\par
\subsection*{P001202 (both; ambiguity): Beta odd-part inequalities}
\noindent\textit{Transcription reads.} Overwritten comparisons may look non-strict\par
\noindent\textit{Critical reading.} Strict decrease is required for the combinatorial interpretation; retain visual uncertainty\par
\noindent\textit{Evidence.} Each product factor contributes at most once, so the parts are distinct. This argument resolves the mathematics, not the uncertain pen strokes.\par
\subsection*{P001203 (both; faithful error): Binary-word indexing}
\noindent\textit{Transcription reads.} All eventually constant words counted by partitions\par
\noindent\textit{Critical reading.} Take words modulo translation, or choose their charge-zero origin\par
\noindent\textit{Evidence.} Without identifying shifts the count is infinite. The manuscript describes words and gaps but leaves the equivalence convention implicit.\par
\subsection*{P001203 (both; ambiguity): Young-boundary coordinate}
\noindent\textit{Transcription reads.} Gap label \(x-y\) with a one drawn upward\par
\noindent\textit{Critical reading.} \(x-y\) for downward-positive y; \(x+y\) for upward-positive y\par
\noindent\textit{Evidence.} The scan fixes step directions but not the ambient coordinate convention clearly enough to call this a transcription error.\par
\subsection*{P001203 (both; faithful error): Levi parenthesis}
\noindent\textit{Transcription reads.} M: semisimple part of P\par
\noindent\textit{Critical reading.} M: reductive Levi factor of P\par
\noindent\textit{Evidence.} The Russian parenthesis is faithfully copied. Unramified-character families require the Levi center; a Levi need not be semisimple.\par
\subsection*{P001203 (both; provenance): Personal name in trace-formula report}
\noindent\textit{Transcription reads.} Jacket / Dzhaket\par
\noindent\textit{Critical reading.} Jacquet is the bibliographic identification; retain the diplomatic spelling\par
\noindent\textit{Evidence.} Trace-formula context identifies H. Jacquet, but the notebook spelling itself is not silently replaced.\par
\subsection*{P001203-P001204 (both; provenance): Trace Paley-Wiener attribution}
\noindent\textit{Transcription reads.} Kazhdan and Bernstein\par
\noindent\textit{Critical reading.} Bernstein--Deligne--Kazhdan in the published attribution\par
\noindent\textit{Evidence.} Full original title page names all three authors; received 26 August 1985. This corrects attribution, not a misread word.\par
\subsection*{P001204 (both; faithful error): Bernstein variety equality}
\noindent\textit{Transcription reads.} \(\coprod\{\rho\chi\}/W=\operatorname{Irr}G\)\par
\noindent\textit{Critical reading.} The left-hand object is the cuspidal-support variety; \(\operatorname{Irr}G\to\Omega(G)\) is finite-to-one\par
\noindent\textit{Evidence.} The scan records an equality. BDK Section 2.1 explicitly gives the surjective finite-to-one map, and reducible induced representations show why equality fails.\par
\subsection*{P001204 (both; faithful error): Paley-Wiener image conditions}
\noindent\textit{Transcription reads.} Algebraicity and only finitely many families\par
\noindent\textit{Critical reading.} Add one compact open subgroup K dominating the trace form, and regularity on all induced families in R(G)\par
\noindent\textit{Evidence.} The scan is abbreviated. BDK Section 1.2 supplies the exact necessary and sufficient hypotheses.\par
\subsection*{P001204 (both; faithful error): Local GSp field}
\noindent\textit{Transcription reads.} \(GSp_4(\mathbb Q)\) beside Weil--Deligne representations\par
\noindent\textit{Critical reading.} \(GSp_4(\mathbb Q_p)\) in this local arrow\par
\noindent\textit{Evidence.} The scan writes Q in a local-packet discussion. A global Q-point representation is not classified by a single local Weil--Deligne group.\par
\subsection*{P001204 (both; faithful error): Hermitian matrix condition}
\noindent\textit{Transcription reads.} \(X=\bar X\)\par
\noindent\textit{Critical reading.} \(X=\bar X^{\,t}\)\par
\noindent\textit{Evidence.} The displayed off-diagonal pair \(x_2\) and \(\overline{x_2}\) is visibly Hermitian; the transpose is missing in the compact condition.\par
\subsection*{P001204 (both; faithful error): Weil representation of similitude groups}
\noindent\textit{Transcription reads.} Weil representation of \(GO_4\times GSp_4\)\par
\noindent\textit{Critical reading.} Specify the compatible-multiplier subgroup and its induced extension\par
\noindent\textit{Evidence.} The full product does not preserve the tensor-product symplectic form under the naive action. Gan--Takeda Section 2 states the exact construction.\par
\subsection*{P001204 (both; provenance): Order of record headings}
\noindent\textit{Transcription reads.} Catalogue record numbers put Feigin 645 before Vigneras 646\par
\noindent\textit{Critical reading.} Physical page order is Vigneras 646, Beilinson intervention in 644, then Feigin 645\par
\noindent\textit{Evidence.} Original PNG headings and separators inspected; Beilinson here is still 16 September, not record 648 of 23 September.\par
\subsection*{P001205 (both; faithful error): First integral in Polyakov action}
\noindent\textit{Transcription reads.} \(\lambda\int\sqrt{\det g}\)\par
\noindent\textit{Critical reading.} \(\lambda\int\sqrt{\det g}\,d\xi_1\,d\xi_2\)\par
\noindent\textit{Evidence.} The scan visibly suppresses the coordinate measure only in the first term; the geometric interpretation is the metric area.\par
\subsection*{P001206 (both; faithful error): p-adic norm unit condition}
\noindent\textit{Transcription reads.} \(a=p^nm,\ (a,m)=1\)\par
\noindent\textit{Critical reading.} \(a=p^nm\) with \(m\) a p-adic unit, equivalently \(v_p(m)=0\)\par
\noindent\textit{Evidence.} The scan confirms the letter a in the gcd condition. For integral m, the intended condition is p not dividing m.\par
\subsection*{P001206 (both; faithful error): Adele definition}
\noindent\textit{Transcription reads.} \(\prod_p\mathbb Q_p\)\par
\noindent\textit{Critical reading.} \(\mathbb R\times\prod^{\prime}_p\mathbb Q_p\), restricted relative to \(\mathbb Z_p\)\par
\noindent\textit{Evidence.} Both the restricted-product condition and archimedean place are absent in the literal manuscript expression.\par
\subsection*{P001206 (both; ambiguity): Inverse-limit quotient}
\noindent\textit{Transcription reads.} General G adelic quotient equals one inverse system of real congruence quotients\par
\noindent\textit{Critical reading.} Valid in the displayed SL2 setting by strong approximation; general groups require component and approximation hypotheses\par
\noindent\textit{Evidence.} The adjacent group is SL2. The formula should not be detached from that context. Arthur 1981 Section 2 verifies the intended case.\par
\subsection*{P001206 (both; faithful error): All congruence curves are spheres}
\noindent\textit{Transcription reads.} \(\mathfrak H/\Gamma_n\) is a sphere with punctures\par
\noindent\textit{Critical reading.} General modular curves have varying genus; level 6 has genus 1 and level 7 genus 3\par
\noindent\textit{Evidence.} The sphere wording is visible in the scan, including an overwritten phrase. Milne Example 2.23 supplies explicit counterexamples.\par
\subsection*{P001206 (both; faithful error): Full modular group quotient}
\noindent\textit{Transcription reads.} \(\mathfrak H/SL_2(\mathbb Z)\) is a sphere with two punctures\par
\noindent\textit{Critical reading.} One cusp is removed; the two elliptic orbifold points are not punctures\par
\noindent\textit{Evidence.} The scan clearly says two. Milne Proposition 2.21 and the standard fundamental domain distinguish the cusp and two elliptic points.\par
\subsection*{P001206 (both; faithful error): Infinite tensor product}
\noindent\textit{Transcription reads.} \(\pi_\infty\otimes\pi_2\otimes\cdots\)\par
\noindent\textit{Critical reading.} Restricted tensor product with specified spherical reference vectors almost everywhere\par
\noindent\textit{Evidence.} The displayed finite-place sequence is correct shorthand only with the restriction; the admissible automorphic factorization supplies it.\par
\subsection*{P001206 (both; faithful error): Sufficiency of local classification}
\noindent\textit{Transcription reads.} Enough to know representations of p-adic groups\par
\noindent\textit{Critical reading.} Local factors are necessary; global occurrence and multiplicity require further conditions and the archimedean factor\par
\noindent\textit{Evidence.} The following notebook question itself says not all products occur. The modern statement separates these two tasks.\par
\subsection*{P001207 (both; faithful error): Euler product after almost-all qualifier}
\noindent\textit{Transcription reads.} \(\prod_p L_p\) using only spherical factors\par
\noindent\textit{Critical reading.} \(L^S=\prod_{p\notin S}L_p\); complete factors require further definitions\par
\noindent\textit{Evidence.} The scan says almost all factors are spherical but prints all primes. Arthur 1981 Section 5 explicitly uses the partial product.\par
\subsection*{P001217 (both; provenance): date, records 652--653}
\noindent\textit{Transcription reads.} Catalogue: 21 September 1985\par
\noindent\textit{Critical reading.} Manuscript date: 21 October 1985\par
\noindent\textit{Evidence.} Heading 21/X/85 visible on scan; the subsequent 28/X and 4/XI sequence agrees. Retain catalogue date in provenance.\par
\subsection*{P001217 (both; provenance): source order}
\noindent\textit{Transcription reads.} Catalogue record order 652 then 653\par
\noindent\textit{Critical reading.} 653 (Ryzhikova, upper page), then 652 (Fuchs, lower page)\par
\noindent\textit{Evidence.} Speaker headings on the same manuscript page determine source order.\par
\subsection*{P001207 lower--P001209 (both; provenance): record 649 boundary}
\noindent\textit{Transcription reads.} Overlapping catalogue ranges\par
\noindent\textit{Critical reading.} P001207 lower--P001209\par
\noindent\textit{Evidence.} Boundary follows handwritten speaker/date/topic changes, checked against Russian transcription and page images.\par
\subsection*{P001210--P001212 (both; provenance): record 650 boundary}
\noindent\textit{Transcription reads.} Overlapping catalogue ranges\par
\noindent\textit{Critical reading.} P001210--P001212\par
\noindent\textit{Evidence.} Boundary follows handwritten speaker/date/topic changes, checked against Russian transcription and page images.\par
\subsection*{P001213--P001216 (both; provenance): record 651 boundary}
\noindent\textit{Transcription reads.} Overlapping catalogue ranges\par
\noindent\textit{Critical reading.} P001213--P001216\par
\noindent\textit{Evidence.} Boundary follows handwritten speaker/date/topic changes, checked against Russian transcription and page images.\par
\subsection*{P001217 upper (both; provenance): record 653 boundary}
\noindent\textit{Transcription reads.} Overlapping catalogue ranges\par
\noindent\textit{Critical reading.} P001217 upper\par
\noindent\textit{Evidence.} Boundary follows handwritten speaker/date/topic changes, checked against Russian transcription and page images.\par
\subsection*{P001217 lower--P001218 (both; provenance): record 652 boundary}
\noindent\textit{Transcription reads.} Overlapping catalogue ranges\par
\noindent\textit{Critical reading.} P001217 lower--P001218\par
\noindent\textit{Evidence.} Boundary follows handwritten speaker/date/topic changes, checked against Russian transcription and page images.\par
\subsection*{P001219--P001221 upper (both; provenance): record 654 boundary}
\noindent\textit{Transcription reads.} Overlapping catalogue ranges\par
\noindent\textit{Critical reading.} P001219--P001221 upper\par
\noindent\textit{Evidence.} Boundary follows handwritten speaker/date/topic changes, checked against Russian transcription and page images.\par
\subsection*{P001221 lower--P001222 (both; provenance): record 655 boundary}
\noindent\textit{Transcription reads.} Overlapping catalogue ranges\par
\noindent\textit{Critical reading.} P001221 lower--P001222\par
\noindent\textit{Evidence.} Boundary follows handwritten speaker/date/topic changes, checked against Russian transcription and page images.\par
\subsection*{P001223--P001224 (both; provenance): record 656 boundary}
\noindent\textit{Transcription reads.} Overlapping catalogue ranges\par
\noindent\textit{Critical reading.} P001223--P001224\par
\noindent\textit{Evidence.} Boundary follows handwritten speaker/date/topic changes, checked against Russian transcription and page images.\par
\subsection*{P001225--P001228 (both; provenance): record 657 boundary}
\noindent\textit{Transcription reads.} Overlapping catalogue ranges\par
\noindent\textit{Critical reading.} P001225--P001228\par
\noindent\textit{Evidence.} Boundary follows handwritten speaker/date/topic changes, checked against Russian transcription and page images.\par
\subsection*{P001207 (both; faithful error): 649, induced metric}
\noindent\textit{Transcription reads.} Euclidean or Minkowski target gives a Riemannian metric\par
\noindent\textit{Critical reading.} Require immersion and spacelike world sheet; use absolute determinant for a timelike sheet\par
\noindent\textit{Evidence.} The original allows both signatures; its blanket Riemannian conclusion needs this qualification.\par
\subsection*{P001208 (both; faithful error): 649, second variation}
\noindent\textit{Transcription reads.} \(\delta^2 A=-\int(\Delta V,V)\)\par
\noindent\textit{Critical reading.} Area uses the Jacobi form, including second fundamental form; fixed-metric Dirichlet energy has the displayed Hessian.\par
\noindent\textit{Evidence.} Scan retains the false simplification. A compactly supported tangential reparametrization of a planar disk keeps area constant but has positive Dirichlet quadratic energy.\par
\subsection*{P001208 (both; faithful error): 649, Taylor expansion}
\noindent\textit{Transcription reads.} Quadratic term written without conventional Hessian factor\par
\noindent\textit{Critical reading.} \(A(x+\varepsilon V)=A(x)+\varepsilon\delta A+\frac12\varepsilon^2\delta^2 A+o(\varepsilon^2)\)\par
\noindent\textit{Evidence.} The conventional second-derivative normalization requires 1/2; do not silently transfer a differently normalized quadratic form.\par
\subsection*{P001208 (both; faithful error): 649, harmonic versus conformal}
\noindent\textit{Transcription reads.} Fixed-metric critical map has conformal induced metric\par
\noindent\textit{Critical reading.} Also vary the metric, imposing the vanishing stress tensor\par
\noindent\textit{Evidence.} Harmonic map (u,v) to (2u,v) is an immediate counterexample; \citet[eq.~(7)]{polyakov1981quantum} contains the additional equation.\par
\subsection*{P001209 (both; faithful error): 649, metric-space norm}
\noindent\textit{Transcription reads.} \(d\xi_1d\xi_k\)\par
\noindent\textit{Critical reading.} \(d\xi_1d\xi_2\)\par
\noindent\textit{Evidence.} The scan has the index error in the area differential.\par
\subsection*{P001210 (both; faithful error): 650, determinant and inverse}
\noindent\textit{Transcription reads.} Unprimed determinant/inverse despite zero modes\par
\noindent\textit{Critical reading.} Use primed determinant and complement of the kernel; include cohomology norms in determinant lines\par
\noindent\textit{Evidence.} Closed-surface scalar and Dolbeault operators have zero modes; \citet[\S\S1--3]{quillen1985determinants}.\par
\subsection*{P001210 (both; faithful error): 650, conformal adjoint}
\noindent\textit{Transcription reads.} \(e^{j\varphi}\bar\partial^*e^{-(j+1)\varphi}\bar\partial\) for \(T^j\)\par
\noindent\textit{Critical reading.} \(e^{-(j+1)\varphi}\bar\partial^*e^{j\varphi}\bar\partial\) for \(T^j\)\par
\noindent\textit{Evidence.} Direct \(L^2\) pairing calculation: zero-form and one-form pairings scale by \(\exp((j+1)\varphi)\) and \(\exp(j\varphi)\). The written expression belongs to \(T^{-j-1}\).\par
\subsection*{P001211 (both; faithful error): 650, scalar heat curvature}
\noindent\textit{Transcription reads.} \(tK/3\), K called scalar curvature\par
\noindent\textit{Critical reading.} \(tR/6=tK_G/3\), where \(R=2K_G\)\par
\noindent\textit{Evidence.} The scalar heat coefficient and the accompanying curvature label conflict; both are clear in the manuscript.\par
\subsection*{P001211 (both; faithful error): 650, finite trace}
\noindent\textit{Transcription reads.} The factor \(1/(4\pi)\) disappears in the next displayed line\par
\noindent\textit{Critical reading.} Retain the overall \(1/(4\pi)\) with the stated Laplacian/curvature normalization\par
\noindent\textit{Evidence.} Algebraic consequence of the immediately preceding heat expansion; faithful omission.\par
\subsection*{P001211 (both; faithful error): 650, coordinate curvature}
\noindent\textit{Transcription reads.} \(K=\partial_z\partial_{\bar z}\log\theta\)\par
\noindent\textit{Critical reading.} For \(ds^2=\theta|dz|^2\), \(K_G=-2\theta^{-1}\partial_z\partial_{\bar z}\log\theta\)\par
\noindent\textit{Evidence.} Requires the inverse conformal factor and a declared real/complex convention.\par
\subsection*{P001211--P001212;P001218 (both; faithful error): 650/652, Liouville claim}
\noindent\textit{Transcription reads.} Computing the functional integral is equivalent to solving the Liouville equation\par
\noindent\textit{Critical reading.} The equation gives classical critical points; the quantum integral needs a measure, regularization, coupling and insertion data\par
\noindent\textit{Evidence.} The same conflation occurs twice; later rigorous sphere theory has additional explicit hypotheses.\par
\subsection*{P001212 (both; faithful error): 650, determinant powers}
\noindent\textit{Transcription reads.} \([\det\Delta_1/(\det\Delta_0)^{d/2}]^{1/2}\)\par
\noindent\textit{Critical reading.} For real fields: \((\det{}' P^*P)^{1/2}(\det{}'\Delta_0)^{-d/2}\)\par
\noindent\textit{Evidence.} Outer square root and scalar exponent mix conventions. Critical dimension 26 remains correct; Polyakov (1981), eqs. (20)--(25).\par
\subsection*{P001212 (both; faithful error): 650, zero-mode line}
\noindent\textit{Transcription reads.} Ratio of exterior powers of H1(T) and H1(O)\par
\noindent\textit{Critical reading.} Include H0 and power 13: \(\lambda(T)\otimes\lambda(\mathcal O)^{-13}\), or the consistently dual line\par
\noindent\textit{Evidence.} The anomaly polynomial has \(a_1=13\); a ratio with both exponents equal to 1 is not the determinant-line string measure.\par
\subsection*{P001214 (both; faithful error): 651, conformal slice}
\noindent\textit{Transcription reads.} One fixed constant-curvature metric times conformal factors and diffeomorphisms covers metrics\par
\noindent\textit{Critical reading.} Include Teichmueller parameters: complex dimension \(3g-3\) for \(g\ge2\)\par
\noindent\textit{Evidence.} One conformal class cannot represent all conformal structures; the next talk explicitly restores moduli.\par
\subsection*{P001215 (both; faithful error): 651, tangent decomposition}
\noindent\textit{Transcription reads.} Metric variations split into two displayed pieces\par
\noindent\textit{Critical reading.} Split into trace, image of trace-free diffeomorphism operator P, and ker(P*) (moduli)\par
\noindent\textit{Evidence.} At genus \(g\ge2\) the omitted transverse-traceless space has real dimension \(6g-6\).\par
\subsection*{P001215 (both; faithful error): 651, infinitesimal diffeomorphisms}
\noindent\textit{Transcription reads.} The full Lie derivative D is identified with d-bar\par
\noindent\textit{Critical reading.} Project to trace-free metric variations before identifying with the real Dolbeault operator\par
\noindent\textit{Evidence.} The full Lie derivative has a trace component \(2\operatorname{div}X\).\par
\subsection*{P001216 (both; faithful error): 651, final operator equality}
\noindent\textit{Transcription reads.} \(D^*D\) equals the ordinary Hodge Laplacian\par
\noindent\textit{Critical reading.} Keep the distinct ghost operator and its normalization\par
\noindent\textit{Evidence.} Flat principal symbol of full Lie derivative squared is \(2(|\xi|^2I+\xi\otimes\xi)\), whereas the Hodge symbol is \(|\xi|^2I\). \citet{polyakov1981quantum} explicitly distinguishes the vector ghost operator.\par
\subsection*{P001218 (both; faithful error): 652, critical dimension}
\noindent\textit{Transcription reads.} At d=26 the whole integrand is constant\par
\noindent\textit{Critical reading.} Only the Weyl-factor anomaly cancels; moduli and zero-mode dependence remain\par
\noindent\textit{Evidence.} The notebook has just fixed moduli. \citet{belavinknizhnik1986complex} describes nontrivial moduli dependence and degeneration poles.\par
\subsection*{P001219 (both; faithful error): 654, characteristic polynomial}
\noindent\textit{Transcription reads.} General p integrated as an Euler number\par
\noindent\textit{Critical reading.} For Euler--Gauss--Bonnet use the Euler Pfaffian in even dimension with orientation conventions\par
\noindent\textit{Evidence.} The introduction of \citet{cheegergromov1985characteristic} supplies the precise characteristic polynomial; general Pontryagin combinations are distinct.\par
\subsection*{P001220 (both; faithful error): 654, conic metric}
\noindent\textit{Transcription reads.} \(g=du^2+u^{-2}g_Y\)\par
\noindent\textit{Critical reading.} \(g=du^2+u^{2}g_Y\) for the exact metric cone\par
\noindent\textit{Evidence.} Negative exponent is visible in scan. \citet[\S1]{cheeger1983spectral} uses the positive exponent; the subsequent radial inverse-square operator agrees with the latter.\par
\subsection*{P001220 (both; ambiguity): 654, angular term}
\noindent\textit{Transcription reads.} A is simultaneously treated as a section-space and as the angular operator\par
\noindent\textit{Critical reading.} Specify a self-adjoint angular differential operator acting in L2 sections; retain its spectral shift\par
\noindent\textit{Evidence.} The terse wording does not determine a unique correction of every symbol.\par
\subsection*{P001220--P001221 (both; ambiguity): 654, Fuchs operator}
\noindent\textit{Transcription reads.} Mixed indices i,j and no summation limits\par
\noindent\textit{Critical reading.} Specify a Fuchs-type expansion, ellipticity, bundle, density and closed domain before asking for an index\par
\noindent\textit{Evidence.} Only the general regular singular index program can be reconstructed; \citet{lesch1997fuchs} is a precise later realization.\par
\subsection*{P001221 (both; ambiguity): 654, endpoint condition}
\noindent\textit{Transcription reads.} \(\lim u^{-1/2}\varphi=\lim u^{1/2}\varphi'\)\par
\noindent\textit{Critical reading.} No unique emendation without the angular mode and extension\par
\noindent\textit{Evidence.} For \(\varphi=a\sqrt u+b\sqrt u\log u\), finite limits force b=0 and their equality then forces a=0. Not the standard exclusion of the logarithmic branch.\par
\subsection*{P001221 (both; faithful error): 654, invariant operators}
\noindent\textit{Transcription reads.} Locally invariant D automatically has finite L2 kernels\par
\noindent\textit{Critical reading.} State ellipticity and the specific controlled cusp operator/realization hypotheses\par
\noindent\textit{Evidence.} The invariant zero operator is a counterexample; The publisher description of \citet{muller1987cusps} concerns generalized Dirac operators.\par
\subsection*{P001221 (both; faithful error): 654, number of cusps}
\noindent\textit{Transcription reads.} Only one cusp for every real quadratic Hilbert modular quotient\par
\noindent\textit{Critical reading.} One cusp requires class number one for the full group; torsion-free covers may have more cusps\par
\noindent\textit{Evidence.} Cusps are ideal classes. Full group may have orbifold points; do not silently invoke a smooth-manifold index.\par
\subsection*{P001221 (both; ambiguity): 654, L-series domain}
\noindent\textit{Transcription reads.} Unexplained final \(N_F\) in the summation domain\par
\noindent\textit{Critical reading.} Use \((M\setminus\{0\})/V\) for the standard unit-orbit series\par
\noindent\textit{Evidence.} Mathematical emendation: without quotient, norm-one units give infinitely repeated terms. Exact L-value coefficient remains unverified.\par
\subsection*{P001221 (both; faithful error): 655, Dolbeault exact sequence}
\noindent\textit{Transcription reads.} Final \(H^1(T)\)\par
\noindent\textit{Critical reading.} Final \(H^1(T^j)\)\par
\noindent\textit{Evidence.} The exponent \(j\) is missing in the scan, while every preceding term is for \(T^j\).\par
\subsection*{P001222 (both; ambiguity): 655, Quillen metric}
\noindent\textit{Transcription reads.} Squared cohomology norm divided by \(\det{}'\Delta_j\)\par
\noindent\textit{Critical reading.} Retain division: this is the dual of Quillen's original operator determinant-line convention\par
\noindent\textit{Evidence.} Quillen (1985), sections 1--3 checked in full. This row prevents an erroneous correction from division to multiplication.\par
\subsection*{P001223 (both; faithful error): 656, lattice heading}
\noindent\textit{Transcription reads.} Definition of a lattice includes the modular law\par
\noindent\textit{Critical reading.} The final displayed implication is the additional modular axiom\par
\noindent\textit{Evidence.} Separates general lattice axioms from modularity.\par
\subsection*{P001224 (both; faithful error): 656, first cubicle list}
\noindent\textit{Transcription reads.} List begins with 0\par
\noindent\textit{Critical reading.} The empty meet is the cubicle top \(1\); its relative bottom is the meet of all \(h_t\)\par
\noindent\textit{Evidence.} \citet[\S1.3]{gelfandponomarev1974modular}. Mapping every generator to V shows the relative cubicle bottom need not map to 0.\par
\subsection*{P001224 (both; faithful error): 656, recursion}
\noindent\textit{Transcription reads.} Second-level exclusion uses \(i_2\); Gamma indices are scrambled\par
\noindent\textit{Critical reading.} Use \(i_1\) in the second-level term and exclude \(i_j,i_{j+1}\) in Gamma; also require adjacent beta indices distinct\par
\noindent\textit{Evidence.} Exact source reconstruction from \citet[\S1.3, eqs.~(1.1)--(1.3)]{gelfandponomarev1974modular}.\par
\subsection*{P001224 (both; faithful error): 656, cubicle size}
\noindent\textit{Transcription reads.} \(2^n\)\par
\noindent\textit{Critical reading.} \(2^r\), r=number of generators\par
\noindent\textit{Evidence.} The theorem has \(2^r\); n is unexplained in these notes.\par
\subsection*{P001225 (both; ambiguity): 657, central convention}
\noindent\textit{Transcription reads.} Bracket factor \((j-i)\) with positive \((i^3-i)C\), followed by standard-c formulas\par
\noindent\textit{Critical reading.} If \(e_i=L_{-i}\), then \(C=-c\); convert later eigenvalues or change the initial bracket convention\par
\noindent\textit{Evidence.} Direct substitution into the usual Virasoro bracket. The algebra itself is valid; the linked formulas need one convention.\par
\subsection*{P001225 (both; ambiguity): 657, category name}
\noindent\textit{Transcription reads.} \ru{Оре}\par
\noindent\textit{Critical reading.} Retain uncertain reading; likely categoryO is not certain enough for a silent correction\par
\noindent\textit{Evidence.} The Russian transcription flags the same ambiguous word. No specific BGG paper is inferred.\par
\subsection*{P001225 (both; faithful error): 657, Kac polynomial}
\noindent\textit{Transcription reads.} \(\frac12(r^2-s^2)^2\)\par
\noindent\textit{Critical reading.} \(\frac1{16}(r^2-s^2)^2\)\par
\noindent\textit{Evidence.} Scan confirms \(1/2\); \citet[\S3.2, printed p.~58]{feiginfuchs1982invariant} has \(1/16\). At \(c=1,r=1,s=2\), the correct polynomial is \((h-1/4)^2\).\par
\subsection*{P001226 (both; faithful error): 657, root quadratic sum}
\noindent\textit{Transcription reads.} Inner \(r\in\Delta\)\par
\noindent\textit{Critical reading.} Inner \(r\in\mathbb Z\)\par
\noindent\textit{Evidence.} Outer sum is over roots; inner variable labels loop modes. The scan itself has Delta.\par
\subsection*{P001226 (both; faithful error): 657, enveloping algebra}
\noindent\textit{Transcription reads.} Infinite S(m) is an element of the ordinary enveloping algebra\par
\noindent\textit{Critical reading.} Use a completed enveloping algebra or its well-defined action on a restricted module\par
\noindent\textit{Evidence.} Normal-ordered infinite sums are not finite PBW sums; \citet[p.~133]{kent1985unitary} supplies the construction.\par
\subsection*{P001227 (both; faithful error): 657, Sugawara denominator}
\noindent\textit{Transcription reads.} \(e_m=S(m)/(2(g+k))\) without restriction\par
\noindent\textit{Critical reading.} Assume \(k\ne-h^\vee\), with invariant form and level normalized consistently\par
\noindent\textit{Evidence.} At the critical level the denominator vanishes; the raw Killing form is not the normalization with long-root length 2.\par
\subsection*{P001227 (both; faithful error): 657, sl2 highest weight}
\noindent\textit{Transcription reads.} \(k=k_1+k_2\), \(h=k_1(k_1+1)/(k+2)\)\par
\noindent\textit{Critical reading.} Affine Dynkin label \(\ell\) gives \(h=\ell(\ell+2)/(4(k+2))\); spin \(j=\ell/2\) gives \(j(j+1)/(k+2)\)\par
\noindent\textit{Evidence.} The same \(k_1\) cannot be both a Dynkin label in level sum and a spin without a factor 2 conversion.\par
\subsection*{P001227 (both; faithful error): 657, commuting subalgebra}
\noindent\textit{Transcription reads.} \(\widehat{\mathfrak g}_i\)\par
\noindent\textit{Critical reading.} \(\widehat{\mathfrak g}_1\)\par
\noindent\textit{Evidence.} Only subalgebra 1 is defined; i is the running Virasoro mode.\par
\subsection*{P001227 (both; faithful error): 657, Selberg integral}
\noindent\textit{Transcription reads.} Vandermonde power \(\rho\), limits 0 and 1, printed Gamma product\par
\noindent\textit{Critical reading.} Use power \(2\rho\) on \(0<t_1<\cdots<t_m<1\); or cube with absolute values and m! times the printed product\par
\noindent\textit{Evidence.} \citet[eqs.~(1.1)--(1.2)]{forresterwarnaar2008selberg}. This is a paired exponent/domain correction; neither alone fixes the formula.\par
\subsection*{P001228 (both; faithful error): 657, mixed-product range}
\noindent\textit{Transcription reads.} Upper bound \(v=m+1\)\par
\noindent\textit{Critical reading.} Upper bound \(v=n+m\)\par
\noindent\textit{Evidence.} Two species have sizes \(n\) and \(m\); the erroneous endpoint is visible in scan.\par
\subsection*{P001228 (both; faithful error): 657, last differential}
\noindent\textit{Transcription reads.} \(dt_1\cdots dt_n\)\par
\noindent\textit{Critical reading.} \(dt_1\cdots dt_{n+m}\)\par
\noindent\textit{Evidence.} The integrand has n+m screening variables.\par
\subsection*{P001228 (both; faithful error): 657, correlation order}
\noindent\textit{Transcription reads.} Two-point correlation function\par
\noindent\textit{Critical reading.} Four-point conformal block before the nonchiral combination\par
\noindent\textit{Evidence.} The same page then calls it four-point and defines the cross-ratio; two points have no such modulus.\par
\subsection*{P001228 (both; faithful error): 657, screening integral}
\noindent\textit{Transcription reads.} Unqualified ordinary real integral from 0 to 1 with mixed double poles\par
\noindent\textit{Critical reading.} Specify disjoint/compatible complex contours or twisted cycles and charge-neutrality/external-charge data\par
\noindent\textit{Evidence.} \citet[\S3]{dotsenkofateev1984conformal}. For both species present, crossing \(t_u=t_v\) gives a nonintegrable real double pole.\par
\subsection*{P001228 (both; ambiguity): 657, continuation claim}
\noindent\textit{Transcription reads.} Analytic continuation from nonnegative integers \(n,m\) gives the real functions\par
\noindent\textit{Critical reading.} State a uniqueness class or independent definition; the discrete data alone do not select a continuation\par
\noindent\textit{Evidence.} Adding \(\sin(\pi n)H(n,m)\) preserves every nonnegative integer pair. This records underdetermination, not a new open problem.\par
\subsection*{P001228-P001232 (both; provenance): Catalogue bounds, record 658}
\noindent\textit{Transcription reads.} P001228-P001232\par
\noindent\textit{Critical reading.} P001229 through the upper part of P001231\par
\noindent\textit{Evidence.} Vershik heading on P001229; Feigin heading midway down P001231. P001228 belongs to preceding record.\par
\subsection*{P001230-P001237 (both; provenance): Catalogue bounds, record 659}
\noindent\textit{Transcription reads.} P001230-P001237\par
\noindent\textit{Critical reading.} Mid-P001231 through P001235\par
\noindent\textit{Evidence.} Feigin heading on P001231; 18 November continuation begins P001236.\par
\subsection*{P001235-P001239 (both; provenance): Catalogue bounds, record 660}
\noindent\textit{Transcription reads.} P001235-P001239\par
\noindent\textit{Critical reading.} P001236 through the upper part of P001238\par
\noindent\textit{Evidence.} 18 November heading on P001236; Zelevinsky heading halfway down P001238.\par
\subsection*{P001237-P001242 (both; provenance): Catalogue bounds, record 661}
\noindent\textit{Transcription reads.} P001237-P001242\par
\noindent\textit{Critical reading.} Mid-P001238 through P001240\par
\noindent\textit{Evidence.} Zelevinsky heading on P001238; 25 November Feigin heading starts P001241.\par
\subsection*{P001240-P001246 (both; provenance): Catalogue bounds, record 662}
\noindent\textit{Transcription reads.} P001240-P001246\par
\noindent\textit{Critical reading.} P001241-P001244\par
\noindent\textit{Evidence.} 25 November heading starts P001241; next dated speaker starts P001245.\par
\subsection*{P001229 (both; ambiguity): Matrix sign-class problem}
\noindent\textit{Transcription reads.} All minors nonzero with the same signs\par
\noindent\textit{Critical reading.} All maximal 3 by 3 minors, if the intended problem is a uniform rank-three oriented matroid\par
\noindent\textit{Evidence.} The wording is not silently changed; minors of smaller orders would impose additional constraints.\par
\subsection*{P001230 (both; ambiguity): Three-polytope realization space}
\noindent\textit{Transcription reads.} The space of fixed-type polytopes is contractible\par
\noindent\textit{Critical reading.} A specified normalized realization space is contractible\par
\noindent\textit{Evidence.} Without affine normalization an affine frame contributes noncontractible group/orientation data. Richter-Gebert--Ziegler section 2 states a normalized result.\par
\subsection*{P001230 (both; ambiguity): Simplicial four-polytope claim}
\noindent\textit{Transcription reads.} More or less any topology\par
\noindent\textit{Critical reading.} Retain the hedge and distinguish the conjectural fixed-dimension claim from proved variable-dimensional universality\par
\noindent\textit{Evidence.} The Russian qualification is explicit. AP Theorem 2 grows dimension; Conjecture 4 retains dimension four.\par
\subsection*{P001230 (both; faithful error): Geometric lattice definition}
\noindent\textit{Transcription reads.} Finite (that is, every element is a finite join of atoms)\par
\noindent\textit{Critical reading.} Finite and atomistic, as separate conditions, together with upper semimodularity\par
\noindent\textit{Evidence.} Finiteness alone does not imply atomisticity. The source parenthesis conflates the two properties.\par
\subsection*{P001230 (both; ambiguity): Fano diagram}
\noindent\textit{Transcription reads.} Seven vertices with six visible line triples\par
\noindent\textit{Critical reading.} Named Fano configuration requires all seven line triples; do not draw the missing stroke into the source\par
\noindent\textit{Evidence.} PNG inspected. The characteristic-two theorem applies to the named Fano matroid, not merely the six visible incidences.\par
\subsection*{P001230-P001231 (both; ambiguity): Faithful lattice representation example}
\noindent\textit{Transcription reads.} Question-marked fragment m=4, E,F in R6\par
\noindent\textit{Critical reading.} Leave the intended dimensional example unresolved\par
\noindent\textit{Evidence.} The scan supplies no complete lattice or proposition to identify; join preservation alone is not a full lattice embedding.\par
\subsection*{P001231 (both; faithful error): Equivalence-relation join}
\noindent\textit{Transcription reads.} Supremum described by union\par
\noindent\textit{Critical reading.} Equivalence closure of the union\par
\noindent\textit{Evidence.} A raw union of equivalence relations need not be transitive.\par
\subsection*{P001231 (both; ambiguity): Scattering frame}
\noindent\textit{Transcription reads.} The total momentum is zero\par
\noindent\textit{Critical reading.} Zero total spatial momentum in the centre-of-mass frame, or zero sum after outgoing four-momenta are signed incoming\par
\noindent\textit{Evidence.} The adjacent equations otherwise use positive-energy two-in/two-out conservation.\par
\subsection*{P001231-P001232 (both; ambiguity): Amplitude axiom}
\noindent\textit{Transcription reads.} Analytic, with simple poles\par
\noindent\textit{Critical reading.} Meromorphic; holomorphic away from its specified poles\par
\noindent\textit{Evidence.} An entire analytic function cannot have the listed finite poles.\par
\subsection*{P001232 (both; faithful error): Veneziano uniqueness}
\noindent\textit{Transcription reads.} The listed conditions determine a unique amplitude\par
\noindent\textit{Critical reading.} They determine neither its scalar normalization nor, without a global pole-expansion prescription, its entire part\par
\noindent\textit{Evidence.} Multiplication by any nonzero constant preserves all listed conditions; suitable symmetric polynomial additions preserve local pole data.\par
\subsection*{P001232 (both; ambiguity): Momentum diagram}
\noindent\textit{Transcription reads.} Two outgoing arrows and sum of all four momenta equals zero\par
\noindent\textit{Critical reading.} Use signed all-incoming momenta for the zero-sum identity\par
\noindent\textit{Evidence.} PNG inspected; an outgoing momentum must be reversed when inserted into an all-incoming sum.\par
\subsection*{P001232 (both; faithful error): Number of oscillators}
\noindent\textit{Transcription reads.} Finitely many oscillators; finitely many do not suffice\par
\noindent\textit{Critical reading.} Finitely many target-space species, countably many frequency modes\par
\noindent\textit{Evidence.} The two source statements contradict each other literally; the immediately following m in Z formulas fix the intended model.\par
\subsection*{P001232-P001233 (both; faithful error): Target index range}
\noindent\textit{Transcription reads.} \(\mu=0,\ldots,d\)\par
\noindent\textit{Critical reading.} \(\mu=0,\ldots,d-1\) for a d-dimensional target\par
\noindent\textit{Evidence.} The printed central charge and metric dimension are d; including both endpoints gives d+1 species.\par
\subsection*{P001233 (both; faithful error): Coordinate Fourier summation}
\noindent\textit{Transcription reads.} \(a_m\) multiplying \(\exp(-in\theta)/n\) in a sum over n\par
\noindent\textit{Critical reading.} \(a_n\) with the same summation index n\par
\noindent\textit{Evidence.} Both the oscillator commutator and derivative Q prime=P require a consistent dummy index; the source has the mismatch.\par
\subsection*{P001233 (both; faithful error): Coordinate-coordinate commutator}
\noindent\textit{Transcription reads.} A function called Heaviside\par
\noindent\textit{Critical reading.} An odd sign kernel on a lifted interval, with zero-mode monodromy\par
\noindent\textit{Evidence.} Antisymmetry of a commutator rules out the plain Heaviside function. Direct Fourier calculation gives -i pi sign(theta-theta prime) for the stated Q convention.\par
\subsection*{P001234 (both; ambiguity): Exponential vertex}
\noindent\textit{Transcription reads.} exp(i k Q) without normal-ordering qualification\par
\noindent\textit{Critical reading.} A normal-ordered exponential, with a specified zero-mode convention\par
\noindent\textit{Evidence.} The factorized creation/annihilation expression is finite formally; the ordinary exponential has a divergent self-contraction.\par
\subsection*{P001235 (both; faithful error): Two-vertex final bra}
\noindent\textit{Transcription reads.} \(p+k_1+k\)\par
\noindent\textit{Critical reading.} \(p+k_1+k_2\)\par
\noindent\textit{Evidence.} The two successive momentum-shifting vertices add k1 and k2. The second subscript is absent in the source.\par
\subsection*{P001236 (both; faithful error): Spurion definition (2)}
\noindent\textit{Transcription reads.} A physical vector equal to one \(L_i v'\) for some \(i<0\)\par
\noindent\textit{Critical reading.} Intersection with the linear sum of negative-mode images\par
\noindent\textit{Evidence.} A union of individual images is not generally a vector subspace and cannot directly be a quotient denominator.\par
\subsection*{P001236 (both; faithful error): Weight equation (5)}
\noindent\textit{Transcription reads.} \(L_{-1}v'=v\) implies \(\operatorname{weight}(v')=2\)\par
\noindent\textit{Critical reading.} \(\operatorname{weight}(v')=0\) when \(\operatorname{weight}(v)=1\) and v is nonzero\par
\noindent\textit{Evidence.} PNG and EvidenceBlock both confirm two. The bracket \([L_0,L_{-1}]=L_{-1}\) raises weight by one.\par
\subsection*{P001236 (both; faithful error): No-ghost parenthesis}
\noindent\textit{Transcription reads.} Absence of tachyons\par
\noindent\textit{Critical reading.} Absence of negative-norm ghosts\par
\noindent\textit{Evidence.} The source clearly reads tachyons. The bosonic no-ghost theorem does not remove a negative-mass-squared ground state.\par
\subsection*{P001236 (both; faithful error): Strict positivity claim}
\noindent\textit{Transcription reads.} Every physical quotient vector has norm >0\par
\noindent\textit{Critical reading.} Every nonzero quotient class at nonzero momentum has positive norm, for the stated critical theory\par
\noindent\textit{Evidence.} Zero has norm zero; zero momentum has an exceptional indefinite level-one sector. Frenkel Theorem 4.1 explicitly removes it.\par
\subsection*{P001237 (both; faithful error): Creation exponential (6)}
\noindent\textit{Transcription reads.} \(\sum a_{-n}z^n\) without \(1/n\)\par
\noindent\textit{Critical reading.} \(\sum a_{-n}z^n/n\)\par
\noindent\textit{Evidence.} PNG confirms missing denominator. Its absence destroys the required Heisenberg field commutator and product contraction.\par
\subsection*{P001237 (both; faithful error): Correlation product (8)}
\noindent\textit{Transcription reads.} Product indexed \(j>1\)\par
\noindent\textit{Critical reading.} Product over \(1\le i<j\le n\)\par
\noindent\textit{Evidence.} PNG confirms the digit 1, leaving i unbound; the contraction is pairwise.\par
\subsection*{P001237 (both; ambiguity): Vertex residues (10)}
\noindent\textit{Transcription reads.} Bare oscillator/shift X integrated with dz/z\par
\noindent\textit{Critical reading.} Use full cylinder field \(X=z^hY\), so for \(h=1\) \(\operatorname{Res}X\,dz/z=\operatorname{Res}Y\,dz\)\par
\noindent\textit{Evidence.} The missing zero-mode powers affect Virasoro covariance. This is a necessary convention qualification, not restoration of invisible text.\par
\subsection*{P001237 (both; faithful error): Cartan condition (11)}
\noindent\textit{Transcription reads.} \(A_{ij}<0\) for unrestricted i,j\par
\noindent\textit{Critical reading.} \(A_{ii}=2\) and \(A_{ij}\le0\) for i not equal to j\par
\noindent\textit{Evidence.} The unrestricted inequality contradicts the preceding diagonal norm two and excludes legitimate zero off-diagonal entries.\par
\subsection*{P001237 (both; faithful error): Gelfand Monster remark}
\noindent\textit{Transcription reads.} Largest finite simple group\par
\noindent\textit{Critical reading.} Largest sporadic finite simple group\par
\noindent\textit{Evidence.} Finite simple groups have unbounded orders; the Monster is the intended sporadic group, with order about \(8.08\times10^{53}\).\par
\subsection*{P001238 (both; ambiguity): Transverse oscillator basis}
\noindent\textit{Transcription reads.} Any basis of c-perp modulo c gives standard oscillators\par
\noindent\textit{Critical reading.} An orthonormal transverse basis and a fixed normalized value of (c,a0)\par
\noindent\textit{Evidence.} Without these conditions the commutator contains a Gram matrix and the central operator (c,a0).\par
\subsection*{P001238-P001239 (both; ambiguity): Schubert homology basis}
\noindent\textit{Transcription reads.} The cells form the ordinary homology basis\par
\noindent\textit{Critical reading.} Fundamental classes of the cell closures form the basis\par
\noindent\textit{Evidence.} An affine cell itself is noncompact; its closure defines the ordinary compact homology cycle.\par
\subsection*{P001239 (both; faithful error): BGG top class}
\noindent\textit{Transcription reads.} Product \((x_i-x_j)\) equals the top monomial modulo J\par
\noindent\textit{Critical reading.} (1/n!) product \((x_i-x_j)\) equals the top monomial modulo J\par
\noindent\textit{Evidence.} PNG confirms no factor. BGG Theorem 3.15 has 1/|W|; \(n=2\) gives \(x_1-x_2=2x_1\) modulo J.\par
\subsection*{P001239 (both; faithful error): Rho exponent and factorial}
\noindent\textit{Transcription reads.} m(n-1)/2\par
\noindent\textit{Critical reading.} n(n-1)/2\par
\noindent\textit{Evidence.} PNG confirms m in both positions, but the flag rank is n. BGG Corollary 3.16 uses the number of positive roots.\par
\subsection*{P001240 (both; faithful error): Divided-difference denominator}
\noindent\textit{Transcription reads.} \(x_{i+1}-x_i\) below \(f-s_i f\)\par
\noindent\textit{Critical reading.} \(x_i-x_{i+1}\) for the listed positive Schubert polynomial conventions\par
\noindent\textit{Evidence.} PNG confirms the order. At \(n=2\) the recorded operator sends \(x_1\) to -1, contradicting \(P_{\mathrm{id}}=1\).\par
\subsection*{P001240 (both; faithful error): Final adjacent transposition and operator}
\noindent\textit{Transcription reads.} Second subscript \(i_{\ell(w)+1}\)\par
\noindent\textit{Critical reading.} \(i_{\ell(w)}+1\)\par
\noindent\textit{Evidence.} The terminal factor must be the adjacent transposition \((i_\ell,i_\ell+1)\); the plus one is at the wrong subscript level in both displays.\par
\subsection*{P001239-P001240 (both; ambiguity): Longest-element operator indexing}
\noindent\textit{Transcription reads.} \(A_{w_0w^{-1}}\) acting on the left\par
\noindent\textit{Critical reading.} Specify action convention; standard left composition uses \(\partial_{w^{-1}w_0}\)\par
\noindent\textit{Evidence.} The two index orders can reflect opposite action conventions. Do not change it as a transcription error without fixing the convention.\par
\subsection*{P001241 (both; faithful error): Sugawara definition (1)}
\noindent\textit{Transcription reads.} \(L_i=:\!\sum a_m a_n\!:\)\par
\noindent\textit{Critical reading.} \(L_i=\tfrac12:\!\sum a_m a_n\!:\)\par
\noindent\textit{Evidence.} PNG confirms absence of the half; the correct normalization is already written on P001233.\par
\subsection*{P001241 (both; ambiguity): Normal-ordering explanation}
\noindent\textit{Transcription reads.} Only adds a constant\par
\noindent\textit{Critical reading.} Finite cutoffs differ by scalars; the infinite scalar requires a prescription, while normal ordering defines locally finite modes\par
\noindent\textit{Evidence.} The unrenormalized zero-mode sum contains divergent scalar contributions.\par
\subsection*{P001241 (both; ambiguity): Classical vector fields}
\noindent\textit{Transcription reads.} \(L_n=e^{int}\,d/dt\) with the usual \((m-n)\) bracket\par
\noindent\textit{Critical reading.} Use \(i e^{int}\,d/dt\), or retain an extra i in the bracket\par
\noindent\textit{Evidence.} Direct vector-field computation distinguishes the conventions.\par
\subsection*{P001241 (both; faithful error): Norm-only Fock-module assertion}
\noindent\textit{Transcription reads.} \(H_p\) is isomorphic to \(H_{p'}\) whenever \(p^2=(p')^2\)\par
\noindent\textit{Critical reading.} Equal isometry orbit suffices; exclude zero versus nonzero null momentum\par
\noindent\textit{Evidence.} At level one the map L1 has rank zero at p=0 and rank one at nonzero null p, an isomorphism invariant.\par
\subsection*{P001242 (both; faithful error): Repeated no-ghost theorem}
\noindent\textit{Transcription reads.} Tilde physical space is positive definite and generated by one vacuum\par
\noindent\textit{Critical reading.} Nonzero-momentum physical quotient by null states is positive definite; a vacuum generates its chosen DDF momentum tower\par
\noindent\textit{Evidence.} The inserted tilde is visible. The prequotient has spurious null states; zero momentum is exceptional.\par
\subsection*{P001242 (both; faithful error): Mode relation (12)}
\noindent\textit{Transcription reads.} \(X(z)=\sum X_i z^i\) and \([a_n,X_i]=pX_{i+n}\)\par
\noindent\textit{Critical reading.} With the given expansion: \([a_n,X_i]=pX_{i-n}\)\par
\noindent\textit{Evidence.} Coefficient comparison in \([a_n,X(z)]=p z^nX(z)\). The plus sign belongs to the negative-exponent convention.\par
\subsection*{P001242 (both; ambiguity): Unimodularity determinant}
\noindent\textit{Transcription reads.} det Gram=1\par
\noindent\textit{Critical reading.} Absolute determinant one in arbitrary nondegenerate signature; +1 in the positive-definite case\par
\noindent\textit{Evidence.} The lecture changes between Lorentzian and positive-definite forms; the signature assumption must accompany this determinant.\par
\subsection*{P001243 (both; ambiguity): Divisibility-by-eight antecedent}
\noindent\textit{Transcription reads.} Such exist only in dimensions divisible by eight, immediately after roots\par
\noindent\textit{Critical reading.} Positive-definite even unimodular lattices have rank divisible by eight\par
\noindent\textit{Evidence.} Rank-one A1 already has norm-two roots. The antecedent must refer to the preceding lattice conditions.\par
\subsection*{P001243 (both; faithful error): Cocycle paragraph}
\noindent\textit{Transcription reads.} Q not necessarily even, with \(\epsilon(a,b)\epsilon(b,a)=(-1)^{(a,b)}\) for all a,b\par
\noindent\textit{Critical reading.} Q must be even, and epsilon must satisfy a cocycle law (bimultiplicativity suffices)\par
\noindent\textit{Evidence.} Put a=b: the left side is one, forcing even norms. The two displayed sign identities alone do not imply associativity/Jacobi.\par
\subsection*{P001243 (both; ambiguity): Finite Lie algebra construction}
\noindent\textit{Transcription reads.} Construction of simple Lie algebras for an arbitrary positive lattice\par
\noindent\textit{Critical reading.} Reductive in general; simple when roots span and form an irreducible system\par
\noindent\textit{Evidence.} The complement of the root span is central; a rootless lattice produces an abelian algebra.\par
\subsection*{P001244 (both; ambiguity): Root-current brackets}
\noindent\textit{Transcription reads.} Only positive and -1 inner products displayed\par
\noindent\textit{Critical reading.} Retain the partial formulas and add the opposite-root Cartan-plus-central case in reconstruction\par
\noindent\textit{Evidence.} The omitted case is needed to explain the projective current representation. It is not an assertion that the partial formulas are false.\par
\subsection*{P001245 (both; faithful error): Equation (1), derivative normalization}
\noindent\textit{Transcription reads.} \(-\partial\bar\partial=-\Delta\)\par
\noindent\textit{Critical reading.} With standard Wirtinger derivatives, \(\partial\bar\partial=\Delta_{\mathbb R^2}/4\); alternatively state a rescaled Laplacian convention.\par
\noindent\textit{Evidence.} The scan and Russian both contain the equality. This is a convention defect, not an OCR correction.\par
\subsection*{P001245 (both; faithful error): Equations (1)--(2), exclusion of Lax pairs}
\noindent\textit{Transcription reads.} \ru{не бывает пар Лакса}; only rigid motions\par
\noindent\textit{Critical reading.} Restrict the intended operator class, or retain as a heuristic; the unrestricted assertion is false.\par
\noindent\textit{Evidence.} The scan confirms the wording. The \(y\)-independent KdV counterexample in 1985-663.1 and \texttt{code/\allowbreak{}check\_\allowbreak{}mathematics.py} refutes the unqualified statement.\par
\subsection*{P001246 (both; faithful error): Equations (9)--(10)}
\noindent\textit{Transcription reads.} \(L=\partial\bar\partial+A\partial+B\), \(A=\bar\partial\log c_2\)\par
\noindent\textit{Critical reading.} For the written asymptotics and \(c_1=1\): \(L=\partial\bar\partial-(\partial\log c_2)\bar\partial+b\).\par
\noindent\textit{Evidence.} Both original derivative placements were inspected in the scan. The order-\(k_1\) and order-\(k_2\) terms at the two punctures give the correction; overall sign and local-parameter rescalings must be kept consistent.\par
\subsection*{P001247 (both; faithful error): Equation (22), auxiliary coefficient and B}
\noindent\textit{Transcription reads.} \(\partial u=3\bar\partial V\), \(B=\partial\bar u+\bar\partial u\)\par
\noindent\textit{Critical reading.} With \(A=\partial^3+u\partial\): \(\bar\partial u=3\partial V\), \(B=\partial u+\bar\partial\bar u\).\par
\noindent\textit{Evidence.} The scan confirms the transcript. Comparing the coefficients of \(\partial^2\), \(\bar\partial^2\), and \(\partial\bar\partial\) in the commutator fixes the placements. Two independent symbolic checks are included.\par
\subsection*{P001248 (both; provenance): Knizhnik heading}
\noindent\textit{Transcription reads.} \ru{Д. Книжник. (Дима)} / D. Knizhnik (Dima)\par
\noindent\textit{Critical reading.} Keep the notebook wording diplomatically; identify catalogue record 664 as V. G. Knizhnik, without silently changing the handwritten initial or nickname.\par
\noindent\textit{Evidence.} The original heading visibly reads Д., and the supplied catalogue names V. G. Knizhnik. The title and content match the latter\textquotesingle s publications. The nickname conflict is unresolved provenance, not evidence for a different mathematician.\par
\subsection*{P001248 (both; faithful error): Equation (5), Gram pairing}
\noindent\textit{Transcription reads.} \(\overline{f^i}f^i\)\par
\noindent\textit{Critical reading.} \(\overline{f^i}f^j\)\par
\noindent\textit{Evidence.} The scan confirms both superscripts are i. The left-hand indices and the Gram-matrix transformation law require j in the second factor.\par
\subsection*{P001249 (both; faithful error): Products of one-forms}
\noindent\textit{Transcription reads.} \ru{Всегда они так получаются}\par
\noindent\textit{Critical reading.} For genus two, or for a nonhyperelliptic curve of genus at least three.\par
\noindent\textit{Evidence.} The scan confirms the universal wording. Max Noether\textquotesingle s theorem has the stated exception; the explicit hyperelliptic basis gives a cokernel of dimension \(g-2\).\par
\subsection*{P001249 (both; faithful error): Equation (6), dual basis}
\noindent\textit{Transcription reads.} \(\int\eta_i f^i=\delta_i^j\)\par
\noindent\textit{Critical reading.} \(\int\eta_i f^j=\delta_i^j\)\par
\noindent\textit{Evidence.} The scan has i on f. The free index j on the right fixes the intended duality statement.\par
\subsection*{P001249 (both; faithful error): Fundamental-domain figure, right vertical axis label}
\noindent\textit{Transcription reads.} \(-1/2\)\par
\noindent\textit{Critical reading.} \(+1/2\)\par
\noindent\textit{Evidence.} The manuscript itself repeats -1/2. The right vertical boundary is the translate of the left by +1.\par
\subsection*{P001250 (both; faithful error): Wedge identity after equation (13)}
\noindent\textit{Transcription reads.} \(d\bar\eta\wedge d\eta=\tfrac14d\tau\wedge d\bar\tau\)\par
\noindent\textit{Critical reading.} \(d\bar\eta\wedge d\eta=-\tfrac14d\tau\wedge d\bar\tau\)\par
\noindent\textit{Evidence.} The displayed fractional-linear formula gives \(d\tau=-2i\,d\eta\). The wedge orders in both languages agree with the scan. Independently checked using real coordinates.\par
\subsection*{P001251 (both; ambiguity): Equation (20), determinant normalization}
\noindent\textit{Transcription reads.} \(\det^{\prime}\Delta=y|\eta|^4\), up to an unstated constant\par
\noindent\textit{Critical reading.} Specify the unit-area metric, or replace the right side by \(C y^2|\eta|^4\) for \(|dz|^2\).\par
\noindent\textit{Evidence.} The scan and translation agree. Area scaling of the spectral determinant distinguishes the two formulas.\par
\subsection*{P001251 (both; faithful error): Stress-tensor definition}
\noindent\textit{Transcription reads.} \(g_{ab}\partial^a\psi\partial^b\psi\)\par
\noindent\textit{Critical reading.} \(g_{ab}g^{cd}\partial_c\psi\partial_d\psi\)\par
\noindent\textit{Evidence.} The scan explicitly has the repeated a,b. The trace contraction needs dummy indices different from the free a,b.\par
\subsection*{P001251 (both; faithful error): Metric-variation identity}
\noindent\textit{Transcription reads.} \(\delta(\sqrt g g^{ab}\partial_a\psi\partial_b\psi)=\delta g^{ab}T_{ab}\)\par
\noindent\textit{Critical reading.} Insert \(\sqrt g\) on the right, or define T as a tensor density.\par
\noindent\textit{Evidence.} Direct variation of det g supplies the density factor; the source does not define a densitized tensor.\par
\subsection*{P001252 (both; faithful error): Odd characteristic and zero orders}
\noindent\textit{Transcription reads.} Any odd characteristic; all zeros of order 2\par
\noindent\textit{Critical reading.} A nonsingular odd characteristic; all zeros of even order.\par
\noindent\textit{Evidence.} Theta gradients may vanish at singular odd characteristics, and even multiplicity need not equal two.\par
\subsection*{P001252 (both; faithful error): Sentence following equation (28)}
\noindent\textit{Transcription reads.} All of this is a quadratic differential\par
\noindent\textit{Critical reading.} The Schwarzian term has an inhomogeneous coordinate transformation; specify a projective connection or a subtraction convention.\par
\noindent\textit{Evidence.} The source gives ordinary z-derivatives and says the omega symbols omit dz. The Schwarzian composition identity prevents the unqualified tensor claim.\par
\subsection*{P001253 (both; faithful error): Equation (3), one-loop density}
\noindent\textit{Transcription reads.} \(|dq/q|^2\Delta(q)^{-2}\log|q|^{-13}\)\par
\noindent\textit{Critical reading.} \(C|dq/q|^2|\Delta(q)|^{-2}(-\log|q|)^{-14}\) in the conventions of P001256.\par
\noindent\textit{Evidence.} The scan confirms exponent -13 and no modulus on Delta. The later correct tau-density and \(q=e^{2\pi i\tau}\) determine the corrected exponent and absolute square.\par
\subsection*{P001253 (both; ambiguity): Parenthesis after the definition of L}
\noindent\textit{Transcription reads.} Holomorphic structure induced by the metric\par
\noindent\textit{Critical reading.} The oriented surface metric induces a complex structure on X; a holomorphic structure on the arbitrary bundle L is additional data.\par
\noindent\textit{Evidence.} The same wording occurs in Russian and the scan. A base metric does not define a Dolbeault operator on an unspecified smooth line bundle.\par
\subsection*{P001254 (both; faithful error): Equation (7), Quillen metric}
\noindent\textit{Transcription reads.} Hodge norm times \(\sqrt{\det^{\prime}\Delta}\) on \(\det H^0/\det H^1\)\par
\noindent\textit{Critical reading.} Hodge norm divided by \(\sqrt{\det^{\prime}\Delta}\), or dualize the determinant line.\par
\noindent\textit{Evidence.} The scan confirms multiplication. Quillen uses the dual line. The independent finite-dimensional acyclic map \(T=a\) sends 1 to \(a^{-1}e/f\) and fixes the reciprocal.\par
\subsection*{P001255 (both; faithful error): Determinant of the one-dimensional map}
\noindent\textit{Transcription reads.} \(\det T\in E_1/E_2\)\par
\noindent\textit{Critical reading.} \(\det T\in E_2/E_1\); its inverse lies in \(E_1/E_2\).\par
\noindent\textit{Evidence.} The source fraction is clear. A map \(E_1\text{ to }E_2\) belongs to \(E_1^*\otimes E_2\).\par
\subsection*{P001255 (both; ambiguity): Bernoulli-polynomial label}
\noindent\textit{Transcription reads.} \(6j^2-6j+1\) is the second Bernoulli polynomial\par
\noindent\textit{Critical reading.} \(6j^2-6j+1=6B_2(j)\) with \(B_2(j)=j^2-j+1/6\).\par
\noindent\textit{Evidence.} This is the standard normalization; the source suppresses the factor six.\par
\subsection*{P001256 (both; faithful error): Equation (18), exponent index}
\noindent\textit{Transcription reads.} \(b_{j_i}\)\par
\noindent\textit{Critical reading.} \(b_i\), in the notation of (14) and (19).\par
\noindent\textit{Evidence.} The scan has the nested index. The preceding cancellation sum and following theorem use coefficients indexed by the list position i.\par
\subsection*{P001256 (both; faithful error): Uniqueness following theorem (19)}
\noindent\textit{Transcription reads.} Zero curvature uniquely determines the canonical metric\par
\noindent\textit{Critical reading.} It can at most determine the metric up to a positive constant after a global constancy argument.\par
\noindent\textit{Evidence.} Constant rescaling preserves zero curvature. P001257 invokes global moduli but still leaves the normalization free.\par
\subsection*{P001256 (en; translation): Label below equation (20)}
\noindent\textit{Transcription reads.} \(\Delta_j\) is the Laplacian on j-forms\par
\noindent\textit{Critical reading.} The Laplacian on j-differentials, i.e. sections of \(K^j\).\par
\noindent\textit{Evidence.} The Russian reads j-differentials; negative j=-1 in the displayed formula excludes ordinary differential-form degree.\par
\subsection*{P001256 (both; ambiguity): Genus-two formula}
\noindent\textit{Transcription reads.} Powers of \(\Im\tau\), \(d^2\tau\), and \(\prod_m\theta_m^2\) with no characteristic restriction\par
\noindent\textit{Critical reading.} Powers of \(\det\Im\tau\), volume in three complex variables, and the product over ten even characteristics.\par
\noindent\textit{Evidence.} The scan omits these qualifications. The dimensional and modular-weight requirements and vanishing of odd theta constants determine them.\par
\subsection*{P001257 (both; faithful error): Last uniqueness conclusion}
\noindent\textit{Transcription reads.} Unique metric and unique holomorphic section\par
\noindent\textit{Critical reading.} Unique up to a positive metric scale and a nonzero section scalar, respectively.\par
\noindent\textit{Evidence.} The source claim of global constancy does not fix a nonzero constant. The precise comparison theorem contains a genus-dependent normalization.\par
\subsection*{P001258-P001259 (both; provenance): Record ownership in the supplied catalogue page spans}
\noindent\textit{Transcription reads.} Pages through P001259 assigned to record 666\par
\noindent\textit{Critical reading.} Record 666 ends on P001257; separately headed records 667 and 668 start on P001258 and P001259.\par
\noindent\textit{Evidence.} Both scan headings name different speakers and dates. The archive all-links list has those adjacent entries. They are included in the adjacent-evidence supplement without changing the prescribed 26-record core.\par
\subsection*{P001258 (both; provenance): Kontsevich date}
\noindent\textit{Transcription reads.} 24 February 1985 in the scan; 24 February 1987 in the archive\par
\noindent\textit{Critical reading.} Preserve both attestations; 24 February 1986 is a plausible contextual inference, not an established reading.\par
\noindent\textit{Evidence.} The scan clearly ends in 85; the next leaf is dated 24 Feb 86. Calendar check: 24 Feb 85 Sunday, 24 Feb 86 Monday, 24 Feb 87 Tuesday. Monday regularity supports but does not prove the inferred year.\par
\subsection*{P001259 (both; provenance): Morozov date}
\noindent\textit{Transcription reads.} Catalogue 24 February 1987\par
\noindent\textit{Critical reading.} Manuscript 24 February 1986, while retaining the catalogue discrepancy.\par
\noindent\textit{Evidence.} The scan clearly writes 24.2.86; that day was Monday and follows Beilinson\textquotesingle s 17 Feb 86 talk.\par
\subsection*{P001258 (both; ambiguity): ADJ-667 moduli object}
\noindent\textit{Transcription reads.} Triple(C,p,x) is called Teichmueller space\par
\noindent\textit{Critical reading.} Formal-parameter pointed moduli, or its pullback to a marked Teichmueller cover\par
\noindent\textit{Evidence.} No topological marking is included in the triple; Kontsevich 1987 section 1.\par
\subsection*{P001258 (both; faithful error): ADJ-667 tangent quotient}
\noindent\textit{Transcription reads.} \(\mathfrak g/\mathfrak k\)\par
\noindent\textit{Critical reading.} \(\mathfrak g/\gamma\) for the full formal-parameter moduli\par
\noindent\textit{Evidence.} The next display gives the correct global stabilizer. Further quotient by k gives unpointed moduli.\par
\subsection*{P001258 (both; ambiguity): ADJ-667 marginal extension}
\noindent\textit{Transcription reads.} \ru{всю C!!}\par
\noindent\textit{Critical reading.} Allow meromorphic pole at \(p\); do not demand holomorphic extension through \(p\)\par
\noindent\textit{Evidence.} Gamma was defined by possible pole at \(p\); whole-curve holomorphic vector fields vanish for genus \(\ge2\).\par
\subsection*{P001259 (both; faithful error): ADJ-668 Euclidean sign}
\noindent\textit{Transcription reads.} Positive exponential of positive kinetic energy\par
\noindent\textit{Critical reading.} With positive kinetic energy use exp(-S), or explicitly declare a different signature/convention\par
\noindent\textit{Evidence.} The subsequent Gaussian determinant computation needs a convergent Gaussian before analytic continuation.\par
\subsection*{P001259 (both; faithful error): ADJ-668 derivative vertex}
\noindent\textit{Transcription reads.} A target delta derivative gives a world-sheet derivative times an exponential\par
\noindent\textit{Critical reading.} Target derivative gives \(-ip_\mu e^{ipX}\) in the stated Fourier convention; world-sheet derivative is another insertion\par
\noindent\textit{Evidence.} Elementary Fourier identity; the incomplete chiral factor is not yet a scalar integrated closed-string vertex.\par
\subsection*{P001259 (both; faithful error): ADJ-668 Gaussian factor}
\noindent\textit{Transcription reads.} Unprimed scalar determinant; all-pairs exponent with unspecified sign/coefficient\par
\noindent\textit{Critical reading.} Use det-prime and isolate the constant mode; fix covariance/action convention for sign and factor 1/2\par
\noindent\textit{Evidence.} Belavin--Knizhnik 1986 section 6 includes momentum conservation.\par
\subsection*{P001259 (both; ambiguity): ADJ-668 Green regularization}
\noindent\textit{Transcription reads.} Smooth the diagonal\par
\noindent\textit{Critical reading.} Choose a specified Wick/finite-part prescription and separately check integrated collisions\par
\noindent\textit{Evidence.} Off-diagonal Gaussian correlations do not imply convergence of the full insertion/moduli integrals.\par
